% Generated by build_chapter.py from the GY1978 files. Do not edit by hand; edit the files and rebuild.
\chapter{The academic year 1978--79}\label{ch:1978}
This chapter follows the seminar from October 1978 to April 1979 through
surreal arithmetic, classical and quantum integrability, functional
integration, geometric representation theory, twistors, superfields,
differential Lie algebras and webs. It covers all fourteen commissioned
catalogue records, including the year divider and reading plan, and adds
one source-established record: D.~B.~Fuchs on 5 March 1979. The latter is
represented only by the supplied pages P000824--P000825 and P000829;
P000826--P000828 are unavailable. The two final supplied pages,
P000841--P000842, continue R.~Kallosh's September 1979 lecture and are
documented as a boundary rather than assigned to Shabat. The resulting
inventory contains 15 records and 216 distinct mathematical items.

\paragraph{What the year contains.}
Kruskal moves from Conway's established field operations to the exponential
construction later developed systematically by Gonshor. Polyakov and
Zakharov connect nonlinear equations, spectral auxiliary systems, loop
variables and variational principles; Schwarz examines the measure and
determinants needed for an instanton approximation. Guichardet's
extensions and cohomology connect representation theory with homogeneous
geometry. Belavin's lectures organize factorized scattering and transfer
matrices, while S.~I.~Gelfand and Beilinson explain the ambitwistor
interpretation of Yang--Mills. Fuchs's supplied fragments, Ogievetsky's
superfields, Lebedev's local Lie structures and Shabat's webs give
contrasting examples of how algebraic structures encode geometry.

The chronology matters. A journal date after a talk does not by itself
make an item class A: earlier announcements or circulating preprints can
make it exposition. Conversely, a contemporary assertion of a theorem
does not establish that its proof was complete. The web algebraization
entry separates the 1978 assertion, the 1981 correction and Tr\'epreau's
complete theorem, first circulated in 2005 and published in 2006
\citep[pp.~78--79]{cherngriffiths1981corrections}
\citep[Thm.~1.1]{trepreau2006algebrisation}.
Likewise, the later anisotropic spin-one model answers a corrected nearby
question, not the literal density printed in the notebook. Surreal
integration illustrates a different relation: the modern resurgent
extension theory implements a broad earlier idea within stated function
classes, without validating the erroneous Gaussian series
\citep[Thms.~78 and 85]{costinehrlich2024integration}.

\begin{center}
\begin{tabular}{cl r}
\toprule
Class & Historical relation & Items\\
\midrule
A & Published soon afterward by the speaker or Gelfand & 10\\
B & Subsequently developed by others & 11\\
C & Precisely identified, partly developed & 0\\
D & Precisely identified, not developed & 0\\
E & Known before the talk & 168\\
F & Intended result or question not identifiable & 27\\
\bottomrule
\end{tabular}
\end{center}

\paragraph{Problems and critical developments.}
The problem index contains 5 definite questions, including solved
problems and partial matches whose differences are retained. No item was
assigned C or D on the evidence checked, and therefore no Stage~2 research
problem met the selection rule. This is a result of this bounded
identification exercise, not a claim that the subjects contain no open
problems. Aesthetic requests, unspecified algebras and broad universality
remarks remain F rather than acquiring invented quantifiers.
The developments are critical annotations: complete local proofs and
independently checked computations concerning a Gaussian remainder,
omnific fractions, projected contour areas, the printed contact bracket
and a necessary condition for a regular spin-chain R-matrix. They carry
no claim of a new solution to an open problem.

\paragraph{Connections and source boundaries.}
Earlier-year entries about instantons, inverse scattering and local Lie
structures provide genuine thematic links. Several apparent recurrences
are instead duplicate indexing of these very scans: 1977-445.1--4,
parts of 1977-451, and the composite 1977-457 block. Their migration map
in the report preserves old identifiers without counting them as earlier
lectures. The field \emph{reconstruction} makes supplied hypotheses
visible; literal correctness, publication match and verification level
remain separate from historical class. Each reference's audit annotation
records whether the result itself or only its metadata was read. Status
statements are dated 10 November 2026.

\GStalk{446}{1978-10-30}{M. D. Kruskal (catalogue: V. S. Kruskal)}{P000786-P000790}{Surreal numbers: arithmetic, exponentiation and asymptotics}
The manuscript names only \ru{Крускал}; the catalogue's initials V.~S. are retained as a discrepant authority field. Identification with Martin David Kruskal is a strong contextual inference: the subject and construction agree with the work explicitly attributed to M.~D.~Kruskal by \citet[Introduction and Chap.~10]{gonshor1986surreal}. The lecture begins with Conway's book and Knuth's mathematical novel (\citet{knuth1974surreal}), then constructs surreal numbers by ordinal sign sequences, orders them, defines arithmetic by simplest cuts, and explains infinite numbers, infinitesimals and natural ordinal arithmetic. The final page sketches exponentiation and proposes asymptotic integration as an application. Basic arithmetic is exposition; the exponential construction belongs to the subsequently published Kruskal--Gonshor development. The Gaussian formula and integer criterion require mathematical emendation. The entries 1977-445.1--4 in the supplied previous index describe the same opening material: these are overlapping-source aliases, not evidence of another earlier lecture. The next dated heading, Polyakov on 13 November, begins on P000791; the catalogue's overlapping page range must not be read as a talk boundary.

\begin{GSitem}{1978-446.1}{E}{Ordinal sign sequences and their numerical order}
\GSfield{Kind; source}{definition; \GSpage{P000786}}
\GSfield{The notes}{A number is a map from a well-ordered set to the two arrows \(\uparrow,\downarrow\); unequal numbers are compared lexicographically. The tree includes \(0,\uparrow,\uparrow\downarrow\).}
\GSfield{Statement}{Work in a class theory such as NBG. A surreal is a function \(s:\alpha\to\{-,+\}\), with \(\alpha\) an ordinal. Identify well-ordered domains by their order type. At the first differing position compare with \(-<\text{undefined}<+\). This gives a total order on the proper class \(\mathbf{No}\); the empty sequence is zero.}
\GSfield{Reconstruction}{The undefined-sign convention and passage from an arbitrary well-order to its ordinal order type are supplied explicitly.}
\GSfield{Fate}{Conway's construction predates the talk (book 1976; Knuth's exposition 1974). The sign-sequence theorem and its order are explicitly recorded in \citet[\S2.1]{kuhlmannmatusinski2015classes}; the early book is \citet{conway1976numbers}. The supplied alias 1977-445.1 is the same source occurrence.}
\GSfield{Present status}{As of 10 November 2026, this is an established part of surreal arithmetic under the hypotheses stated above; no open problem is asserted here.}
\GSfield{Assessment}{literal statement: correct with missing hypotheses; publication match: exact; verification: read the result; development: none; confidence: high}
\end{GSitem}
\begin{GSitem}{1978-446.2}{E}{Simplicity, truncations and the canonical cut}
\GSfield{Kind; source}{definition; \GSpage{P000786-P000787}}
\GSfield{The notes}{\(x\) is earlier than \(y\) when it lies on the path from zero to \(y\). The symbols \(y]\), \(y\rfloor\), \(y\rceil\) denote earlier numbers, respectively all, smaller, and greater. Theorem~2 writes \(x=\{x\rfloor:x\rceil\}\).}
\GSfield{Statement}{Write \(a<_s x\) when \(a\) is a proper initial segment of \(x\). Set \(L_s(x)=\{a<_s x:a<x\}\) and \(R_s(x)=\{a<_s x:a>x\}\). Then \(x=\{L_s(x)\mid R_s(x)\}\), the simplest number strictly between these two sets. Both option sets are sets, although \(\mathbf{No}\) is a proper class.}
\GSfield{Fate}{This is the canonical-cut theorem, already in Conway; \citet[Thm.~2.1(3), Def.~2.2]{kuhlmannmatusinski2015classes} gives the precise modern statement. Alias 1977-445.2 represents duplicated material.}
\GSfield{Present status}{As of 10 November 2026, this is an established part of surreal arithmetic under the hypotheses stated above; no open problem is asserted here.}
\GSfield{Assessment}{literal statement: correct; publication match: exact; verification: read the result; development: none; confidence: high}
\end{GSitem}
\begin{GSitem}{1978-446.3}{E}{The simplest separator of two sets}
\GSfield{Kind; source}{theorem; \GSpage{P000787}}
\GSfield{The notes}{For sets \(A,B\) with \(a<b\) for every \(a\in A,b\in B\), Theorem~1 asserts a number between them and a unique earliest such number, proved by transfinite induction.}
\GSfield{Statement}{For sets \(A,B\subset\mathbf{No}\), possibly empty, with \(A<B\), the class \(\{x:A<x<B\}\) has a unique simplest member \(\{A\mid B\}\). In particular it has unique minimal sign-sequence length. The claim is not extended to arbitrary proper-class cuts.}
\GSfield{Fate}{\citet[Thm.~2.1(1)]{kuhlmannmatusinski2015classes} states the exact existence theorem; it is part of the pre-1978 Conway theory. The previous index's 1977-445.3 is a source overlap.}
\GSfield{Present status}{As of 10 November 2026, this is an established part of surreal arithmetic under the hypotheses stated above; no open problem is asserted here.}
\GSfield{Assessment}{literal statement: correct; publication match: exact; verification: read the result; development: none; confidence: high}
\end{GSitem}
\begin{GSitem}{1978-446.4}{E}{Recursive addition}
\GSfield{Kind; source}{construction; \GSpage{P000787}}
\GSfield{The notes}{The definition is \(x+y=\{x\rfloor+y,x+y\rfloor:x\rceil+y,x+y\rceil\}\).}
\GSfield{Statement}{For canonical options of \(x,y\), define recursively
\[x+y=\{x^L+y,\ x+y^L\mid x^R+y,\ x+y^R\}.\]
Each family ranges over its entire option set. Induction on the well-founded product of the two simplicity orders makes the recursion well-defined.}
\GSfield{Fate}{The ordered-group construction is in \citet[Chap.~1]{conway1976numbers}, and the formula is reproduced in \citet[Thm.~2.3]{kuhlmannmatusinski2015classes}. Alias 1977-445.4 covers this same page.}
\GSfield{Present status}{As of 10 November 2026, this is an established part of surreal arithmetic under the hypotheses stated above; no open problem is asserted here.}
\GSfield{Assessment}{literal statement: correct; publication match: exact; verification: read the result; development: none; confidence: high}
\end{GSitem}
\begin{GSitem}{1978-446.5}{E}{Additive identities, order and sign reversal}
\GSfield{Kind; source}{theorem; \GSpage{P000787}}
\GSfield{The notes}{The notes prove \(0+0=0\), \(x+0=x\), assert commutativity and associativity, and obtain the opposite by reversing arrows. A sentence calls the sum the simplest operation monotone in each argument.}
\GSfield{Statement}{The operation makes \((\mathbf{No},+,0,<)\) an ordered abelian group: \(x+0=x\), \(x+y=y+x\), \((x+y)+z=x+(y+z)\), and \(x+(-x)=0\). Negation interchanges the two signs at every position. Translation preserves strict order. Simplicity characterizes the recursive cut, not an arbitrary pointwise monotone binary function.}
\GSfield{Reconstruction}{The dubious handwritten symbol in the associative-law display is read mathematically as a plus. The informal monotonicity sentence is restricted to its recursive context.}
\GSfield{Fate}{These are the standard field-operation theorems of \citet[Chap.~1]{conway1976numbers}; compare \citet[Thm.~2.3 and the following uniformity paragraph]{kuhlmannmatusinski2015classes}.}
\GSfield{Present status}{As of 10 November 2026, this is an established part of surreal arithmetic under the hypotheses stated above; no open problem is asserted here.}
\GSfield{Assessment}{literal statement: correct with missing hypotheses; publication match: exact; verification: read the result; development: annotated; confidence: high}
\end{GSitem}
\begin{GSitem}{1978-446.6}{E}{Two arithmetic exercises}
\GSfield{Kind; source}{question; \GSpage{P000788}}
\GSfield{The notes}{The exercises ask \(\uparrow\downarrow+\uparrow\downarrow=?\) and \(\widehat{\uparrow}+\downarrow\downarrow\downarrow=?\).}
\GSfield{Statement}{In sign notation \(+-=1/2\), hence \((+-)+(+-)=1\). With \(+^\omega=\omega\) and \(---=-3\), the second answer is \(\omega-3=+^\omega---\). These are exercise questions, not unresolved conjectures.}
\GSfield{Fate}{They follow directly from the recursive group law and the ordinal embedding, already established in \citet[Chaps.~1--2]{conway1976numbers}; \citet[\S2.1]{kuhlmannmatusinski2015classes} gives the conventions used here.}
\GSfield{Present status}{As of 10 November 2026, this is an established part of surreal arithmetic under the hypotheses stated above; no open problem is asserted here.}
\GSfield{Assessment}{literal statement: correct; publication match: exact; verification: read the result; development: annotated; confidence: high}
\end{GSitem}
\begin{GSitem}{1978-446.7}{E}{The basic multiplication inequality}
\GSfield{Kind; source}{claim; \GSpage{P000788}}
\GSfield{The notes}{For \(r<s\), \(u<v\), the overwritten sign in \(ru+sv\mathrel{?}rv+su\) resembles \(>\).}
\GSfield{Statement}{In an ordered field, if \(r<s\) and \(u<v\), then
\[ru+sv-rv-su=(s-r)(v-u)>0.\]
Thus the correct strict sign is \(>\). Positivity of \(r,s,u,v\) is not needed.}
\GSfield{Fate}{This is an elementary ordered-field identity underlying Conway's multiplication. Its role in the recursive product is described in \citet[\S2.2]{mantovamatusinski2017survey}; the displayed factorization independently checks the sign.}
\GSfield{Present status}{As of 10 November 2026, this is an established part of surreal arithmetic under the hypotheses stated above; no open problem is asserted here.}
\GSfield{Assessment}{literal statement: not determinable; publication match: exact; verification: read the result; development: annotated; confidence: high}
\end{GSitem}
\begin{GSitem}{1978-446.8}{E}{The four-family product cut}
\GSfield{Kind; source}{construction; \GSpage{P000788}}
\GSfield{The notes}{The notebook gives four expressions: same-side options on the left of the cut and opposite-side options on the right.}
\GSfield{Statement}{For canonical surreal options, let \(P(a,b)=ay+xb-ab\). Then
\[xy=\{P(x^L,y^L),P(x^R,y^R)\mid P(x^L,y^R),P(x^R,y^L)\}.\]
All indicated pairs of options are included. The induction uses smaller arguments, and the inequalities follow from the signs of \((x-a)(y-b)\).}
\GSfield{Fate}{This is precisely Conway's product, known before the talk; see \citet[Chap.~1]{conway1976numbers} and \citet[Thm.~2.3]{kuhlmannmatusinski2015classes}.}
\GSfield{Present status}{As of 10 November 2026, this is an established part of surreal arithmetic under the hypotheses stated above; no open problem is asserted here.}
\GSfield{Assessment}{literal statement: correct; publication match: exact; verification: read the result; development: none; confidence: high}
\end{GSitem}
\begin{GSitem}{1978-446.9}{E}{Ring laws and the multiplicative identities}
\GSfield{Kind; source}{theorem; \GSpage{P000788}}
\GSfield{The notes}{The notes derive \(x\cdot0=0\) and \(x\cdot1=x\), and assign associativity, commutativity and distributivity as homework.}
\GSfield{Statement}{For all surreals \(x,y,z\), the recursive operations satisfy \(xy=yx\), \((xy)z=x(yz)\), \(x(y+z)=xy+xz\), \(x0=0\) and \(x1=x\), with \(1=\{0\mid\varnothing\}\).}
\GSfield{Fate}{The established ordered-field theorem is the exact match: \citet[Chap.~1]{conway1976numbers}, restated in \citet[Thm.~2.3]{kuhlmannmatusinski2015classes}. A request to prove these facts at home is not a research problem.}
\GSfield{Present status}{As of 10 November 2026, this is an established part of surreal arithmetic under the hypotheses stated above; no open problem is asserted here.}
\GSfield{Assessment}{literal statement: correct; publication match: exact; verification: read the result; development: none; confidence: high}
\end{GSitem}
\begin{GSitem}{1978-446.10}{E}{Uniformity: independence of a cut representation}
\GSfield{Kind; source}{method; \GSpage{P000788}}
\GSfield{The notes}{The advice is to prove addition is ``bracket indifferent'': if \(x=\{a:b\}\), \(y=\{c:d\}\), use \(\{a+y,x+c:b+y,x+d\}\). The same property is claimed for multiplication.}
\GSfield{Statement}{For any set presentations \(x=\{A\mid B\}\), \(y=\{C\mid D\}\) as simplest cuts, the usual addition and multiplication formulas remain valid with options from \(A,B,C,D\), even when these are not the canonical options. The value is independent of the presentation.}
\GSfield{Reconstruction}{The singleton-looking letters in the notes are treated as generic members of option sets; the claim is not invariance under arbitrary changes of inequalities.}
\GSfield{Fate}{This is uniformity of the Conway operations, explicitly distinguished from associativity in \citet[paragraph after Thm.~2.3]{kuhlmannmatusinski2015classes}; it was already part of \citet[Chap.~1]{conway1976numbers}.}
\GSfield{Present status}{As of 10 November 2026, this is an established part of surreal arithmetic under the hypotheses stated above; no open problem is asserted here.}
\GSfield{Assessment}{literal statement: correct; publication match: exact; verification: read the result; development: none; confidence: high}
\end{GSitem}
\begin{GSitem}{1978-446.11}{E}{An infinite number and its reciprocal infinitesimal}
\GSfield{Kind; source}{construction; \GSpage{P000788-P000789}}
\GSfield{The notes}{\(\omega=\widehat{\uparrow}\), \(\varepsilon=\uparrow\widehat{\downarrow}\), \(\omega\varepsilon=1\), and \(\omega+\varepsilon=\widehat{\uparrow}\uparrow\widehat{\downarrow}\).}
\GSfield{Statement}{The sign sequences \(+^\omega\) and \(+-^\omega\) represent \(\omega\) and \(\omega^{-1}=\{0\mid2^{-n}:n\in\mathbb N\}\). Their product is one. The sum has \(\omega+1\) plus signs followed by \(\omega\) minus signs; equivalently its normal form is \(\omega+\omega^{-1}\). Here hats indicate an \(\omega\)-long block, not an ordinal exponentiation rule.}
\GSfield{Reconstruction}{The lengths of the repeated sign blocks are made explicit; the notation is interpreted as concatenation.}
\GSfield{Fate}{Conway's infinite arithmetic predates 1978. The product calculation is written out in \citet[\S2.2, example before Thm.~2.6]{mantovamatusinski2017survey}; the normal-form theorem is also in \citet[Thm.~2.4]{kuhlmannmatusinski2015classes}.}
\GSfield{Present status}{As of 10 November 2026, this is an established part of surreal arithmetic under the hypotheses stated above; no open problem is asserted here.}
\GSfield{Assessment}{literal statement: correct; publication match: exact; verification: read the result; development: none; confidence: high}
\end{GSitem}
\begin{GSitem}{1978-446.12}{E}{Ordinals and Hessenberg natural arithmetic}
\GSfield{Kind; source}{theorem; \GSpage{P000789}}
\GSfield{The notes}{Ordinals have only up arrows. Ordinary ordinal addition is noncommutative; the note names ``Hessenburg'' and says his natural addition and multiplication coincide with surreal arithmetic.}
\GSfield{Statement}{The embedding \(\alpha\mapsto+^\alpha\) sends ordinal natural sum and natural product to surreal addition and multiplication. For Cantor normal forms, natural addition adds coefficients of equal powers and natural multiplication distributes with natural sums of exponents. Ordinal concatenation instead satisfies \(1+_{\rm ord}\omega=\omega\ne\omega+_{\rm ord}1\).}
\GSfield{Fate}{\citet[Chap.~2]{conway1976numbers} treats the ordinal embedding; \citet[\S2.1, end]{kuhlmannmatusinski2015classes} explicitly identifies the induced operations as natural sum and product. The surname is Hessenberg.}
\GSfield{Present status}{As of 10 November 2026, this is an established part of surreal arithmetic under the hypotheses stated above; no open problem is asserted here.}
\GSfield{Assessment}{literal statement: correct; publication match: exact; verification: read the result; development: annotated; confidence: high}
\end{GSitem}
\begin{GSitem}{1978-446.13}{E}{Division and positive square roots}
\GSfield{Kind; source}{theorem; \GSpage{P000789}}
\GSfield{The notes}{Every nonzero number has \(1/x\); \(y/x\) exists; every \(x>0\) has \(\sqrt{x}\), by transfinite induction.}
\GSfield{Statement}{For every nonzero surreal \(x\), there is a unique \(y\) with \(xy=1\). Every positive surreal has a unique positive square root. These are assertions about a proper-class real-closed ordered field, with the usual finite algebraic formulas.}
\GSfield{Fate}{The field and real-closedness results were already established by Conway; see \citet[Chaps.~1, 4]{conway1976numbers}. \citet[Thm.~2.3]{kuhlmannmatusinski2015classes} states the combined theorem precisely.}
\GSfield{Present status}{As of 10 November 2026, this is an established part of surreal arithmetic under the hypotheses stated above; no open problem is asserted here.}
\GSfield{Assessment}{literal statement: correct; publication match: exact; verification: read the result; development: none; confidence: high}
\end{GSitem}
\begin{GSitem}{1978-446.14}{E}{Polynomial intermediate values do not eliminate gaps}
\GSfield{Kind; source}{claim; \GSpage{P000789}}
\GSfield{The notes}{A polynomial taking opposite signs at two points has a zero between them; nevertheless there are gaps, for example between finite and infinite numbers. A class-long arrow is said no longer to be a surreal.}
\GSfield{Statement}{If \(p\in\mathbf{No}[T]\) has finite degree and \(a<b\) satisfy \(p(a)p(b)<0\), some \(c\in(a,b)\) satisfies \(p(c)=0\). This does not imply Dedekind completeness. The cut with all finite surreals on the left and all positive infinite surreals on the right is unfilled: a finite candidate can be increased by one, and a positive infinite candidate can be halved. There is no ordinal-length sign sequence with proper-class domain.}
\GSfield{Reconstruction}{Finite degree, the coefficient field, and the precise finite/infinite cut are supplied. The sentence is not a denial of polynomial order-continuity.}
\GSfield{Fate}{Polynomial intermediate values follow from the pre-existing real-closedness theorem, \citet[Chap.~4]{conway1976numbers}; \citet[Thm.~2.3]{kuhlmannmatusinski2015classes} verifies the relevant algebraic hypothesis. The gap argument is supplied here to disambiguate ``without continuity''.}
\GSfield{Present status}{As of 10 November 2026, this is an established part of surreal arithmetic under the hypotheses stated above; no open problem is asserted here.}
\GSfield{Assessment}{literal statement: correct with missing hypotheses; publication match: exact; verification: read the result; development: annotated; confidence: high}
\end{GSitem}
\begin{GSitem}{1978-446.15}{E}{Periodic signs and two rational examples}
\GSfield{Kind; source}{claim; \GSpage{P000789}}
\GSfield{The notes}{The notebook records \(\widehat{\uparrow\downarrow}=2/3\) and \(\uparrow\downarrow\uparrow=3/4\).}
\GSfield{Statement}{The finite sequence \(+-+\) denotes \(3/4\). The \(\omega\)-long alternating sequence \(+-+-\cdots\) denotes \(2/3\): its finite truncations alternate about \(1-1/2+1/4-\cdots=2/3\), and their cuts determine that real number.}
\GSfield{Fate}{These are examples of the real embedding in Conway's theory, \citet[Chap.~2]{conway1976numbers}; compare the sign representation of real numbers in \citet[\S2.1]{kuhlmannmatusinski2015classes}. No new irrationality or summation question is posed.}
\GSfield{Present status}{As of 10 November 2026, this is an established part of surreal arithmetic under the hypotheses stated above; no open problem is asserted here.}
\GSfield{Assessment}{literal statement: correct; publication match: exact; verification: read the result; development: none; confidence: high}
\end{GSitem}
\begin{GSitem}{1978-446.16}{B}{Exponentiation, logarithms and real-base powers}
\GSfield{Kind; source}{construction; \GSpage{P000789-P000790}}
\GSfield{The notes}{The notes say that \(\exp x\), \(x^y\) and \(\log\) exist, then begin a recursive construction.}
\GSfield{Statement}{The Kruskal--Gonshor exponential is an order-preserving group isomorphism \(\exp:(\mathbf{No},+)\to(\mathbf{No}_{>0},\cdot)\), agreeing with the real exponential. Its inverse is \(\log\). For \(x>0\) define \(x^y=\exp(y\log x)\). This power operation must be distinguished from the normal-form monomial map \(a\mapsto\omega^a\).}
\GSfield{Reconstruction}{Positivity of the base and the distinction between ordinary powers and the monomial map are supplied.}
\GSfield{Fate}{\citet[Introduction]{gonshor1986surreal} dates his first exposure to Kruskal's ideas to January 1977; Chap.~10, published in 1986, develops exponentiation from Kruskal's hints. Thus B records subsequent development by Gonshor, not discovery in this particular 1978 talk or publication soon afterward. The exact map is stated in \citet[Thm.~2.7]{kuhlmannmatusinski2015classes}.}
\GSfield{Present status}{As of 10 November 2026, the exponential ordered-field structure is established. The notebook does not specify a new open problem about it.}
\GSfield{Assessment}{literal statement: correct with missing hypotheses; publication match: exact; verification: read the result; development: none; confidence: high}
\end{GSitem}
\begin{GSitem}{1978-446.17}{B}{Exponential growth beyond every fixed polynomial}
\GSfield{Kind; source}{claim; \GSpage{P000790}}
\GSfield{The notes}{For positive infinite \(x\), \(2^x>x^n\) ``for any \(n\).}
\GSfield{Statement}{If \(x\in\mathbf{No}\) exceeds every real number and \(n\in\mathbb N\) is fixed, then \(\exp(x\log2)>x^n\). The quantifier is over ordinary finite integers, not arbitrary surreal exponents: taking \(n=x\) would reverse the inequality when \(x>2\).}
\GSfield{Reconstruction}{The range of n is restricted to ordinary finite integers.}
\GSfield{Fate}{This is an established property of the Kruskal--Gonshor exponential; \citet[Chap.~10]{gonshor1986surreal} is the definitive later treatment. \citet[\S2.3]{costinehrlich2024integration} states that \((\mathbf{No},\exp)\) is an elementary extension of the real exponential field, which yields the stated fixed-integer inequality.}
\GSfield{Present status}{As of 10 November 2026, established for finite polynomial degree. The unrestricted exponent reading is false.}
\GSfield{Assessment}{literal statement: correct with missing hypotheses; publication match: exact; verification: read the result; development: none; confidence: high}
\end{GSitem}
\begin{GSitem}{1978-446.18}{E}{The exponential series at an infinitesimal}
\GSfield{Kind; source}{theorem; \GSpage{P000790}}
\GSfield{The notes}{For infinitesimal \(x\), \(\exp x=\sum_{n\in\text{finite}}x^n/n!\).}
\GSfield{Statement}{For \(|\varepsilon|<1/n\) for every positive integer \(n\),
\[\exp(\varepsilon)=\sum_{n=0}^{\infty}\varepsilon^n/n!,\]
where the sum is Hahn/Conway summation of normal forms. It is not a claim that ordinary sequence convergence in the full order topology is sufficient.}
\GSfield{Reconstruction}{The summation sense and starting index zero are supplied.}
\GSfield{Fate}{The local analytic construction is part of Conway's pre-1978 power-series calculus; its agreement with the global exponential is \citet[Chap.~10]{gonshor1986surreal} and \citet[Prop.~8(ii)]{costinehrlich2024integration}.}
\GSfield{Present status}{As of 10 November 2026, this is an established part of surreal arithmetic under the hypotheses stated above; no open problem is asserted here.}
\GSfield{Assessment}{literal statement: correct with missing hypotheses; publication match: exact; verification: read the result; development: none; confidence: high}
\end{GSitem}
\begin{GSitem}{1978-446.19}{E}{Arbitrary real-coefficient power series at infinitesimals}
\GSfield{Kind; source}{claim; \GSpage{P000790}}
\GSfield{The notes}{The line reads ``\(x\) infinitesimal, \(n\) real'' followed by \(\sum x^n a_n\) convergent for all \(a_n\); the placement of the coefficient qualifier is doubtful.}
\GSfield{Statement}{For a sequence \((a_n)_{n\ge0}\) of real coefficients and a surreal infinitesimal \(\varepsilon\), the family \((a_n\varepsilon^n)_{n\ge0}\) is summable in Conway's normal-form sense. No analytic radius of convergence is required. Neither arbitrary real-indexed sums nor arbitrary surreal coefficients satisfy this assertion without extra summability hypotheses.}
\GSfield{Reconstruction}{The intended reading is \(a_n\) real and \(n\) a nonnegative integer; this is a mathematical reconstruction, not a claimed certain paleographic correction.}
\GSfield{Fate}{\citet[Chap.~4, p.~40]{conway1976numbers} is the pre-talk result. It is stated verbatim mathematically as a formal-power-series theorem in \citet[Prop.~4]{costinehrlich2024integration}. For a counterexample to unrestricted coefficients, choose nonzero \(\varepsilon\) and \(a_n=\varepsilon^{-n}\): every term is one.}
\GSfield{Present status}{As of 10 November 2026, this is an established part of surreal arithmetic under the hypotheses stated above; no open problem is asserted here.}
\GSfield{Assessment}{literal statement: correct with missing hypotheses; publication match: exact; verification: read the result; development: annotated; confidence: high}
\end{GSitem}
\begin{GSitem}{1978-446.20}{B}{Asymptotic integration and an incorrect Gaussian formula}
\GSfield{Kind; source}{method; \GSpage{P000790}}
\GSfield{The notes}{A possible application is written \(\int_x^{\text{class-long }\uparrow}e^{-\xi^2}\,d\xi=\sum n!/x^n\). The scan confirms the Gaussian exponent and the displayed factorial series; no scale factor or alternating signs are visible.}
\GSfield{Statement}{For real \(x>0\), the corrected classical expansion is
\[\int_x^\infty e^{-t^2}\,dt\sim\frac{e^{-x^2}}{2x}\sum_{n\ge0}\frac{(-1)^n(2n-1)!!}{2^n x^{2n}}.\]
At a positive infinite surreal argument, substitution of \(x^{-1}\) makes the formal power series summable; an equality to a surreal integral additionally requires a specified extension and integration operator. The factorial series in the notebook is not that Gaussian expansion.}
\GSfield{Reconstruction}{The corrected Gaussian asymptotic series is our explicit mathematical repair. It does not recover an unwritten 1978 definition of the integral.}
\GSfield{Fate}{The broader asymptotic-integration method was subsequently developed by \citet[Defs.~68, 76; Thms.~78, 85]{costinehrlich2024integration}: preprint 30 August 2022, revised July 2024, journal 2024. Their result concerns a specified class of resurgent functions and their extensions, not all smooth functions; \S13 gives foundational obstructions to broader constructive operators. This is a partial match to the proposed method, not a proof of the literal formula.}
\GSfield{Present status}{As of 10 November 2026, the specified resurgent integration theory is developed. The notebook leaves no precise function-class or axiomatic open problem; no new problem-index entry is manufactured. A complete finite-remainder check is supplied in the critical annotations.}
\GSfield{Assessment}{literal statement: false as written; publication match: partial; verification: read the result; development: annotated; confidence: high}
\end{GSitem}
\begin{GSitem}{1978-446.21}{E}{Omnific integers and an insufficient truncation criterion}
\GSfield{Kind; source}{definition; \GSpage{P000790}}
\GSfield{The notes}{The notes write \(\widehat{\uparrow}\downarrow=\omega-1\) and \(|x]-x|=1\) as the definition of an integer.}
\GSfield{Statement}{The relevant integers are the omnific integers \(\mathbf{Oz}\), equivalently the surreals \(z=\{z-1\mid z+1\}\). In normal form they have no negative exponents and the coefficient at exponent zero is an ordinary integer. The formula with an arbitrary earlier truncation is not an equivalent criterion: for \(x=3\), its truncation zero has distance three.}
\GSfield{Reconstruction}{The Oz interpretation is required by the adjacent quotient theorem. The notes do not say which truncation, if any, should be selected.}
\GSfield{Fate}{Omnific integers were already treated in \citet[Chap.~5]{conway1976numbers}. The book's opened excerpt is insufficient to recheck this chapter directly; the quotient assertion is independently verified from the stated normal-form characterization in the critical annotations, and the historical attribution remains indirect.}
\GSfield{Present status}{As of 10 November 2026, this is an established part of surreal arithmetic under the hypotheses stated above; no open problem is asserted here.}
\GSfield{Assessment}{literal statement: false as written; publication match: partial; verification: indirect; development: annotated; confidence: medium}
\end{GSitem}
\begin{GSitem}{1978-446.22}{E}{Every surreal is a quotient of omnific integers}
\GSfield{Kind; source}{theorem; \GSpage{P000790}}
\GSfield{The notes}{\ru{Всякое число} is a ratio \(m/n\) of two integers.}
\GSfield{Statement}{For every \(x\in\mathbf{No}\), there exist \(m,n\in\mathbf{Oz}\), \(n\ne0\), with \(x=m/n\). Indeed choose a positive surreal \(a\) exceeding the negatives of every exponent in the set-sized support of the normal form of \(x\). Then \(n=\omega^a\) and \(m=\omega^a x\) have only positive exponents and are omnific integers. For zero use \(0/1\). The assertion with ordinary integers is false.}
\GSfield{Reconstruction}{The qualifier omnific is supplied from the theory; it is essential. The short proof is an independent verification of the intended assertion.}
\GSfield{Fate}{This is the field-of-fractions result of Conway's \emph{Number Theory in the Land of Oz}, \citet[Chap.~5]{conway1976numbers}; the normal-form theorem needed for the displayed verification is read in \citet[Thm.~2.4]{kuhlmannmatusinski2015classes}. No claim of novelty is attached to the proof.}
\GSfield{Present status}{As of 10 November 2026, this is an established part of surreal arithmetic under the hypotheses stated above; no open problem is asserted here.}
\GSfield{Assessment}{literal statement: correct with missing hypotheses; publication match: exact; verification: read the result; development: annotated; confidence: high}
\end{GSitem}
\begin{GSitem}{1978-446.23}{B}{The unfinished recursive exponential cut}
\GSfield{Kind; source}{construction; \GSpage{P000790}}
\GSfield{The notes}{The last formula starts
\[\exp x=\{0,\exp x]\sum_{m=1}^n(x-x])^m/m!:\exp x]/\sum_{m=0}^n(x]-x)^m/m!\]
with no final brace and undifferentiated earlier options.}
\GSfield{Statement}{A valid full recursive definition must separate left and right options, include the constant term in Taylor polynomials \(E_n(t)=\sum_{m=0}^nt^m/m!\), and retain only positive reciprocal denominators. The Kruskal--Gonshor construction uses two families on each side of the cut, besides zero, and yields the exponential of item 1978-446.16. The printed two-family fragment is not a standalone total definition.}
\GSfield{Reconstruction}{Completion requires more than closing delimiters. The canonical modern construction is identified separately, without silently inserting it into the diplomatic transcript.}
\GSfield{Fate}{The complete formula is given in \citet[Chap.~10]{gonshor1986surreal} and \citet[\S2.3, recursive exponential definition]{mantovamatusinski2017survey}. This is a partial mathematical match: no claim is made that the missing terms were actually spoken. For example, at \(x=1\), the earlier option zero and \(n=1\) give a reciprocal denominator \(E_1(-1)=0\) in the unqualified notebook rule.}
\GSfield{Present status}{As of 10 November 2026, the completed Kruskal--Gonshor construction is established; the literal notebook recursion remains incomplete.}
\GSfield{Assessment}{literal statement: false as written; publication match: partial; verification: read the result; development: annotated; confidence: high}
\end{GSitem}

\GStalk{447}{1978-11-13}{A. M. Polyakov}{P000791-P000798}{Nonabelian Toda chains, loop-space gauge fields, and a proposed area law}
The lecture begins on P000791, after the preceding meeting’s material, and ends on P000798. Its first part passes from a lattice principal chiral model to the matrix Toda chain, a difference-operator Lax equation, formal transfer matrices, and conserved coefficients. The second part treats lattice curvature and holonomy as a route to a chiral-field description of Yang--Mills fields on loop space. A functional flatness equation, conservation equation, prospective higher currents, and a spectral linear problem are written, but the measure in the proposed loop-space action is expressly left undefined. The final part seeks a string analogue of the Klein--Gordon Green function and an area-law asymptotic, then reduces the contour dependence to its six projected areas in four dimensions. The proposed self-dual area expression must be distinguished from the minimum area of an arbitrary spanning surface.\par The matrix Toda discussion was subsequently developed by Krichever and others; a modern primary survey explicitly attributes its unpublished origin to Polyakov \citep[\S1, equation (1.3), and \S9]{takebezabrodin2025multicomponent}. Polyakov’s 1979 publication identifies the loop-space programme, but the available author-hosted publication record verifies its abstract and chronology, not every formula in this notebook \citep[abstract]{polyakov1979string}. The projected-area proposal therefore receives a mathematical audit instead of being silently promoted to a general Plateau theorem.

\begin{GSitem}{1978-447.1}{E}{A two-dimensional matrix lattice action}
\GSfield{Kind; source}{construction; \GSpage{P000791}}
\GSfield{The notes}{The notes take \(g_x\in GL(n,\mathbb R)\), \(S=\operatorname{Tr}\sum_{x,\delta}g_x^{-1}g_{x+\delta}\), and write \(\sum_\delta(g_{x+\delta}g_x^{-1}-g_xg_{x-\delta}^{-1})=0\).}
\GSfield{Statement}{On a finite two-dimensional lattice with periodic boundary conditions, or with variations supported away from its boundary, vary each invertible matrix by \(\delta g_x=\epsilon_xg_x\). Stationarity of the displayed oriented nearest-neighbour action is equivalent to \(\sum_\delta(A_{x\delta}-A_{x-\delta,\delta})=0\), where \(A_{x\delta}=g_{x+\delta}g_x^{-1}\). The two positive lattice directions are summed. This is an algebraic Euler--Lagrange identity; positivity of the action is not asserted for \(GL(n,\mathbb R)\).}
\GSfield{Reconstruction}{The finite-lattice or compact-support condition makes the variation legitimate. It is not stated in the notebook. The action and its equation can be verified directly without assuming integrability.}
\GSfield{Fate}{The continuum principal-chiral-field setting was already in print in \citet[\S2]{zakharovmikhailov1978relativistic}, submitted 4 January 1978. That is a partial model match, not a citation proving the exact fully discrete action integrable.}
\GSfield{Present status}{As of 10 November 2026: The variational identity is valid. The stronger integrability assertion is assessed separately.}
\GSfield{Assessment}{literal statement: correct with missing hypotheses; publication match: partial; verification: read the result; development: annotated; confidence: high}
\end{GSitem}
\begin{GSitem}{1978-447.2}{F}{Complete integrability and higher discrete currents}
\GSfield{Kind; source}{claim; \GSpage{P000791}}
\GSfield{The notes}{The system is called completely integrable, with an \(L,A\)-pair and infinitely many currents satisfying \(\sum_\delta(J^{(k)}_{x\delta}-J^{(k)}_{x-\delta,\delta})=0\). The first current is \(A\); the second ends after \(\varepsilon_{\delta\gamma}A_{x\gamma}+\).}
\GSfield{Statement}{The first conservation law is exactly the preceding lattice field equation. A complete theorem for higher currents would have to specify the remainder in \(J^{(2)}\), the recursion, boundary conditions, and the meaning of independence or completeness of the conserved quantities. Those data are absent.}
\GSfield{Reconstruction}{No missing second-current term is invented. The semi-discrete matrix Toda result below and the continuum principal chiral field do not by themselves prove the fully discrete assertion.}
\GSfield{Fate}{\citet[\S2]{zakharovmikhailov1978relativistic} treats the continuum model, while \citet[\S1 and \S9]{takebezabrodin2025multicomponent} treats matrix Toda hierarchies. Neither checked passage identifies the unfinished current for precisely this lattice action.}
\GSfield{Present status}{As of 10 November 2026: The first current is established; the claimed complete hierarchy is not identifiable from the surviving formula. This is an unresolved identification, not evidence that a precise open problem remains.}
\GSfield{Assessment}{literal statement: not determinable; publication match: thematic only; verification: unresolved lead; development: annotated; confidence: medium}
\end{GSitem}
\begin{GSitem}{1978-447.3}{B}{The semi-discrete nonabelian Toda chain}
\GSfield{Kind; source}{construction; \GSpage{P000791}}
\GSfield{The notes}{With \(A_s=g_s^{-1}\dot g_s\), \(B_s=g_s^{-1}g_{s+1}\), the equations are \(\dot B_s=-A_sB_s+B_sA_{s+1}\), \(\dot A_s=B_s-B_{s-1}\). A kinetic term is written \(\dot g_s^{-1}\dot g_s\), without parentheses or a factor \(1/2\).}
\GSfield{Statement}{Let \(g_s(t)\in GL(n,\mathbb C)\) be twice differentiable, with a finite periodic chain or locally finite variations. The first equation follows identically from the definitions. The second is the nonabelian Toda equation \(\frac{d}{dt}(g_s^{-1}\dot g_s)=g_s^{-1}g_{s+1}-g_{s-1}^{-1}g_s\). A compatible Lagrangian is \(\sum_s\operatorname{Tr}(\tfrac12 A_s^2-B_s)\). For \(n=1\), \(g_s=e^{q_s}\) gives \(\ddot q_s=e^{q_{s+1}-q_s}-e^{q_s-q_{s-1}}\).}
\GSfield{Reconstruction}{The normalized Lagrangian is a reconstruction. Interpreting the notebook literally as \(\frac{d}{dt}(g_s^{-1})\dot g_s\) gives \(-A_s^2\), not the stated kinetic normalization.}
\GSfield{Fate}{\citet[\S1, equation (1.3), references 34 and 45, and \S9]{takebezabrodin2025multicomponent} explicitly attributes the unpublished matrix Toda equation and its integrals to Polyakov, and records Krichever’s 1981 periodic/quasiperiodic solution theory and Korepin’s 1981 Hamiltonian quantization.}
\GSfield{Present status}{As of 10 November 2026: The matrix Toda system is an established integrable hierarchy. The scalar reduction and the displayed first-order system are precise; the notebook Lagrangian needs normalization.}
\GSfield{Assessment}{literal statement: false as written; publication match: exact; verification: read the result; development: annotated; confidence: high}
\end{GSitem}
\begin{GSitem}{1978-447.4}{B}{A spectral zero-curvature representation for matrix Toda}
\GSfield{Kind; source}{problem; \GSpage{P000791-P000792}}
\GSfield{The notes}{The explicit request is to find \(L_s(A,B,\lambda)\), \(M_s(A,B,\lambda)\) such that \(\dot L_s+M_sL_s-L_sM_{s+1}=0\) is equivalent to the Toda equations.}
\GSfield{Statement}{For the matrix Toda system, construct a spectral auxiliary linear system whose compatibility yields both first-order equations. One concrete realization uses the spectral equation for the block Jacobi operator \(\mathcal L=E^{-1}+A+BE\), with \((Ef)_s=f_{s+1}\), and time evolution generated by \(\mathcal M=BE\); compatibility is \(\dot{\mathcal L}=[\mathcal M,\mathcal L]\). Rewriting its second-order difference equation as a first-order block system gives a transfer-matrix representation, with a convention-dependent sign in the time equation.}
\GSfield{Reconstruction}{The distinction between transfer matrices and difference operators is added. The notebook uses \(L_s\) for both, although their zero-curvature equations are not literally identical.}
\GSfield{Fate}{The matrix Toda solution theory is explicitly traced to Krichever’s 1981 work in \citet[\S1 and \S9]{takebezabrodin2025multicomponent}. The operator calculation here also verifies directly the requested compatibility at the formal algebraic level.}
\GSfield{Present status}{As of 10 November 2026: The Lax-representation problem is resolved; analytic spectral theory depends on finite, periodic, or infinite-chain boundary conditions.}
\GSfield{Assessment}{literal statement: correct with missing hypotheses; publication match: exact; verification: read the result; development: annotated; confidence: high; problem \hyperref[prob:GP-1978-01]{GP-1978-01}}
\end{GSitem}
\begin{GSitem}{1978-447.5}{E}{Conservation of a transfer product and its endpoint terms}
\GSfield{Kind; source}{claim; \GSpage{P000792}}
\GSfield{The notes}{The notebook writes \(\frac{d}{dt}\prod_{s=-\infty}^{\infty}L_s=0\) and calls the product a generating function of integrals; its diagram omits endpoint transport.}
\GSfield{Statement}{For finitely many ordered matrices satisfying \(\dot L_s=L_sM_{s+1}-M_sL_s\), the product \(T_{a,b}=L_a\cdots L_b\) satisfies \(\dot T_{a,b}=T_{a,b}M_{b+1}-M_aT_{a,b}\). If the endpoint matrices vanish, \(T\) is constant. If they agree, \(T\) evolves by conjugation and \(\operatorname{Tr}(T^k)\) is constant. An infinite product additionally needs convergence, a normalization, and control of the endpoint limit.}
\GSfield{Reconstruction}{The telescoping identity supplies the missing hypotheses; equality of endpoint matrices alone does not make the matrix product constant.}
\GSfield{Fate}{The analogous monodromy mechanism already appears in the inverse-scattering framework of \citet[\S\S1--2]{zakharovmikhailov1978relativistic}. The particular endpoint identity is checked algebraically here.}
\GSfield{Present status}{As of 10 November 2026: The unconditional equality is false; its conserved-spectrum version is standard and valid under the stated conditions.}
\GSfield{Assessment}{literal statement: false as written; publication match: partial; verification: read the result; development: annotated; confidence: high}
\end{GSitem}
\begin{GSitem}{1978-447.6}{B}{A formal matrix Riccati recursion and conserved coefficients}
\GSfield{Kind; source}{method; \GSpage{P000792}}
\GSfield{The notes}{The proposed recursion is \(L_s=1+\lambda A_s-\lambda^2B_sL_{s+1}^{-1}\). Expanding its ordered product produces \(\sum_s A_s\) and \(\sum_{s<s^{\prime}}A_sA_{s^{\prime}}-\sum_sB_s\); the latter is called a nontrivial integral.}
\GSfield{Statement}{In formal power series with constant term the identity, the recursion determines coefficients successively: \(L_s=I+\lambda A_s-\lambda^2B_s+O(\lambda^3)\). For a finite product in increasing \(s\), its coefficient of \(\lambda^2\) is exactly the displayed ordered expression. Conservation of a coefficient follows only when the associated product is conserved, or after taking conjugacy-invariant functions in the periodic case. An infinite-chain coefficient requires summability or subtraction of a background.}
\GSfield{Reconstruction}{Formal-power-series interpretation, remainder terms, ordering, and the distinction between matrix coefficients and scalar invariants are supplied.}
\GSfield{Fate}{This belongs to the matrix Toda construction subsequently developed by others; see \citet[\S1 and \S9]{takebezabrodin2025multicomponent}. The two coefficients are directly checkable, but the checked source does not establish the notebook’s exact recursion with its boundary normalization.}
\GSfield{Present status}{As of 10 November 2026: The displayed low-order expansion is correct after reinstating remainder terms; unqualified infinite-product conservation is not.}
\GSfield{Assessment}{literal statement: correct with missing hypotheses; publication match: partial; verification: read the result; development: annotated; confidence: high}
\end{GSitem}
\begin{GSitem}{1978-447.7}{B}{The block Jacobi Lax equation}
\GSfield{Kind; source}{theorem; \GSpage{P000792}}
\GSfield{The notes}{Equation (14) gives \(L=E^{-1}+A_s+B_sE\), \(M=B_sE\), where \(E\) shifts by one site.}
\GSfield{Statement}{In the algebra of matrix-valued difference operators, use \(EC_s=C_{s+1}E\). Then \([BE,E^{-1}+A+BE]=(B_s-B_{s-1})+(B_sA_{s+1}-A_sB_s)E\). Thus \(\dot L=[M,L]\) is equivalent to the two matrix Toda equations. For a finite periodic block matrix, all coefficients of its characteristic polynomial are conserved.}
\GSfield{Reconstruction}{The commutator equation and finite periodic consequence are made explicit. For an infinite operator, isospectrality requires a suitable operator domain and well-posed evolution.}
\GSfield{Fate}{The identification with nonabelian Toda is supported by \citet[\S1 and \S9]{takebezabrodin2025multicomponent}; the complete formal commutator calculation is elementary and included in the statement.}
\GSfield{Present status}{As of 10 November 2026: The algebraic equivalence is established. No infinite-dimensional spectral theorem is inferred from the formal identity alone.}
\GSfield{Assessment}{literal statement: correct with missing hypotheses; publication match: exact; verification: read the result; development: annotated; confidence: high}
\end{GSitem}
\begin{GSitem}{1978-447.8}{E}{Lattice curvature and the continuum Yang--Mills action}
\GSfield{Kind; source}{construction; \GSpage{P000792-P000793}}
\GSfield{The notes}{A link matrix \(B_{x\alpha}\) is attached to each positive edge. The plaquette is \(F_{x\alpha\beta}=B_{x\alpha}B_{x+\alpha,\beta}B^{-1}_{x+\beta,\alpha}B^{-1}_{x\beta}\); the trace sum is said to tend to \(F^2\).}
\GSfield{Statement}{For a smooth connection sampled on a mesh of size \(a\), oriented parallel transport around a plaquette has expansion \(I+a^2F_{\alpha\beta}+O(a^3)\). For a compact unitary gauge group, a real Wilson action using \(\operatorname{Re}\operatorname{Tr}(I-F_{x\alpha\beta})\), with the appropriate power of \(a\) and normalization, tends to the positive curvature-square action. A raw trace sum differs by a constant, sign, and normalization and is not automatically real.}
\GSfield{Reconstruction}{The mesh scaling and the compact-unitary reality convention repair the shorthand. The exponential expansion on the page is a truncation, not an exact identity.}
\GSfield{Fate}{\citet[\S\S II--III]{wilson1974confinement} introduced the gauge-invariant plaquette lattice action in 1974; the original paper and the continuum-limit discussion were checked.}
\GSfield{Present status}{As of 10 November 2026: The construction is standard. The notebook records the structural continuum limit, with its normalization suppressed.}
\GSfield{Assessment}{literal statement: correct with missing hypotheses; publication match: exact; verification: read the result; development: annotated; confidence: high}
\end{GSitem}
\begin{GSitem}{1978-447.9}{E}{The lattice Yang--Mills equation as a sum of plaquette differences}
\GSfield{Kind; source}{claim; \GSpage{P000793}}
\GSfield{The notes}{For fixed \(\alpha\), neighbouring plaquettes in the \(\alpha\beta\) planes are subtracted and summed over \(\beta\). In the expanded second plaquette, its first inverse has index \(x,\beta\).}
\GSfield{Statement}{Varying a link in a trace plaquette action gives a sum of neighbouring oriented plaquettes, all based at the same vertex; for a constrained gauge group the result is projected to its Lie algebra. The inverse edge leaving \(x\) in the negative \(\beta\) direction is \(B^{-1}_{x-\beta,\beta}\), not \(B^{-1}_{x,\beta}\). With this correction the second product closes and the continuum equation is \(D_\beta F_{\alpha\beta}=0\), up to the chosen sign of curvature.}
\GSfield{Reconstruction}{The source diagram determines the corrected edge index. Lie-algebra projection is necessary for groups such as \(SU(n)\), rather than unrestricted matrix variations.}
\GSfield{Fate}{The underlying lattice gauge variational framework is prior work \citep[\S III]{wilson1974confinement}; the erroneous edge can be diagnosed independently by following the four vertices.}
\GSfield{Present status}{As of 10 November 2026: The diagrammatic equation has the standard meaning. Its expanded matrix formula contains a faithful index error.}
\GSfield{Assessment}{literal statement: false as written; publication match: partial; verification: read the result; development: annotated; confidence: high}
\end{GSitem}
\begin{GSitem}{1978-447.10}{E}{Based holonomy and curvature insertions}
\GSfield{Kind; source}{construction; \GSpage{P000793}}
\GSfield{The notes}{\(\Psi(C_y)\) is the ordered product of links along a based contour. Attaching a small square compares \(\Psi(C+\Pi)\Psi(C)^{-1}\) with a curvature transported to the base point.}
\GSfield{Statement}{For a smooth matrix connection and a piecewise smooth based loop, holonomy transforms by conjugation at the base point. Inserting a sufficiently small oriented plaquette at a point reached by transport \(U\) changes the holonomy, to first order in its area, by \(UF_{\alpha\beta}U^{-1}\) in the corresponding left or right logarithmic derivative. Taking a trace removes dependence on the base-point frame. Matrix holonomy itself is not a gauge-invariant scalar.}
\GSfield{Reconstruction}{The small-area qualification, transport convention, and distinction between a holonomy and its trace are supplied.}
\GSfield{Fate}{The loop observable and ordered lattice transport are already present in \citet[\S\S II--III]{wilson1974confinement}. Polyakov’s subsequent abstract identifies their loop-space reinterpretation \citep[abstract]{polyakov1979string}.}
\GSfield{Present status}{As of 10 November 2026: The curvature-insertion construction is established for smooth connections and finite lattice paths.}
\GSfield{Assessment}{literal statement: correct with missing hypotheses; publication match: exact; verification: read the result; development: annotated; confidence: high}
\end{GSitem}
\begin{GSitem}{1978-447.11}{A}{The logarithmic loop field: transversality and flatness}
\GSfield{Kind; source}{construction; \GSpage{P000793-P000794}}
\GSfield{The notes}{The loop field is \(\mathcal F_\mu(s,C)=(\delta\Psi/\delta x_\mu(s))\Psi^{-1}\). The notes impose \(x^{\prime}_\mu\delta\Psi/\delta x_\mu=0\) and a Maurer--Cartan equation.}
\GSfield{Statement}{Where a twice differentiable map \(\Psi\) from a smooth loop chart to a matrix group is defined, its right logarithmic derivative satisfies \(\delta_\nu\mathcal F_\mu-\delta_\mu\mathcal F_\nu+[\mathcal F_\mu,\mathcal F_\nu]=0\), with the point labels on functional derivatives retained. Reparametrization invariance implies tangential variations vanish, hence \(x^{\prime}_\mu(s)\mathcal F_\mu(s,C)=0\). These kinematic identities alone do not impose the Yang--Mills equation.}
\GSfield{Reconstruction}{The statement separates a legitimate Maurer--Cartan identity from undefined coincident-point functional derivatives. Smoothness is a local assumption, not a construction of a measure on loop space.}
\GSfield{Fate}{Polyakov’s 1979 paper and 1980 sequel identify this chiral-field description \citep[abstracts]{polyakov1979string,polyakov1980rings}. Their author-hosted records were opened; the complete printed derivations were not accessible.}
\GSfield{Present status}{As of 10 November 2026: The ordinary differential-geometric identities are valid in a smooth loop chart. Distributional contact terms require a separately defined calculus.}
\GSfield{Assessment}{literal statement: correct with missing hypotheses; publication match: partial; verification: publisher metadata checked; development: annotated; confidence: high}
\end{GSitem}
\begin{GSitem}{1978-447.12}{A}{Yang--Mills as a functional conservation equation}
\GSfield{Kind; source}{claim; \GSpage{P000794}}
\GSfield{The notes}{The divergence equation \(\delta\mathcal F_\mu(s,C)/\delta x_\mu(s)=0\) is proposed as the loop-space counterpart of \(D_\mu F_{\mu\nu}=0\). The notes say three dimensions were checked carefully and four dimensions not completely.}
\GSfield{Statement}{For a holonomy induced by a smooth connection, variations of the loop insert curvature, and the appropriate local part of a second variation inserts its covariant divergence. Identifying that local coefficient with zero can encode the Yang--Mills equation. The coincident functional derivative written on the page is not an ordinary trace of the Hessian on an unspecified infinite-dimensional space.}
\GSfield{Reconstruction}{The local-coefficient interpretation states the formal content without choosing an unrecorded regulator or claiming an unrestricted loop-space PDE theorem.}
\GSfield{Fate}{The same programme is explicitly described in \citet[abstracts]{polyakov1979string,polyakov1980rings}. The 1980 abstract still distinguishes three-dimensional higher currents from the unclear four-dimensional hidden symmetry.}
\GSfield{Present status}{As of 10 November 2026: Identified as Polyakov’s published loop-space formulation, with verification limited to abstracts and metadata. The notebook is not a rigorous general existence result for this functional PDE.}
\GSfield{Assessment}{literal statement: correct with missing hypotheses; publication match: partial; verification: publisher metadata checked; development: annotated; confidence: medium}
\end{GSitem}
\begin{GSitem}{1978-447.13}{F}{An undefined chiral action on loop space}
\GSfield{Kind; source}{method; \GSpage{P000794}}
\GSfield{The notes}{The notes formally write an integral over all contours of \(\delta\Psi^{-1}\delta\Psi\), followed by the explicit reservation \ru{Неясно, как определять}.}
\GSfield{Statement}{To define the proposed action one must specify a space of loops, a measure or substitute for it, domains for the derivatives, and a finite or renormalized value of the integral. None is given. The expression can only be retained as a formal analogy with a two-dimensional chiral action.}
\GSfield{Reconstruction}{No preferred Wiener measure, Sobolev completion, or regularization is attributed to the speaker.}
\GSfield{Fate}{The loop-space viewpoint matches \citet[abstracts]{polyakov1979string,polyakov1980rings}; neither checked abstract resolves the notebook’s expressly stated definitional gap.}
\GSfield{Present status}{As of 10 November 2026: No precise variational problem is identifiable here. An absent construction is not classified as a presently open theorem.}
\GSfield{Assessment}{literal statement: not determinable; publication match: thematic only; verification: unresolved lead; development: annotated; confidence: medium}
\end{GSitem}
\begin{GSitem}{1978-447.14}{A}{Recovering a connection from loop data}
\GSfield{Kind; source}{claim; \GSpage{P000794}}
\GSfield{The notes}{The discussion says that when \(\Psi\) comes from a connection, or is a homomorphism from loops to the group, the loop equations give Yang--Mills fields.}
\GSfield{Statement}{A connection determines smooth parallel transport respecting concatenation, inversion, reparametrization, and thin homotopy. Conversely, reconstruction requires smooth transport data of this kind, with compatible base points or a path-groupoid formulation; an arbitrary homomorphism of an unspecified loop group is insufficient. Under a valid reconstruction theorem, the local curvature and its covariant divergence recover the original Yang--Mills equation.}
\GSfield{Reconstruction}{The listed transport properties explain what the abbreviated homomorphism condition suppresses; no specific global reconstruction theorem is claimed to have been stated in the lecture.}
\GSfield{Fate}{The published identification is supported at programme level by \citet[abstract]{polyakov1979string}. The exact global hypotheses are not recoverable from the notebook or the accessible abstract.}
\GSfield{Present status}{As of 10 November 2026: The forward holonomy construction is clear; the unrestricted converse requires added hypotheses and is only a partial publication match.}
\GSfield{Assessment}{literal statement: correct with missing hypotheses; publication match: partial; verification: publisher metadata checked; development: annotated; confidence: medium}
\end{GSitem}
\begin{GSitem}{1978-447.15}{A}{A prospective second loop current in three dimensions}
\GSfield{Kind; source}{method; \GSpage{P000795}}
\GSfield{The notes}{The second current is written using \(\varepsilon_{\mu\nu\lambda}x^{\prime}_\nu\mathcal F_\lambda+[\mathcal F_\mu,\chi]\), with \(\delta_\mu\chi=-\varepsilon_{\mu\nu\lambda}x^{\prime}_\nu\mathcal F_\lambda\); an inverse relation and a tail-integral construction are sketched.}
\GSfield{Statement}{For a unit tangent \(t\) and a transverse vector \(\mathcal F\), the algebraic identity \(\delta\chi=-t\times\mathcal F\) implies \(\mathcal F=t\times\delta\chi\). For an unnormalized tangent, the inverse needs \(|t|^{-2}\). Existence of \(\chi\) also requires compatibility of its functional derivatives. A claimed conservation law for the resulting nonlocal current depends on the loop equation and a consistent treatment of coincident-point terms.}
\GSfield{Reconstruction}{The unit-speed and compatibility conditions are supplied. The half-factor implicit in antisymmetrizing functional derivatives is not fixed by the notebook’s sketch.}
\GSfield{Fate}{Three-dimensional higher currents are expressly reported in \citet[abstract]{polyakov1980rings}, following the 1979 announcement \citep[abstract]{polyakov1979string}. Only those abstracts and bibliographic records were checked.}
\GSfield{Present status}{As of 10 November 2026: The proposal has an identified published continuation. The displayed sketch is not a complete proof and its unnormalized inverse is incomplete.}
\GSfield{Assessment}{literal statement: correct with missing hypotheses; publication match: partial; verification: publisher metadata checked; development: annotated; confidence: medium}
\end{GSitem}
\begin{GSitem}{1978-447.16}{A}{A spectral linear system on loop space}
\GSfield{Kind; source}{construction; \GSpage{P000795-P000796}}
\GSfield{The notes}{The auxiliary equation has the form \((\delta_\mu+z\varepsilon_{\mu\nu\lambda}x^{\prime}_\nu\delta_\lambda)\phi=\mathcal F_\mu\phi\); a special spectral value is related to \(\Psi\). Compatibility is said to produce the loop equations.}
\GSfield{Statement}{The recorded expression is a formal spectral deformation of the loop-space first-order system. To obtain an actual compatibility theorem one must define the functional differential operators, include their action on the varying tangent, choose the parameter at which holonomy is recovered, and fix multiplicative normalization. The notes do not provide these data.}
\GSfield{Reconstruction}{The two-dimensional auxiliary problem that follows is not substituted as a proof of this infinite-dimensional one.}
\GSfield{Fate}{The hidden-symmetry and conserved-current claims of \citet[abstracts]{polyakov1979string,polyakov1980rings} identify the programme, but not an exact printed formula verified in this study.}
\GSfield{Present status}{As of 10 November 2026: A partial publication match is established. Analytic solvability and full compatibility of this precise formal operator remain unverified here.}
\GSfield{Assessment}{literal statement: not determinable; publication match: partial; verification: publisher metadata checked; development: annotated; confidence: medium}
\end{GSitem}
\begin{GSitem}{1978-447.17}{E}{The principal chiral field and a spectral flat connection}
\GSfield{Kind; source}{theorem; \GSpage{P000795-P000796}}
\GSfield{The notes}{As a finite-dimensional analogue, \(A_\mu=g^{-1}\partial_\mu g\) and \((\partial_\mu+z\varepsilon_{\mu\nu}\partial_\nu)\phi=A_\mu\phi\) lead to a formula with denominator \(1-z^2\).}
\GSfield{Statement}{On a simply connected two-dimensional Lorentzian domain, a smooth group-valued principal chiral field satisfies \(dA+A\wedge A=0\) and \(d{*}A=0\). These are equivalent to flatness, for all admissible spectral values, of a rational family built from \(A\) and \({*}A\), whose denominator can be written \(1-z^2\) when \({*}^2=1\) on one-forms. In Euclidean signature \({*}^2=-1\), the corresponding inversion uses \(1+z^2\) after a consistent parameter convention.}
\GSfield{Reconstruction}{Signature and left/right multiplication conventions are essential additions; the notebook does not identify them, so its formula is not asserted in both signatures.}
\GSfield{Fate}{\citet[\S2]{zakharovmikhailov1978relativistic} had already established the inverse-scattering formulation of the two-dimensional chiral field before this meeting.}
\GSfield{Present status}{As of 10 November 2026: The chiral-model Lax construction is established background, with the signature ambiguity explicitly recorded.}
\GSfield{Assessment}{literal statement: correct with missing hypotheses; publication match: exact; verification: read the result; development: annotated; confidence: high}
\end{GSitem}
\begin{GSitem}{1978-447.18}{E}{Flat holonomy and spectral integrals of the chiral model}
\GSfield{Kind; source}{method; \GSpage{P000796}}
\GSfield{The notes}{Flatness is used to write a closed-contour ordered exponential as 1 and to infer conservation of the transport along a spatial line.}
\GSfield{Statement}{A flat connection has trivial holonomy around contractible loops in its domain; noncontractible loops can have nontrivial monodromy. Transport along a spatial interval evolves by endpoint multiplication. On a periodic spatial circle its conjugacy class is conserved; on a line, a normalized scattering matrix is constant under suitable decay and endpoint normalization. Expanding conjugacy-invariant spectral functions then gives conserved quantities.}
\GSfield{Reconstruction}{Contractibility, endpoint behavior, and normalization are restored. Flatness alone does not set every closed-loop holonomy to the identity.}
\GSfield{Fate}{This is the transport mechanism underlying \citet[\S\S1--2]{zakharovmikhailov1978relativistic}, already published in 1978.}
\GSfield{Present status}{As of 10 November 2026: Valid with the stated topology and boundary hypotheses; false if read as trivial holonomy of every flat connection.}
\GSfield{Assessment}{literal statement: false as written; publication match: partial; verification: read the result; development: annotated; confidence: high}
\end{GSitem}
\begin{GSitem}{1978-447.19}{E}{The massive Green-function analogy}
\GSfield{Kind; source}{remark; \GSpage{P000796}}
\GSfield{The notes}{The scalar equation \((\Delta-m^2)G=\delta\) and decay \(G\sim e^{-mr}\) motivate an exponential in the area of a spanning surface.}
\GSfield{Statement}{For \(m>0\), the Euclidean fundamental solution of \(-\Delta+m^2\) on \(\mathbb R^d\) has large-distance behavior \(c_d r^{-(d-1)/2}e^{-mr}(1+O(r^{-1}))\). Changing the operator sign changes the source normalization. The notebook records the exponential rate and suppresses the dimension-dependent prefactor.}
\GSfield{Reconstruction}{The prefactor is needed if \(\sim\) means asymptotic equivalence. Nothing in this finite-dimensional example proves a corresponding contour Green-function formula.}
\GSfield{Fate}{The use of massive propagators is standard field-theoretic background; the inverse-scattering/field-theory setting is already explicit in \citet[Introduction]{zakharovmikhailov1978relativistic}. This citation is thematic, while the asymptotic follows directly from the modified-Bessel fundamental solution.}
\GSfield{Present status}{As of 10 November 2026: The exponential rate is correct. Literal asymptotic equivalence without the power is generally false for \(d>1\).}
\GSfield{Assessment}{literal statement: false as written; publication match: thematic only; verification: read the result; development: annotated; confidence: high}
\end{GSitem}
\begin{GSitem}{1978-447.20}{F}{A proposed string Green function and area-law asymptotic}
\GSfield{Kind; source}{problem; \GSpage{P000796}}
\GSfield{The notes}{The notes ask for a Green functional satisfying
\[(\delta^2/\delta x_\mu(s)^2-m^4|x^{\prime}(s)|^2)G(C)=0,\]
with \(G(C)\) decaying like \(e^{-m^2A_{\min}(C)}\), and identify it heuristically with an expectation of \(\Psi(C)\).}
\GSfield{Statement}{This is a proposed functional differential problem, not a defined Green-function problem: its loop space, operator regularization, source, boundary conditions, measure defining the expectation, and asymptotic regime are absent. For a gauge-invariant observable, a representation and trace or equivalent invariant functional must also be specified. The expected leading exponential is a physical area-law ansatz.}
\GSfield{Reconstruction}{No single accepted continuum quantum field theory or rigorous loop Laplacian is silently inserted into the statement.}
\GSfield{Fate}{The relation to later loop-space gauge/string equations is thematic in \citet[abstract]{polyakov1980rings}. Wilson’s earlier area-law mechanism concerns a regulated strong-coupling lattice theory, not this continuum functional equation \citep[Introduction and \S III]{wilson1974confinement}.}
\GSfield{Present status}{As of 10 November 2026: The specific mathematical boundary-value problem is not identifiable. This item is not equated with the Yang--Mills mass-gap problem and receives no problem-index identifier.}
\GSfield{Assessment}{literal statement: not determinable; publication match: thematic only; verification: unresolved lead; development: annotated; confidence: medium}
\end{GSitem}
\begin{GSitem}{1978-447.21}{E}{Projected areas as contour coordinates}
\GSfield{Kind; source}{definition; \GSpage{P000797}}
\GSfield{The notes}{The contour tensor is \(\sigma_{\alpha\beta}=\oint_C x_\alpha\,dx_\beta\). Its functional derivative is written using \(\delta_{\mu\alpha}x^{\prime}_\beta-\delta_{\mu\beta}x^{\prime}_\alpha\).}
\GSfield{Statement}{For a closed \(C^1\) loop in \(\mathbb R^4\), \(\sigma_{\beta\alpha}=-\sigma_{\alpha\beta}\). For a smooth periodic variation \(h\), \(D\sigma_{\alpha\beta}[h]=\int(h_\alpha x^{\prime}_\beta-h_\beta x^{\prime}_\alpha)\,ds\). These six oriented projected areas are translation invariant and insensitive to reparametrization preserving orientation. By Stokes’ theorem they are integrals of the constant two-forms \(dx_\alpha\wedge dx_\beta\) over any spanning surface.}
\GSfield{Reconstruction}{The differentiability and closure hypotheses, independent antisymmetric components, and orientation convention are explicit additions.}
\GSfield{Fate}{This is a classical Stokes/first-variation identity. Its role in Polyakov’s proposed one-dimensional approximation is thematically identified by \citet[abstract]{polyakov1979string}; no novelty is assigned to the identity itself.}
\GSfield{Present status}{As of 10 November 2026: The identity is correct; six projected areas do not determine a loop or its Plateau area.}
\GSfield{Assessment}{literal statement: correct; publication match: thematic only; verification: publisher metadata checked; development: annotated; confidence: high}
\end{GSitem}
\begin{GSitem}{1978-447.22}{E}{Reduction to area variables and self-dual momentum}
\GSfield{Kind; source}{method; \GSpage{P000797}}
\GSfield{The notes}{Assuming \(G\) depends only on \(\sigma\), the notes replace contour derivatives by derivatives \(\pi_{\alpha\beta}=\partial/\partial\sigma_{\alpha\beta}\), require their quadratic contraction to be scalar, and invoke \(\pi=\pm{}^{*}\pi\).}
\GSfield{Statement}{For a real skew-symmetric \(4\times4\) matrix \(P\), the algebraic condition \(P^2=-c^2I\) is equivalent, for \(c>0\), to \(P/c\) being an orthogonal complex structure. With a fixed orientation, its two-form is either self-dual or anti-self-dual. The Hodge convention is \(({*}P)_{\alpha\beta}=\tfrac12\varepsilon_{\alpha\beta\gamma\delta}P_{\gamma\delta}\). This finite-dimensional fact does not by itself classify solutions of a differential system in six area variables.}
\GSfield{Reconstruction}{The missing Hodge factor and free-index correction in \(\pi_{\alpha\beta}\pi_{\beta\gamma}\propto\delta_{\alpha\gamma}\) are supplied. Fourier-symbol algebra must be distinguished from an operator identity, and the contour Hessian can contain contact terms.}
\GSfield{Fate}{The self-dual decomposition is elementary four-dimensional linear algebra. The relation to the published approximation is only thematic \citep[abstract]{polyakov1979string}; the notebook’s unrestricted differential conclusion is not verified by that source.}
\GSfield{Present status}{As of 10 November 2026: The algebraic symbol statement is valid. The purported reduction of all functional solutions to self-dual branches has not been established by the recorded reasoning.}
\GSfield{Assessment}{literal statement: correct with missing hypotheses; publication match: thematic only; verification: publisher metadata checked; development: annotated; confidence: medium}
\end{GSitem}
\begin{GSitem}{1978-447.23}{F}{An unspecified Bessel-function solution in two area invariants}
\GSfield{Kind; source}{claim; \GSpage{P000797}}
\GSfield{The notes}{After defining \(U_\pm=(\sigma^2\pm\sigma{*}\sigma)^{1/2}\), the notes say that \(G(U_+,U_-)\) is found through Bessel functions, without giving an order, equation, or solution.}
\GSfield{Statement}{The two invariants encode the lengths of the self-dual and anti-self-dual parts of the area tensor, once the contraction normalization is fixed. A Bessel solution would require a precise reduced differential operator, branch, source, and boundary/asymptotic data. None is recorded, so a specific formula cannot be reconstructed uniquely.}
\GSfield{Reconstruction}{No modified-Bessel order is guessed from a finite-dimensional radial analogy.}
\GSfield{Fate}{The one-dimensional approximation in \citet[abstract]{polyakov1979string} is a thematic lead only; the full paper was not obtained, and its detailed match to this Bessel assertion remains unchecked.}
\GSfield{Present status}{As of 10 November 2026: The invariants are identifiable, but the alleged Green function is not. No claim of mathematical openness follows.}
\GSfield{Assessment}{literal statement: not determinable; publication match: thematic only; verification: unresolved lead; development: annotated; confidence: low}
\end{GSitem}
\begin{GSitem}{1978-447.24}{E}{Self-dual area profiles and the Hamilton--Jacobi identity}
\GSfield{Kind; source}{claim; \GSpage{P000797-P000798}}
\GSfield{The notes}{The proposed area expression is \((\sigma^2\pm\sigma{*}\sigma)^{1/2}\), and the notes say it solves \(|\delta A/\delta x(s)|^2=|x^{\prime}(s)|^2\).}
\GSfield{Statement}{Use independent projected areas and the Euclidean norm on two-forms. Define \(A_\pm(C)=\sqrt2\,|\sigma_\pm(C)|\), where \(\sigma_\pm=(\sigma\pm{*}\sigma)/2\). At \(\sigma_\pm\ne0\), the derivative is pairing with the constant two-form \(\omega=\sqrt2\,\sigma_\pm/|\sigma_\pm|\), which has comass one and corresponds to an orthogonal complex structure \(J\). Consequently the contour gradient is \(Jx^{\prime}\), up to orientation convention, and its norm equals \(|x^{\prime}|\). The functional is nondifferentiable at \(\sigma_\pm=0\).}
\GSfield{Reconstruction}{This normalization makes a positively oriented unit circle have value \(\pi\). The derivative includes complementary-plane components even for a curve contained in \(\mathbb R^3\). Satisfaction of a Hamilton--Jacobi equation does not select the Plateau value function.}
\GSfield{Fate}{The identity is checked by self-dual linear algebra and contour variation; it is classical variational background rather than evidence for a new Plateau solution. Polyakov’s publication is a thematic historical link only \citep[abstract]{polyakov1979string}.}
\GSfield{Present status}{As of 10 November 2026: The normalized identity is valid away from the zero self-dual component. Its implication for true minimum area is a separate claim.}
\GSfield{Assessment}{literal statement: correct with missing hypotheses; publication match: thematic only; verification: publisher metadata checked; development: annotated; confidence: high}
\end{GSitem}
\begin{GSitem}{1978-447.25}{F}{Can projected areas give the minimum spanning area?}
\GSfield{Kind; source}{claim; \GSpage{P000797-P000798}}
\GSfield{The notes}{The expression in the preceding item is labelled \(A_{\min}\), followed by the qualification \ru{Для не слишком закрученных контуров} and Sinai’s recollection of something in Courant.}
\GSfield{Statement}{For every oriented smooth spanning surface, \(\operatorname{Area}(S)\ge A_\pm(C)\), by the comass-one bound for the associated constant two-form. Equality requires tangent planes calibrated by that same form almost everywhere. It therefore requires a surface holomorphic for one fixed compatible orthogonal complex structure. A generic contour has no such filling. The phrase about insufficient twisting specifies neither a topology nor a quantitative condition and cannot define a theorem.}
\GSfield{Reconstruction}{The calibration inequality is a precise reconstruction of what the projected-area expression can establish. The unrestricted equality is false; the unspecified qualified version cannot be assigned a definitive intended scope. A contour’s being close to a planar circle is not by itself a valid equality condition.}
\GSfield{Fate}{The checked Polyakov publication records do not verify this exact area claim \citep[abstracts]{polyakov1979string,polyakov1980rings}. A checked primary account of the comass argument is \citet[Introduction, pages 48--49]{harveylawson1982calibrated}. The elementary audit distinguishes a false unrestricted equality from an unspecified restricted question.}
\GSfield{Present status}{As of 10 November 2026: The general equality fails; a mathematically meaningful equality class is supplied by calibrated fillings. The notebook’s intended restricted family remains unidentifiable.}
\GSfield{Assessment}{literal statement: false as written; publication match: none found; verification: unresolved lead; development: annotated; confidence: medium}
\end{GSitem}

\GStalk{448}{1978-11-20}{A. S. Schwarz}{P000799-P000803}{Gauge quotients, instanton determinants, and the sigma-model Coulomb gas}
The meeting begins on P000799 with palindromes and Krichever’s isolated phrase about instantons without bundles; neither supplies an additional developed argument. Schwarz explains regularized functional integrals, gauge degeneracy, the Faddeev--Popov prescription, and its finite-dimensional quotient-measure analogue. He then describes the semiclassical contribution of rational-map instantons in the two-dimensional sphere-valued sigma model. The key calculation replaces a difficult determinant by its metric variation, separates local heat coefficients from zero modes, and cancels the latter against the Gram determinant on instanton moduli. Its final zeros-and-poles measure becomes a two-component Coulomb gas and, at the indicated coupling, a free-fermion calculation.\par Most of the method is identifiable in work already circulating before this meeting. The detailed 1979 Fateev--Frolov--Schwarz paper says that its central measure formula had been sketched in Frolov and Schwarz’s 1978 announcement \citep[\S1, equation (1.9), and reference 13]{fateevfrolovschwarz1979quantum}. The notebook’s missing zero-mode exclusions, determinant sign, and dropped heat derivative are recorded as faithful mathematical errors. The resulting one-loop instanton-sector calculation is not a proof of the full interacting continuum field theory.

\begin{GSitem}{1978-448.1}{F}{Krichever’s isolated remark about instantons}
\GSfield{Kind; source}{remark; \GSpage{P000799}; I. M. Krichever}
\GSfield{The notes}{Before Schwarz begins, Krichever says \ru{инстантоны без расслоений}. No formula, construction, or precise claim follows.}
\GSfield{Statement}{The phrase can be translated as instantons without bundles, but it does not specify the equation, gauge group, boundary condition, or proposed replacement for bundle language. It cannot be identified as a theorem from these notes.}
\GSfield{Reconstruction}{No ADHM construction or later spectral construction is assigned to the phrase merely because it concerns instantons.}
\GSfield{Fate}{The neighbouring lecture has a definite match in \citet[\S1]{fateevfrolovschwarz1979quantum}; that match does not identify this independent intervention.}
\GSfield{Present status}{As of 10 November 2026: The remark remains an unresolved identification, without an inferred open problem.}
\GSfield{Assessment}{literal statement: not determinable; publication match: none found; verification: unresolved lead; development: annotated; confidence: low}
\end{GSitem}
\begin{GSitem}{1978-448.2}{E}{Normalized functional integrals and cutoff-dependent parameters}
\GSfield{Kind; source}{method; \GSpage{P000799}}
\GSfield{The notes}{\(Z=\int e^{-S(A)}dA\) is approximated on a lattice. The notes say the partition-function limit does not exist but ratios giving correlations have a meaning after changing \(\lambda_a\) with the mesh size.}
\GSfield{Statement}{At finite volume and finite lattice spacing, a normalized expectation is \(\langle\Phi\rangle_a=Z_a^{-1}\int\Phi e^{-S_a}\,dA\), whenever \(0<Z_a<\infty\). Multiplying the measure by a cutoff-dependent constant cancels in this ratio. Existence of a continuum limit still requires a specified model, parameter tuning, observable renormalization, and control of the limit; taking a ratio does not itself ensure convergence.}
\GSfield{Reconstruction}{The unqualified assertion about ratios is restricted to a regularization strategy. No general construction of four-dimensional interacting quantum fields is inferred.}
\GSfield{Fate}{Lattice regularization is prior background \citep[\S\S II--III]{wilson1974confinement}. The sigma-model calculation later written out in \citet[\S2]{fateevfrolovschwarz1979quantum} keeps both ultraviolet and infrared regularization explicit.}
\GSfield{Present status}{As of 10 November 2026: Valid as a method with model-specific hypotheses, not as a universal convergence theorem.}
\GSfield{Assessment}{literal statement: correct with missing hypotheses; publication match: partial; verification: read the result; development: annotated; confidence: high}
\end{GSitem}
\begin{GSitem}{1978-448.3}{E}{Gauge degeneracy of the Yang--Mills action}
\GSfield{Kind; source}{definition; \GSpage{P000799-P000800}}
\GSfield{The notes}{The action is \(S(A)=\int\langle F_{\mu\nu},F^{\mu\nu}\rangle dV\), with \(F=dA+A\wedge A\), and gauge transformation \(A^g=gAg^{-1}+g\,dg^{-1}\). The Lorenz condition \(\partial_\mu A_\mu=0\) is given as a gauge condition.}
\GSfield{Statement}{For a connection on a principal bundle with an invariant inner product on the Lie algebra, curvature transforms by \(F^g=gFg^{-1}\), so the curvature-square action is gauge invariant. Its Hessian at a critical point has infinitesimal gauge directions in its kernel. A local gauge slice can remove these directions under appropriate ellipticity and stabilizer assumptions; the Lorenz condition is not a globally unique representative of every nonabelian orbit.}
\GSfield{Reconstruction}{The bundle and invariant-inner-product hypotheses, critical-point qualification for Hessian degeneracy, and local character of gauge fixing are supplied.}
\GSfield{Fate}{Gauge invariance and its lattice counterpart predate the lecture \citep[\S III]{wilson1974confinement}. Schwarz’s treatment of degenerate quadratic functionals was already published in 1978 \citep[title and publication metadata]{schwarz1978partition}.}
\GSfield{Present status}{As of 10 November 2026: Established gauge-theoretic background. A global gauge section is not asserted.}
\GSfield{Assessment}{literal statement: correct with missing hypotheses; publication match: exact; verification: read the result; development: annotated; confidence: high}
\end{GSitem}
\begin{GSitem}{1978-448.4}{E}{The Faddeev--Popov slice and Jacobian}
\GSfield{Kind; source}{method; \GSpage{P000800}}
\GSfield{The notes}{The notes propose inserting \(\Delta_f(A)\delta(f(A))\) and integrating on a transverse gauge condition. They call this a theorem in finite dimensions and explicitly distinguish physical rigor in field theory.}
\GSfield{Statement}{Locally, let a Lie group act freely and properly on a finite-dimensional manifold, and let \(f=0\) be a smooth transverse slice meeting each relevant orbit once. The change of variables from an orbit parameter and slice coordinates produces a determinant \(\Delta_f=|\det(df\circ T_x)|\), with chosen measures. If multiple slice intersections occur, multiplicities must be accounted for. Infinite-dimensional gauge fixing additionally requires a regularized determinant and a defined integration theory.}
\GSfield{Reconstruction}{The one-intersection and measure conditions make the finite-dimensional theorem precise. Absolute values and orientations depend on whether one integrates positive densities or signed forms.}
\GSfield{Fate}{The finite-dimensional quotient mechanism is the background for \citet[title and metadata]{schwarz1978partition}; the detailed one-loop application is read in \citet[\S2]{fateevfrolovschwarz1979quantum}. The notebook itself labels the method as earlier Faddeev--Popov work.}
\GSfield{Present status}{As of 10 November 2026: A correct local change-of-variables method. The global slice and functional-measure problems are not solved by the insertion.}
\GSfield{Assessment}{literal statement: correct with missing hypotheses; publication match: partial; verification: read the result; development: annotated; confidence: high}
\end{GSitem}
\begin{GSitem}{1978-448.5}{E}{The orbit-volume determinant in a finite-dimensional quotient}
\GSfield{Kind; source}{claim; \GSpage{P000800}}
\GSfield{The notes}{For a free isometric action the notebook writes
\[\int_M e^{-S}=\int_{M/G}e^{-S}\det(TT^*)^{1/2}dV,\]
with \(T:\mathfrak g\to T_xM\). No prime or group-volume factor is visible.}
\GSfield{Statement}{Let a compact Lie group act freely and isometrically on a Riemannian manifold, and choose an Ad-invariant inner product on \(\mathfrak g\). For an invariant integrable function \(h\), the quotient integration formula is \(\int_Mh\,dV_M=\operatorname{vol}(G)\int_{M/G}h\sqrt{\det(T_x^*T_x)}\,dV_{M/G}\), with Haar volume defined by that inner product. Equivalently use the nonzero determinant of \(T_xT_x^*\). The full latter operator is singular whenever the quotient has positive dimension.}
\GSfield{Reconstruction}{The determinant is moved to the injective operator on \(\mathfrak g\), or replaced by a nonzero determinant. Normalizing Haar volume to one absorbs the explicit group-volume constant.}
\GSfield{Fate}{This is the finite-dimensional quotient construction underlying the already published degenerate-Gaussian discussion \citep[metadata]{schwarz1978partition}. The correction follows from singular-value decomposition and the coarea formula; it is not an attribution to an unread theorem in that paper.}
\GSfield{Present status}{As of 10 November 2026: The notebook formula is false with the ordinary determinant on \(T_xM\); the corrected density is standard.}
\GSfield{Assessment}{literal statement: false as written; publication match: partial; verification: publisher metadata checked; development: annotated; confidence: high}
\end{GSitem}
\begin{GSitem}{1978-448.6}{E}{Degenerate Gaussian integrals and analytic torsion}
\GSfield{Kind; source}{method; \GSpage{P000800-P000801}}
\GSfield{The notes}{The mathematical theory for quadratic actions and an application to Ray--Singer invariants are announced, without a torsion formula.}
\GSfield{Statement}{For a finite-dimensional quadratic form nondegenerate on a complement to its kernel, Gaussian integration over that complement gives a nonzero determinant to the power \(-1/2\); quotienting symmetries supplies additional Jacobians. In an elliptic complex on a closed manifold, zeta-regularized alternating Laplacian determinants define analytic torsion, with treatment of cohomology required if zero modes are present. The notebook records this connection, not a new proof of metric independence.}
\GSfield{Reconstruction}{No particular elliptic complex, power of torsion, or acyclicity assumption is attributed to the notes. Those choices distinguish precise partition functions.}
\GSfield{Fate}{\citet[Definition 1.6 and Theorem 2.1]{raysinger1971torsion} already defined torsion and proved metric independence in the acyclic closed-manifold setting. Schwarz’s partition-function interpretation appeared in January 1978 \citep[metadata]{schwarz1978partition}.}
\GSfield{Present status}{As of 10 November 2026: Established prior work; the exact Gaussian model advertised here is not specified.}
\GSfield{Assessment}{literal statement: correct with missing hypotheses; publication match: partial; verification: read the result; development: annotated; confidence: high}
\end{GSitem}
\begin{GSitem}{1978-448.7}{E}{Semiclassical integration over a critical manifold}
\GSfield{Kind; source}{method; \GSpage{P000801}}
\GSfield{The notes}{The problem is the small-parameter asymptotic of \(\int e^{-S/\alpha}\), with a manifold of stationary points and strong degeneracy. The proposed answer integrates a determinant over instantons.}
\GSfield{Statement}{On a compact Riemannian manifold without boundary, suppose the minimum set \(N\) is a compact smooth submanifold and the Hessian of \(S\) is positive definite in its normal bundle. For Riemannian volume, the leading Laplace term is \(e^{-S(N)/\alpha}(2\pi\alpha)^{(\dim M-\dim N)/2}\int_N(\det H_\perp)^{-1/2}dV_N\). Gauge directions must first be quotiented or fixed. The field-theory version replaces the determinant by a regularized nonzero determinant and introduces collective-coordinate zero modes.}
\GSfield{Reconstruction}{The finite-dimensional Morse--Bott hypotheses clarify the heuristic. An arbitrary stationary point, saddle, or noncompact critical family does not satisfy this formula without further analysis.}
\GSfield{Fate}{The instanton-sector realization is given in \citet[\S2, equation (2.5)]{fateevfrolovschwarz1979quantum}; that paper explicitly identifies its central result as announced in 1978.}
\GSfield{Present status}{As of 10 November 2026: The finite-dimensional formula is standard. Its sigma-model use is a controlled formal one-loop calculation after the specified regularizations, not a general path-integral theorem.}
\GSfield{Assessment}{literal statement: correct with missing hypotheses; publication match: exact; verification: read the result; development: annotated; confidence: high}
\end{GSitem}
\begin{GSitem}{1978-448.8}{E}{Dimension of the SU(2) instanton family}
\GSfield{Kind; source}{theorem; \GSpage{P000801}}
\GSfield{The notes}{The shorthand says that for every \(q\) there is a family of \(8q-3\) instantons, besides infinite-dimensional gauge symmetry.}
\GSfield{Statement}{For charge \(q\ge1\), the unframed moduli space of smooth anti-self-dual \(SU(2)\) connections on the round four-sphere, modulo gauge, has real dimension \(8q-3\). The framed problem on \(\mathbb R^4\) has dimension \(8q\). The integer describes parameters, not a count of solutions. Orientation conventions interchange self-duality and anti-self-duality.}
\GSfield{Reconstruction}{The base, framing, positivity of charge, and meaning of the dimension are supplied; the zero-charge case is not covered by the written expression.}
\GSfield{Fate}{\citet[abstract and theorem, pp. 2662--2663]{atiyahhitchinsinger1977deformations} established the dimension in 1977. The primary abstract was read, but the scanned-paper link yielded only a download-preparation page. The earlier cumulative entry reports a reading of both pages; that reading was not repeated here.}
\GSfield{Present status}{As of 10 November 2026: This is established prior instanton moduli theory, not a new 1978 assertion.}
\GSfield{Assessment}{literal statement: correct with missing hypotheses; publication match: exact; verification: publisher metadata checked; development: annotated; confidence: high}
\end{GSitem}
\begin{GSitem}{1978-448.9}{E}{Rational maps as sphere-valued instantons}
\GSfield{Kind; source}{construction; \GSpage{P000801}}
\GSfield{The notes}{For \(n:\mathbb R^2\to S^2\), stereographic projection gives \(w(z)\), and the notes say extremals correspond to rational functions, crediting Belavin and Polyakov.}
\GSfield{Statement}{For smooth finite-action maps extending across infinity to \(S^2\), the positive-degree energy-minimizing instantons are holomorphic maps \(S^2\to\mathbb{CP}^1\), hence rational functions of degree \(q\). Negative-degree anti-instantons are antiholomorphic; constant maps give degree zero. Thus the unqualified phrase all extremals are rational in \(z\) omits the opposite orientation. No claim is made here about arbitrary noncompact-domain or singular stationary maps.}
\GSfield{Reconstruction}{The claim is restricted to its instanton setting, and the antiholomorphic sector is restored. Finite action and extension at infinity are essential.}
\GSfield{Fate}{The prior Belavin--Polyakov construction is explicitly restated and credited in \citet[\S1, equations (1.6)--(1.8)]{fateevfrolovschwarz1979quantum}. The original 1975 PDF link could not be opened, so the attribution is verified indirectly through this primary paper.}
\GSfield{Present status}{As of 10 November 2026: The rational-map instanton classification is established; the notebook suppresses the anti-instanton sector.}
\GSfield{Assessment}{literal statement: correct with missing hypotheses; publication match: exact; verification: read the result; development: annotated; confidence: high}
\end{GSitem}
\begin{GSitem}{1978-448.10}{E}{The sigma-model topological bound}
\GSfield{Kind; source}{theorem; \GSpage{P000801}}
\GSfield{The notes}{The action splits into \(c_1N+c_2\int |\partial_{\bar z}w|^2/(1+|w|^2)^2\); the measure and constants are suppressed.}
\GSfield{Statement}{For \(n:S^2\to S^2\), or a smooth finite-energy map on \(\mathbb R^2\) with constant value at infinity, define \(q=(4\pi)^{-1}\int n\cdot(\partial_1n\times\partial_2n)\,d^2x\). Then \(\int|dn|^2\,d^2x\ge8\pi|q|\). Equality holds for the appropriate first-order conformality equation, equivalently a holomorphic or antiholomorphic stereographic map. With the notebook action \(S=f_0^{-1}\int|dn|^2\), the bound is \(8\pi|q|/f_0\).}
\GSfield{Reconstruction}{The factor reflects the absence of \(1/2\) in the notebook action. The signed decomposition uses the sign of \(q\), and \(d^2x\) must be included.}
\GSfield{Fate}{The completed-square argument is read in \citet[\S1, equation (1.6)]{fateevfrolovschwarz1979quantum}, where the earlier Belavin--Polyakov work is cited.}
\GSfield{Present status}{As of 10 November 2026: The topological bound is standard and sharp under the stated finite-action conditions.}
\GSfield{Assessment}{literal statement: correct with missing hypotheses; publication match: exact; verification: read the result; development: annotated; confidence: high}
\end{GSitem}
\begin{GSitem}{1978-448.11}{E}{The instanton Hessian as a Dolbeault Laplacian}
\GSfield{Kind; source}{construction; \GSpage{P000801-P000802}}
\GSfield{The notes}{The fluctuation bundle is pulled back from \(T\mathbb{CP}^1\), and \(\Delta=\bar\partial^*\bar\partial\). For fixed degree, \(\bar\partial\) is said to be the same and the metric depends on \(w\).}
\GSfield{Statement}{For a degree-\(q\) holomorphic map \(w:S^2\to\mathbb{CP}^1\), the complex tangent fluctuation line bundle is \(w^*T\mathbb{CP}^1\cong\mathcal O(2q)\). The second variation at the instanton is, up to normalization, the nonnegative Dolbeault Laplacian on its sections. Holomorphic line bundles of the same degree on \(S^2\) are isomorphic, so one can fix a holomorphic model while varying its Hermitian metric. Its complex kernel has dimension \(2q+1\), or real dimension \(4q+2\).}
\GSfield{Reconstruction}{The compactified domain and infrared regulator are restored. A bare-plane differential operator does not automatically have the compact discrete spectrum used in the next step.}
\GSfield{Fate}{\citet[\S2, including equations (2.5)--(2.8), and \S3]{fateevfrolovschwarz1979quantum} gives the regulated fluctuation calculation and counts its zero modes. The calculation extends work already announced in 1978.}
\GSfield{Present status}{As of 10 November 2026: Established elliptic and instanton-moduli description. The identification of holomorphic structures requires a choice of isomorphism.}
\GSfield{Assessment}{literal statement: correct with missing hypotheses; publication match: exact; verification: read the result; development: annotated; confidence: high}
\end{GSitem}
\begin{GSitem}{1978-448.12}{E}{Zeta-regularized nonzero determinants}
\GSfield{Kind; source}{definition; \GSpage{P000802}}
\GSfield{The notes}{The notes define \(\zeta(s)=\sum_{\lambda_i\ne0}\lambda_i^{-s}\) and \(\log\det=-\zeta^{\prime}(0)\), saying that zero is a regular point.}
\GSfield{Statement}{For a nonnegative self-adjoint elliptic differential operator of positive order on a smooth closed compact manifold, remove its finite-dimensional kernel. Its spectral zeta function is initially defined for large real part, continues meromorphically, and is regular at zero; set \(\det{}^{\prime}_{\zeta}\Delta=\exp[-\zeta_\Delta^{\prime}(0)]\). For a noncompact domain, a relative determinant, boundary/infrared regularization, or other definition is required.}
\GSfield{Reconstruction}{The ellipticity, positivity, compactness, and zero-mode exclusion are stated explicitly. The determinant is of the positive fluctuation operator, not an arbitrary first-order operator without a spectral cut.}
\GSfield{Fate}{The construction predates the lecture: \citet[\S1, equations (1.5)--(1.6)]{raysinger1971torsion} uses the same zeta regularity and determinant principle. \citet[Appendix B]{fateevfrolovschwarz1979quantum} applies it to the instanton computation.}
\GSfield{Present status}{As of 10 November 2026: An established regularization with precise spectral hypotheses.}
\GSfield{Assessment}{literal statement: correct with missing hypotheses; publication match: exact; verification: read the result; development: annotated; confidence: high}
\end{GSitem}
\begin{GSitem}{1978-448.13}{E}{Metric variation and the Frullani determinant formula}
\GSfield{Kind; source}{method; \GSpage{P000802}}
\GSfield{The notes}{For a fixed differential operator \(A\) with varying adjoint, the notes posit \(\dot{(A^*A)}=BA^*A+A^*CA\), quote Frullani’s identity, then write a positive heat integral for \(\log\det A\).}
\GSfield{Statement}{For positive real \(a,b\), \(\int_0^\infty(e^{-ta}-e^{-tb})\,dt/t=\log(b/a)\). Hence a determinant ratio for positive finite-dimensional operators uses a minus sign in front of the corresponding difference of heat traces. For \(\Delta=A^*A\), a regularized heat representation of \(\log\det{}^{\prime}_{\zeta}\Delta\) also has this sign, removes zero modes, and includes the finite-part convention dictated by the Mellin transform. A determinant of \(|A|\) is half that logarithm.}
\GSfield{Reconstruction}{The signs of \(B,C\) depend on the convention for the metric derivative; the displayed variation identity can be taken as their definition. The later positive integral is a separate faithful error and also conflates \(A\) with \(A^*A\).}
\GSfield{Fate}{The exact regulated counterpart is read in \citet[\S2, equation (2.10), and Appendix B]{fateevfrolovschwarz1979quantum}; the heat determinant method was already standard before the lecture.}
\GSfield{Present status}{As of 10 November 2026: Frullani’s identity is correct. The notebook heat-integral sign and unspecified zero-mode subtraction require correction.}
\GSfield{Assessment}{literal statement: false as written; publication match: exact; verification: read the result; development: annotated; confidence: high}
\end{GSitem}
\begin{GSitem}{1978-448.14}{E}{Variation of the heat trace and local coefficients}
\GSfield{Kind; source}{claim; \GSpage{P000803}}
\GSfield{The notes}{Equation (22) puts \(t\partial_t\) only on the \(B\)-trace and adds an un-differentiated \(C\)-trace. The constant heat coefficient is then singled out.}
\GSfield{Statement}{Whenever the relevant compositions are trace class and differentiation is justified, \(\dot\Delta=B\Delta+A^*CA\) implies \(\partial_u\operatorname{Tr}e^{-t\Delta}=t\partial_t\{\operatorname{Tr}(Be^{-tA^*A})+\operatorname{Tr}(Ce^{-tAA^*})\}\). Both summands carry \(t\partial_t\). After dividing by \(t\) and integrating with the determinant sign and zero-mode subtraction, the small-time local coefficients and the large-time kernel projections control the variation.}
\GSfield{Reconstruction}{The missing derivative is confirmed by the source PNG and by cyclicity of trace using \(Ae^{-tA^*A}=e^{-tAA^*}A\). No full determinant is determined by a single local coefficient without normalization and zero-mode data.}
\GSfield{Fate}{\citet[\S3, equations (3.1)--(3.3)]{fateevfrolovschwarz1979quantum} provides the appropriate metric-variation calculation; the general heat-variation machinery is already in \citet[\S6]{raysinger1971torsion}.}
\GSfield{Present status}{As of 10 November 2026: The corrected trace identity is valid. Its literal notebook version is false.}
\GSfield{Assessment}{literal statement: false as written; publication match: exact; verification: read the result; development: annotated; confidence: high}
\end{GSitem}
\begin{GSitem}{1978-448.15}{E}{Zero modes and the instanton moduli measure}
\GSfield{Kind; source}{claim; \GSpage{P000803}}
\GSfield{The notes}{The \(4q+2\)-dimensional metric determinant on rational maps must be included, and the notes say that its unresolved part cancels the unresolved zero-mode contribution to the fluctuation determinant.}
\GSfield{Statement}{In the regulated degree-\(q\) sector, choose real collective coordinates \(\gamma^a\) for the rational maps and their zero modes \(v_a=\partial n/\partial\gamma^a\). The induced volume has density \(\sqrt{\det(\langle v_a,v_b\rangle)}\,d\gamma\). Combined with \((\det{}^{\prime}\Delta)^{-1/2}\), the variation of this Gram determinant cancels the kernel-projector contribution in the determinant variation, leaving the local anomaly and an overall normalization.}
\GSfield{Reconstruction}{Cancellation is a statement about the combined density and a specified regulator, not permission to drop all zero modes from both calculations.}
\GSfield{Fate}{\citet[\S2, equation (2.7), and \S3]{fateevfrolovschwarz1979quantum} gives the calculation; its central measure result was announced in Frolov--Schwarz 1978, as explicitly recorded in \S1 and reference 13.}
\GSfield{Present status}{As of 10 November 2026: The one-loop instanton measure mechanism is identified and established within its regularized semiclassical setting.}
\GSfield{Assessment}{literal statement: correct with missing hypotheses; publication match: exact; verification: read the result; development: annotated; confidence: high}
\end{GSitem}
\begin{GSitem}{1978-448.16}{E}{The zeros-and-poles instanton measure}
\GSfield{Kind; source}{construction; \GSpage{P000803}}
\GSfield{The notes}{Writing \(w=c\prod_i(z-a_i)/\prod_i(z-b_i)\), the notebook displays same-sign squared distances above unlike-sign squared distances, and \((1+|c^2|)^{-2}\). The lower product has a free \(j\).}
\GSfield{Statement}{On the chart of degree-\(q\) rational maps with distinct finite zeros and poles and finite nonzero value \(c\) at infinity, the one-loop measure has, up to its regulator-dependent fugacity and overall constant, density \[\frac{1}{(q!)^2}\frac{\prod_{i<j}|a_i-a_j|^2|b_i-b_j|^2}{\prod_{i,j}|a_i-b_j|^2}\frac{d^2c}{(1+|c|^2)^2}\prod_i d^2a_i\,d^2b_i.\] Short-distance collisions require the stated ultraviolet prescription; other rational-map charts cover infinity and degenerations.}
\GSfield{Reconstruction}{The missing product ranges, coordinate differentials, permutation factors, and sector fugacity are supplied from the published formula. Since \(|c^2|=|c|^2\), moving the square outside the modulus does not constitute a mathematical correction.}
\GSfield{Fate}{This exact structural density is \citet[\S1, equation (1.9), and \S3, equation (3.10)]{fateevfrolovschwarz1979quantum}. The authors state that its proof had been sketched in Frolov--Schwarz 1978; the full paper was received 24 January 1979 and published in July 1979.}
\GSfield{Present status}{As of 10 November 2026: The formula is an established one-loop density on the instanton sector, with its ultraviolet and infrared restrictions.}
\GSfield{Assessment}{literal statement: correct with missing hypotheses; publication match: exact; verification: read the result; development: annotated; confidence: high}
\end{GSitem}
\begin{GSitem}{1978-448.17}{E}{The neutral two-component Coulomb gas}
\GSfield{Kind; source}{construction; \GSpage{P000803}}
\GSfield{The notes}{The denominator of the normalized instanton sum is identified with a Coulomb-gas partition function: poles carry positive charges, zeros negative charges.}
\GSfield{Statement}{Taking logarithms of the zeros-and-poles density gives pair interactions proportional to \(-q_iq_j\log|z_i-z_j|\), with \(q_i=\pm1\), equal numbers of each sign, and the inverse-temperature/coupling product fixed by the exponent 2. A sector sum introduces a fugacity. The collision singularities and infinite-volume limit need a cutoff or a specified renormalized prescription; the measure is not a finite gas partition function without them.}
\GSfield{Reconstruction}{The sign labelling is conventional; exchanging both species leaves the density unchanged. The mapping applies to the semiclassical instanton ensemble described here.}
\GSfield{Fate}{The Coulomb-gas interpretation is developed in \citet[\S4]{fateevfrolovschwarz1979quantum}, whose central instanton measure was already announced in 1978.}
\GSfield{Present status}{As of 10 November 2026: The algebraic mapping is established. It does not identify the entire interacting sigma model with this gas.}
\GSfield{Assessment}{literal statement: correct with missing hypotheses; publication match: exact; verification: read the result; development: annotated; confidence: high}
\end{GSitem}
\begin{GSitem}{1978-448.18}{E}{The free-fermion value of the Coulomb-gas coupling}
\GSfield{Kind; source}{claim; \GSpage{P000803}}
\GSfield{The notes}{At one temperature, identified parenthetically as the temperature occurring here, the gas is said to be equivalent to free fermions.}
\GSfield{Statement}{For the two-dimensional neutral logarithmic gas at the coupling encoded by squared Cauchy determinants, the identity \(\det(1/(a_i-b_j))\) converts the same-sign/unlike-sign product into a fermionic determinant. With consistent regulators and fugacity, the resulting grand partition function and suitable correlations admit a free-fermion representation. The claim concerns this gas representation of a one-loop instanton ensemble, not an exact equivalence of all sigma-model observables to a free field.}
\GSfield{Reconstruction}{The Cauchy-determinant condition supplies the missing normalization of temperature. A coupling quoted numerically would depend on the convention for logarithmic potential.}
\GSfield{Fate}{The free-fermion treatment is explicit in \citet[\S4]{fateevfrolovschwarz1979quantum}; its references identify the earlier bosonization/free-fermion machinery.}
\GSfield{Present status}{As of 10 November 2026: Established with the stated sector and regularization qualifications.}
\GSfield{Assessment}{literal statement: correct with missing hypotheses; publication match: exact; verification: read the result; development: annotated; confidence: high}
\end{GSitem}

\GStalk{449}{1978-11-27}{V. E. Zakharov; L. N. Lipatov}{P000804-P000807}{A variational principle for rational zero-curvature systems; large-order perturbation theory}
The catalogue names Zakharov, but the dated page P000804 first announces Lev Nikolaevich Lipatov, and P000806 explicitly changes speaker to Lipatov. Both contributions belong to this meeting. Zakharov expands a rational zero-curvature equation with disjoint simple poles, conjugates its residues to prescribed Jordan forms, and supplies an action in those orbit variables. The conservation-law remark and a conjecture about multiple Hamiltonian structures follow. The corresponding Zakharov--Mikhailov paper was received on 3 October 1978 and published in 1980; receipt proves the existence of an earlier manuscript, but is not by itself evidence of public circulation before this lecture \citep[equations (1)--(10) and receipt line]{zakharovmikhailov1980variational}.\par Lipatov explicitly cites his 1977 paper. He motivates the asymptotics of perturbative coefficients through renormalized scalar field theory, invariant charge, and the Landau leading-logarithm formula, then explains the saddle-point method and its sharp four-dimensional Sobolev profile \citep[\S\S2--4]{lipatov1977divergence}. His portion is therefore primarily exposition of prior work. The notebook’s claimed infinite limit of the displayed running coupling contradicts its own positive denominator. A final, separately attributed Zakharov intervention connects the critical cubic equation with Yang--Mills equations without recording the required ansatz.

\begin{GSitem}{1978-449.1}{A}{Rational auxiliary systems and zero curvature}
\GSfield{Kind; source}{construction; \GSpage{P000804}; V. E. Zakharov}
\GSfield{The notes}{The auxiliary equations are \(\Psi_\xi=U\Psi\), \(\Psi_\eta=V\Psi\), with \(U=U_0+\sum_{n=1}^{N_1}U_n/(\lambda-\lambda_n)\), \(V=V_0+\sum_{m=1}^{N_2}V_m/(\lambda-\mu_m)\). All poles are fixed, simple, pairwise distinct, and the two pole sets are disjoint. Equation (3) instead prints \(U_\xi-V_\eta+[U,V]=0\).}
\GSfield{Statement}{For smooth matrix coefficients on a simply connected coordinate domain and an invertible fundamental solution away from the poles, compatibility of the stated auxiliary equations is \(U_\eta-V_\xi+[U,V]=0\). Conversely, this flatness identity gives a local fundamental solution at each regular spectral parameter. The distinct-pole assumptions define the general-position family considered in the lecture; they are not a definition of every integrable system.}
\GSfield{Reconstruction}{The derivatives in equation (3) are interchanged. The corrected convention is fixed by differentiating the two auxiliary equations, rather than by guessing from later formulas.}
\GSfield{Fate}{The rational family appears in \citet[equations (1)--(2)]{zakharovmikhailov1980variational}. Its receipt line is 3 October 1978, preceding the 27 November lecture; this establishes an earlier manuscript, not the date of first public circulation. The rational family was published by the speakers in 1980.}
\GSfield{Present status}{As of 10 November 2026: The corrected compatibility identity is established and the publication match is exact at the level of the rational ansatz.}
\GSfield{Assessment}{literal statement: false as written; publication match: exact; verification: read the result; development: annotated; confidence: high}
\end{GSitem}
\begin{GSitem}{1978-449.2}{E}{A common pole and polynomial Lax pairs}
\GSfield{Kind; source}{remark; \GSpage{P000804-P000805}; V. E. Zakharov}
\GSfield{The notes}{For a Lax pair, the notes describe a simple pole of \(U\) and an arbitrary-order pole of \(V\) at the same point, then write \(U=U_0+U\lambda\), \(V=P_n(\lambda)\).}
\GSfield{Statement}{If rational matrix functions on the spectral sphere have no pole except a common point, a fractional-linear spectral change taking that point to infinity makes them polynomials. Their degrees are the corresponding pole orders. Thus a simple pole for \(U\) and order \(n\) for \(V\) give degree at most one and at most \(n\). This is a subclass of Lax representations, lying outside the disjoint-pole hypothesis above.}
\GSfield{Reconstruction}{The spectral-sphere and no-other-poles conditions are supplied. The statement is not interpreted as a classification of all Lax pairs.}
\GSfield{Fate}{\citet[opening discussion following equations (1)--(2)]{zakharovmikhailov1980variational} separates the generic rational family from familiar degenerations. The polynomial assertion is elementary complex analysis.}
\GSfield{Present status}{As of 10 November 2026: Valid under the stated pole hypotheses; the notebook suppresses the change of spectral coordinate.}
\GSfield{Assessment}{literal statement: correct with missing hypotheses; publication match: partial; verification: read the result; development: annotated; confidence: high}
\end{GSitem}
\begin{GSitem}{1978-449.3}{A}{Coefficient equations at the poles}
\GSfield{Kind; source}{theorem; \GSpage{P000805}; V. E. Zakharov}
\GSfield{The notes}{Equation (4) gives the flatness of \((U_0,V_0)\), followed by \(U_{n\eta}=[U_n,A_n]\), \(V_{m\xi}=[V_m,B_m]\), where \(A_n=V(\lambda_n)\), \(B_m=U(\mu_m)\).}
\GSfield{Statement}{With the auxiliary-equation convention fixed above, equality of the constant term and all residues is equivalent to \[U_{0\eta}-V_{0\xi}+[U_0,V_0]=0,\qquad U_{n\eta}=[A_n,U_n],\qquad V_{m\xi}=[B_m,V_m].\] The evaluations are finite because the pole sets are disjoint. The two residue commutators in the notebook have the opposite order.}
\GSfield{Reconstruction}{The signs are determined by residue extraction from \(U_\eta-V_\xi+[U,V]=0\). No change of commutator convention is made between equations.}
\GSfield{Fate}{The coefficient system is the starting reduction in \citet[equations (2)--(4)]{zakharovmikhailov1980variational}, received before the lecture and published in 1980; the first public circulation date was not verified. The residue calculation independently checks the signs.}
\GSfield{Present status}{As of 10 November 2026: The corrected system is equivalent to zero curvature; the literal notebook signs are inconsistent with its own auxiliary equations.}
\GSfield{Assessment}{literal statement: false as written; publication match: exact; verification: read the result; development: annotated; confidence: high}
\end{GSitem}
\begin{GSitem}{1978-449.4}{A}{Conjugacy orbits and one-variable residue data}
\GSfield{Kind; source}{construction; \GSpage{P000805}; V. E. Zakharov}
\GSfield{The notes}{The residues are written \(U_n=\varphi_nU_n^{(0)}\varphi_n^{-1}\), \(V_m=\psi_mV_m^{(0)}\psi_m^{-1}\), with \(U_n^{(0)}\) depending only on \(\xi\) and \(V_m^{(0)}\) only on \(\eta\). These matrices are regarded as prescribed Jordan forms.}
\GSfield{Statement}{An equation \(U_\eta=[A,U]\) transports \(U\) by conjugation along \(\eta\); its characteristic polynomial and Jordan type are constant on each such curve. Locally on a stratum where smooth conjugating frames exist, choose representatives \(a_n(\xi)\), \(b_m(\eta)\) and write \(U_n=\varphi_na_n\varphi_n^{-1}\), \(V_m=\psi_mb_m\psi_m^{-1}\). Right multiplication of a frame by an element centralizing its representative changes no residue.}
\GSfield{Reconstruction}{Local smooth frames and constant Jordan type are required; a globally smooth choice of Jordan normal form is not automatic through eigenvalue collisions.}
\GSfield{Fate}{The orbit parametrization and its stabilizer freedom are explicit in \citet[equations (6)--(8)]{zakharovmikhailov1980variational}; its manuscript receipt preceded the lecture, while first public circulation was not established.}
\GSfield{Present status}{As of 10 November 2026: The local conjugacy reduction is established. Global gauge charts and singular orbit strata require additional treatment.}
\GSfield{Assessment}{literal statement: correct with missing hypotheses; publication match: exact; verification: read the result; development: annotated; confidence: high}
\end{GSitem}
\begin{GSitem}{1978-449.5}{A}{Gauge freedom and equations for the frames}
\GSfield{Kind; source}{construction; \GSpage{P000805}; V. E. Zakharov}
\GSfield{The notes}{With \(\nabla_\eta=\partial_\eta-V_0\), \(\nabla_\xi=\partial_\xi-U_0\), equation (9) gives linear-looking equations for \(\varphi_n,\psi_m\), with the residues themselves expressed through these frames.}
\GSfield{Statement}{Flatness of \((U_0,V_0)\) permits a local gauge in which \(U_0=V_0=0\). The residue equations then allow representatives satisfying \[\partial_\eta\varphi_n=\left(\sum_m\frac{V_m}{\lambda_n-\mu_m}\right)\varphi_n,\qquad \partial_\xi\psi_m=\left(\sum_n\frac{U_n}{\mu_m-\lambda_n}\right)\psi_m.\] For arbitrary frames there are extra right-hand terms in their centralizers; these can be removed locally by the allowed frame changes.}
\GSfield{Reconstruction}{The passage from commutator equations to frame equations fixes stabilizer gauge. The resulting coupled equations are nonlinear, even though each equation is linear in its own frame when the other residues are held fixed.}
\GSfield{Fate}{\citet[equations (5)--(9)]{zakharovmikhailov1980variational} records the gauge and stabilizer construction in the manuscript received 3 October 1978 and published in 1980; receipt alone does not establish a public preprint date.}
\GSfield{Present status}{As of 10 November 2026: The local gauge construction is established; no global triviality is inferred for a multiply connected base.}
\GSfield{Assessment}{literal statement: correct with missing hypotheses; publication match: exact; verification: read the result; development: annotated; confidence: high}
\end{GSitem}
\begin{GSitem}{1978-449.6}{A}{The orbit-variable variational principle}
\GSfield{Kind; source}{theorem; \GSpage{P000805}; V. E. Zakharov}
\GSfield{The notes}{Equation (10) displays a trace action with two first-order frame terms and a pairwise interaction \(-U_mV_n/(\lambda_m-\mu_n)\), but the summation letters and the displayed equality sign do not form a consistent expression.}
\GSfield{Statement}{In the gauge \(U_0=V_0=0\), with fixed smooth \(a_n(\xi),b_m(\eta)\), set \(U_n=\varphi_na_n\varphi_n^{-1}\), \(V_m=\psi_mb_m\psi_m^{-1}\). Compactly supported variations of the invertible frames in \[S=\int\operatorname{Tr}\left(\sum_n a_n\varphi_n^{-1}\partial_\eta\varphi_n-\sum_m b_m\psi_m^{-1}\partial_\xi\psi_m-\sum_{n,m}\frac{U_nV_m}{\lambda_n-\mu_m}\right)d\xi\,d\eta\] give exactly \(U_{n\eta}=[V(\lambda_n),U_n]\) and \(V_{m\xi}=[U(\mu_m),V_m]\). Frame equations follow after the centralizer gauge choice. On real forms one must also impose the reality conditions making the action real.}
\GSfield{Reconstruction}{This is an explicitly gauge-fixed reconstruction with corrected dummy indices. Varying the frame kinetic term and integrating by parts verifies the commutator signs; the printed equality before the last term is read as a broken continuation, not an additional identity.}
\GSfield{Fate}{The matching action is \citet[equation (10), with equations (5)--(9)]{zakharovmikhailov1980variational}. Both the published result and the receipt date were checked; the latter is not treated as a public-circulation date; OCR does not reliably preserve every symbol in the original formula, so the displayed normalized action is also checked by direct variation.}
\GSfield{Present status}{As of 10 November 2026: The variational principle is established for the specified generic rational family. It does not assert that every integrable equation admits this exact action without a limiting or reduction argument.}
\GSfield{Assessment}{literal statement: false as written; publication match: exact; verification: read the result; development: annotated; confidence: high}
\end{GSitem}
\begin{GSitem}{1978-449.7}{E}{Formal conservation-law expansions at the poles}
\GSfield{Kind; source}{claim; \GSpage{P000806}; V. E. Zakharov}
\GSfield{The notes}{The talk claims \(N_1+N_2\) series of conservation laws, one for each pole, and mentions sine--Gordon with poles at \(0\) and \(\infty\).}
\GSfield{Statement}{For a zero-curvature system admitting local formal diagonalization near a spectral pole, the diagonal scalar compatibility equations generate formal conserved densities and fluxes order by order. Distinct, semisimple leading eigenvalues are a standard sufficient genericity condition. Each of the \(N_1+N_2\) poles offers such an expansion. This counts formal expansion sites, not proven independent global Hamiltonians; integrating the densities requires boundary conditions. Familiar sine--Gordon pairs are reductions with singularities at zero and infinity on the spectral sphere.}
\GSfield{Reconstruction}{The notebook supplies neither the densities nor hypotheses for diagonalization. These qualifications make its mechanism precise without claiming a complete set of independent integrals.}
\GSfield{Fate}{The inverse-scattering setting and sine--Gordon reductions predate this meeting \citep[\S\S1--3]{zakharovmikhailov1978relativistic}. The variational rational family was already in the October 1978 manuscript \citep[equations (1)--(2)]{zakharovmikhailov1980variational}.}
\GSfield{Present status}{As of 10 November 2026: The formal mechanism is standard. The exact independence and analytic convergence of the asserted series are not established in the notes.}
\GSfield{Assessment}{literal statement: correct with missing hypotheses; publication match: partial; verification: read the result; development: annotated; confidence: high}
\end{GSitem}
\begin{GSitem}{1978-449.8}{F}{Many Hamiltonian structures for the rational family}
\GSfield{Kind; source}{conjecture; \GSpage{P000806}; V. E. Zakharov}
\GSfield{The notes}{The explicit conjecture is that there are many Hamiltonian structures, as for KdV, which is said to have infinitely many.}
\GSfield{Statement}{A definite conjecture would need to specify a phase space, allowable local or nonlocal Poisson brackets, equivalence of brackets, compatibility, and which flows must be Hamiltonian. None of these choices is recorded. Merely multiplying one Poisson bracket by constants would already give an uninformative meaning of many.}
\GSfield{Reconstruction}{The historical motivation is retained, but no quantifier or modern conjecture is inserted in place of the missing specification.}
\GSfield{Fate}{The final paragraph of \citet{zakharovmikhailov1980variational} discusses the Hamiltonian interpretation under reality conditions and leaves uniqueness of the constructed structure as a question at that time. This is a related, more specific formulation, not evidence for a present unsolved problem matching every possible reading of the lecture.}
\GSfield{Present status}{As of 10 November 2026: The precise conjecture cannot be identified. The 1980 question is not promoted to an open claim in 2026 without a defined category and a further literature match.}
\GSfield{Assessment}{literal statement: not determinable; publication match: partial; verification: unresolved lead; development: annotated; confidence: medium}
\end{GSitem}
\begin{GSitem}{1978-449.9}{E}{Euclidean scalar correlators and renormalization}
\GSfield{Kind; source}{construction; \GSpage{P000806}; L. N. Lipatov}
\GSfield{The notes}{Lipatov motivates a small coupling by the approximate fine-structure value \(\alpha=e^2/(4\pi\hbar c)\simeq1/137\), then writes the four-dimensional scalar action with \(\lambda_0\varphi^4/4!\), the normalized Euclidean integral of \(\prod_i\varphi(x_i)\), and proposes cutoff-dependent bare mass and coupling to make Green functions finite. His cited source is ZhETF 72, 411 (1977).}
\GSfield{Statement}{At finite volume and an ultraviolet lattice or mode cutoff, a positive quartic coupling defines an ordinary finite-dimensional normalized integral. Its moments are full Schwinger functions; connected functions are their cumulants. Perturbative continuum renormalization in four dimensions requires mass, coupling, and field-strength counterterms, plus a vacuum normalization. Order-by-order renormalized coefficients do not by themselves establish an interacting nonperturbative continuum limit.}
\GSfield{Reconstruction}{The finite cutoff and field-strength counterterm are supplied. The ambiguous subscript \((c)\) cannot turn the displayed full moment into a connected correlator without a cumulant subtraction.}
\GSfield{Fate}{\citet[\S2]{lipatov1977divergence}, received 14 July 1976 and published February 1977, is explicitly named in the notebook and sets up the renormalized perturbative calculation.}
\GSfield{Present status}{As of 10 November 2026: The perturbative framework is established. The literal unqualified assertion about finite continuum Green functions is stronger than the calculation and is not adopted.}
\GSfield{Assessment}{literal statement: correct with missing hypotheses; publication match: exact; verification: read the result; development: annotated; confidence: high}
\end{GSitem}
\begin{GSitem}{1978-449.10}{E}{Momentum invariants, amputation, and the physical mass}
\GSfield{Kind; source}{construction; \GSpage{P000806}; L. N. Lipatov}
\GSfield{The notes}{The four-point function carries \(\delta(p_1+p_2-p_3-p_4)\), depends on scalar products, and is divided by four propagators to define \(\Gamma\). The free propagator is \((p^2+m^2)^{-1}\), while the physical pole is called \(m_c^2\).}
\GSfield{Statement}{In a translation- and rotation-invariant Euclidean scalar theory, strip the overall momentum delta distribution from the connected four-point function and amputate its four external full propagators. The resulting vertex depends on invariant scalar products, subject to momentum conservation. At zero coupling the two-point function is \((p_E^2+m^2)^{-1}\). A stable one-particle mass produces, after the appropriate analytic continuation, a Euclidean pole at \(p_E^2=-m_{\rm phys}^2\), or a Minkowski pole at \(p_M^2=m_{\rm phys}^2\).}
\GSfield{Reconstruction}{Connectedness and the Euclidean/Minkowski sign distinction are explicit repairs. Amputation alone does not in general define a one-particle-irreducible vertex unless the relevant disconnected and reducible contributions have also been specified.}
\GSfield{Fate}{The Green-function and vertex conventions underlying the calculation are treated in \citet[\S2]{lipatov1977divergence}; the lecture is exposition of that prior perturbative setup.}
\GSfield{Present status}{As of 10 November 2026: The kinematic construction is established under the stated invariance and analyticity conditions. The notebook changes signature conventions without saying so.}
\GSfield{Assessment}{literal statement: correct with missing hypotheses; publication match: partial; verification: read the result; development: annotated; confidence: high}
\end{GSitem}
\begin{GSitem}{1978-449.11}{E}{Invariant charge and its normalization}
\GSfield{Kind; source}{construction; \GSpage{P000807}; L. N. Lipatov}
\GSfield{The notes}{The invariant charge is written \(\lambda(p_i)=\lambda_0\Gamma(p_i)\prod_i[G(p_i)/G_0(p_i)]^{1/2}\). The preceding definition of \(\Gamma\) does not divide out its tree coupling.}
\GSfield{Statement}{A renormalized effective quartic coupling is obtained by evaluating an amputated connected vertex at a specified nonexceptional subtraction configuration and including the chosen external-field normalization. If a factor \(\lambda_0\) is written separately, the vertex factor must be normalized to one at tree level; if \(\Gamma\) already includes the tree quartic coupling, that extra factor is omitted. These two conventions cannot be used simultaneously.}
\GSfield{Reconstruction}{The missing tree normalization and subtraction scheme are supplied as alternatives, not silently selected from the incomplete formula. The cancellation of a constant field rescaling is the invariant-charge mechanism.}
\GSfield{Fate}{\citet[\S2, definitions of renormalized coupling and renormalization-group function]{lipatov1977divergence} is the expressly cited prior calculation. The notebook’s exact normalization is only a partial match.}
\GSfield{Present status}{As of 10 November 2026: The invariant-charge construction is standard; the literal product has an unresolved extra coupling unless the vertex is renormalized by its tree value.}
\GSfield{Assessment}{literal statement: correct with missing hypotheses; publication match: partial; verification: read the result; development: annotated; confidence: high}
\end{GSitem}
\begin{GSitem}{1978-449.12}{E}{The leading-logarithm formula and the Landau limit}
\GSfield{Kind; source}{claim; \GSpage{P000807}; L. N. Lipatov}
\GSfield{The notes}{The displayed expression is \(\lambda_\mu=\lambda_0/[1+(3\lambda_0/32\pi^2)\log(\Lambda/\mu^2)]\), followed by the claim \(\lambda_\mu\to\infty\) for fixed positive \(\lambda_0\) as \(\Lambda\to\infty\). The same paradox is mentioned for electrodynamics.}
\GSfield{Statement}{Taking the formula on its own terms, with \(\Lambda\) a squared ultraviolet scale, \(\mu>0\), and fixed \(\lambda_0>0\), the limit is zero. Conversely, keeping a positive target \(\lambda_\mu\) fixed gives \(\lambda_0=\lambda_\mu/[1-(3\lambda_\mu/32\pi^2)\log(\Lambda/\mu^2)]\), which becomes singular at a finite cutoff. This is the leading-logarithm zero-charge/Landau-pole tension, not a nonperturbative theorem about all possible ultraviolet completions.}
\GSfield{Reconstruction}{The claimed infinite limit is a faithful notebook error, confirmed against the scan. The inverse relation distinguishes removing the cutoff at fixed bare coupling from fixing the renormalized charge.}
\GSfield{Fate}{This is prior renormalization-group background for the coefficient problem in \citet[\S1--2]{lipatov1977divergence}. Its quoted limiting conclusion is disproved directly by the displayed rational function.}
\GSfield{Present status}{As of 10 November 2026: The sign and limit audit is decisive. Neither this scalar leading-log formula nor the brief QED analogy proves a general constructive-field-theory result.}
\GSfield{Assessment}{literal statement: false as written; publication match: partial; verification: read the result; development: annotated; confidence: high}
\end{GSitem}
\begin{GSitem}{1978-449.13}{E}{Large-order coefficients from a functional saddle point}
\GSfield{Kind; source}{method; \GSpage{P000807}; L. N. Lipatov}
\GSfield{The notes}{The question is the \(k\to\infty\) behaviour of coefficients in \(\psi(\lambda)=a_2\lambda^2+a_3\lambda^3+\cdots\). Four coefficients are said to be known; the proposed method is Laplace’s method for the continuum integral.}
\GSfield{Statement}{In Lipatov’s regulated, renormalized scalar \(\varphi^4\) calculation with a fixed scheme, extract the order-\(k\) coefficient by a contour integral in the coupling and combine it with the regulated field integral. Rescaling the field exposes a large parameter \(k\); saddles at analytically continued negative quartic coupling yield formal factorial large-order contributions, with collective coordinates and fluctuation determinants determining powers and prefactors. This is not a general saddle-dominance theorem. In the exceptional four-dimensional quartic case, one-loop insertions make contributions comparable to the skeleton diagrams and must be included. The notebook does not state a numerical asymptotic constant or error estimate. With \(t=\log(\Lambda/\mu^2)\), the derivative \(\psi=d\lambda_\mu/dt\) and the physical-scale derivative \(\beta=d\lambda_\mu/d\log\mu\) obey \(\psi=-\beta/2\) when the cutoff and bare parameters used to define \(\beta\) are fixed.}
\GSfield{Reconstruction}{The coefficient-extraction step, scheme dependence, and sign convention are supplied from the cited method. No asymptotic formula for an unspecified beta function is invented, and a formal infinite-dimensional saddle calculation is not called a constructive proof.}
\GSfield{Fate}{\citet[\S\S2--4]{lipatov1977divergence} develops this very large-order method; section 4, p.~221, records the exceptional quartic one-loop contributions. It was submitted July 1976 and published in 1977, well before the meeting.}
\GSfield{Present status}{As of 10 November 2026: The historical method is identifiable and influential. The notes supply too little normalization to verify a particular numerical \(a_k\) formula, and the claim about the first four coefficients is a dated report rather than a result proved here.}
\GSfield{Assessment}{literal statement: correct with missing hypotheses; publication match: exact; verification: read the result; development: annotated; confidence: high}
\end{GSitem}
\begin{GSitem}{1978-449.14}{E}{The sharp four-dimensional Sobolev inequality}
\GSfield{Kind; source}{theorem; \GSpage{P000807}; L. N. Lipatov}
\GSfield{The notes}{The inequality is \([\int_{\mathbb R^4}(\partial_\sigma\varphi)^2]^2\ge(32\pi^2/3)\int_{\mathbb R^4}\varphi^4\), with equality at \(c_1/(|x-x_0|^2+\lambda^2)\).}
\GSfield{Statement}{For real \(\varphi\in\dot H^1(\mathbb R^4)\), \[\left(\int_{\mathbb R^4}|\nabla\varphi|^2dx\right)^2\ge\frac{32\pi^2}{3}\int_{\mathbb R^4}|\varphi|^4dx.\] The constant is sharp. Nonzero equality cases are \(c/(|x-x_0|^2+\rho^2)\), with \(c\ne0\), \(\rho>0\), and arbitrary \(x_0\); zero is also an equality case. The homogeneous energy space is essential, since these nonzero profiles are not in \(L^2(\mathbb R^4)\). For \(u=(1+|x|^2)^{-1}\), direct radial integration gives \(\int u^4=\pi^2/6\), \(\int|\nabla u|^2=4\pi^2/3\), confirming the quoted constant.}
\GSfield{Reconstruction}{The domain, parameter restrictions, zero case, and distinction from inhomogeneous \(H^1\) are supplied. The numerical normalization is independently checked.}
\GSfield{Fate}{The sharp inequality and extremal family are the \(p=2,m=4\) specialization of \citet[main theorem, equations (1)--(3), and \S2]{talenti1976best}, received 16 July 1975 and published in 1976. Its role in the field-theory saddle is explained in \citet[\S3]{lipatov1977divergence}.}
\GSfield{Present status}{As of 10 November 2026: The inequality and equality family are established prior results, correctly quoted once the function space is fixed.}
\GSfield{Assessment}{literal statement: correct with missing hypotheses; publication match: exact; verification: read the result; development: annotated; confidence: high}
\end{GSitem}
\begin{GSitem}{1978-449.15}{F}{A critical cubic equation and a missing Yang--Mills ansatz}
\GSfield{Kind; source}{claim; \GSpage{P000807}; V. E. Zakharov}
\GSfield{The notes}{A separately labelled Zakharov insertion says that \(\Delta\varphi+\varphi^3=0\) is obtained from the Yang--Mills equations.}
\GSfield{Statement}{With \(\Delta=\sum_{j=1}^4\partial_j^2\), the scalar equation is the Euler--Lagrange equation of \(\int_{\mathbb R^4}(\tfrac12|\nabla\varphi|^2-\tfrac14\varphi^4)dx\). The smooth finite-energy bubble \(\varphi(x)=\sqrt8\,\rho/(|x-x_0|^2+\rho^2)\) solves it. To identify a Yang--Mills reduction one must specify the gauge group, connection ansatz, spacetime signature, and reduction constraints; none is recorded.}
\GSfield{Reconstruction}{The scalar variational equation and bubble are independently checkable reconstructions. They do not verify the missing reduction of the full gauge equations.}
\GSfield{Fate}{\citet[\S3]{lipatov1977divergence} supplies the scalar saddle-point context, and \citet[main theorem]{talenti1976best} the critical Sobolev profile. Neither checked passage identifies the particular Yang--Mills ansatz intended by Zakharov’s one-line intervention.}
\GSfield{Present status}{As of 10 November 2026: The scalar equation is precise and solved by the stated bubble; the asserted gauge-theory provenance remains unidentified, not an established open problem.}
\GSfield{Assessment}{literal statement: not determinable; publication match: thematic only; verification: unresolved lead; development: annotated; confidence: medium}
\end{GSitem}

\GStalk{450}{1978-12-04}{Alain Guichardet}{P000808-P000811}{Extensions of representations and first cohomology}
Guichardet starts with extensions split as underlying abelian groups and the cocycle construction of their first cohomology classes. He then treats Lie-algebra extensions, separation by infinitesimal characters, and the relative cohomology describing extensions split over a subalgebra. The topological setting introduces a substantive choice of category: continuous representations, smooth vectors, and Harish--Chandra modules need not have identical extension groups. The notes distinguish ordinary cohomology from the quotient by the closure of coboundaries through the affine-group example, and mention the established semisimple classification. The meeting ends on P000811 immediately after introducing dual orbits for a semidirect product; no extension theorem for those orbits is recorded.\par Delorme’s 1977 paper supplies a checked primary source for the continuous unitary cohomology statements \citep[\S0(a), \S V(C), \S V(D)]{delorme1977cohomologie}. Guichardet’s related article appeared in 1979, but a contemporary primary bibliography lists a June 1978 preprint; its later publication date alone cannot establish that this lecture contained a new theorem \citep[reference 20]{charfi1980extensions}. Several compressed statements on P000811 lack the representations or category needed to identify a theorem and are retained as unresolved identifications.

\begin{GSitem}{1978-450.1}{E}{Strong extensions and equivalence}
\GSfield{Kind; source}{definition; \GSpage{P000808}}
\GSfield{The notes}{For two \(G\)-modules \(E,F\), an extension is an exact sequence \(0\to F\to H\to E\to0\). Equivalence fixes the end modules. It is called strong if a section exists as a homomorphism of abelian groups, without requiring \(G\)-equivariance.}
\GSfield{Statement}{For modules over \(\mathbb Z[G]\), restrict attention to extensions split after forgetting the \(G\)-action. A strong extension has a \(\mathbb Z\)-linear section, so \(H\cong F\oplus E\) as abelian groups. Equivalence means a \(G\)-module isomorphism commuting with injection and projection and inducing the identity on \(F\) and \(E\). An equivariant section makes the extension trivial. Over a field every short exact sequence splits as vector spaces, but this is false for general abelian groups.}
\GSfield{Reconstruction}{The relative, underlying-module splitting condition is distinguished from unrestricted \(\operatorname{Ext}^1_{\mathbb Z[G]}\). The field case does not justify omitting the condition over \(\mathbb Z\).}
\GSfield{Fate}{This is classical extension theory \citep[Chapters X and XIV]{cartaneilenberg1956homological}. Publisher chapter metadata were checked; the specific degree-one block construction is independently verified in the next item rather than attributed to an unread page.}
\GSfield{Present status}{As of 10 November 2026: The definition is standard and unambiguous once the underlying category is fixed.}
\GSfield{Assessment}{literal statement: correct; publication match: partial; verification: publisher metadata checked; development: annotated; confidence: high}
\end{GSitem}
\begin{GSitem}{1978-450.2}{E}{Strong group extensions and first cohomology}
\GSfield{Kind; source}{theorem; \GSpage{P000808-P000809}}
\GSfield{The notes}{The coefficient module is \(\operatorname{Hom}_{\mathbb Z}(E,F)\), with \(g\cdot u=V_guU_g^{-1}\). A section produces a cocycle; conversely the notebook gives the block matrix with upper-right entry \(f(g)U_g\).}
\GSfield{Statement}{Let \(M=\operatorname{Hom}_{\mathbb Z}(E,F)\) with the stated action. A cocycle satisfies \(f(gh)=f(g)+g\cdot f(h)\), and a coboundary has the form \(f(g)=g\cdot u-u\). The formula \[\rho_f(g)=\begin{pmatrix}V_g&f(g)U_g\\0&U_g\end{pmatrix}\] defines a group action on \(F\oplus E\) exactly when \(f\) is a cocycle. Changing the section changes \(f\) by a coboundary. Thus equivalence classes of strong extensions are naturally \(H^1(G,M)\), with Baer addition agreeing with cocycle addition.}
\GSfield{Reconstruction}{The normalization of the off-diagonal block and the group action on Hom are made explicit; they can be checked by multiplying two block matrices.}
\GSfield{Fate}{The construction belongs to the classical cohomological extension framework \citep[Chapters X and XIV]{cartaneilenberg1956homological}. The checked publisher record confirms the book and chapters; the exact degree-one equivalence is reconstructed directly from the source formulas.}
\GSfield{Present status}{As of 10 November 2026: Established. It classifies all vector-space extensions over a field, and only underlying-split extensions for arbitrary \(\mathbb Z\)-modules.}
\GSfield{Assessment}{literal statement: correct; publication match: partial; verification: indirect; development: annotated; confidence: high}
\end{GSitem}
\begin{GSitem}{1978-450.3}{E}{Lie-algebra extensions and the Chevalley--Eilenberg differential}
\GSfield{Kind; source}{theorem; \GSpage{P000809-P000810}}
\GSfield{The notes}{For representations of \(\mathfrak g\) on \(E,F\), the off-diagonal map is \(f(X)\); the notes give its cocycle equation and write \((df)(X,Y)=f([X,Y])-Xf(Y)+Yf(X)\).}
\GSfield{Statement}{Over \(\mathbb C\), put \(M=\operatorname{Hom}_{\mathbb C}(E,F)\) and \(X\cdot u=V_Xu-uU_X\). The block map \(X\mapsto\begin{pmatrix}V_X&f(X)\\0&U_X\end{pmatrix}\) is a representation exactly when \(f([X,Y])=X\cdot f(Y)-Y\cdot f(X)\). Changes of splitting add \(X\mapsto X\cdot u\), giving \(\operatorname{Ext}^1_{U(\mathfrak g)}(E,F)\cong H^1(\mathfrak g,M)\). The notebook differential is the negative of the usual degree-one Chevalley--Eilenberg differential, which has the same kernel.}
\GSfield{Reconstruction}{The action on Hom replaces the ambiguous word commutator when \(E\) and \(F\) differ. Algebraic vector-space splitting imposes no finite-dimensionality assumption on \(E,F\).}
\GSfield{Fate}{The enveloping-algebra interpretation and cochain construction are prior results \citep[Chapter XIII, \S\S1,5,8]{cartaneilenberg1956homological}. The chapter text was opened, including its algebraic setup; the displayed degree-one block verification is elementary.}
\GSfield{Present status}{As of 10 November 2026: Established. The sign convention for \(d\) is harmless when used consistently.}
\GSfield{Assessment}{literal statement: correct; publication match: exact; verification: read the result; development: annotated; confidence: high}
\end{GSitem}
\begin{GSitem}{1978-450.4}{E}{Different infinitesimal characters split an algebraic extension}
\GSfield{Kind; source}{theorem; \GSpage{P000809-P000810}}
\GSfield{The notes}{The lemma states that if \(E\) and \(F\) have different infinitesimal characters, then \(H^1(\mathfrak g,\operatorname{Hom}(E,F))=0\). The next page defines the character by scalar action of the centre of the enveloping algebra.}
\GSfield{Statement}{Let \(E,F\) be \(U(\mathfrak g)\)-modules on which every element of \(Z(U(\mathfrak g))\) acts through algebra characters \(\chi_E,\chi_F\). If these characters differ, choose a central \(z\) with values \(\alpha\ne\beta\). In an extension \(H\), \((z-\alpha)(z-\beta)=0\); the two scalar roots are distinct, so its eigenspaces split \(H\) into invariant submodules. The eigenspace mapping to \(E\) is an equivariant complement to \(F\). Therefore \(\operatorname{Ext}^1(E,F)=0\).}
\GSfield{Reconstruction}{The argument makes the scalar-central-action hypothesis explicit. It uses the centre of the universal enveloping algebra, not merely the centre of the Lie algebra, and does not require \(\mathfrak g\) semisimple.}
\GSfield{Fate}{This is an algebraic separation-of-central-characters lemma in the classical module framework \citep[Chapter XIII, \S1]{cartaneilenberg1956homological}. The splitting argument above verifies it directly; later topological-category caveats do not refute it.}
\GSfield{Present status}{As of 10 November 2026: Established for the algebraic category under the stated hypotheses. Additional continuity is needed when its polynomial projection is used on a topological representation.}
\GSfield{Assessment}{literal statement: correct with missing hypotheses; publication match: partial; verification: read the result; development: annotated; confidence: high}
\end{GSitem}
\begin{GSitem}{1978-450.5}{E}{Relative first cohomology and extensions split over a subalgebra}
\GSfield{Kind; source}{theorem; \GSpage{P000810}}
\GSfield{The notes}{The notes say that relative cocycles vanish on \(\mathfrak h\subset\mathfrak g\) and describe \(\mathfrak h\)-trivial extensions. They first write that absolute \(H^1\) is relative \(H^1\), and subsequently place an element-of sign between the two groups.}
\GSfield{Statement}{For a \(\mathfrak g\)-module \(M\), relative degree-one cocycles are ordinary cocycles \(f\) with \(f|_{\mathfrak h}=0\); their \(\mathfrak h\)-equivariance follows from the cocycle identity. Relative coboundaries are \(du\) with \(u\in M^{\mathfrak h}\). The natural map \(H^1(\mathfrak g,\mathfrak h;M)\to H^1(\mathfrak g;M)\) is injective, with image the kernel of restriction to \(H^1(\mathfrak h;M)\). For \(M=\operatorname{Hom}(E,F)\), this image consists precisely of extensions admitting an \(\mathfrak h\)-equivariant section.}
\GSfield{Reconstruction}{The image and relative coboundaries replace the two erroneous literal relations. Injectivity is special to degree one here: if a relative cocycle equals \(du\), its vanishing on \(\mathfrak h\) forces \(u\in M^{\mathfrak h}\).}
\GSfield{Fate}{The relative cochain complex is standard prior theory; a checked modern primary formulation is \citet[\S3, opening definition and equation (10)]{wockel2014topological}. Its warning that higher-degree maps need not be injective reinforces the need to retain the degree-one hypothesis.}
\GSfield{Present status}{As of 10 November 2026: Established as stated. Equality with absolute cohomology requires additional hypotheses and is false in general, for example when \(\mathfrak h=\mathfrak g\) has nonzero absolute \(H^1\).}
\GSfield{Assessment}{literal statement: false as written; publication match: partial; verification: read the result; development: annotated; confidence: high}
\end{GSitem}
\begin{GSitem}{1978-450.6}{E}{Continuous cocycles and topologically split extensions}
\GSfield{Kind; source}{construction; \GSpage{P000810}}
\GSfield{The notes}{A \(G\)-module is now a topological vector space with continuous action. The first idea uses \(f:G\to\operatorname{Hom}(E,F)\), requiring \((g,e)\mapsto f(g)e\) to be jointly continuous.}
\GSfield{Statement}{For continuous representations on locally convex spaces, restrict to short exact sequences with a continuous linear section and the resulting topological identification \(H\cong F\oplus E\). A crossed homomorphism with jointly continuous evaluation gives a continuous block action, and changes by continuous linear sections give coboundaries. Thus the relevant continuous \(H^1\) classifies these topologically split extensions. An arbitrary topological exact sequence need not admit such a section, and the topology on the Hom coefficient space must be compatible with the chosen evaluation condition.}
\GSfield{Reconstruction}{The continuous-section condition is essential. One may specialize to Banach modules and bounded linear maps, but strong continuity must not silently be replaced by operator-norm continuity.}
\GSfield{Fate}{This is the continuous version of the classical cocycle mechanism, used in the prior unitary cohomology theory \citep[\S0(a)]{delorme1977cohomologie}. The general extension-category distinctions are also documented by \citet[Introduction and \S4]{charfi1980extensions}; its bibliography traces the relevant preprints to before this meeting.}
\GSfield{Present status}{As of 10 November 2026: Established for the specified split category. The notation \(H^1(G,\operatorname{Hom}(E,F))\) alone does not choose that category.}
\GSfield{Assessment}{literal statement: correct with missing hypotheses; publication match: partial; verification: read the result; development: annotated; confidence: high}
\end{GSitem}
\begin{GSitem}{1978-450.7}{E}{Smooth vectors and their topology}
\GSfield{Kind; source}{construction; \GSpage{P000810}}
\GSfield{The notes}{The second idea replaces \(E,F\) by their differentiable-vector spaces \(E_\infty,F_\infty\). The word describing their imperfection is unexplained; the next theorem specializes to \(E=\mathbb C\).}
\GSfield{Statement}{For a continuous representation of a finite-dimensional Lie group on a Fréchet space \(E\), let \(E_\infty\) consist of vectors whose orbit maps are smooth. Give it the Fréchet topology induced by the embedding \(v\mapsto(g\mapsto gv)\) into \(C^\infty(G,E)\), equivalently by seminorms controlling all differentiated orbit maps on compact sets. This is a smooth representation with continuous derived action. Replacing both arguments in Hom by smooth vectors changes the extension category and is not merely a change of notation.}
\GSfield{Reconstruction}{The natural smooth-vector topology is supplied; using the original norm on a dense smooth subspace would lose the stated completeness and continuity properties.}
\GSfield{Fate}{\citet[\S5.1, pages 152--153]{blanc1979cohomologie} defines precisely this space and topology. The smooth-vector construction is prior background. The stronger regularization theorem is assessed separately: the manuscript was received 14 December 1977, revised 27 October 1978, and published in 1979, but no earlier public-release date was verified.}
\GSfield{Present status}{As of 10 November 2026: The construction is established. The notebook’s isolated qualitative criticism has no independent mathematical content without a specified property.}
\GSfield{Assessment}{literal statement: correct with missing hypotheses; publication match: exact; verification: read the result; development: annotated; confidence: high}
\end{GSitem}
\begin{GSitem}{1978-450.8}{B}{Continuous, smooth-vector, and relative Lie-algebra cohomology}
\GSfield{Kind; source}{theorem; \GSpage{P000810}}
\GSfield{The notes}{For \(E=\mathbb C\), the displayed theorem is \(H^1(G,F)\cong H^1(G,F_\infty)\cong H^1(\mathfrak g,\mathfrak k;F_\infty)\), with arrows attributed respectively to Delorme--Blanc and van Est.}
\GSfield{Statement}{Let \(G\) be a connected finite-dimensional Lie group, \(K\) a maximal compact subgroup, and \(F\) a continuous Fréchet \(G\)-module. With continuous cochains on the group, the natural smooth-vector topology, and the derived action on \(F_\infty\), regularization and differentiation give \[H^1_{\rm cont}(G,F)\cong H^1_{\rm cont}(G,F_\infty)\cong H^1(\mathfrak g,\mathfrak k;F_\infty).\] Equivalently one may use smooth group cochains after passing to smooth vectors. If disconnected groups are admitted, \(K\)-equivariance must replace merely \(\mathfrak k\)-equivariance and component-group effects must be retained.}
\GSfield{Reconstruction}{Connectedness, Fréchet completeness, and the cochain topologies are supplied. The assertion concerns first cohomology of a single coefficient module; it does not automatically identify all topological Ext groups between two arbitrary representations.}
\GSfield{Fate}{The regularization theorem is read in \citet[Theorem 5.2, page 153]{blanc1979cohomologie}; the notebook attributes the result to Delorme--Blanc. The paper was received 14 December 1977 and revised 27 October 1978, but these are submission dates, not verified public-release dates. The relative differentiation statement is classical van Est theory, checked in the modern primary treatment \citet[\S3, hypotheses and Theorem 3.3]{wockel2014topological}. That source distinguishes \((\mathfrak g,K)\) from \((\mathfrak g,\mathfrak k)\).}
\GSfield{Present status}{As of 10 November 2026: Established under the explicit assumptions. It is an attributed theorem in the lecture, subsequently documented in Blanc’s 1979 publication; its first public circulation before this meeting is unverified.}
\GSfield{Assessment}{literal statement: correct with missing hypotheses; publication match: exact; verification: read the result; development: annotated; confidence: high}
\end{GSitem}
\begin{GSitem}{1978-450.9}{E}{The unipotent extension of the additive group}
\GSfield{Kind; source}{construction; \GSpage{P000810}}
\GSfield{The notes}{The nontrivial extension of the trivial \(\mathbb R\)-module by itself is \(t\mapsto\begin{pmatrix}1&at\\0&1\end{pmatrix}\). A compressed following sentence says there is no correspondence of group and Lie-algebra representations, and mentions triviality on compact subgroups.}
\GSfield{Statement}{For the additive real group and trivial complex coefficients, continuous cocycles are \(f(t)=at\), and coboundaries vanish. Thus \(H^1_{\rm cont}(\mathbb R,\mathbb C)\cong\mathbb C\); the block representation is nonsplit if \(a\ne0\). Its derived nilpotent matrix integrates exactly to this representation of \(\mathbb R\). For the circle instead, every continuous homomorphism to the additive group \(\mathbb C\) vanishes, so \(H^1(S^1,\mathbb C)=0\), whereas its absolute Lie-algebra \(H^1\) is nonzero.}
\GSfield{Reconstruction}{The circle comparison supplies a valid explanation of the compact-subgroup caveat. It is explicitly a reconstruction, since the page writes \(\mathbb R\), whose displayed representation does integrate.}
\GSfield{Fate}{The example follows directly from the cocycle definition and illustrates the role of relative cohomology in van Est comparison \citep[\S3]{wockel2014topological}. No unidentified counterexample to integration over the simply connected additive group is asserted.}
\GSfield{Present status}{As of 10 November 2026: The displayed additive-group example is correct; the following compressed sentence cannot literally mean that this representation fails to integrate.}
\GSfield{Assessment}{literal statement: correct with missing hypotheses; publication match: partial; verification: read the result; development: annotated; confidence: high}
\end{GSitem}
\begin{GSitem}{1978-450.10}{E}{The affine group: large ordinary cohomology and zero reduced cohomology}
\GSfield{Kind; source}{theorem; \GSpage{P000811}}
\GSfield{The notes}{For an irreducible unitary representation of the affine group of the line, the notes assert that \(B^1\) is dense in \(Z^1\), while \(Z^1/B^1\) is infinite-dimensional, and link this to weak containment of the trivial representation.}
\GSfield{Statement}{Use the connected affine group \(G=\mathbb R\rtimes\mathbb R_{>0}\) with \((a,b)(a^{\prime},b^{\prime})=(aa^{\prime},b+ab^{\prime})\), and its irreducible unitary representation on \(L^2(\mathbb R_{>0},dt/t)\), \((\pi(a,b)u)(t)=e^{ibt}u(at)\). Its ordinary continuous \(H^1\) is infinite-dimensional, but the reduced quotient \(Z^1/\overline{B^1}\), using uniform convergence on compact subsets of \(G\), is zero. Nonclosed coboundaries and failure of the ordinary quotient to be Hausdorff are the central distinction; weak containment alone is not a universal proof of infinite-dimensionality for every group.}
\GSfield{Reconstruction}{The specific irreducible representation and topology are supplied. Infinite-dimensionality can be checked using \(u_\alpha(t)=\chi(t)(\log(1/t))^{-\alpha}\), \(0<\alpha\le1/2\), with \(\chi=1\) near zero and supported below \(e^{-1}\): each \(\pi(g)u_\alpha-u_\alpha\) is in \(L^2(dt/t)\), while no nonzero finite linear combination of the \(u_\alpha\) is. These yield independent cocycle classes; an invariant correction cannot occur because the translation action multiplies by \(e^{ibt}\).}
\GSfield{Fate}{Density of coboundaries for this non-character irreducible representation is a direct specialization of \citet[Theorem V.6, page 323]{delorme1977cohomologie}, published in 1977. The elementary family in the reconstruction checks the notebook’s stronger dimensional assertion; the paper’s general density theorem alone is not cited as proving that dimension.}
\GSfield{Present status}{As of 10 November 2026: The stated example is established and exhibits why reduced and ordinary first cohomology must be kept distinct.}
\GSfield{Assessment}{literal statement: correct with missing hypotheses; publication match: partial; verification: read the result; development: annotated; confidence: high}
\end{GSitem}
\begin{GSitem}{1978-450.11}{E}{First cohomology of irreducible unitary representations of simple groups}
\GSfield{Kind; source}{theorem; \GSpage{P000811}}
\GSfield{The notes}{For semisimple groups and irreducible unitary representations, the lecture says it is known when \(H^1\) vanishes, without listing the representations.}
\GSfield{Statement}{For a connected real simple Lie group with finite centre, nontrivial irreducible unitary representations with nonzero continuous first cohomology occur only in the real- and complex-hyperbolic families. For Lie algebra \(\mathfrak{so}(2,1)\) or \(\mathfrak{su}(n,1)\), \(n\ge2\), there are two mutually conjugate complex irreducibles; for \(\mathfrak{so}(n,1)\), \(n\ge3\), there is one. Their \(H^1\) is one-dimensional. Other real simple families have vanishing first cohomology for all unitary representations. Product groups require the corresponding factor analysis rather than a literal transfer of this simple-group count.}
\GSfield{Reconstruction}{The count is an explicit reconstruction of the known result alluded to, restricted to connected simple groups to keep the hypotheses precise. It is not presented as a numerical theorem written in the notebook.}
\GSfield{Fate}{The exceptional counts and dimensions are read in \citet[Theorem V.5, page 317; see also \S V(A) and \S V(C)]{delorme1977cohomologie}. The paper was published in 1977, so this is prior exposition.}
\GSfield{Present status}{As of 10 November 2026: The classification cited here is established. The distinction between complex unitary and real orthogonal irreducibles explains why some formulations give one exceptional real representation instead of two conjugate complex ones.}
\GSfield{Assessment}{literal statement: correct with missing hypotheses; publication match: exact; verification: read the result; development: annotated; confidence: high}
\end{GSitem}
\begin{GSitem}{1978-450.12}{F}{An inconsistent orthogonal-group vanishing sentence}
\GSfield{Kind; source}{claim; \GSpage{P000811}}
\GSfield{The notes}{Under \ru{Для неполупростых групп} the notes write \(SO_0(m,n)\) and say there are no \ru{неунит.} representations with nonzero cohomology. The English translation preserves the contradiction.}
\GSfield{Statement}{No determinate theorem follows: most indicated orthogonal groups are semisimple; the dimensions and exceptional signatures are unspecified; and the abbreviation nonunitary cannot be silently changed to unitary. Possible intended statements differ substantially in their representation categories and in the exceptional rank-one cases.}
\GSfield{Reconstruction}{The scan confirms that the contradictory labels occur in the notebook. The nearby correct unitary classification is treated separately and does not authorize rewriting this sentence.}
\GSfield{Fate}{\citet[\S V(C), especially Theorem V.5]{delorme1977cohomologie} identifies the nearby unitary theory, but does not justify the literal nonsemisimple/nonunitary assertion.}
\GSfield{Present status}{As of 10 November 2026: Unidentified because the recorded category and group description conflict. No new vanishing theorem or open problem is extracted.}
\GSfield{Assessment}{literal statement: not determinable; publication match: thematic only; verification: unresolved lead; development: annotated; confidence: low}
\end{GSitem}
\begin{GSitem}{1978-450.13}{F}{Failure of central-character splitting in a continuous category}
\GSfield{Kind; source}{claim; \GSpage{P000811}}
\GSfield{The notes}{The first approach is said to violate the earlier infinitesimal-character lemma, with a counterexample for \(SL(2,\mathbb R)\). No representations, topology, or extension are written.}
\GSfield{Statement}{A counterexample in the continuous category would consist of two specified continuous representations with distinct scalar infinitesimal characters on their smooth vectors and a nonsplit continuous extension, despite splitting after passage to a smoother algebraic category. Differentiated central operators need not act as everywhere-defined bounded operators on the original spaces, so the algebraic polynomial-projection proof cannot automatically be used there. The notes omit the data needed to identify the claimed example.}
\GSfield{Reconstruction}{The category obstruction explains how the claim can coexist with the correct algebraic lemma; it is not a reconstruction of the missing \(SL(2,\mathbb R)\) example.}
\GSfield{Fate}{The category distinctions and relevant Pinczon--Simon/Guichardet work are documented in \citet[Introduction, \S4, and references 17--20]{charfi1980extensions}. Searches found Pinczon--Simon’s 1979 article, but its full text was inaccessible. Neither the retrieved abstract nor Charfi’s discussion identifies this exact counterexample.}
\GSfield{Present status}{As of 10 November 2026: The particular counterexample remains an unresolved identification. The general warning about transferring algebraic splitting to a continuous completion is valid.}
\GSfield{Assessment}{literal statement: not determinable; publication match: thematic only; verification: unresolved lead; development: annotated; confidence: medium}
\end{GSitem}
\begin{GSitem}{1978-450.14}{E}{Central-character splitting in the smooth category}
\GSfield{Kind; source}{theorem; \GSpage{P000811}}
\GSfield{The notes}{The second, smooth-vector approach is said to restore the infinitesimal-character lemma, although \(H^1(G,\operatorname{Hom}(E_\infty,F_\infty))\) is difficult to compute.}
\GSfield{Statement}{For an extension in a category of smooth Fréchet representations of a connected Lie group, suppose the derived central operators are continuous and have distinct scalar characters on the end modules. A separating central operator satisfies a polynomial with two distinct roots on the extension. Its polynomial spectral projections are continuous and commute with the group action, and yield a topological equivariant splitting. This argument uses the smooth category and cannot infer boundedness on the original Hilbert or Banach completions.}
\GSfield{Reconstruction}{The continuity and connectedness conditions make the earlier algebraic proof a valid topological proof. The qualitative difficulty of computation is not treated as a separate precise open problem.}
\GSfield{Fate}{The natural continuity of derived operators in smooth-vector topology is part of \citet[\S5.1]{blanc1979cohomologie}. The extension setting and central-character separation also appear in \citet[\S4, Remark 4.1 and reference 18]{charfi1980extensions}. The degree-one projection argument is explicitly reconstructed here.}
\GSfield{Present status}{As of 10 November 2026: Established for the stated smooth category. No equality with continuous extension groups between arbitrary completions is claimed.}
\GSfield{Assessment}{literal statement: correct with missing hypotheses; publication match: partial; verification: read the result; development: annotated; confidence: high}
\end{GSitem}
\begin{GSitem}{1978-450.15}{E}{K-finite vectors and relative algebraic extensions}
\GSfield{Kind; source}{method; \GSpage{P000811}}
\GSfield{The notes}{The third approach replaces \(E,F\) by \(E_K,F_K\), the vectors lying in finite-dimensional \(K\)-subrepresentations, and writes \(H^1(\mathfrak g,\mathfrak k;\operatorname{Hom}(E_K,F_K))\), algebra without topology. It says that for \(SL(2,\mathbb R)\) the groups can be computed for all irreducible Harish--Chandra modules.}
\GSfield{Statement}{For Harish--Chandra \((\mathfrak g,K)\)-modules, compact \(K\)-types are semisimple, so extensions split as algebraic \(K\)-modules. Their classes are computed by relative degree-one cocycles with coefficients in \(\operatorname{Hom}_{\mathbb C}(E_K,F_K)\), requiring \(K\)-equivariance. When \(K\) is connected, as \(SO(2)\) is for \(SL(2,\mathbb R)\), the corresponding \(\mathfrak k\)-condition is equivalent. The notebook gives no table of dimensions or labels for irreducibles, so no particular complete Ext table is asserted here.}
\GSfield{Reconstruction}{The Harish--Chandra category and \(K\)-equivariance are supplied. Merely taking finite-dimensional \(K\)-orbits in an arbitrary nonsmooth representation is insufficient to define a \(\mathfrak g\)-module; one uses the standard admissible smooth setting.}
\GSfield{Fate}{The relative cochain framework is checked in \citet[\S3]{wockel2014topological}, and the passage between analytic and \(K\)-finite vectors in the pertinent rank-one example is used in \citet[proof of Theorem V.5, pages 318--321]{delorme1977cohomologie}. These identify the method; an exact match for the unrecorded complete Ext table was not verified.}
\GSfield{Present status}{As of 10 November 2026: The algebraic computation method is established. The lecture’s unprinted table remains only a bibliographic lead, not a stated family of new numerical results.}
\GSfield{Assessment}{literal statement: correct with missing hypotheses; publication match: partial; verification: read the result; development: annotated; confidence: high}
\end{GSitem}
\begin{GSitem}{1978-450.16}{E}{Semidirect products and dual orbits}
\GSfield{Kind; source}{construction; \GSpage{P000811}}
\GSfield{The notes}{The meeting ends by defining \(G=B\rtimes A\), \(B=\mathbb R^n\), the duality with \(B^*\), and an \(A\)-orbit \(X\subset B^*\). No ensuing cohomology formula is present.}
\GSfield{Statement}{For a continuous linear action of a Lie group \(A\) on \(B=\mathbb R^n\), define the semidirect product and contragredient action \((a\cdot x)(b)=x(a^{-1}b)\) on \(B^*\). A point \(x\) gives the character \(b\mapsto e^{i\langle b,x\rangle}\), and its stabilizer \(A_x\) is the little group used in induced-representation constructions. A full classification or extension theorem needs further regularity and representation hypotheses; none is stated on this page.}
\GSfield{Reconstruction}{The dual-action convention and stabilizer are supplied to explain the setup. No theorem from the subsequent literature is retroactively inserted into the missing continuation.}
\GSfield{Fate}{This is the standard orbit/induction framework. \citet[Introduction, \S3, and reference 20]{charfi1980extensions} applies it and explicitly cites Guichardet’s preprint M.357-06-78, establishing June 1978 circulation of his related work. Guichardet’s later 1979 article was located bibliographically, but its text was not accessible in this research pass.}
\GSfield{Present status}{As of 10 November 2026: The setup is identified as prior representation theory. The surviving record contains no additional semidirect-product extension theorem to classify.}
\GSfield{Assessment}{literal statement: correct with missing hypotheses; publication match: partial; verification: read the result; development: annotated; confidence: high}
\end{GSitem}

\GStalk{451}{}{}{P000810-P000812 (catalogue range)}{Calendar-year divider: 1979}
This is a catalogue divider, \emph{YEAR 1979}, rather than a separate lecture. Its linked page range overlaps the conclusion of Guichardet's December lecture and the reading programme and Belavin material beginning with the dated heading 12 February 1979 on P000812. No independent mathematical occurrence is created from the divider. In particular, the earlier index's 1977-451.1--13 must not be interpreted as a talk by this record: those titles describe the overlapping Guichardet and Belavin pages. The actual occurrences are reassigned by the handwritten headings in the present chapter. The catalogue is verified at \citet{mccme2026gelfand}; its grouping is retained as provenance, not taken as evidence of authorship.


\GStalk{452}{1979-02-12}{I. M. Gelfand (catalogue attribution; reading list credited to A. A. Belavin)}{P000812 (upper part); catalogue P000810-P000813}{A reading programme for the physics discussions}
The heading \ru{Что нужно знать? (по А. А. Белавину)} introduces a reading programme: Feynman's lectures, statistical physics and path integrals; Huang on Onsager's solution; Klein on nineteenth-century mathematics; Dirac on monopoles, quantum mechanics and relativity; Akhiezer--Berestetskii; Yang--Mills; and Konopleva--Popov. The catalogue attributes the record to I.~M. Gelfand, but the page explicitly credits Belavin for the programme. No mathematical assertion is manufactured from this bibliography. The separate underlined \ru{Белавин} introduces record 453 on the same page. The notebook's Feynman coauthor \ru{Гиббсом} refers to Hibbs, not Gibbs; the identification is independently supported by the emended edition's editor \citep{feynmanhibbs1965path}.


\GStalk{453}{1979-02-12}{A. A. Belavin}{P000812 (from Belavin heading)-P000815; catalogue P000811-P000817}{Instanton matrices, lattice fermions and factorized scattering}
An instanton parametrization challenge precedes the main exposition of exactly solvable two-dimensional quantum field theories. The talk connects the finite XYZ chain with lattice fermions, the massive Thirring and sine-Gordon models, and Zamolodchikov's exchange algebra. Baxter's elliptic weights and commuting transfer matrices organize the hidden symmetries. Several displayed formulas have normalization or spectral-origin errors, while the suggestion that gauge averaging alone proves confinement overstates its conclusion. The ADHM paragraph and beginning of the gauge argument already appear in the supplied prior index as 1977-451.11--13; their source is this same dated page, so this is a provenance duplication, not an earlier historical occurrence.

\begin{GSitem}{1978-453.1}{E}{Quaternionic ADHM reality equations and quotient}
\GSfield{Kind; source}{construction; \GSpage{P000812}}
\GSfield{The notes}{For quaternionic \(N\times(N+1)\) matrices \(A,B\), put \(M(x)=xA+B\); require \(MM^+\) real for every quaternion \(x\), modulo \(M\mapsto GMV\), with \(VV^+=1\).}
\GSfield{Statement}{The nondegenerate quaternionic ADHM data for charge \(N\ge1\) consist, in the row-matrix convention, of an affine matrix \(M(x)\) such that \(M(x)M(x)^*\) is real and positive definite for all finite \(x\in\mathbb H\), with the corresponding full-rank condition at infinity. Quotient by \(GL(N,\mathbb R)\) on the left and \(Sp(N+1)\) on the right. Its one-dimensional quaternionic kernel defines an \(SU(2)\) instanton bundle with the orthogonally projected connection; gauge classes of smooth instantons on the round \(S^4\) arise this way, with self-duality sign depending on orientation.}
\GSfield{Reconstruction}{The full-rank condition and the precise groups are supplied from ADHM; quaternion multiplication conventions may be transposed/conjugated when changing row and column conventions. Here \(xA\) means left multiplication of each entry by x. It is compatible with the stated action: \(G(xA+B)V=x(GAV)+GBV\), since G is real and quaternion multiplication is associative. No commutation of x with entries of V is assumed. Passing to the adjoint column convention conjugates x and can reverse the self-duality orientation.}
\GSfield{Fate}{Already announced by \citet[pp.~185--187, reality condition (2) and equivalence paragraph]{atiyahhitchindrinfeldmanin1978instantons}, received 16 December 1977 and published 6 March 1978. The notebook omits the nondegeneracy and the real nature of the left change-of-basis group. Prior-index 1977-451.11 is the same source paragraph under an incorrect year assignment.}
\GSfield{Present status}{As of 10 November 2026: The ADHM classification is established. Reality alone also includes degenerate data and must not be equated with the smooth-instanton moduli space.}
\GSfield{Assessment}{literal statement: correct with missing hypotheses; publication match: partial; verification: read the result; development: annotated; confidence: high}
\end{GSitem}
\begin{GSitem}{1978-453.2}{E}{The unframed charge-N instanton dimension}
\GSfield{Kind; source}{theorem; \GSpage{P000812}}
\GSfield{The notes}{\ru{Всего должно быть \(8N-3\) параметров.}}
\GSfield{Statement}{The moduli space of irreducible charge-\(N\) \(SU(2)\) self-dual connections on the round \(S^4\), modulo the full gauge group, has real dimension \(8N-3\), for \(N\ge1\). Keeping a framing at one point instead gives dimension \(8N\).}
\GSfield{Reconstruction}{Irreducibility, the round sphere, positive charge, the full gauge quotient and the distinction from a framed moduli space are supplied hypotheses.}
\GSfield{Fate}{The dimension theorem predates the talk: \citet[p.~2662, theorem]{atiyahhitchinsinger1977deformations}; ADHM's parameter count is on p.~186 of \citet{atiyahhitchindrinfeldmanin1978instantons}. The former was contributed 20 April 1977 and published July 1977.}
\GSfield{Present status}{As of 10 November 2026: Established for the stated smooth unframed moduli problem; this is not a count of arbitrary solutions of the unqualified matrix equation.}
\GSfield{Assessment}{literal statement: correct with missing hypotheses; publication match: exact; verification: publisher metadata checked; development: annotated; confidence: high}
\end{GSitem}
\begin{GSitem}{1978-453.3}{F}{A beautiful parametrization of instanton matrices}
\GSfield{Kind; source}{problem; \GSpage{P000812}}
\GSfield{The notes}{The competition question asks what remains unsolved about instantons and requests a \ru{красивую параметризацию} of the pairs above.}
\GSfield{Statement}{The mathematical target identifiable from the page is the ADHM quotient. No meaning of ``beautiful'' is specified: the page does not demand global coordinates, rational coordinates, an atlas, an algorithm, or a topology theorem. Consequently no definite additional conjecture is recoverable.}
\GSfield{Fate}{ADHM already supplied a complete algebraic construction \citep[pp.~186--187]{atiyahhitchindrinfeldmanin1978instantons}. Searches for later parametrizations do not establish a solution to an unspecified aesthetic criterion. Prior-index 1977-451.12 is this same paragraph; GP-1977-05 concerns classification and connectedness and must not silently acquire new rationality/global-coordinate quantifiers.}
\GSfield{Present status}{As of 10 November 2026: Insufficiently specified for a separate open-problem entry or a bounded research development.}
\GSfield{Assessment}{literal statement: not determinable; publication match: thematic only; verification: read the result; development: annotated; confidence: high}
\end{GSitem}
\begin{GSitem}{1978-453.4}{E}{Gauge averaging forces a charged bilocal expectation to vanish}
\GSfield{Kind; source}{claim; \GSpage{P000812-P000813}}
\GSfield{The notes}{Local invariance gives \(\langle\psi(x)\overline\psi(y)\rangle=e^{i(\alpha(x)-\alpha(y))}\langle\psi(x)\overline\psi(y)\rangle\), hence zero.}
\GSfield{Statement}{In a finite-volume lattice gauge theory with compact local gauge group, gauge-invariant measure and action, no gauge fixing, and two distinct sites \(x,y\), an undressed bilocal transforming nontrivially at either site has zero expectation. For an Abelian unit-charge field, averaging an independent phase at \(x\) gives the displayed conclusion. The same argument does not apply to \(x=y\) or to a bilocal completed by a parallel transporter.}
\GSfield{Reconstruction}{The finite lattice, compact gauge group, invariant measure and action, and independent transformations at distinct sites specify a setting in which the averaging argument is valid. These hypotheses are not all stated in the notebook.}
\GSfield{Fate}{This is a gauge-averaging identity of the kind analysed by \citet[Sec.~I]{elitzur1975local}, received 29 July and published 15 December 1975. It is established independently of any confinement phase.}
\GSfield{Present status}{As of 10 November 2026: Correct with gauge-invariant integration and distinct sites; it supplies no confinement criterion.}
\GSfield{Assessment}{literal statement: correct with missing hypotheses; publication match: partial; verification: read the result; development: annotated; confidence: high}
\end{GSitem}
\begin{GSitem}{1978-453.5}{F}{The proposed confinement inference}
\GSfield{Kind; source}{claim; \GSpage{P000812-P000813}}
\GSfield{The notes}{The vanishing correlation is followed by \ru{невозможность ... для заряж. объектов}; the argument is labelled naive and attributed to Wilson.}
\GSfield{Statement}{The implication ``vanishing of an undressed charged correlator implies physical confinement'' is false without a gauge-invariant confinement observable and a dynamical hypothesis. Gauge averaging also vanishes in a Coulomb regime. A meaningful confinement problem must specify Wilson loops, charge sectors or long-range potentials and the theory and limit involved.}
\GSfield{Reconstruction}{The distinction between a gauge-dependent correlator and a gauge-invariant confinement criterion is an editorial correction; no area-law theorem is reconstructed from the paragraph.}
\GSfield{Fate}{\citet[Introduction and Sec.~IV]{elitzur1975local} explicitly separates local-gauge noninvariant expectation values from physical phase diagnostics. No definite four-dimensional confinement conjecture is formulated on this page; none is inserted here.}
\GSfield{Present status}{As of 10 November 2026: The inference fails; the incomplete source supplies no exact unsolved problem.}
\GSfield{Assessment}{literal statement: false as written; publication match: thematic only; verification: read the result; development: annotated; confidence: medium}
\end{GSitem}
\begin{GSitem}{1978-453.6}{F}{A finite cutoff can destroy a hidden symmetry}
\GSfield{Kind; source}{remark; \GSpage{P000813}}
\GSfield{The notes}{Passing from a functional integral to a finite multiple integral may break symmetry; the talk is introduced as a discussion of hidden symmetry.}
\GSfield{Statement}{A regulator need not preserve a symmetry of a formal continuum action. To turn this into a theorem one must specify the regulated measure, the transformations and convergence. No regulator or obstruction is identified here.}
\GSfield{Reconstruction}{No choice of ultraviolet regulator or precise hidden symmetry can be recovered. The modern statement only identifies the information needed to turn the sentence into a claim.}
\GSfield{Fate}{\citet[Introduction]{faddeev1996bethe} explains integrability-preserving lattice regularizations in the same programme, but this is only a thematic match to the notebook's general warning; no particular anomaly is asserted.}
\GSfield{Present status}{As of 10 November 2026: A methodological observation, not an identifiable open anomaly problem.}
\GSfield{Assessment}{literal statement: not determinable; publication match: thematic only; verification: read the result; development: annotated; confidence: medium}
\end{GSitem}
\begin{GSitem}{1978-453.7}{E}{A semiclassical sine-Gordon spectrum}
\GSfield{Kind; source}{claim; \GSpage{P000813}}
\GSfield{The notes}{The equation \(\Box\varphi=\sin\varphi\) is written and the notes say that a quasiclassical spectrum has been obtained.}
\GSfield{Statement}{For the quantized sine-Gordon model in \(1+1\) dimensions, with normal-ordered cosine interaction and a specified coupling/renormalization convention, semiclassical quantization treats solitons and their bound states by fluctuations about classical finite-energy solutions. The page records existence of such a calculation, not a numerical mass formula or an exact remainder estimate.}
\GSfield{Reconstruction}{The distinction between semiclassical asymptotics and an exact spectrum is supplied; no numerical mass formula absent from the notes is attributed to the speaker.}
\GSfield{Fate}{Pre-talk work is surveyed in \citet[abstract and soliton-quantization discussion]{faddeevkorepin1978quantum}, published June 1978; the exact scattering description was already announced by \citet{zamolodchikov1977exact}. These are not later solutions of a problem posed in February 1979.}
\GSfield{Present status}{As of 10 November 2026: Established semiclassical programme; exact quantum statements require a chosen normalization absent here.}
\GSfield{Assessment}{literal statement: correct with missing hypotheses; publication match: partial; verification: publisher metadata checked; development: annotated; confidence: high}
\end{GSitem}
\begin{GSitem}{1978-453.8}{E}{The finite periodic XYZ spectral problem}
\GSfield{Kind; source}{problem; \GSpage{P000813}}
\GSfield{The notes}{\(H=\sum_{n=1}^N\sum_{a=1}^3J_a\sigma_n^a\sigma_{n+1}^a\), \(\sigma_{N+1}=\sigma_1\), \(J_a\in\mathbb R\); find the eigenvectors and eigenvalues. The notes immediately say Baxter solved it.}
\GSfield{Statement}{For \(N\ge2\), the indicated Hermitian operator acts on \((\mathbb C^2)^{\otimes N}\). The problem is to diagonalize this finite matrix, including multiplicities and exceptional couplings, using the commuting eight-vertex transfer family. Algebraic diagonalizability is elementary; the substantive issue is an exact integrable description uniform in \(N\).}
\GSfield{Reconstruction}{The finite Hilbert space, multiplicity issue and exceptional-coupling qualification make the spectral question explicit; the notes themselves already call Baxter's problem solved.}
\GSfield{Fate}{Baxter's eight-vertex solution was announced in 1971 \citep{baxter1971eight}, with the eigenvector programme published in the 1973 series \citep{baxter1973eigenvectors}. This is explicitly presented as solved exposition, not a new conjecture. Publication matching for exceptional degeneracies is not claimed from the accessible metadata.}
\GSfield{Present status}{As of 10 November 2026: An established exactly solvable family. Singular parametrizations and completeness subtleties must be treated separately from the elementary existence of a spectral decomposition.}
\GSfield{Assessment}{literal statement: correct with missing hypotheses; publication match: partial; verification: publisher metadata checked; development: annotated; confidence: high; problem \hyperref[prob:GP-1978-11]{GP-1978-11}}
\end{GSitem}
\begin{GSitem}{1978-453.9}{E}{The chain-to-Thirring scaling limit}
\GSfield{Kind; source}{method; \GSpage{P000813-P000814}}
\GSfield{The notes}{Luther reduces the chain to the Thirring model; even and odd sites supply two fermion components, and the continuum theory is related to sine-Gordon.}
\GSfield{Statement}{A continuum low-energy limit of a suitably tuned anisotropic spin chain, after Jordan--Wigner transformation, staggered decomposition, lattice-spacing rescaling and coupling renormalization, yields a massive interacting Dirac theory of Thirring type. This is a scaling limit, not an equality of a fixed finite-dimensional XYZ Hamiltonian with a continuum field operator.}
\GSfield{Reconstruction}{The low-energy regime, tuning, cutoff and renormalization qualifications are supplied from the identified scaling-limit interpretation. No specified topology of convergence is claimed.}
\GSfield{Fate}{\citet[abstract, pp.~2153--2159]{luther1976eigenvalue} computes the excitation spectrum and identifies the corresponding massive Thirring/sine-Gordon couplings; received 19 April and published 1 September 1976.}
\GSfield{Present status}{As of 10 November 2026: Established correspondence in the stated scaling regime; the notebook omits the scaling prescription and cannot support convergence in a specified operator topology.}
\GSfield{Assessment}{literal statement: correct with missing hypotheses; publication match: partial; verification: publisher metadata checked; development: annotated; confidence: high}
\end{GSitem}
\begin{GSitem}{1978-453.10}{E}{Pauli spin algebra}
\GSfield{Kind; source}{definition; \GSpage{P000813}}
\GSfield{The notes}{The three standard Pauli matrices are displayed, but (2) writes \([\sigma_n^a,\sigma_n^b]=\varepsilon_{abc}\sigma_n^c\).}
\GSfield{Statement}{With exactly the displayed matrices, \([\sigma_n^a,\sigma_m^b]=2i\delta_{nm}\sum_c\varepsilon_{abc}\sigma_n^c\) and \(\{\sigma_n^a,\sigma_n^b\}=2\delta_{ab}I\). Tensor factors at distinct sites commute.}
\GSfield{Reconstruction}{The corrected factor and different-site commutator are obtained directly from the matrices printed in the notes.}
\GSfield{Fate}{This is the standard local spin algebra used in \citet[Sec.~3]{faddeev1996bethe}; the factor \(2i\) is also checked directly from the displayed \(2\times2\) matrices. The scan faithfully lacks it.}
\GSfield{Present status}{As of 10 November 2026: The corrected identities are exact; the error is mathematical shorthand in the notebook, not an English translation error.}
\GSfield{Assessment}{literal statement: false as written; publication match: partial; verification: read the result; development: annotated; confidence: high}
\end{GSitem}
\begin{GSitem}{1978-453.11}{E}{Jordan--Wigner fermions and their normalization}
\GSfield{Kind; source}{construction; \GSpage{P000813}}
\GSfield{The notes}{\(\psi_n^\pm=(\sigma_n^1\pm i\sigma_n^2)\prod_{m<n}\sigma_m^3\) is claimed to satisfy \(\{\psi_n^-,\psi_m^+\}=\delta_{nm}\).}
\GSfield{Statement}{For a finite ordered spin chain, define \(c_n=\frac12(\sigma_n^1-i\sigma_n^2)\prod_{m<n}\sigma_m^3\). Then \(\{c_n,c_m\}=0\), \(\{c_n,c_m^\dagger\}=\delta_{nm}I\). Without the factor \(1/2\), the latter right side is \(4\delta_{nm}I\). A periodic spin chain produces a fermion-parity-dependent boundary term.}
\GSfield{Reconstruction}{The factor one-half and periodic-boundary parity qualification are supplied; the normalization is checked directly from the printed matrices.}
\GSfield{Fate}{The construction and fermionic solution of chain models are pre-talk material \citep{liebschultzmattis1961soluble}. The normalization follows directly from the displayed Pauli matrices; the missing periodic parity condition is essential for a finite-chain spectrum.}
\GSfield{Present status}{As of 10 November 2026: The canonical anticommutation relations hold exactly after correction, with parity sectors retained at periodic boundaries.}
\GSfield{Assessment}{literal statement: false as written; publication match: partial; verification: publisher metadata checked; development: annotated; confidence: high}
\end{GSitem}
\begin{GSitem}{1978-453.12}{E}{Quadratic fermion diagonalization of the transverse Ising chain}
\GSfield{Kind; source}{method; \GSpage{P000813-P000814}}
\GSfield{The notes}{Expressing spins through fermions makes the Hamiltonian quadratic; the concrete example is \(H=\sum_n(\sigma_n^1\sigma_{n+1}^1+\Delta\sigma_n^3)\). A Fourier transform is proposed.}
\GSfield{Statement}{For the transverse Ising chain the Jordan--Wigner image is quadratic in creation and annihilation operators, and a Fourier--Bogoliubov transformation diagonalizes each fixed parity sector. For a general XYZ chain the \(J_3\sigma_n^3\sigma_{n+1}^3\) term yields quartic density interactions, so the quadratic assertion is restricted to the Ising/free-fermion example.}
\GSfield{Reconstruction}{Quadratic fermions, the absence of a longitudinal interaction and the treatment of the boundary parity sector specify the solvable specialization.}
\GSfield{Fate}{This is part of the pre-1979 free-fermion method \citep{liebschultzmattis1961soluble}. Onsager's 1944 solution uses a related transfer/Clifford structure \citep[Secs.~II--VII]{onsager1944crystal}; the notebook compresses these related implementations into one historical sentence.}
\GSfield{Present status}{As of 10 November 2026: Exact in the specified free-fermion model, not for arbitrary XYZ couplings.}
\GSfield{Assessment}{literal statement: correct with missing hypotheses; publication match: partial; verification: read the result; development: annotated; confidence: high}
\end{GSitem}
\begin{GSitem}{1978-453.13}{E}{Heisenberg evolution preserves the algebra, but the displayed equation is schematic}
\GSfield{Kind; source}{claim; \GSpage{P000814}}
\GSfield{The notes}{Equation (7) is \(\dot\sigma=i[\sigma,H]\); after fermionization (8) becomes \(\dot\psi_n=\psi_{n+1}-\psi_{n-1}\).}
\GSfield{Statement}{Unitary conjugation preserves every polynomial operator relation. With the usual convention \(U(t)=e^{-itH}\), Heisenberg evolution satisfies \(\dot O=i[H,O]\). For the transverse Ising chain, the fermion equations are linear in both \(c\) and \(c^\dagger\), and depend on \(\Delta\); the printed single-component, \(\Delta\)-independent equation is not its general equation of motion.}
\GSfield{Reconstruction}{The time-evolution convention is stated explicitly. Dependence on the transverse field and coupling to adjoint fermions are corrections to the schematic displayed equation.}
\GSfield{Fate}{The quadratic solution framework is already in \citet{liebschultzmattis1961soluble}. The missing \(\Delta\) term and adjoint coupling follow by direct commutation with (6), so no later theorem repairs (8) as a literal statement.}
\GSfield{Present status}{As of 10 November 2026: Algebra preservation is exact. Equation (8) requires a different specialization or variables not supplied by the notes.}
\GSfield{Assessment}{literal statement: false as written; publication match: partial; verification: publisher metadata checked; development: annotated; confidence: high}
\end{GSitem}
\begin{GSitem}{1978-453.14}{E}{Interacting fermions and the continuum quartic term}
\GSfield{Kind; source}{construction; \GSpage{P000814}}
\GSfield{The notes}{The original chain gives a schematic quartic Hamiltonian (9), followed by a continuum expression with \(\psi_1^+\psi_2\psi_1^+\psi_2^+\).}
\GSfield{Statement}{A nonzero longitudinal XYZ coupling produces a density-density term \(n_n n_{n+1}\), in addition to quadratic hopping and, for transverse anisotropy, pairing. A continuum Thirring interaction is a regularized current-current product \(\frac g2 j_\mu j^\mu\), \(j^\mu=\overline\psi\gamma^\mu\psi\). The printed repeated \(\psi_1^+\) expression does not specify that product and cannot be repaired uniquely without conventions and point splitting.}
\GSfield{Reconstruction}{The regularized current-current expression is the identified continuum interaction, not a unique emendation of the malformed monomial in the notebook.}
\GSfield{Fate}{The lattice/continuum identification predates the talk \citep{luther1976eigenvalue,coleman1975equivalence}. The quartic monomial is retained as a faithful source error or unresolved shorthand, rather than silently replaced with a guessed sequence of indices.}
\GSfield{Present status}{As of 10 November 2026: The current-current interaction is established; the exact literal continuum expression is not a usable Hamiltonian.}
\GSfield{Assessment}{literal statement: false as written; publication match: partial; verification: publisher metadata checked; development: annotated; confidence: high}
\end{GSitem}
\begin{GSitem}{1978-453.15}{E}{Massive Thirring and quantum sine-Gordon equivalence}
\GSfield{Kind; source}{theorem; \GSpage{P000814}}
\GSfield{The notes}{The notes attribute the 1975 equivalence to Mandelstam and an omitted name.}
\GSfield{Statement}{In Coleman's normal-ordering convention, the massive Thirring and sine-Gordon theories have matching neutral-sector correlation functions with \(4\pi/\beta^2=1+g/\pi\), after matching the mass parameters and operator normalization. The equivalence is a quantum field-theoretic statement with regularization; it is not established by equating their unregularized classical equations.}
\GSfield{Reconstruction}{The neutral sector, normal ordering, operator normalization and coupling convention are supplied from Coleman's formulation.}
\GSfield{Fate}{The omitted name is Sidney Coleman: \citet[abstract, pp.~2088--2097]{coleman1975equivalence}, received 6 January and published 15 April 1975; \citet{mandelstam1975soliton}, received 24 February and published 15 May 1975, constructs the soliton fields.}
\GSfield{Present status}{As of 10 November 2026: A standard equivalence with explicitly specified operator and coupling conventions; no open problem was posed here.}
\GSfield{Assessment}{literal statement: correct with missing hypotheses; publication match: exact; verification: publisher metadata checked; development: annotated; confidence: high}
\end{GSitem}
\begin{GSitem}{1978-453.16}{E}{Mandelstam soliton operators and canonical factors}
\GSfield{Kind; source}{construction; \GSpage{P000814}}
\GSfield{The notes}{Equations (12)--(14) write \([\dot\varphi(x),\varphi(y)]=\delta(x-y)\), an unnormal-ordered exponential, and \(\alpha\beta=4\pi\).}
\GSfield{Statement}{For a canonically normalized real boson, \([\varphi(t,x),\Pi(t,y)]=i\delta(x-y)\), hence \([\Pi,\varphi]=-i\delta\). Fermion fields can be represented by regulated normal-ordered vertex/disorder operators of the schematic form \(C:\!\exp[-i(\frac{2\pi}{\beta}\int_{-\infty}^x\Pi\,dy\pm\frac\beta2\varphi(x))]\!:\), with component/Klein factors and cutoff normalization. Their exchange signs and delta-function normalization depend on this full prescription.}
\GSfield{Reconstruction}{This regulated vertex operator is a schematic reconstruction. The notebook does not supply its component signs, cutoff or Klein factors. No unique repair of its alpha-beta relation is asserted.}
\GSfield{Fate}{\citet[abstract and operator construction]{mandelstam1975soliton} explicitly constructs operators satisfying the Thirring equations and anticommutation relations. The scan confirms the missing imaginary factors and gives no normal-ordering prescription; its \(\alpha\beta=4\pi\) cannot be accepted with the unqualified displayed conventions.}
\GSfield{Present status}{As of 10 November 2026: Established regulated bosonization; the raw notebook exponential and CCR are not literal quantum-operator identities.}
\GSfield{Assessment}{literal statement: false as written; publication match: partial; verification: publisher metadata checked; development: annotated; confidence: high}
\end{GSitem}
\begin{GSitem}{1978-453.17}{E}{The Zamolodchikov exchange algebra}
\GSfield{Kind; source}{definition; \GSpage{P000814-P000815}}
\GSfield{The notes}{Generators \(A_i(\theta)\) exchange through \(S(\theta_1-\theta_2)\); the printed range \(1\le i<M\) is inconsistent with the following \(M=2\) example using \(A_2\).}
\GSfield{Statement}{For \(V=\mathbb C^M\), use indices \(1\le i\le M\). Introduce generators \(A_i(\theta)\) with quadratic exchange rules determined by a meromorphic two-body matrix \(S(u)\in\operatorname{End}(V\otimes V)\), where the output generators have interchanged rapidities. This defines an associative quotient algebra; requiring noncollapsed, consistently ordered scattering states imposes further conditions on \(S\).}
\GSfield{Reconstruction}{The inclusive index range, fixed permutation convention and distinction between an associative quotient and independent ordered states are editorial clarifications.}
\GSfield{Fate}{The factorized scattering framework predates the talk \citep[Secs.~1--3]{zamolodchikovzamolodchikov1978on}; the sine-Gordon two-particle solution appeared in 1977 \citep{zamolodchikov1977exact}.}
\GSfield{Present status}{As of 10 November 2026: Established algebraic framework after the index-range correction and the distinction between quotient associativity and nondegenerate exchange consistency.}
\GSfield{Assessment}{literal statement: false as written; publication match: partial; verification: read the result; development: annotated; confidence: high}
\end{GSitem}
\begin{GSitem}{1978-453.18}{E}{Three-particle consistency and factorization}
\GSfield{Kind; source}{claim; \GSpage{P000815}}
\GSfield{The notes}{Associativity of a three-generator product is said to impose the condition on \(S\) that leads to factorized multiparticle scattering.}
\GSfield{Statement}{If three-particle ordered states with distinct rapidities are linearly independent and pair exchanges act by invertible two-body matrices, consistency of the two reorderings gives \(R_{12}(u)R_{13}(u+v)R_{23}(v)=R_{23}(v)R_{13}(u+v)R_{12}(u)\), with \(R\) defined according to a fixed permutation convention. The Yang--Baxter equation alone does not establish absence of particle production in a given QFT.}
\GSfield{Reconstruction}{Invertibility, linear independence and absence of particle production are additional hypotheses; associativity of an arbitrary quotient alone is insufficient.}
\GSfield{Fate}{The consistency equations are explicit in \citet[Sec.~2]{zamolodchikovzamolodchikov1978on}; a later precise matrix presentation is \citet[eqs.~(1)--(4)]{belavin1980discrete}. These sources distinguish the algebraic constraint from the physical identification of a model.}
\GSfield{Present status}{As of 10 November 2026: Established under the stated factorized-scattering hypotheses; merely forming an associative quotient does not force a Yang--Baxter matrix.}
\GSfield{Assessment}{literal statement: correct with missing hypotheses; publication match: partial; verification: read the result; development: annotated; confidence: high}
\end{GSitem}
\begin{GSitem}{1978-453.19}{E}{Baxter elliptic weights and their limited completeness claim}
\GSfield{Kind; source}{construction; \GSpage{P000815}}
\GSfield{The notes}{The four weights are \(a=\operatorname{sn}(\theta+\eta,k)\), \(b=\operatorname{sn}(\theta-\eta,k)\), \(c=\operatorname{sn}(2\eta,k)\), \(d=kabc\); these are said to be all solutions up to a factor.}
\GSfield{Statement}{In the zero-field eight-vertex ansatz \(R(u)=\begin{psmallmatrix}a&0&0&d\\0&b&c&0\\0&c&b&0\\d&0&0&a\end{psmallmatrix}\), a standard regular parametrization is \(a=\operatorname{sn}(u+\eta_0,k)\), \(b=\operatorname{sn}(u,k)\), \(c=\operatorname{sn}(\eta_0,k)\), \(d=kabc\), divided by \(c\ne0\). Then \(R(0)=P\). The notebook form is obtained with \(u=\theta-\eta\), \(\eta_0=2\eta\); that spectral-origin change must also be respected in the addition law. No unrestricted classification of all \(4\times4\) Yang--Baxter solutions follows.}
\GSfield{Reconstruction}{The regular spectral origin and nonzero normalizing denominator are supplied explicitly. The restricted eight-vertex ansatz does not support the literal unrestricted completeness claim.}
\GSfield{Fate}{The eight-vertex solution is pre-talk Baxter material \citep{baxter1971eight}. Belavin later constructed higher-dimensional solutions with additional discrete symmetry \citep[pp.~189--200]{belavin1981dynamical}; this is not evidence for the notebook's unrestricted ``all solutions'' sentence.}
\GSfield{Present status}{As of 10 November 2026: The elliptic eight-vertex family is established; origin, degenerations, gauge transformations and imposed symmetries delimit any completeness assertion.}
\GSfield{Assessment}{literal statement: false as written; publication match: partial; verification: publisher metadata checked; development: annotated; confidence: high}
\end{GSitem}
\begin{GSitem}{1978-453.20}{E}{Transfer matrices and commuting Hamiltonians}
\GSfield{Kind; source}{construction; \GSpage{P000815}}
\GSfield{The notes}{The cyclic contraction of \(N\) copies of \(S\) defines \(T(\theta)\), a \(2^N\times2^N\) matrix, and satisfies \([T(\theta),H]=[T(\theta_1),T(\theta_2)]=0\).}
\GSfield{Statement}{For a meromorphic, generically invertible difference-form \(R\) satisfying Yang--Baxter, define \(t(u)=\operatorname{tr}_a R_{aN}(u)\cdots R_{a1}(u)\). Then \([t(u),t(v)]=0\). If \(R(0)=P\), the logarithmic derivative \(t(0)^{-1}t'(0)\) is a periodic nearest-neighbour Hamiltonian, up to scalar normalization/additive constants, and commutes with every \(t(u)\). At fixed \(N\), an infinite Taylor list need not be algebraically independent.}
\GSfield{Reconstruction}{Generic invertibility, regularity, the auxiliary trace and the restriction on independence at fixed chain length are supplied to make the standard construction precise.}
\GSfield{Fate}{The pre-talk transfer construction is recalled explicitly by \citet[eqs.~(3)--(6)]{belavin1980discrete}; the complete RTT/trace derivation and Hamiltonian derivative are given by \citet[Sec.~3, eqs.~(44)--(65)]{faddeev1996bethe}, based on his 1995 lectures, first posted May 1996.}
\GSfield{Present status}{As of 10 November 2026: Established finite-dimensional identities. The phrase ``infinite set of integrals'' describes a generating series or varying-chain hierarchy, not infinitely many independent operators on a fixed finite chain.}
\GSfield{Assessment}{literal statement: correct with missing hypotheses; publication match: partial; verification: read the result; development: annotated; confidence: high}
\end{GSitem}

\GStalk{454}{1979-02-19}{A. A. Belavin}{P000816-P000820; catalogue P000815-P000820}{Exchange algebras, transfer matrices and the spin-1 question}
The continuation develops the Yang--Baxter consistency condition, commuting transfer matrices, the eight-vertex ansatz and monodromy multiplication. It asks for anisotropic parameters in a displayed spin-one Hamiltonian and ends with an unfinished Ising transfer-matrix assertion. Literal errors matter: the sample functional equation is incompatible with the stated regularity unless its elliptic weight vanishes, and the spin-one density fails a necessary regular difference-form Yang--Baxter condition under the commutator normalization explicitly given on P000819. The later Fateev--Zamolodchikov model supplies a nearby correctly normalized construction; it is not a solution of the printed density. The chapter's separately marked algebraic audit makes that distinction testable.

\begin{GSitem}{1978-454.1}{E}{Exchange generators and associativity}
\GSfield{Kind; source}{definition; \GSpage{P000816}}
\GSfield{The notes}{The \ru{Алгебра Замолодчикова} has generators \(A_i(\theta)\), \(1\le i\le M\), quadratic exchange rules and a requirement of associativity.}
\GSfield{Statement}{For \(V=\mathbb C^M\), an exchange matrix \(S(u)\) specifies relations between generators with interchanged spectral parameters. Their quotient of the tensor algebra is associative by construction. A nondegenerate ordered-state interpretation requires that the exchange relations do not introduce inconsistent additional dependencies.}
\GSfield{Reconstruction}{The tensor-algebra and nondegeneracy distinction clarifies the notebook's shorthand.}
\GSfield{Fate}{The construction is pre-talk factorized-scattering material \citep[Sec.~2]{zamolodchikovzamolodchikov1978on}. The distinction between quotient associativity and a basis of ordered states is essential to the next assertion.}
\GSfield{Present status}{As of 10 November 2026: A standard construction; associativity alone is not a test for a Yang--Baxter matrix.}
\GSfield{Assessment}{literal statement: correct with missing hypotheses; publication match: partial; verification: read the result; development: annotated; confidence: high}
\end{GSitem}
\begin{GSitem}{1978-454.2}{E}{Rapidity ordering and collision diagrams}
\GSfield{Kind; source}{method; \GSpage{P000816}}
\GSfield{The notes}{The parameters \(\theta_i\) are called velocities; two collision histories are identified because only the initial and final order matters.}
\GSfield{Statement}{For massive relativistic particles in \(1+1\) dimensions, rapidity satisfies \(E=m\cosh\theta\), \(p=m\sinh\theta\), \(v=\tanh\theta\). A difference-form two-body amplitude depends on rapidity differences. Equating the two three-particle collision orderings presupposes elastic factorized scattering with well-defined asymptotic states.}
\GSfield{Reconstruction}{Rapidity replaces literal velocity; the scattering assumptions are made explicit.}
\GSfield{Fate}{These are the physical hypotheses of \citet[Secs.~1--2]{zamolodchikovzamolodchikov1978on}, published February 1978. Its authors expressly distinguish a consistent proposed S-matrix from identification with a particular field theory.}
\GSfield{Present status}{As of 10 November 2026: Established kinematics and conditional factorization argument, not a proof that every two-dimensional theory factorizes.}
\GSfield{Assessment}{literal statement: correct with missing hypotheses; publication match: partial; verification: read the result; development: annotated; confidence: high}
\end{GSitem}
\begin{GSitem}{1978-454.3}{E}{Three-body reordering and the Yang--Baxter equation}
\GSfield{Kind; source}{claim; \GSpage{P000816-P000817}}
\GSfield{The notes}{Two reorderings of \(A_i(\theta_1)A_j(\theta_2)A_k(\theta_3)\) lead to equation (2); the last generator on P000816 is written \(A_i(\theta)\).}
\GSfield{Statement}{For linearly independent ordered three-particle states and a fixed unbraided convention, consistency is \(R_{12}(u)R_{13}(u+v)R_{23}(v)=R_{23}(v)R_{13}(u+v)R_{12}(u)\). The last factor in the displayed reordering must carry the final contracted index and original rapidity, \(A_{i_2}(\theta_1)\).}
\GSfield{Reconstruction}{The leg convention, independence and repaired dummy indices are supplied; the notebook's component indices are not silently treated as reliable.}
\GSfield{Fate}{The matrix relation is already present in the factorized-scattering programme \citep[Sec.~2]{zamolodchikovzamolodchikov1978on}; \citet[eqs.~(1)--(4)]{belavin1980discrete} gives an explicit later matrix formulation.}
\GSfield{Present status}{As of 10 November 2026: The corrected Yang--Baxter identity is standard; a relation in a collapsed quotient would not imply it.}
\GSfield{Assessment}{literal statement: false as written; publication match: partial; verification: read the result; development: annotated; confidence: high}
\end{GSitem}
\begin{GSitem}{1978-454.4}{E}{The component count of the consistency equations}
\GSfield{Kind; source}{remark; \GSpage{P000817}}
\GSfield{The notes}{There are \(M^6\) equations for \(M^4\) functions \(S_{ik}^{\alpha\beta}(\theta)\).}
\GSfield{Statement}{An endomorphism of \(V^{\otimes3}\) has \(M^6\) matrix entries, whereas one of \(V^{\otimes2}\) has \(M^4\). Thus the Yang--Baxter identity supplies this many displayed component equations for this many scalar matrix functions. This is a raw count, with no claim that the equations are independent.}
\GSfield{Reconstruction}{The distinction between a raw component count and independent constraints is supplied.}
\GSfield{Fate}{The count follows directly from the matrix formulation \citep[eqs.~(1)--(4)]{belavin1980discrete}; symmetries and identities can substantially reduce the system.}
\GSfield{Present status}{As of 10 November 2026: An exact dimension count with no classification consequence by itself.}
\GSfield{Assessment}{literal statement: correct; publication match: exact; verification: read the result; development: annotated; confidence: high}
\end{GSitem}
\begin{GSitem}{1978-454.5}{E}{Cyclic contraction and commuting transfer matrices}
\GSfield{Kind; source}{construction; \GSpage{P000817}}
\GSfield{The notes}{Equation (3) cyclically contracts \(N\) copies of \(S\) into an \(M^N\times M^N\) matrix; (4) asserts pairwise commutation.}
\GSfield{Statement}{For a generically invertible meromorphic Yang--Baxter matrix on \(V\otimes V\), the periodic transfer matrix \(t(u)=\operatorname{tr}_aR_{aN}(u)\cdots R_{a1}(u)\) satisfies \([t(u),t(v)]=0\). Invertibility is used in the auxiliary-space intertwining argument; the identity then continues meromorphically.}
\GSfield{Reconstruction}{The auxiliary trace, analytic domain and invertibility hypothesis supply the proof's missing structure.}
\GSfield{Fate}{This is pre-talk transfer-matrix theory. A complete RTT and trace proof appears in \citet[Sec.~3, eqs.~(44)--(65)]{faddeev1996bethe}, lectures of 1995 posted in May 1996; Belavin recalls it in \citet[eqs.~(3)--(6)]{belavin1980discrete}.}
\GSfield{Present status}{As of 10 November 2026: Established, including inhomogeneous variants not needed for the literal notes.}
\GSfield{Assessment}{literal statement: correct with missing hypotheses; publication match: partial; verification: read the result; development: annotated; confidence: high}
\end{GSitem}
\begin{GSitem}{1978-454.6}{E}{Conservation laws as a route to factorization}
\GSfield{Kind; source}{claim; \GSpage{P000817}}
\GSfield{The notes}{\ru{Факториз. \(S\)-матр. --- следствие кон. числа законов сохр.}}
\GSfield{Statement}{In a massive relativistic scattering theory, suitable nontrivial higher-spin local conserved charges, acting additively and with nondegenerate momentum dependence on asymptotic particles, can constrain scattering to elastic factorized processes. Merely having a finite number of conserved quantities, for example energy and momentum alone, is not sufficient.}
\GSfield{Reconstruction}{Higher-spin, locality, additivity and nondegeneracy are necessary qualifications to the identified argument; the exact minimal list of charges is not recoverable here.}
\GSfield{Fate}{\citet[Sec.~1]{zamolodchikovzamolodchikov1978on} states the physical higher-conservation-law setting. The notebook's unrestricted finite-number assertion has no matched theorem with its literal quantifiers.}
\GSfield{Present status}{As of 10 November 2026: The conditional mechanism is established. The bare finite-number assertion is false.}
\GSfield{Assessment}{literal statement: false as written; publication match: partial; verification: read the result; development: annotated; confidence: high}
\end{GSitem}
\begin{GSitem}{1978-454.7}{E}{The four-weight eight-vertex ansatz}
\GSfield{Kind; source}{construction; \GSpage{P000817}}
\GSfield{The notes}{For \(M=2\), four functions \(a,b,c,d\) are retained from sixteen possible entries; the note mentions five or six more general functions and parenthetically Maxwell's equations.}
\GSfield{Statement}{The zero-field eight-vertex ansatz is \(R=\begin{pmatrix}a&0&0&d\\0&b&c&0\\0&c&b&0\\d&0&0&a\end{pmatrix}\). Homogeneity of Yang--Baxter permits division by a nonzero scalar function, so \(c=1\) is a local normalization wherever \(c\ne0\). Neither a five-or-six-function classification nor a Maxwell equation identification is specified.}
\GSfield{Reconstruction}{The symmetry-restricted matrix and the domain of scalar normalization are made explicit. The Maxwell parenthesis is retained without interpretation.}
\GSfield{Fate}{The restricted ansatz is Baxter's pre-talk construction \citep{baxter1971eight}. The unelaborated parenthesis and numerical generality claim have no exact publication match in the sources checked.}
\GSfield{Present status}{As of 10 November 2026: Established eight-vertex specialization; the omitted general ansatz cannot be reconstructed uniquely.}
\GSfield{Assessment}{literal statement: correct with missing hypotheses; publication match: partial; verification: publisher metadata checked; development: annotated; confidence: high}
\end{GSitem}
\begin{GSitem}{1978-454.8}{F}{An inconsistent sample functional equation}
\GSfield{Kind; source}{claim; \GSpage{P000817}}
\GSfield{The notes}{Equation (6), after setting \(c=1\), is \(a(u)a(u+v)d(v)=d(u+v)a(v)+b(u)b(u+v)d(v)\). Equation (7) requires \(R(0)=P\).}
\GSfield{Statement}{In the displayed four-weight convention regularity means \(a(0)=c(0)=1\), \(b(0)=d(0)=0\). Substituting \(v=0\) in the printed (6) yields \(0=d(u)\). Thus the printed equation cannot describe a regular elliptic eight-vertex family with nonzero \(d\).}
\GSfield{Reconstruction}{The contradiction is an elementary substitution using the notebook's own initial condition; only the diagnosis, not a unique intended equation, is reconstructed.}
\GSfield{Fate}{Baxter's elliptic family exists \citep{baxter1971eight}; the conflict is internal to the notebook. The precise missing factor or argument in (6) has not been uniquely recovered, so no guessed correction replaces it.}
\GSfield{Present status}{As of 10 November 2026: Literal inconsistency established. The exact intended component equation remains unidentified.}
\GSfield{Assessment}{literal statement: false as written; publication match: none found; verification: unresolved lead; development: annotated; confidence: high}
\end{GSitem}
\begin{GSitem}{1978-454.9}{E}{Differential reduction and elliptic solutions}
\GSfield{Kind; source}{method; \GSpage{P000817-P000818}}
\GSfield{The notes}{Differentiating the functional system is said to yield four displayed ODEs and two omitted ones, with \(a'(0)=\alpha\), \(b'(0)=\beta\), \(d'(0)=\gamma\). Two independent algebraic constraints then lead to elliptic functions.}
\GSfield{Statement}{For a differentiable regular Yang--Baxter solution, differentiating the full functional equations at one parameter zero gives necessary ODEs, with first derivatives encoding local couplings. Integrating an incomplete ODE list is not sufficient for the original functional equation. The notebook's (6) already fails the regularity test, and the two unprinted ODEs and algebraic constraints cannot be recovered solely from this page.}
\GSfield{Reconstruction}{The necessary-versus-sufficient distinction and requirement to check the full functional equation are supplied. No omitted equation is invented.}
\GSfield{Fate}{The elliptic eight-vertex result is known before the talk \citep{baxter1971eight}; the precise displayed differential system is only a partial, faulty transcription of that method. It is not used here to assert a new classification.}
\GSfield{Present status}{As of 10 November 2026: The differential method is standard; the notebook's incomplete system does not identify an independently verifiable theorem.}
\GSfield{Assessment}{literal statement: not determinable; publication match: partial; verification: publisher metadata checked; development: annotated; confidence: high}
\end{GSitem}
\begin{GSitem}{1978-454.10}{F}{Representing an unspecified algebra of four functions}
\GSfield{Kind; source}{question; \GSpage{P000818}}
\GSfield{The notes}{A noncommutative, nonassociative algebra of quadruples is said to arise. Find an explicit realization, for example by matrices. Manin suggests torsion points of an elliptic curve.}
\GSfield{Statement}{The page specifies neither a multiplication of quadruples nor the relations to be represented. Ordinary matrix multiplication is associative, so a faithful realization of a genuinely nonassociative algebra would require a different operation. Torsion points alone do not define such an operation.}
\GSfield{Reconstruction}{Only the missing structural data are identified. No later Sklyanin or elliptic algebra is substituted for the unspecified object.}
\GSfield{Fate}{Belavin's later elliptic constructions use discrete groups \citep[eqs.~(7)--(17)]{belavin1980discrete}; his 1981 paper concerns dynamical symmetry \citep{belavin1981dynamical}. These are thematic leads, not an exact identification of this missing algebra.}
\GSfield{Present status}{As of 10 November 2026: No definite representability problem can be extracted. It is not classified as an unresolved modern conjecture.}
\GSfield{Assessment}{literal statement: not determinable; publication match: thematic only; verification: read the result; development: annotated; confidence: medium}
\end{GSitem}
\begin{GSitem}{1978-454.11}{E}{Commuting Taylor coefficients}
\GSfield{Kind; source}{theorem; \GSpage{P000818}}
\GSfield{The notes}{From \(T(u)=\sum_{\ell\ge0}I_\ell u^\ell\) the notes infer \([I_n,I_m]=0\).}
\GSfield{Statement}{If a matrix-valued function \(T\) is analytic near zero and \([T(u),T(v)]=0\) for independent \(u,v\), coefficient comparison gives \([I_n,I_m]=0\) for all \(n,m\). The same holds formally for a commuting formal power series.}
\GSfield{Reconstruction}{Analyticity or formal-series interpretation is supplied.}
\GSfield{Fate}{This is a direct consequence of the established transfer-matrix relation \citep[Sec.~3]{faddeev1996bethe}, with no later mathematical problem attached.}
\GSfield{Present status}{As of 10 November 2026: Exact coefficient identity; infinitely many coefficients need not be independent for fixed finite N.}
\GSfield{Assessment}{literal statement: correct; publication match: exact; verification: read the result; development: annotated; confidence: high}
\end{GSitem}
\begin{GSitem}{1978-454.12}{E}{The XYZ Hamiltonian as a logarithmic derivative}
\GSfield{Kind; source}{construction; \GSpage{P000818}}
\GSfield{The notes}{Equation (11) gives \(I_0^{-1}I_1=T(0)^{-1}T'(0)=\sum_{n,a}J_a\sigma_n^a\sigma_{n+1}^a\), and all \(I_\ell\) commute with it.}
\GSfield{Statement}{For an analytic regular eight-vertex matrix with \(R(0)=P\), the transfer matrix at zero is the cyclic shift, hence invertible. Its first logarithmic derivative is a sum of translated copies of \(h=PR'(0)\), which has XYZ form up to a scalar identity and convention-dependent coupling normalization.}
\GSfield{Reconstruction}{Regularity, the shift and the scalar normalization are supplied.}
\GSfield{Fate}{This pre-talk construction is recorded in \citet[eqs.~(3)--(6)]{belavin1980discrete}; the regular-point derivation is spelled out in \citet[Sec.~3]{faddeev1996bethe}.}
\GSfield{Present status}{As of 10 November 2026: Established local Hamiltonian construction from the regular transfer family.}
\GSfield{Assessment}{literal statement: correct with missing hypotheses; publication match: partial; verification: read the result; development: annotated; confidence: high}
\end{GSitem}
\begin{GSitem}{1978-454.13}{E}{Local three-site charges versus raw Taylor coefficients}
\GSfield{Kind; source}{claim; \GSpage{P000818}}
\GSfield{The notes}{The notes say that \(I_2\) contains three spins, and similarly for higher coefficients.}
\GSfield{Statement}{Local higher charges are obtained from derivatives of \(\log(t(u)t(0)^{-1})\), with conventions for their index and factorial. The second such derivative is \(t(0)^{-1}t''(0)-[t(0)^{-1}t'(0)]^2\). The subtraction cancels disconnected products; a raw Taylor coefficient \(I_2\) generally contains nonlocal products as well.}
\GSfield{Reconstruction}{The logarithmic generating function is the required correction; it is not written in the notebook for the higher coefficients.}
\GSfield{Fate}{This is the local-charge organization of the transfer hierarchy \citep[Sec.~3]{faddeev1996bethe}. The notebook's raw-coefficient locality assertion is not literally correct.}
\GSfield{Present status}{As of 10 November 2026: The local hierarchy is established; raw Taylor coefficients must not be described as single three-site densities.}
\GSfield{Assessment}{literal statement: false as written; publication match: partial; verification: read the result; development: annotated; confidence: high}
\end{GSitem}
\begin{GSitem}{1978-454.14}{F}{Gelfand suggests SL(2,R)}
\GSfield{Kind; source}{remark; \GSpage{P000818}; I. M. Gelfand}
\GSfield{The notes}{\ru{И. М. Г. \quad a,b,c,d --- наверное, \(SL(2,R)\)}.}
\GSfield{Statement}{Four scalar functions do not specify an element, representation or coordinate system of \(SL(2,\mathbb R)\) without a determinant relation, real form and multiplication law. None is stated.}
\GSfield{Reconstruction}{No group identification is reconstructed from the tentative wording.}
\GSfield{Fate}{The nearby eight-vertex and elliptic-symmetry literature \citep{baxter1971eight,belavin1981dynamical} supplies context but not a proposition equivalent to this tentative remark.}
\GSfield{Present status}{As of 10 November 2026: An unformulated analogy, not a surviving precise research problem.}
\GSfield{Assessment}{literal statement: not determinable; publication match: thematic only; verification: unresolved lead; development: annotated; confidence: low}
\end{GSitem}
\begin{GSitem}{1978-454.15}{E}{The six-vertex trigonometric degeneration}
\GSfield{Kind; source}{construction; \GSpage{P000818}}
\GSfield{The notes}{Setting \(d=0\) gives ordinary trigonometric rather than elliptic functions.}
\GSfield{Statement}{Within Baxter's family, sending the elliptic modulus to zero gives \(\operatorname{sn}(u,0)=\sin u\) and \(d=0\), hence the six-vertex/XXZ trigonometric family, with a rational limit after further scaling. Merely imposing \(d=0\) does not classify all degenerate solutions.}
\GSfield{Reconstruction}{The modulus limit, domain and rational degeneration delimit the intended claim.}
\GSfield{Fate}{This is a degeneration of the pre-talk eight-vertex solution \citep{baxter1971eight}; the trigonometric higher-spin framework is explicitly developed in \citet{bytsko2003xxz}.}
\GSfield{Present status}{As of 10 November 2026: Established restricted degeneration, without an unrestricted completeness assertion.}
\GSfield{Assessment}{literal statement: correct with missing hypotheses; publication match: partial; verification: read the result; development: annotated; confidence: high}
\end{GSitem}
\begin{GSitem}{1978-454.16}{E}{Arbitrary auxiliary dimension and a nearest-neighbour density}
\GSfield{Kind; source}{theorem; \GSpage{P000818-P000819}}
\GSfield{The notes}{Equations (13) and (17) replace Pauli matrices by matrix units and form \(H=I_0^{-1}I_1\) for general \(M\).}
\GSfield{Statement}{For \(V=\mathbb C^M\) and analytic regular Yang--Baxter \(R\), the periodic Hamiltonian is \(H=\sum_n h_{n,n+1}\), \(h=PR'(0)\), with the contraction order fixed by the chosen convention. Matrix units \(E_{ij}\) span every two-site operator. The spanning statement does not imply that every two-site operator is the derivative of a Yang--Baxter family.}
\GSfield{Reconstruction}{Analyticity, regularity and the distinction between arbitrary density and integrable density are supplied.}
\GSfield{Fate}{The regular transfer-matrix derivative is standard \citep[Sec.~3]{faddeev1996bethe}; its converse has nontrivial obstructions, exemplified by the audited spin-one density.}
\GSfield{Present status}{As of 10 November 2026: Established forward construction. Its unrestricted converse is false.}
\GSfield{Assessment}{literal statement: correct with missing hypotheses; publication match: partial; verification: read the result; development: annotated; confidence: high}
\end{GSitem}
\begin{GSitem}{1978-454.17}{F}{The literal rotationally invariant spin-1 density}
\GSfield{Kind; source}{claim; \GSpage{P000818-P000819}}
\GSfield{The notes}{The displayed density is \(h=X^2-2X\), \(X=\sum_aS^a\otimes S^a\), while the next page fixes \([S^a,S^b]=\varepsilon_{abc}S^c\) in the real standard/adjoint representation of \(SO(3)\). It is attributed to a Zamolodchikov solution.}
\GSfield{Statement}{With exactly that normalization, \(X=K-P\), where \(P\) flips the factors and \(K(e_i\otimes e_j)=\delta_{ij}\sum_a e_a\otimes e_a\). For \(h_t=X^2+tX\), the matrix entry \(\langle001|[h_{12}+h_{23},[h_{12},h_{23}]]|212\rangle=t(t-1)(t+1)\). It is \(-6\) at the printed \(t=-2\), whereas every \(Q_{23}-Q_{12}\) has zero in that entry. Therefore this density, even up to nonzero scale and added identity, cannot arise from an analytic regular difference-form braided R-matrix.}
\GSfield{Reconstruction}{The analytic regular difference-form scope is explicitly supplied. The exact obstruction is this chapter's critical algebraic audit, not a claim about every possible notion of integrability. In Hermitian generators \(J^a=iS^a\), the familiar isotropic density is proportional to \(J\cdot J-(J\cdot J)^2\), not to the notebook expression.}
\GSfield{Fate}{The necessary condition is reviewed by \citet[Sec.~2, eqs.~(1)--(6)]{zhang2025bootstrap}, first posted April 2025, checked version February 2026. The chapter derives it and computes the obstruction independently twice. The intended published density is not identifiable from the printed coefficient; the known spin-one model has a different bilinear/biquadratic ratio \citep[Sec.~5]{bytsko2003xxz}.}
\GSfield{Present status}{As of 10 November 2026: The literal regular difference-form claim is excluded by an exact necessary-condition calculation. The historical intended coefficient remains uncertain.}
\GSfield{Assessment}{literal statement: false as written; publication match: partial; verification: read the result; development: developed in this chapter; confidence: high; problem \hyperref[prob:GP-1978-10]{GP-1978-10}}
\end{GSitem}
\begin{GSitem}{1978-454.18}{E}{Uniqueness of an invariant two-site spin-one-half interaction}
\GSfield{Kind; source}{claim; \GSpage{P000818}}
\GSfield{The notes}{For \(M=2\), \(H=\sum_n\vec\sigma_n\cdot\vec\sigma_{n+1}\); the notes say there is nothing else invariant for spin one-half.}
\GSfield{Statement}{On \(\mathbb C^2\otimes\mathbb C^2\), the diagonal \(SU(2)\) representation splits as one spin-zero and one spin-one summand. Its commutant is two-dimensional, spanned by \(I\) and \(P\), equivalently \(I\) and \(\sum_a\sigma^a\otimes\sigma^a\). Thus a rotationally invariant two-site density has one nontrivial coupling after adding constants and scaling. This does not classify interactions involving more sites.}
\GSfield{Reconstruction}{Two-site support and the quotient by additive constants are the omitted scope restrictions.}
\GSfield{Fate}{This elementary representation-theoretic fact underlies the Heisenberg example \citep[Sec.~3]{faddeev1996bethe}; the signs of ferromagnetic versus antiferromagnetic couplings require an energy convention.}
\GSfield{Present status}{As of 10 November 2026: Exact elementary classification in the stated two-site setting.}
\GSfield{Assessment}{literal statement: correct with missing hypotheses; publication match: partial; verification: read the result; development: annotated; confidence: high}
\end{GSitem}
\begin{GSitem}{1978-454.19}{B}{Adding anisotropic parameters to the spin-1 construction}
\GSfield{Kind; source}{question; \GSpage{P000819}}
\GSfield{The notes}{\ru{Неизвестно, как к гам. (14) добавить параметры, чтобы уложиться в сх. Замолодчикова.}}
\GSfield{Statement}{The literal question asks for a family of regular factorizing matrices whose local Hamiltonian reduces to (14). Under the analytic difference-form interpretation used in the transfer construction, no such family can pass through the printed density, by item 1978-454.17. A separate corrected question asks for anisotropic integrable deformations of the isotropic spin-one density \(J\cdot J-(J\cdot J)^2\).}
\GSfield{Reconstruction}{The regular analytic difference-form scope comes from the preceding construction. The correctly normalized isotropic reference Hamiltonian is an explicitly separate reconstruction, not an emendation asserted to be the speaker's intention.}
\GSfield{Fate}{The nearby corrected problem was answered by \citet[pp.~298--303]{zamolodchikovfateev1980spin}, published in 1980; an earlier circulation date was not verified. \citet[Sec.~5, eqs.~(52)--(55)]{bytsko2003xxz}, preprint December 2001, explicitly identifies the resulting spin-one XXZ Hamiltonian and its isotropic limit. That is a partial match because the literal coefficient in (14) differs.}
\GSfield{Present status}{As of 10 November 2026: The literal regular-family extension is excluded; the corrected spin-one anisotropic construction is established. No unresolved remainder is manufactured from the coefficient error.}
\GSfield{Assessment}{literal statement: false as written; publication match: partial; verification: read the result; development: developed in this chapter; confidence: high; problem \hyperref[prob:GP-1978-10]{GP-1978-10}}
\end{GSitem}
\begin{GSitem}{1978-454.20}{F}{The proposed universality of the factorization scheme}
\GSfield{Kind; source}{claim; \GSpage{P000819}}
\GSfield{The notes}{The scheme is said to contain all completely integrable two-dimensional quantum systems, and classical systems after a semiclassical limit.}
\GSfield{Statement}{No definition of complete integrability, class of field theories, scattering assumptions, limiting parameter or convergence of a classical limit is given. Consequently an all-systems classification theorem cannot be extracted. A regular finite-dimensional difference R-matrix construction is a precise sufficient framework, not this unspecified universal quantifier.}
\GSfield{Reconstruction}{The statement records exactly which quantifiers and definitions are absent; no universal theorem is substituted.}
\GSfield{Fate}{The checked original factorized-scattering paper \citep[Secs.~1--3]{zamolodchikovzamolodchikov1978on} treats particular massive models and identification issues, not the asserted universal classification.}
\GSfield{Present status}{As of 10 November 2026: Too broad to identify as a definite theorem or open problem.}
\GSfield{Assessment}{literal statement: not determinable; publication match: none found; verification: read the result; development: annotated; confidence: high}
\end{GSitem}
\begin{GSitem}{1978-454.21}{F}{Gelfand asks for a Maurer-Cartan interpretation}
\GSfield{Kind; source}{question; \GSpage{P000819}; I. M. Gelfand}
\GSfield{The notes}{\ru{Вопрос И.М.Г.: соотн. на S --- типа Маурера--Картана}.}
\GSfield{Statement}{A precise Maurer--Cartan interpretation requires a specified differential graded or higher algebra, a degree-one object representing S and an equation equivalent to the given Yang--Baxter relation. The notebook supplies none of these.}
\GSfield{Reconstruction}{No differential, graded bracket or deformation complex is reconstructed.}
\GSfield{Fate}{Later algebraic interpretations of Yang--Baxter equations are a broad thematic continuation. The checked Belavin constructions \citep{belavin1980discrete,belavin1981dynamical} do not identify an exact theorem with this unformulated question.}
\GSfield{Present status}{As of 10 November 2026: A suggestive but unquantified analogy; not eligible for an open-problem or development claim.}
\GSfield{Assessment}{literal statement: not determinable; publication match: thematic only; verification: unresolved lead; development: annotated; confidence: medium}
\end{GSitem}
\begin{GSitem}{1978-454.22}{F}{Second-order extension does not make every local Hamiltonian integrable}
\GSfield{Kind; source}{claim; \GSpage{P000819}}
\GSfield{The notes}{The most general nearest-neighbour Hamiltonian is said always to have a nontrivial integral because \(S''(0)\) can be reconstructed from \(S'(0)\).}
\GSfield{Statement}{Even when a regular Yang--Baxter expansion can be normalized so that \(\check R''(0)=h^2\), its cubic coefficient requires \([h_{12}+h_{23},[h_{12},h_{23}]]=Q_{23}-Q_{12}\). Arbitrary densities do not satisfy this condition. A second-order formal jet alone gives neither a full R-matrix nor a nontrivial independent local conserved charge.}
\GSfield{Reconstruction}{A precise regular expansion and cubic obstruction replace the unsupported inference. The undefined word nontrivial is not interpreted as a new general integrability theorem.}
\GSfield{Fate}{The obstruction is the standard necessary Reshetikhin condition \citep[Sec.~2]{zhang2025bootstrap}. The exact spin-one calculation in this chapter supplies an explicit failure of the notebook's proposed automatic extension.}
\GSfield{Present status}{As of 10 November 2026: False as an unrestricted assertion; a finite Taylor jet is insufficient.}
\GSfield{Assessment}{literal statement: false as written; publication match: partial; verification: read the result; development: developed in this chapter; confidence: high; problem \hyperref[prob:GP-1978-10]{GP-1978-10}}
\end{GSitem}
\begin{GSitem}{1978-454.23}{E}{Matrix associativity is not enough for exchange consistency}
\GSfield{Kind; source}{remark; \GSpage{P000819}}
\GSfield{The notes}{If the \(A_i(\theta)\) are matrices, associativity is automatic.}
\GSfield{Statement}{Matrix multiplication is associative. However, a finite-dimensional representation of exchange relations can have dependent or vanishing products, so associativity of that representation alone does not imply an equality of abstract coefficient tensors such as Yang--Baxter.}
\GSfield{Reconstruction}{The failure of coefficient comparison without independence is supplied.}
\GSfield{Fate}{This clarifies the ordered-state hypothesis in the original construction \citep[Sec.~2]{zamolodchikovzamolodchikov1978on}. It is an elementary logical distinction, not a missing later solution.}
\GSfield{Present status}{As of 10 November 2026: The literal associativity observation is true, but its suggested sufficiency for solving the coefficient equations is not.}
\GSfield{Assessment}{literal statement: correct; publication match: partial; verification: read the result; development: annotated; confidence: high}
\end{GSitem}
\begin{GSitem}{1978-454.24}{E}{Monodromy multiplication and RTT representations}
\GSfield{Kind; source}{construction; \GSpage{P000819-P000820}}
\GSfield{The notes}{An untraced product of local S-matrices gives matrix-valued blocks \(L^{ik}\); equation (20) asserts the same exchange relation for L. Equation (21) calls these blocks a new higher-order algebra.}
\GSfield{Statement}{If each local L-operator satisfies
\[R_{ab}(u-v)L_{an}(u)L_{bn}(v)=L_{bn}(v)L_{an}(u)R_{ab}(u-v),\]
and distinct quantum sites commute, their ordered product satisfies the same RTT relation. Its entries act on \(V^{\otimes N}\). Thus tensor products of local representations generate larger representations; the blocks need not literally obey the one-index Zamolodchikov algebra without an additional identification.}
\GSfield{Reconstruction}{Independent quantum sites and the RTT rather than one-index exchange convention are supplied. The broken internal chain of indices in (18) is repaired by the unambiguous ordered-product definition.}
\GSfield{Fate}{The multiplication identity is an algebraic consequence of the pre-talk Yang--Baxter method. The quantum inverse-scattering programme was announced by \citet[Secs.~1--3]{sklyaninfaddeev1978quantum}, submitted July and published December 1978; that short paper is not claimed to contain every later RTT formulation. A complete multiplication and trace derivation is read in \citet[Sec.~3, eqs.~(44)--(65)]{faddeev1996bethe}. Prior-index 1977-436.16 is a thematic link to the quantum inverse-scattering programme, not the same finite spectral or spin-one deformation question.}
\GSfield{Present status}{As of 10 November 2026: Established representation multiplication; the literal index string contains a faithful source error.}
\GSfield{Assessment}{literal statement: false as written; publication match: partial; verification: read the result; development: annotated; confidence: high}
\end{GSitem}
\begin{GSitem}{1978-454.25}{E}{The zero-field Ising partition sum and trace}
\GSfield{Kind; source}{construction; \GSpage{P000820}}
\GSfield{The notes}{For periodic spins \(\sigma_{n,m}=\pm1\), (22) sums horizontal and vertical nearest-neighbour weights; (23) writes \(Z=\operatorname{Tr}T^N\).}
\GSfield{Statement}{For an \(N\times M\) torus with dimensionless couplings \(K_1,K_2\) and zero magnetic field, let \(T_{s,t}=\exp\{\frac{K_1}{2}\sum_j(s_js_{j+1}+t_jt_{j+1})+K_2\sum_js_jt_j\}\), with periodic j. Then expanding the trace gives the stated finite partition sum exactly: \(Z=\operatorname{Tr}T^N\).}
\GSfield{Reconstruction}{The explicit symmetric row-transfer convention and lattice dimensions are supplied as one realization of the trace identity; the notebook later sketches a diagonal transfer instead.}
\GSfield{Fate}{The transfer method and anisotropic zero-field solution are already in \citet[Secs.~II--VII]{onsager1944crystal}, received October 1943 and published 1944. The notebook's separate unexplained 1965 attribution of commuting transfer matrices is not used as an established priority claim.}
\GSfield{Present status}{As of 10 November 2026: Exact finite combinatorial identity, and an established zero-field solvable model.}
\GSfield{Assessment}{literal statement: correct; publication match: exact; verification: read the result; development: annotated; confidence: high}
\end{GSitem}
\begin{GSitem}{1978-454.26}{F}{Diagonal transfer weights and an unfinished commutation assertion}
\GSfield{Kind; source}{method; \GSpage{P000820}}
\GSfield{The notes}{The diagonal drawing sums spins on intermediate diagonals, writes \(T_{\vec\alpha\vec\beta}=S_1\cdots S_n\), and ends \ru{\(T(J_1,J_2)\) коммут. пр}, apparently the beginning of a conditional assertion.}
\GSfield{Statement}{Diagonal transfer matrices can be assembled by contracting local Boltzmann weights, with torus dimensions and boundary identifications specified. Commutation of two transfer matrices is an additional relation between their couplings or spectral parameters. The page does not state that relation, so no literal two-parameter commutation theorem is recoverable.}
\GSfield{Reconstruction}{The local contraction is identifiable; the missing commutation condition is deliberately not supplied as though present in the notes.}
\GSfield{Fate}{Onsager's transfer-operator analysis \citep{onsager1944crystal} and Baxter's subsequent commuting families \citep{baxter1971eight} provide context. They do not identify the missing words or determine which transfer convention the final assertion used.}
\GSfield{Present status}{As of 10 November 2026: Incomplete source statement. No open question is inferred from a line that simply stops.}
\GSfield{Assessment}{literal statement: not determinable; publication match: thematic only; verification: unresolved lead; development: annotated; confidence: high}
\end{GSitem}

\GStalk{455}{1979-02-26}{S. I. Gelfand, reporting a lecture by Yu. I. Manin}{P000821-P000822; catalogue P000821-P000824}{Complex Minkowski space and the full Yang-Mills ambitwistor correspondence}
S.~I. Gelfand reports Manin's explanation of the 1978 Witten and Isenberg--Yasskin--Green construction for general Yang--Mills fields, without a self-duality assumption. The geometry passes from the Grassmannian and real null twistors to bundles on the incidence quadric and their third formal neighbourhood. The notebook's conformal stabilizer, group equality and Yang--Mills commutator sign require correction. Beilinson's separate intervention begins on P000823 and supplies the missing triviality condition on the quadrics of null directions. The broad catalogue range must not absorb Fuchs's separately dated lecture beginning on P000824.

\begin{GSitem}{1978-455.1}{E}{The Grassmannian compactification of complex Minkowski space}
\GSfield{Kind; source}{definition; \GSpage{P000821-P000822}}
\GSfield{The notes}{\(\widehat{\mathbb CM}=G(4,2)\), the two-planes in \(\mathbb C^4\), contains an open subset isomorphic to \(\mathbb C^4\).}
\GSfield{Statement}{Put \(X=\operatorname{Gr}(2,\mathbb C^4)\). For a splitting \(\mathbb C^4=U\oplus W\), planes transverse to W are graphs of \(Z\in\operatorname{Hom}(U,W)\cong M_2(\mathbb C)\). Their tangent space carries the conformal quadratic form \(q(dZ)=\det(dZ)\), well-defined up to scale under changes of chart. Thus X is compactified complex conformal four-space, with affine Minkowski charts.}
\GSfield{Reconstruction}{The graph chart and determinant conformal structure spell out the compactification; they are not fully written on the page.}
\GSfield{Fate}{This is classical twistor geometry before 1979, reviewed by \citet[Secs.~1--2]{penrosemaccallum1973twistor}; the graph coordinates and determinant metric are explicit in \citet[Sec.~3.1]{woit2021euclidean}. Only that paper's classical geometric review is used here.}
\GSfield{Present status}{As of 10 November 2026: Established finite-dimensional geometry; compactification and affine complexification must be distinguished.}
\GSfield{Assessment}{literal statement: correct; publication match: exact; verification: read the result; development: annotated; confidence: high}
\end{GSitem}
\begin{GSitem}{1978-455.2}{E}{The conformal stabilizer and an impossible quotient}
\GSfield{Kind; source}{claim; \GSpage{P000821-P000822}}
\GSfield{The notes}{The notes write \(\widehat M=SU(2,2)/SO(3,1)\); the next page tentatively suggests \(GL(2,\mathbb C)\times\mathbb R^3\) as stabilizer.}
\GSfield{Statement}{The stabilizer P of a totally null complex two-plane in \(SU(2,2)\) is a real parabolic subgroup of dimension 11. Its Lie algebra is \((\mathfrak{so}(3,1)\oplus\mathbb R)\ltimes\mathbb R^{3,1}\). The compactified real conformal space is \(SU(2,2)/P\), of dimension four. The literal quotient by \(SO(3,1)\) has dimension \(15-6=9\), and the suggested direct product has the wrong group structure.}
\GSfield{Reconstruction}{The parabolic subgroup and its Lie algebra are editorial corrections. Finite central quotients and disconnected components depend on the chosen global conformal action.}
\GSfield{Fate}{The fifteen infinitesimal generators, including the dilation and four special conformal generators missing from the notebook's stabilizer, are explicit in \citet[Sec.~1.1, pp.~245--247]{penrosemaccallum1973twistor}. The dimension and stabilizer corrections follow directly from their action.}
\GSfield{Present status}{As of 10 November 2026: The printed homogeneous quotient is false; the corrected parabolic model is standard.}
\GSfield{Assessment}{literal statement: false as written; publication match: partial; verification: read the result; development: annotated; confidence: high}
\end{GSitem}
\begin{GSitem}{1978-455.3}{E}{SU(2,2) and the conformal spin cover}
\GSfield{Kind; source}{claim; \GSpage{P000821}}
\GSfield{The notes}{Equation (2) writes \(SU(2,2)=SO(4,2)\).}
\GSfield{Statement}{Their Lie algebras are isomorphic, but the groups are not equal: \(SU(2,2)\cong\operatorname{Spin}_0(4,2)\) maps two-to-one onto \(SO_0(4,2)\). The effective projective conformal action has additional central identifications. A literal group equality suppresses this covering information.}
\GSfield{Reconstruction}{Connected components, the spin cover and the effective projective quotient replace the informal equality.}
\GSfield{Fate}{The consecutive two-to-one covering relations are stated in \citet[Sec.~1.1, p.~247]{penrosemaccallum1973twistor}; \citet[Sec.~3.3.2]{woit2021euclidean} uses the spin-group identification.}
\GSfield{Present status}{As of 10 November 2026: Established covering-group relation; equality of Lie algebras must not be confused with equality of groups.}
\GSfield{Assessment}{literal statement: false as written; publication match: exact; verification: read the result; development: annotated; confidence: high}
\end{GSitem}
\begin{GSitem}{1978-455.4}{E}{The projective null-twistor hypersurface}
\GSfield{Kind; source}{construction; \GSpage{P000821}}
\GSfield{The notes}{\(R=\{|z_1|^2+|z_2|^2-|z_3|^2-|z_4|^2=0\}\subset\mathbb CP^3\) is real five-dimensional and preserved by \(SU(2,2)\).}
\GSfield{Statement}{For a nondegenerate Hermitian form h of signature (2,2) on \(\mathbb C^4\), its nonzero null vectors modulo \(\mathbb C^*\) form a smooth real hypersurface \(R\subset\mathbb CP^3\), of real dimension five. Its projective lines are projectivizations of totally h-null two-planes.}
\GSfield{Reconstruction}{Nonzero vectors, projectivization and total nullity clarify the displayed homogeneous equation.}
\GSfield{Fate}{This is the standard real twistor geometry \citep[Sec.~3.3.2]{woit2021euclidean}, already part of the pre-talk Penrose framework \citep{penrosemaccallum1973twistor}. Smoothness and the dimension follow from nondegeneracy of h and projectivization.}
\GSfield{Present status}{As of 10 November 2026: An exact geometric construction; the ambient projective space is complex but this hypersurface is real.}
\GSfield{Assessment}{literal statement: correct; publication match: exact; verification: read the result; development: annotated; confidence: high}
\end{GSitem}
\begin{GSitem}{1978-455.5}{E}{Null two-planes and the U(2) compactification}
\GSfield{Kind; source}{theorem; \GSpage{P000821}}
\GSfield{The notes}{The manifold of lines entirely contained in R is identified with Minkowski space; the lines are graphs \((z_3,z_4)=U(z_1,z_2)\), \(U\in U(2)\).}
\GSfield{Statement}{In the positive-negative splitting \(\mathbb C^4=\mathbb C^2\oplus\mathbb C^2\), projection of a totally null two-plane to the first summand is injective, hence an isomorphism. The plane is the graph of U, and nullity is exactly \(U^*U=I\). Thus the full manifold of these lines is \(U(2)\cong(S^1\times S^3)/\mathbb Z_2\), the conformal compactification of real Minkowski space, not the affine space itself.}
\GSfield{Reconstruction}{The compactification qualification and elementary graph argument make precise the page's compressed theorem.}
\GSfield{Fate}{The identification with U(2) is explicit in \citet[Sec.~3.3.2, p.~21]{woit2021euclidean}. The distinction between Minkowski space and its compactification is already emphasized by \citet[Sec.~1.1]{penrosemaccallum1973twistor}.}
\GSfield{Present status}{As of 10 November 2026: Established identification. The uncertain initial abbreviation in Theorem 1 is not used to infer a different space.}
\GSfield{Assessment}{literal statement: correct with missing hypotheses; publication match: exact; verification: read the result; development: annotated; confidence: high}
\end{GSitem}
\begin{GSitem}{1978-455.6}{E}{Incidence of twistor lines gives the light cone}
\GSfield{Kind; source}{construction; \GSpage{P000821}}
\GSfield{The notes}{Lines in R meeting a fixed line \(L_0\) form a three-dimensional light cone; their intersection points give its generators, and points of R represent light rays.}
\GSfield{Statement}{In a Minkowski graph chart, two twistor lines corresponding to matrices Z and \(Z_0\) meet exactly when \(\det(Z-Z_0)=0\). On the Hermitian real slice this is the ordinary null-separation equation. The fixed null-twistor point determines a null geodesic; varying it over \(L_0\cong\mathbb CP^1\) parametrizes the real celestial two-sphere of null directions through the associated spacetime point.}
\GSfield{Reconstruction}{A local Minkowski chart and the Hermitian real structure distinguish a real light cone from a complex incidence locus.}
\GSfield{Fate}{The determinant/incidence calculation is explicit in \citet[Sec.~3.1, eq.~(3.2) and following paragraphs]{woit2021euclidean}; the real-slice interpretation is classical twistor geometry.}
\GSfield{Present status}{As of 10 November 2026: Established local conformal geometry, with global compactified identifications at infinity.}
\GSfield{Assessment}{literal statement: correct with missing hypotheses; publication match: partial; verification: read the result; development: annotated; confidence: high}
\end{GSitem}
\begin{GSitem}{1978-455.7}{E}{Conformal inversion and its domain}
\GSfield{Kind; source}{construction; \GSpage{P000822}}
\GSfield{The notes}{The conformal group is obtained by adding an inversion \(1/\tau\), with \(\tau\) the Minkowski distance, to Poincare motions.}
\GSfield{Statement}{For the quadratic Minkowski form q, the vector inversion is \(x\mapsto x/q(x)\), defined off \(q(x)=0\) in an affine chart. Together with translations and Lorentz transformations it generates the corresponding extended conformal transformations; compositions yield special conformal transformations and dilations. Its global interpretation uses compactified spacetime, and a single inversion can change components or orientation.}
\GSfield{Reconstruction}{The vector formula, null-cone exclusion and global component qualifications replace the scalar shorthand.}
\GSfield{Fate}{\citet[Sec.~1.1, eq.~(1.2)]{penrosemaccallum1973twistor} gives this inversion and explains why compactification is needed when its null cone goes to infinity.}
\GSfield{Present status}{As of 10 November 2026: Established conformal transformation; the notebook's reciprocal notation is schematic.}
\GSfield{Assessment}{literal statement: correct with missing hypotheses; publication match: partial; verification: read the result; development: annotated; confidence: high}
\end{GSitem}
\begin{GSitem}{1978-455.8}{E}{Holomorphic connection and curvature}
\GSfield{Kind; source}{definition; \GSpage{P000822}}
\GSfield{The notes}{On an open subset of \(\mathbb C^4\), \(A_\mu(z)\in\mathfrak g\) are holomorphic and \(F_{\mu\nu}=\partial_\mu A_\nu-\partial_\nu A_\mu+[A_\mu,A_\nu]\).}
\GSfield{Statement}{For a holomorphic principal G-bundle and a local holomorphic trivialization, \(A=\sum_\mu A_\mu dz^\mu\) is a connection form and \(F=dA+A\wedge A\) its curvature. On overlaps the usual gauge transition rule is required. The displayed coefficients describe one trivializing patch, not an automatic global trivialization of the bundle.}
\GSfield{Reconstruction}{The bundle, patch and transition-function qualifications are supplied.}
\GSfield{Fate}{These are the connection data in the 1978 ambitwistor correspondence, summarized precisely by \citet[Sec.~2]{masonskinner2006ambitwistor}.}
\GSfield{Present status}{As of 10 November 2026: Standard holomorphic gauge geometry.}
\GSfield{Assessment}{literal statement: correct; publication match: exact; verification: read the result; development: annotated; confidence: high}
\end{GSitem}
\begin{GSitem}{1978-455.9}{E}{The Yang-Mills commutator sign}
\GSfield{Kind; source}{claim; \GSpage{P000822}}
\GSfield{The notes}{Equation (6), with the curvature convention in (5), writes \(\sum_\mu(\partial_\mu F_{\mu\nu}+[F_{\mu\nu},A_\mu])=0\).}
\GSfield{Statement}{For \(F=dA+A\wedge A\), the covariant derivative is \(D=d+[A,\cdot]\). The source-free equation is \(D^\mu F_{\mu\nu}=\partial^\mu F_{\mu\nu}+[A^\mu,F_{\mu\nu}]=0\), with indices raised using the chosen flat complex conformal representative. The printed plus sign and reversed commutator are incompatible with this curvature convention.}
\GSfield{Reconstruction}{The commutator order and metric contraction are corrected explicitly; no alternative sign convention is silently imposed.}
\GSfield{Fate}{The Yang--Mills equation in the geometric correspondence is the conventional covariant equation \citep[Secs.~1--2]{masonskinner2006ambitwistor}. The scan explicitly confirms the reversed order, so this is a faithful notebook error.}
\GSfield{Present status}{As of 10 November 2026: The corrected equation is standard and gauge covariant.}
\GSfield{Assessment}{literal statement: false as written; publication match: partial; verification: read the result; development: annotated; confidence: high}
\end{GSitem}
\begin{GSitem}{1978-455.10}{F}{The attributed Landau preference for the S-matrix}
\GSfield{Kind; source}{remark; \GSpage{P000822}; I. M. Gelfand (attributing a remark to L. D. Landau)}
\GSfield{The notes}{Gelfand attributes to Landau the phrase \ru{лагранжиан --- труп} and the view that knowing the S-matrix suffices.}
\GSfield{Statement}{This is a recorded methodological comment, without a defined reconstruction theorem for a Lagrangian from an S-matrix or a uniqueness assertion. Its attribution is the notebook's, not independently authenticated here.}
\GSfield{Reconstruction}{No mathematical quantifiers or historical authentication are supplied beyond the notebook.}
\GSfield{Fate}{No exact mathematical publication is identified with this remark. Searches of the surrounding gauge and factorized-scattering literature give thematic context only; the statement is not made into a technical conjecture.}
\GSfield{Present status}{As of 10 November 2026: A historical-methodological remark, not an identifiable open problem.}
\GSfield{Assessment}{literal statement: not determinable; publication match: none found; verification: unresolved lead; development: annotated; confidence: low}
\end{GSitem}
\begin{GSitem}{1978-455.11}{E}{The Yang-Mills action and Euler-Lagrange equation}
\GSfield{Kind; source}{theorem; \GSpage{P000822}}
\GSfield{The notes}{The real equations are said to follow from \(L=\operatorname{tr}\int F_{\mu\nu}F^{\mu\nu}\,d^4x\).}
\GSfield{Statement}{On a real oriented flat four-dimensional domain with a nondegenerate invariant bilinear form on \(\mathfrak g\), compactly supported variations of \(S[A]=\int\langle F_{\mu\nu},F^{\mu\nu}\rangle\,d^4x\) give \(D_\mu F^{\mu\nu}=0\), up to a nonzero overall normalization. An admissible matrix representation supplies the trace pairing; boundary conditions are needed if variations are not compactly supported.}
\GSfield{Reconstruction}{The invariant pairing, nondegeneracy and boundary/variation conditions are supplied.}
\GSfield{Fate}{This is the classical Yang--Mills variational equation, also used in \citet[Sec.~1, eq.~(1) and surrounding discussion]{masonskinner2006ambitwistor}. It yields the corrected commutator order of item 1978-455.9.}
\GSfield{Present status}{As of 10 November 2026: Established variational identity; it does not validate the reversed commutator printed in (6).}
\GSfield{Assessment}{literal statement: correct with missing hypotheses; publication match: partial; verification: read the result; development: annotated; confidence: high}
\end{GSitem}
\begin{GSitem}{1978-455.12}{E}{The five-dimensional ambitwistor incidence quadric}
\GSfield{Kind; source}{construction; \GSpage{P000822}}
\GSfield{The notes}{The notes form \(Q=\{(\ell,\pi):\ell\subset\pi\}\subset\mathbb CP^3\times\mathbb CP^3\), with a dimension glyph transcribed as \(\dim_{\mathbb C}G=5\).}
\GSfield{Statement}{The second factor is the dual projective space: \(Q=F(1,3;\mathbb C^4)=\{([z],[w])\in\mathbb P^3\times\mathbb P^{3*}:w(z)=0\}\). It is a smooth hypersurface of bidegree (1,1), hence complex dimension five. It parametrizes complex null geodesics in compactified complex Minkowski space, with local domains chosen for the connection transform.}
\GSfield{Reconstruction}{The dual projective space and the corrected dimension variable Q are supplied; the dimension glyph in the scan remains palaeographically ambiguous.}
\GSfield{Fate}{The incidence equation and dual space are explicit in \citet[Sec.~2]{masonskinner2006ambitwistor}. Beilinson's next-page diagram itself restores the dual star.}
\GSfield{Present status}{As of 10 November 2026: Established incidence geometry; the arbitrary gauge group G is not asserted to have dimension five.}
\GSfield{Assessment}{literal statement: correct with missing hypotheses; publication match: exact; verification: read the result; development: annotated; confidence: high}
\end{GSitem}
\begin{GSitem}{1978-455.13}{E}{The third formal neighbourhood criterion for full Yang-Mills}
\GSfield{Kind; source}{theorem; \GSpage{P000821-P000822}}
\GSfield{The notes}{The talk credits Witten and independently Yasskin, Isenberg and Green (1978). It identifies solutions with holomorphic G-bundles on Q extending to its third infinitesimal neighbourhood.}
\GSfield{Statement}{Let U be a sufficiently small Stein domain in complex Minkowski space, with every complex null geodesic intersection connected and simply connected. Let \(Q_U\) be the space of null geodesics meeting U. Under the local ambitwistor correspondence, a holomorphic connection solves the full Yang--Mills equations precisely when its associated bundle E, holomorphically trivial on every quadric \(Q_x\cong\mathbb P^1\times\mathbb P^1\), extends to the third formal neighbourhood of \(Q_U\) in \(\mathbb P^3\times\mathbb P^{3*}\). With ideal sheaf I, this means the ringed space with structure sheaf \(\mathcal O/I^4\).}
\GSfield{Reconstruction}{The small-domain hypotheses, quadric triviality, dual projective factor and \(I^4\) convention are explicit additions. Triviality is supported by Beilinson's continuation on P000823; it is absent from this page's first statement. Vector bundles or faithful linear realizations of principal bundles are understood with the relevant structure preserved.}
\GSfield{Fate}{Published before the talk: \citet{witten1978interpretation} in August 1978 and \citet{isenbergyasskingreen1978nonselfdual} in October 1978; an earlier first-circulation date was not established here. The exact third-order criterion and domain qualifications are read in \citet[Sec.~2, Theorem~2 and footnote~1]{masonskinner2006ambitwistor}, first posted October 2005, published 2006.}
\GSfield{Present status}{As of 10 November 2026: An established local holomorphic correspondence for full Yang--Mills fields, not a self-duality-only construction or a new 1979 result.}
\GSfield{Assessment}{literal statement: correct with missing hypotheses; publication match: exact; verification: read the result; development: annotated; confidence: high}
\end{GSitem}
\begin{GSitem}{1978-455.14}{E}{Arbitrary holomorphic connections and bundles before the field equation}
\GSfield{Kind; source}{claim; \GSpage{P000822}}
\GSfield{The notes}{The last sentence says all connections correspond to all bundles, without necessarily extending to the third neighbourhood.}
\GSfield{Statement}{On the same admissible local domains, holomorphic connections up to gauge correspond to holomorphic bundles on \(Q_U\) trivial on every \(Q_x\), up to isomorphism preserving the relevant structure. No Yang--Mills equation is imposed at this stage. Arbitrary bundles without the quadric-triviality condition do not satisfy this correspondence.}
\GSfield{Reconstruction}{The local-domain and fibrewise holomorphic-triviality conditions are supplied. Mere vanishing of a Chern class is not substituted for explicit holomorphic triviality. Editorial check: the checked Mason--Skinner review uses a Chern-class shorthand, but \(\mathcal O(1,0)\oplus\mathcal O(-1,0)\) on \(\mathbb P^1\times\mathbb P^1\) has \(c_1=c_2=0\) and restricts to \(\mathcal O(1)\oplus\mathcal O(-1)\) on a ruling, so it is not holomorphically trivial. The precise fibrewise holomorphic-triviality condition is therefore retained, exactly as P000823 requires; the review's broader shorthand is not used.}
\GSfield{Fate}{This is the pre-equation transform of \citet[Sec.~2, Theorem~1]{masonskinner2006ambitwistor}, attributed there to the 1978 papers. Beilinson's intervention supplies the missing condition and construction.}
\GSfield{Present status}{As of 10 November 2026: Established conditional correspondence. The unrestricted phrase all bundles is overbroad.}
\GSfield{Assessment}{literal statement: correct with missing hypotheses; publication match: partial; verification: read the result; development: annotated; confidence: high}
\end{GSitem}

\GStalk{456}{1979-02-26}{A. A. Beilinson}{P000823 only; catalogue P000822-P000825}{The complete-flag explanation of the ambitwistor transform}
Beilinson explains the preceding construction by the complete-flag double fibration: pull a connection from the Grassmannian to the correspondence space, take horizontal sections along null curves, and descend to the incidence quadric. The final line explicitly requires triviality on \(\mathbb P^1\times\mathbb P^1\). This is an explanation of the 1978 full-connection ambitwistor correspondence, not the unrelated derived-category theorem bearing Beilinson's name. The supplied prior index places this same page under 1977-457.1--3; these are source aliases, not an earlier seminar occurrence. P000824 begins the separately dated Fuchs lecture. The page itself has no new date: 26 February 1979 is the catalogue date and is consistent with its position immediately after the dated S.~I. Gelfand report.

\begin{GSitem}{1978-456.1}{E}{The complete-flag double fibration}
\GSfield{Kind; source}{construction; \GSpage{P000823}}
\GSfield{The notes}{Let F parametrize line, plane and hyperplane flags in \(\mathbb C^4\); one projection forgets the outer terms and the other forgets the middle plane.}
\GSfield{Statement}{For \(F=F(1,2,3;\mathbb C^4)\), let \(p:F\to X=\operatorname{Gr}(2,4)\) and \(q:F\to Q=F(1,3;4)\). Then \(p^{-1}(x)=\mathbb P(x)\times\mathbb P((\mathbb C^4/x)^*)\cong\mathbb P^1\times\mathbb P^1\), whereas \(q^{-1}(\ell,\pi)=\mathbb P(\pi/\ell)\cong\mathbb P^1\). The p-image of a q-fibre is a complex null geodesic.}
\GSfield{Reconstruction}{The two fibre calculations and null-geodesic identification expand the diagram. The local version restricts p to the chosen Minkowski domain.}
\GSfield{Fate}{This is the incidence geometry underlying the pre-talk 1978 transform, reviewed in \citet[Sec.~2]{masonskinner2006ambitwistor}. Prior-index 1977-457.1 refers to this same source page and must be an alias, not a prior chronological link.}
\GSfield{Present status}{As of 10 November 2026: Established flag geometry, with the two different fibre types distinguished.}
\GSfield{Assessment}{literal statement: correct; publication match: exact; verification: read the result; development: annotated; confidence: high}
\end{GSitem}
\begin{GSitem}{1978-456.2}{E}{Flatness and horizontal sections along a null curve}
\GSfield{Kind; source}{claim; \GSpage{P000823}}
\GSfield{The notes}{A holomorphic bundle with connection becomes canonically trivial on \(\mathbb CP^1\): curvature vanishes by dimension and the fundamental group is zero.}
\GSfield{Statement}{A holomorphic connection restricted to a complex curve is flat, because a holomorphic two-form restricts to zero. On a connected simply connected curve C, evaluation identifies its horizontal sections with every fibre, so \(E|_C\cong\mathcal O_C\otimes H^0_\nabla(C,E)\) canonically. No preferred basis of that vector space is selected. In the local transform C is the null-geodesic intersection with U, not necessarily the full projective line.}
\GSfield{Reconstruction}{Holomorphic rather than arbitrary real connection, connected simply connected domain and the basis-free sense of canonical triviality are supplied.}
\GSfield{Fate}{The horizontal-section construction and simply connected intersections are explicit in \citet[Sec.~2, Theorem~1 and footnote~1]{masonskinner2006ambitwistor}. This is part of the same paragraph aliased by prior-index 1977-457.2.}
\GSfield{Present status}{As of 10 November 2026: Correct in the stated holomorphic local setting. It holds for arbitrary such connections and does not impose the Yang--Mills equation.}
\GSfield{Assessment}{literal statement: correct with missing hypotheses; publication match: partial; verification: read the result; development: annotated; confidence: high}
\end{GSitem}
\begin{GSitem}{1978-456.3}{E}{Pullback and descent to the space of null geodesics}
\GSfield{Kind; source}{construction; \GSpage{P000823}}
\GSfield{The notes}{Lift \((E,\nabla)\) from \(G(4,2)\) to F and descend it to a holomorphic bundle \(\widetilde E\) on Q because it is trivial on every right-hand fibre.}
\GSfield{Statement}{On an admissible domain U, define \(\widetilde E_l\) to be the vector space of \(\nabla\)-horizontal sections along the null-geodesic intersection \(l\cap U\). Local holomorphic dependence and the absence of monodromy give a holomorphic bundle on \(Q_U\). Equivalently this is descent of \(p^*E\) with its relative flat connection along q.}
\GSfield{Reconstruction}{The admissible local domain, relative connection and holomorphic family qualification are supplied; fibrewise topological triviality by itself is not used as a descent theorem.}
\GSfield{Fate}{This is the forward transform in \citet[Sec.~2, preceding Theorem~1]{masonskinner2006ambitwistor}, attributed to Witten and Isenberg--Yasskin--Green in 1978. Prior-index 1977-457.2 is the same page's descent argument.}
\GSfield{Present status}{As of 10 November 2026: Established local transform for arbitrary holomorphic connections; no field equation is needed for this step.}
\GSfield{Assessment}{literal statement: correct with missing hypotheses; publication match: partial; verification: read the result; development: annotated; confidence: high}
\end{GSitem}
\begin{GSitem}{1978-456.4}{E}{Reverse transform and the quadric-triviality condition}
\GSfield{Kind; source}{theorem; \GSpage{P000823}}
\GSfield{The notes}{The final line gives \((E,\nabla)\longleftrightarrow\widetilde E\), with \(\widetilde E\) trivial on \(\mathbb CP^1\times\mathbb CP^1\).}
\GSfield{Statement}{For each \(x\in U\), let \(Q_x\) be its quadric of null directions. A holomorphic bundle \(\widetilde E\) trivial on every \(Q_x\), with the same local-domain conditions, reconstructs \(E_x=H^0(Q_x,\widetilde E|_{Q_x})\); incidence of null curves reconstructs the connection. To obtain Yang--Mills, one separately imposes extension to \(\mathcal O/I^4\), as in item 1978-455.13.}
\GSfield{Reconstruction}{The meaning of each quadric, reconstruction fibre and separation of the connection transform from the third-neighbourhood criterion are supplied.}
\GSfield{Fate}{The reverse connection construction and the separate field-equation criterion are \citet[Sec.~2, Theorems~1--2]{masonskinner2006ambitwistor}. Prior-index 1977-457.3 is a duplicate of this final source line, not evidence of an earlier occurrence.}
\GSfield{Present status}{As of 10 November 2026: An established correspondence. Flatness along one-dimensional null curves alone cannot prove the full Yang--Mills equations.}
\GSfield{Assessment}{literal statement: correct with missing hypotheses; publication match: exact; verification: read the result; development: annotated; confidence: high}
\end{GSitem}

\GStalk{457}{1979-03-05}{D. B. Fuchs}{P000824-P000825 and P000829; P000826-P000828 not supplied}{Graded Lie algebras: a partially supplied lecture omitted from the record list}
The dated heading \ru{5/III-79, Д. Б. Фукс} on P000824 establishes a separate lecture between Beilinson and Ogievetsky. It is omitted from the supplied fourteen-record list and is added here explicitly to avoid suppressing mathematics present in the attached corpus. P000824--P000825 cover transitivity, Cartan vector-field families and prolongation; P000829 resumes with twisted loops and a period-three bracket. P000826--P000828 are not supplied, so this is not a complete account of the lecture. The same available segments were previously placed under 1977-457.4--7 and 1977-457.17--19, although their visible date is March 1979. The prior index's 1977-457.1--3 instead describe Beilinson's P000823, and 1977-457.20--25 describe Ogievetsky's P000830--P000831. These are provenance aliases, not historical recurrences.

\begin{GSitem}{1978-457.1}{E}{Grading and the two transitivity conditions}
\GSfield{Kind; source}{definition; \GSpage{P000824}}
\GSfield{The notes}{A \(\mathbb Z\)-graded algebra has \([\mathfrak g_i,\mathfrak g_j]\subset\mathfrak g_{i+j}\). Bracketing with \(\mathfrak g_{-1}\), respectively \(\mathfrak g_1\), is injective in nonnegative, respectively nonpositive, degree.}
\GSfield{Statement}{For a graded Lie algebra over \(\mathbb C\), each homogeneous piece is a \(\mathfrak g_0\)-module. The displayed maps are \(x\mapsto(v\mapsto[x,v])\); their injectivity means that no nonzero element in the indicated degree centralizes the specified neighbouring component. Classification results require additional assumptions such as finite-dimensional graded pieces and generation by the local part; these are not consequences of the displayed bracket inclusion.}
\GSfield{Reconstruction}{The adjoint maps and missing classification hypotheses are made explicit; the notebook's non-strict degree ranges include degree zero.}
\GSfield{Fate}{\citet[Chapter~I, Sec.~1, Definitions~1--3]{kac1968graded} states these conventions together with finite-dimensionality and local generation. The paper was received 13 February 1968 and published in 1968. Prior alias: 1977-457.4, from this same page.}
\GSfield{Present status}{As of 10 November 2026: Established graded-algebra definitions; transitivity alone is not the asserted classification theorem.}
\GSfield{Assessment}{literal statement: correct; publication match: exact; verification: read the result; development: annotated; confidence: high}
\end{GSitem}
\begin{GSitem}{1978-457.2}{E}{Polynomial vector fields and the degree-zero matrix algebra}
\GSfield{Kind; source}{definition; \GSpage{P000824}}
\GSfield{The notes}{\(W_n=\{\sum_i p_i(x)\partial_{x_i}\}\), with \(\deg x_i=1\), \(\deg\partial_{x_i}=-1\); degree zero is the matrix algebra.}
\GSfield{Statement}{For \(n\ge1\), \(W_n=\operatorname{Der}_{\mathbb C}\mathbb C[x_1,\ldots,x_n]\). Its degree-k part consists of vector fields with coefficient polynomials homogeneous of degree \(k+1\), so \((W_n)_{-1}\cong\mathbb C^n\) and \((W_n)_0\cong\mathfrak{gl}_n\), with the usual choice of sign in identifying linear fields with matrices.}
\GSfield{Reconstruction}{The field, derivation algebra and homogeneous-degree convention are expanded.}
\GSfield{Fate}{The polynomial family and grading are explicitly described in \citet[Chapter~III, Sec.~1, (40), type~(1)]{kac1968graded}. Prior alias: the W-family portion of 1977-457.5.}
\GSfield{Present status}{As of 10 November 2026: Classical Cartan-type example, already known before the lecture.}
\GSfield{Assessment}{literal statement: correct; publication match: exact; verification: read the result; development: annotated; confidence: high}
\end{GSitem}
\begin{GSitem}{1978-457.3}{E}{Divergence-free polynomial vector fields}
\GSfield{Kind; source}{definition; \GSpage{P000824}}
\GSfield{The notes}{The second series is the divergence-free vector fields \(S_n\); a possible prime after S is uncertain.}
\GSfield{Statement}{For \(n\ge2\), define \(S_n=\{\sum_i p_i\partial_{x_i}:\sum_i\partial_i p_i=0\}\). Equivalently these vector fields preserve the standard volume form. The inherited degree-zero component is \(\mathfrak{sl}_n\). The one-variable divergence-free space consists only of translations and is not an infinite simple Cartan family.}
\GSfield{Reconstruction}{The divergence formula, volume form and exclusion of the degenerate n=1 case are supplied. The uncertain prime is not resolved as a distinct derived algebra without evidence.}
\GSfield{Fate}{This is \citet[Chapter~III, Sec.~1, type~(2)]{kac1968graded}. Prior alias: the S-family portion of 1977-457.5.}
\GSfield{Present status}{As of 10 November 2026: Established family with the dimension restriction; the possible prime remains a transcription ambiguity.}
\GSfield{Assessment}{literal statement: correct with missing hypotheses; publication match: partial; verification: read the result; development: annotated; confidence: high}
\end{GSitem}
\begin{GSitem}{1978-457.4}{E}{Hamiltonian polynomial vector fields}
\GSfield{Kind; source}{definition; \GSpage{P000824}}
\GSfield{The notes}{The third family is simply named \(H_n\), Hamiltonian vector fields.}
\GSfield{Statement}{On \(\mathbb C^{2n}\), \(n\ge1\), fix \(\omega=\sum_i dq_i\wedge dp_i\). The Hamiltonian vector fields are \(X_f=\sum_i(\partial_{p_i}f\,\partial_{q_i}-\partial_{q_i}f\,\partial_{p_i})\) for polynomial f; constants give the zero field. Equivalently they preserve \(\omega\). A homogeneous polynomial of degree \(k+2\) gives a degree-k field.}
\GSfield{Reconstruction}{The even-dimensional space, symplectic form, quotient by constant Hamiltonians and grading are supplied from the identified standard family.}
\GSfield{Fate}{This is \citet[Chapter~III, Sec.~1, type~(4)]{kac1968graded}. Prior alias: the H-family portion of 1977-457.5.}
\GSfield{Present status}{As of 10 November 2026: Established polynomial Hamiltonian Lie algebra; the notebook itself does not give this defining formula.}
\GSfield{Assessment}{literal statement: correct with missing hypotheses; publication match: partial; verification: read the result; development: annotated; confidence: high}
\end{GSitem}
\begin{GSitem}{1978-457.5}{E}{Contact polynomial vector fields}
\GSfield{Kind; source}{definition; \GSpage{P000824}}
\GSfield{The notes}{The fourth family \(K_n\) preserves a contact structure on \(\mathbb C^{2n+1}\); its grading is explicitly omitted.}
\GSfield{Statement}{For \(n\ge1\), with \(\alpha=dz-\sum_i p_i\,dq_i\), the polynomial contact algebra is \(K_n=\{X:\mathcal L_X\alpha=f\alpha\text{ for some polynomial }f\}\). It preserves the contact distribution, not necessarily the one-form itself. Giving q and p weight one and z weight two yields the usual depth-two contact grading.}
\GSfield{Reconstruction}{The form, conformal preservation condition and weighted grading are supplied. The later depth-one simplification cannot include the full contact grading unchanged.}
\GSfield{Fate}{The contact family and its special filtration appear in \citet[Introduction, p.~1272; Chapter~III, Sec.~1, type~(6), Proposition~18]{kac1968graded}. Prior alias: the K-family portion of 1977-457.5.}
\GSfield{Present status}{As of 10 November 2026: Established Cartan-type family; omitted grading matters to the scope of the prolongation argument.}
\GSfield{Assessment}{literal statement: correct with missing hypotheses; publication match: partial; verification: read the result; development: annotated; confidence: high}
\end{GSitem}
\begin{GSitem}{1978-457.6}{E}{The four-series classification needs more hypotheses}
\GSfield{Kind; source}{claim; \GSpage{P000824}}
\GSfield{The notes}{After the lower grading bound and transitivity conditions, the notes say Cartan's classification gives four series.}
\GSfield{Statement}{An identifiable corrected statement is the classical classification of infinite-dimensional simple Cartan-type algebras under the relevant irreducibility, finite-piece and generation/primitive assumptions, giving vector, divergence-free, Hamiltonian and contact types. The displayed conditions alone neither exclude finite-dimensional examples nor enforce simplicity or irreducibility. Even the general transitive classification includes volume- and symplectic-conformal extensions absent from the four-series shorthand.}
\GSfield{Reconstruction}{The additional assumptions and the distinction between simple and merely transitive classification are explicitly reconstructed from the cited theorem. The exact version intended by Fuchs is not uniquely specified.}
\GSfield{Fate}{\citet[Chapter~III, Theorem~4; Chapter~IV, Theorem~5]{kac1968graded} separates the simple four-family statement from the broader six Cartan types. Prior alias: the classification portion of 1977-457.5. This is known exposition with omitted hypotheses, not an open classification problem.}
\GSfield{Present status}{As of 10 November 2026: The corrected classical classifications are established; the literal sufficiency of the displayed assumptions is false.}
\GSfield{Assessment}{literal statement: false as written; publication match: partial; verification: read the result; development: annotated; confidence: high}
\end{GSitem}
\begin{GSitem}{1978-457.7}{E}{The maximal Cartan prolongation and Jacobi symmetry}
\GSfield{Kind; source}{construction; \GSpage{P000824-P000825}}
\GSfield{The notes}{At depth one, everything is said to be recovered from \(\mathfrak g_{-1},\mathfrak g_0\). Equations (8)--(13) impose symmetry on multilinear maps. The inputs may be arbitrary, simplicity is extra, and infinite dimension is called good luck.}
\GSfield{Statement}{Let V be a finite-dimensional complex vector space and \(\mathfrak h\subset\mathfrak{gl}(V)\) a specified faithful Lie subalgebra. Define \(\mathfrak h^{(k)}=(S^{k+1}V^*\otimes V)\cap(S^kV^*\otimes\mathfrak h)\), with \(\mathfrak h^{(0)}=\mathfrak h\). Then \(V\oplus\mathfrak h\oplus\bigoplus_{k\ge1}\mathfrak h^{(k)}\) is the maximal depth-one transitive prolongation. Any other such graded algebra embeds in it; equality is an additional maximality condition. Finite versus infinite type is a precise property of this pair, whereas good luck or usually finite has no specified genericity meaning.}
\GSfield{Reconstruction}{Finite-dimensionality, a faithful specified action and maximality are supplied. Recoverability of an arbitrary graded algebra from its nonpositive part is corrected to an embedding into the maximal prolongation.}
\GSfield{Fate}{The classical prolongation framework predates the talk; \citet{guilleminquillensternberg1967infinite} studies infinite-type irreducible inputs. The exact tensor-intersection definition is read in \citet[Definition~2.1]{fuhwang2010prolongation}. Prior alias: 1977-457.6. No unseen rank-one criterion from P000826 is attributed to the available pages.}
\GSfield{Present status}{As of 10 November 2026: Established construction. The notebook's unrestricted uniqueness and informal genericity wording require these qualifications.}
\GSfield{Assessment}{literal statement: false as written; publication match: partial; verification: read the result; development: annotated; confidence: high}
\end{GSitem}
\begin{GSitem}{1978-457.8}{E}{The inductive bracket in the prolongation}
\GSfield{Kind; source}{construction; \GSpage{P000825}}
\GSfield{The notes}{To define \([\alpha,\beta]\), specify \([[\alpha,\beta],\delta]\) for \(\delta\in\mathfrak g_{-1}\); Jacobi reduces it to lower degrees.}
\GSfield{Statement}{For homogeneous positive-degree elements X,Y in the maximal prolongation and \(v\in V\), define \([[X,Y],v]=[X,[Y,v]]-[Y,[X,v]]\). Transitivity gives uniqueness. Existence and Jacobi follow by realizing the symmetric tensors as homogeneous polynomial vector fields and using their ordinary commutator; the symmetry conditions are precisely those needed to remain in the prolongation.}
\GSfield{Reconstruction}{The explicit Jacobi identity and the polynomial-vector-field realization justify existence as well as uniqueness. The latter is explanatory reconstruction, not the additional construction on the missing following page.}
\GSfield{Fate}{This is the standard prolongation bracket associated to the tensor definition \citep[Definition~2.1 and discussion of symmetry algebras]{fuhwang2010prolongation}, within the pre-1979 Cartan framework. Prior alias: 1977-457.7.}
\GSfield{Present status}{As of 10 November 2026: Established algebraic construction; no unsolved bracket-definition problem remains.}
\GSfield{Assessment}{literal statement: correct with missing hypotheses; publication match: partial; verification: read the result; development: annotated; confidence: high}
\end{GSitem}
\begin{GSitem}{1978-457.9}{E}{The loop algebra twisted by an involution}
\GSfield{Kind; source}{construction; \GSpage{P000829}}
\GSfield{The notes}{An involution \(\tau\) gives \(\mathfrak g=\mathfrak g_+\oplus\mathfrak g_-\), and the algebra is \(\cdots\oplus\mathfrak g_-t^{-1}\oplus\mathfrak g_+\oplus\mathfrak g_-t\oplus\mathfrak g_+t^2\oplus\cdots\).}
\GSfield{Statement}{For a complex Lie algebra \(\mathfrak g\) and an automorphism \(\tau^2=1\), let \(\mathfrak g_\pm\) be its eigenspaces. Then \([\mathfrak g_\epsilon,\mathfrak g_\delta]\subset\mathfrak g_{\epsilon\delta}\), and \(L(\mathfrak g,\tau)=\bigoplus_{k\in\mathbb Z}\mathfrak g_{(-1)^k}t^k\) is a Lie subalgebra of \(\mathfrak g\otimes\mathbb C[t,t^{-1}]\). This is the centreless twisted loop algebra; adding an affine central extension and a derivation is separate data.}
\GSfield{Reconstruction}{The automorphism rather than vector-space involution hypothesis, eigenspace brackets and distinction from the full affine algebra are supplied.}
\GSfield{Fate}{\citet[Chapter~II, Sec.~6, Definition~9 and following construction, pp.~1292--1293]{kac1968graded} gives the general finite-order construction. Prior alias: the construction in 1977-457.17.}
\GSfield{Present status}{As of 10 November 2026: Established loop construction; no affine central generator or derivation is printed in this source segment.}
\GSfield{Assessment}{literal statement: correct with missing hypotheses; publication match: exact; verification: read the result; development: annotated; confidence: high}
\end{GSitem}
\begin{GSitem}{1978-457.10}{E}{Negative transpose on sl(3)}
\GSfield{Kind; source}{construction; \GSpage{P000829}}
\GSfield{The notes}{The example is \(\mathfrak g=\mathfrak{sl}_3\), with \(\tau\) described as reflection across the diagonal, negative transposition.}
\GSfield{Statement}{The map \(\tau(X)=-X^{\mathsf T}\) is an involutive Lie algebra automorphism of \(\mathfrak{sl}_3(\mathbb C)\). Its fixed algebra is \(\mathfrak{so}_3(\mathbb C)\), dimension three; its negative eigenspace consists of symmetric traceless matrices, dimension five. Alternating these spaces gives the standard twisted \(A_2\) loop construction.}
\GSfield{Reconstruction}{The minus sign is essential: ordinary transpose is an anti-automorphism. Eigenspaces and the centreless qualification are supplied.}
\GSfield{Fate}{This is the order-two outer-automorphism case of \citet[Chapter~II, Sec.~6, Lemmas~21--22]{kac1968graded}. Prior alias: the explicit example within 1977-457.17.}
\GSfield{Present status}{As of 10 November 2026: Established example; the phrase reflection is interpreted using the page's explicit negative-transpose parenthesis.}
\GSfield{Assessment}{literal statement: correct; publication match: exact; verification: read the result; development: annotated; confidence: high}
\end{GSitem}
\begin{GSitem}{1978-457.11}{E}{A period-three one-dimensional graded bracket}
\GSfield{Kind; source}{construction; \GSpage{P000829}}
\GSfield{The notes}{For basis elements \(\varepsilon_i\), \(i\in\mathbb Z\), the bracket is \([\varepsilon_i,\varepsilon_j]=a_{ij}\varepsilon_{i+j}\), where \(a_{ij}\in\{-1,0,1\}\) is congruent to \(j-i\) modulo three.}
\GSfield{Statement}{The literal rule defines a Lie algebra isomorphic to \(\mathfrak{sl}_2\otimes\mathbb C[z,z^{-1}]\). With \([H,E]=2E\), \([H,F]=-2F\), \([E,F]=H\), an explicit isomorphism is \(\varepsilon_{3k}\mapsto\frac12Hz^k\), \(\varepsilon_{3k+1}\mapsto Ez^k\), \(\varepsilon_{3k-1}\mapsto\frac12Fz^k\). Checking the nine residue pairs verifies the structure constants and Jacobi. This algebra is graded-simple for this grading, but is not simple as an ungraded Lie algebra: evaluation at z=1 has a proper nonzero kernel.}
\GSfield{Reconstruction}{The basis isomorphism and the graded-versus-ungraded simplicity distinction are supplied. This particular grading is not generated by degrees -1,0,1, so Kac's local-generation convention cannot be silently applied to it.}
\GSfield{Fate}{The loop-algebra setting and the distinction between homogeneous simplicity and arbitrary ideals are explicit in \citet[Introduction; Chapter~II, Sec.~6, Proposition~10]{kac1968graded}. The displayed basis identification is an elementary check of the notes' exact coefficients, not a claimed later discovery. Prior alias: 1977-457.18.}
\GSfield{Present status}{As of 10 November 2026: An identifiable centreless loop algebra with an unusual grading. Calling it simple without specifying homogeneous ideals is false.}
\GSfield{Assessment}{literal statement: correct with missing hypotheses; publication match: partial; verification: read the result; development: annotated; confidence: high}
\end{GSitem}
\begin{GSitem}{1978-457.12}{F}{The unspecified three-algebra and two-construction exhaustion claims}
\GSfield{Kind; source}{claim; \GSpage{P000829}}
\GSfield{The notes}{There are said to be three simple algebras with the displayed one-dimensional pieces, but only one is written. The final line says these two constructions exhaust everything up to grading.}
\GSfield{Statement}{No complete list of three, definition of simplicity, finite-growth bound, support condition, local-generation hypothesis or intended equivalence of gradings is supplied. The literal universal claim about all graded Lie algebras is false. A precise classical finite-growth classification exists under additional hypotheses, but the absent pages prevent identification of the exact quantified version intended here.}
\GSfield{Reconstruction}{The missing quantifiers are identified rather than supplied as historical facts; the unseen P000826--P000828 could have contained qualifications, but that is not evidence of their content.}
\GSfield{Fate}{\citet[Introduction and Chapter~III, Theorem~4]{kac1968graded} classifies irreducible graded-simple finite-growth algebras with its explicit conventions. It does not justify dropping those assumptions. Prior alias: 1977-457.19. No definite open classification problem is inferred from the incomplete list.}
\GSfield{Present status}{As of 10 November 2026: The literal scope is unidentified. Established classifications provide context, not proof of an unrestricted exhaustion claim.}
\GSfield{Assessment}{literal statement: not determinable; publication match: partial; verification: read the result; development: annotated; confidence: high}
\end{GSitem}

\GStalk{458}{1979-03-19}{V. I. Ogievetskii}{P000829-P000835 (catalogue); actual P000830-P000833}{Supersymmetry algebra, superspace, and chiral fields}
Although headed \ru{Супергравитация}, the surviving talk develops the kinematics of flat four-dimensional supersymmetry: the super-Poincar\'e brackets, the spinor description of Lorentz vectors, odd coordinates, superfields, supertranslations, and chiral coordinates. It contains no gravitational action, local supersymmetry constraints, or quantum-finiteness theorem. The first two pages overlap items 1977-457.20--25 in the supplied previous index; the dated heading on P000830 shows that these are duplicate indexing of the same manuscript material, not evidence of an earlier seminar occurrence. The substantive continuation on P000832--P000833 supplies the transformation and chirality formulas. P000829 belongs to the preceding Fuchs material, and P000834 begins Lebedev. Most of this introduction was already printed in the speaker's 1975 review, whose English translation carries a 1976 copyright \citep[Secs.~2--3]{ogievetskymezincescu1975superfields}.

\begin{GSitem}{1978-458.1}{E}{The real N=1 super-Poincare algebra}
\GSfield{Kind; source}{construction; \GSpage{P000830}}
\GSfield{The notes}{The notes list four \(P_m\), six \(L_{mn}\), and two-component \(Q_\alpha,\bar Q_{\dot\alpha}\), with \(\{Q_\alpha,\bar Q_{\dot\beta}\}=2\sigma^m_{\alpha\dot\beta}P_m\), zero same-chirality anticommutators and \([P,Q]=0\); the dimension is \((10,4)\).}
\GSfield{Statement}{Let \(\eta=\operatorname{diag}(1,-1,-1,-1)\), \(\sigma^m=(I,\sigma^1,\sigma^2,\sigma^3)\). Extend \(\mathfrak{iso}(1,3)\) by a real Majorana odd module, written as conjugate Weyl modules. The only nonzero odd brackets are the displayed mixed anticommutators; Lorentz generators act in \((\tfrac12,0)\oplus(0,\tfrac12)\). Together with the Poincar\'e brackets this is a real Lie superalgebra of superdimension \(10|4\). The complexification has the corresponding two independent complex Weyl modules.}
\GSfield{Reconstruction}{The Majorana reality condition, metric convention and Lorentz action complete shorthand in the notes. A bar/dot lost in the second bracket display is restored mathematically, not claimed as visible handwriting. The total 14 counts real dimensions.}
\GSfield{Fate}{Already standard in \citet[Sec.~2, Eqs.~(2.2)--(2.3)]{ogievetskymezincescu1975superfields}; publication predates the talk.}
\GSfield{Present status}{As of 10 November 2026: established in the stated scope; this passage poses no additional open question.}
\GSfield{Assessment}{literal statement: correct with missing hypotheses; publication match: exact; verification: read the result; development: none; confidence: high}
\end{GSitem}
\begin{GSitem}{1978-458.2}{E}{Hermitian matrices and the Lorentz double cover}
\GSfield{Kind; source}{construction; \GSpage{P000830}}
\GSfield{The notes}{The vector \(x\) is written as \(X=x^m\sigma_m=\left(\begin{smallmatrix}x_0+x_3&x_1-ix_2\\x_1+ix_2&x_0-x_3\end{smallmatrix}\right)\). The transformation \(A\mapsto X^+AX\), \(X\in SL(2,\mathbb C)\), preserves the determinant.}
\GSfield{Statement}{The real vector space \(\operatorname{Herm}_2(\mathbb C)\) has quadratic form \(\det X=(x^0)^2-\sum_{i=1}^3(x^i)^2\). The action \(X\mapsto gXg^\dagger\) induces the surjection \(SL(2,\mathbb C)\to SO^+(1,3)\) with kernel \(\{\pm I\}\). The fundamental and conjugate Weyl modules carry the invariant alternating spinor forms \(\varepsilon_{12}=-\varepsilon_{21}=1\).}
\GSfield{Reconstruction}{The notes reuse \(X\) for both a vector matrix and a group element. We separate these roles. The identity component and double-cover qualification are supplied.}
\GSfield{Fate}{The spinor conventions are reviewed in \citet[Appendix]{ogievetskymezincescu1975superfields}; this is classical Lorentz geometry.}
\GSfield{Present status}{As of 10 November 2026: established in the stated scope; this passage poses no additional open question.}
\GSfield{Assessment}{literal statement: correct with missing hypotheses; publication match: exact; verification: read the result; development: none; confidence: high}
\end{GSitem}
\begin{GSitem}{1978-458.3}{E}{Constructing a superalgebra from its even algebra and odd module}
\GSfield{Kind; source}{method; \GSpage{P000830-P000831}}
\GSfield{The notes}{At the end of P000830 the speaker starts with \(\mathfrak g_0\) and its module \(\mathfrak g_1\), seeks a bracket on their sum, and on P000831 says it is unique in the present case.}
\GSfield{Statement}{Over a field of characteristic zero, a Lie algebra \(\mathfrak g_0\), a \(\mathfrak g_0\)-module \(V\), and an equivariant symmetric map \(b:S^2V\to\mathfrak g_0\) define a Lie superalgebra precisely when \(b(v,w)\cdot u+b(w,u)\cdot v+b(u,v)\cdot w=0\). For the minimal supertranslation algebra, require translation-invariant odd generators and a nonzero Lorentz-equivariant mixed Weyl bracket into the translation module. That map is unique up to scale; fixing its coefficient as \(2\) fixes the stated normalization.}
\GSfield{Reconstruction}{We supply the odd Jacobi condition and limit uniqueness to the stated modules and target. Unrestricted extensions, including central charges, are not unique.}
\GSfield{Fate}{The required graded Jacobi framework appears in \citet[Sec.~2]{ogievetskymezincescu1975superfields}.}
\GSfield{Present status}{As of 10 November 2026: established in the stated scope; this passage poses no additional open question.}
\GSfield{Assessment}{literal statement: correct with missing hypotheses; publication match: exact; verification: read the result; development: none; confidence: high}
\end{GSitem}
\begin{GSitem}{1978-458.4}{E}{Extended supersymmetry and internal indices}
\GSfield{Kind; source}{construction; \GSpage{P000831}}
\GSfield{The notes}{Equation (2) introduces \(Q^i_\alpha\) and \(\{Q^i_\alpha,\bar Q^j_{\dot\beta}\}=2\sigma^m_{\alpha\dot\beta}P_m\delta^{ij}\), with \(i,j\) internal-symmetry indices.}
\GSfield{Statement}{For \(N\) Weyl supercharges and their conjugates, the displayed relation defines the central-charge-free extended supertranslation algebra; same-chirality brackets vanish. Internal unitary transformations act on the multiplicity index. More general \(N>1\) supersymmetry algebras may have \(\{Q^i_\alpha,Q^j_\beta\}=\varepsilon_{\alpha\beta}Z^{ij}\), where \(Z^{ij}=-Z^{ji}\) are central.}
\GSfield{Reconstruction}{The central-charge-free restriction is essential; the last clause records the scope of the classification, not an extra assertion made in the notebook.}
\GSfield{Fate}{\citet[Sec.~1, Eqs.~(1.1)--(1.7)]{haaglopuszanskisohnius1975supersymmetries} already included central charges in 1975 (received 15 November 1974).}
\GSfield{Present status}{As of 10 November 2026: established in the stated scope; this passage poses no additional open question.}
\GSfield{Assessment}{literal statement: correct with missing hypotheses; publication match: exact; verification: read the result; development: none; confidence: high}
\end{GSitem}
\begin{GSitem}{1978-458.5}{E}{The no-go theorem is a statement about scattering symmetries}
\GSfield{Kind; source}{remark; \GSpage{P000831}}
\GSfield{The notes}{The notes call this a nontrivial mixing of internal and spatial symmetry, write \ru{протон + антипротон = импульс}, and say ordinary Lie algebras cannot do it because of a theorem.}
\GSfield{Statement}{Under the hypotheses of a nontrivial relativistic scattering theory with a mass gap, finitely many particle types at each mass, and sufficiently regular additive symmetry generators, even scattering symmetries split into Poincar\'e and commuting internal parts; graded extensions can possess the displayed spinorial generators. This is not a theorem forbidding arbitrary abstract Lie-algebra extensions or semidirect products.}
\GSfield{Reconstruction}{The scattering hypotheses are restored from the relevant no-go/supersymmetry theorem. The proton phrase is an analogy for the spinor anticommutator, not a particle-reaction identity.}
\GSfield{Fate}{\citet[Secs.~1--2, especially limitations on p.~258]{haaglopuszanskisohnius1975supersymmetries} gives the precise scope.}
\GSfield{Present status}{As of 10 November 2026: established in the stated scope; this passage poses no additional open question.}
\GSfield{Assessment}{literal statement: correct with missing hypotheses; publication match: partial; verification: read the result; development: none; confidence: high}
\end{GSitem}
\begin{GSitem}{1978-458.6}{E}{Superspace and the finite superfield expansion}
\GSfield{Kind; source}{construction; \GSpage{P000831}}
\GSfield{The notes}{The notes adjoin \(\theta^\alpha,\bar\theta^{\dot\alpha}\) to \(x^m\), impose anticommutation, and expand \(\varphi=A+\theta\psi+\bar\theta\bar\psi+\theta\theta B+\bar\theta\bar\theta C+\theta\sigma\bar\theta h+\cdots+\theta\theta\bar\theta\bar\theta D\).}
\GSfield{Statement}{The complexified function sheaf of flat superspace is \(C^\infty(\mathbb R^{1,3};\mathbb C)\otimes\Lambda(\theta^1,\theta^2,\bar\theta^{\dot1},\bar\theta^{\dot2})\). Every superfunction has a unique expansion in its \(2^4=16\) exterior monomials; its total odd degree is at most four. The quadratic mixed term is \(\theta\sigma^m\bar\theta\,h_m(x)\), with a vector index \(m\), and the cubic terms are included in the ellipsis. Translations act by \(P_m=-i\partial_m\); Lorentz transformations also act on the spinor coordinates.}
\GSfield{Reconstruction}{We distinguish a sheaf of superfunctions from ordinary points with numeric anticommuting coordinates. The ambiguous manuscript \(\alpha\) on \(\sigma,h\) must be a vector index in the mixed term.}
\GSfield{Fate}{See \citet[Sec.~3(a)--(b), Eq.~(3.12)]{ogievetskymezincescu1975superfields} for this pre-talk construction.}
\GSfield{Present status}{As of 10 November 2026: established in the stated scope; this passage poses no additional open question.}
\GSfield{Assessment}{literal statement: correct with missing hypotheses; publication match: exact; verification: read the result; development: none; confidence: high}
\end{GSitem}
\begin{GSitem}{1978-458.7}{E}{Odd differential operators realize the supertranslation brackets}
\GSfield{Kind; source}{construction; \GSpage{P000831-P000832}}
\GSfield{The notes}{Equation (5) gives \(Q_\alpha=-i\partial_\alpha-(\sigma^m\bar\theta)_\alpha\partial_m\) and \(\bar Q_{\dot\alpha}=i\bar\partial_{\dot\alpha}+(\theta\sigma^m)_{\dot\alpha}\partial_m\). The next page says these realize (1).}
\GSfield{Statement}{Use left odd derivatives on the exterior algebra and \(P_m=-i\partial_m\). The displayed odd operators satisfy \(\{Q_\alpha,\bar Q_{\dot\beta}\}=-2i\sigma^m_{\alpha\dot\beta}\partial_m=2\sigma^m_{\alpha\dot\beta}P_m\), while \(\{Q,Q\}=\{\bar Q,\bar Q\}=0\) and they commute with all translations.}
\GSfield{Reconstruction}{The side of odd differentiation is specified; changing it without changing signs changes the formulas. The assertion follows directly by the graded Leibniz rule.}
\GSfield{Fate}{This superspace realization predates 1979; see \citet[Sec.~3]{ogievetskymezincescu1975superfields}.}
\GSfield{Present status}{As of 10 November 2026: established in the stated scope; this passage poses no additional open question.}
\GSfield{Assessment}{literal statement: correct with missing hypotheses; publication match: exact; verification: read the result; development: none; confidence: high}
\end{GSitem}
\begin{GSitem}{1978-458.8}{E}{Supertranslations on coordinates}
\GSfield{Kind; source}{construction; \GSpage{P000832}}
\GSfield{The notes}{Equation (6) shifts \(\theta\) by \(\lambda\), \(\bar\theta\) by \(\bar\lambda\), and \(x^m\) by \(i(\theta\sigma^m\bar\lambda-\lambda\sigma^m\bar\theta)\).}
\GSfield{Statement}{For constant odd parameters \(\lambda,\bar\lambda\) anticommuting with all odd coordinates, this polynomial transformation, together with ordinary translations, gives the supertranslation group law. In infinitesimal form \(\delta_\lambda\theta=\lambda\), \(\delta_\lambda\bar\theta=\bar\lambda\), and \(\delta_\lambda x^m=i(\theta\sigma^m\bar\lambda-\lambda\sigma^m\bar\theta)\).}
\GSfield{Reconstruction}{All odd parameters anticommute mutually. Formula (6) is interpreted in the functor-of-points sense; its body is an ordinary spacetime translation.}
\GSfield{Fate}{Compare \citet[Sec.~3(a), Eq.~(3.4)]{ogievetskymezincescu1975superfields}, with the chosen sign convention.}
\GSfield{Present status}{As of 10 November 2026: established in the stated scope; this passage poses no additional open question.}
\GSfield{Assessment}{literal statement: correct with missing hypotheses; publication match: exact; verification: read the result; development: none; confidence: high}
\end{GSitem}
\begin{GSitem}{1978-458.9}{E}{The commutator closes on translations, with a factor correction}
\GSfield{Kind; source}{claim; \GSpage{P000832}}
\GSfield{The notes}{After saying that doing two transformations gives a pure translation, equation (7) writes \((\delta_\lambda\delta_\mu-\delta_\mu\delta_\lambda)x^m=i(\mu\sigma^m\bar\lambda-\lambda\sigma^m\bar\mu)\). The scan has no factor \(2\).}
\GSfield{Statement}{For the preceding infinitesimal variations, treated as even derivations with fixed parameters, \([\delta_\lambda,\delta_\mu]x^m=2i(\lambda\sigma^m\bar\mu-\mu\sigma^m\bar\lambda)\) and \([\delta_\lambda,\delta_\mu]\theta=[\delta_\lambda,\delta_\mu]\bar\theta=0\). Reversing commutator order reverses the sign but never removes the factor two. An ordinary composition of two supertranslations still shifts the odd coordinates; it is their commutator that is purely translational.}
\GSfield{Reconstruction}{This is an explicit algebraic correction of the notebook, independently checked by applying each variation to equation (6). It is not a transcription error.}
\GSfield{Fate}{Closure with coefficient two is already displayed in \citet[Eq.~(1.2)]{ogievetskymezincescu1975superfields}.}
\GSfield{Present status}{As of 10 November 2026: established in the stated scope; this passage poses no additional open question.}
\GSfield{Assessment}{literal statement: false as written; publication match: partial; verification: read the result; development: annotated; confidence: high}
\end{GSitem}
\begin{GSitem}{1978-458.10}{E}{A scalar superfield transforms by pullback}
\GSfield{Kind; source}{definition; \GSpage{P000832}}
\GSfield{The notes}{Equation (8) is \(\varphi'(x',\theta',\bar\theta')=\varphi(x,\theta,\bar\theta)\); the component transformations are described as complicated.}
\GSfield{Statement}{For a superspace automorphism \(T\), a scalar superfield has the passive transformation \(\varphi'=\varphi\circ T^{-1}\). Expanding in the exterior generators determines its component transformations uniquely and mixes even and odd component fields. This is a transformation law, not an invariance assertion \(\varphi' =\varphi\) at fixed coordinates.}
\GSfield{Reconstruction}{Passive and active conventions are distinguished to avoid a sign ambiguity in component variations.}
\GSfield{Fate}{Compare \citet[Sec.~3(b), Eqs.~(3.13)--(3.14)]{ogievetskymezincescu1975superfields}.}
\GSfield{Present status}{As of 10 November 2026: established in the stated scope; this passage poses no additional open question.}
\GSfield{Assessment}{literal statement: correct; publication match: exact; verification: read the result; development: none; confidence: high}
\end{GSitem}
\begin{GSitem}{1978-458.11}{E}{Chiral coordinates and left-right decomposition}
\GSfield{Kind; source}{construction; \GSpage{P000832-P000833}}
\GSfield{The notes}{Equations (9), (11) define \(x_L=x+i\theta\sigma\bar\theta\), \(x_R=x-i\theta\sigma\bar\theta\), with \(\delta x_L=2i\theta\sigma\bar\lambda\) and \(\delta x_R=-2i\lambda\sigma\bar\theta\). P000833 records their sum and difference.}
\GSfield{Statement}{In the same Grassmann conventions, the infinitesimal transformation of \((x_L,\theta)\) depends only on \((x_L,\theta)\) and the parameters; likewise for \((x_R,\bar\theta)\). Moreover \(x=(x_L+x_R)/2\) and \(x_L-x_R=2i\theta\sigma\bar\theta\). These are coordinate identities on complexified superspace. For finite odd shifts the induced \(x_L,x_R\) transformations also contain terms quadratic in the shift parameters.}
\GSfield{Reconstruction}{The formulas for \(\delta\) are explicitly infinitesimal, as the notebook says; they must not be read as complete finite transformation formulas.}
\GSfield{Fate}{The equivalent chiral changes of realization are in \citet[Sec.~3(a), Eqs.~(3.8)--(3.9)]{ogievetskymezincescu1975superfields}.}
\GSfield{Present status}{As of 10 November 2026: established in the stated scope; this passage poses no additional open question.}
\GSfield{Assessment}{literal statement: correct with missing hypotheses; publication match: exact; verification: read the result; development: none; confidence: high}
\end{GSitem}
\begin{GSitem}{1978-458.12}{E}{The chiral multiplet and its auxiliary component}
\GSfield{Kind; source}{construction; \GSpage{P000832}}
\GSfield{The notes}{The special \ru{киральные} fields are \(\varphi_L(x_L,\theta)=A(x_L)+\theta^\alpha\psi_\alpha(x_L)+\theta^\alpha\theta_\alpha F(x_L)\).}
\GSfield{Statement}{A chiral scalar superfield is a superfunction independent of \(\bar\theta\) in the coordinates \((x_L,\theta,\bar\theta)\), equivalently annihilated by the corresponding covariant odd derivatives. Its expansion has exactly a complex scalar \(A\), a Weyl spinor \(\psi\), and a complex scalar auxiliary coefficient \(F\). This off-shell construction is closed under the global supersymmetry action; no field equation or gravitational coupling is asserted.}
\GSfield{Reconstruction}{The normalization of \(\psi\) and \(F\) is conventional; factors such as \(\sqrt2\) can be absorbed into them. The covariant-derivative formulation is a modern equivalent definition.}
\GSfield{Fate}{Already in \citet[Sec.~3(c)]{ogievetskymezincescu1975superfields}.}
\GSfield{Present status}{As of 10 November 2026: established in the stated scope; this passage poses no additional open question.}
\GSfield{Assessment}{literal statement: correct; publication match: exact; verification: read the result; development: none; confidence: high}
\end{GSitem}

\GStalk{459}{1979-03-26}{D. R. Lebedev}{P000832-P000837 (catalogue); actual P000834-P000836}{Local Lie algebras and formal differential groups}
The dated heading is \ru{Лебедев (Митя). О результатах Ритта и Кириллова по дифференциальной алгебре}. Three strands are juxtaposed: continuous support-local brackets on sections of bundles, Kirillov's rank-one classification, and Ritt's formal differential groups with one derivation. A formal pseudodifferential-symbol example connects these to the speaker's Volterra-group work. The source does not justify identifying the smooth, formal and differential-algebraic categories. In particular its ``ten Ritt algebras'' count omits three finite types, and its universal extension assertion omits even the product of two substitution groups. The contact formula also has a genuine sign inconsistency, confirmed from the scan and by two independent calculations. The question marks at higher ranks specify no classification conjecture.

\begin{GSitem}{1978-459.1}{E}{A continuous local Lie bracket on sections}
\GSfield{Kind; source}{definition; \GSpage{P000834}}
\GSfield{The notes}{For a smooth manifold \(M\) and vector bundle \(E\), the notes give a bilinear skew bracket on \(\Gamma^\infty(E)\), satisfying Jacobi, continuous for uniform convergence of every derivative on compact sets, and \(\operatorname{supp}[s,t]\subset\operatorname{supp}s\cap\operatorname{supp}t\).}
\GSfield{Statement}{Let \(M\) be a finite-dimensional smooth manifold and \(E\to M\) a finite-rank smooth real vector bundle. A local Lie algebra structure is a jointly continuous \(\mathbb R\)-bilinear Lie bracket on the Fr\'echet space \(\Gamma^\infty(E)\) with the displayed support property. Restriction gives a sheaf of Lie algebras. No \(C^\infty(M)\)-bilinearity or fixed global differential order is assumed.}
\GSfield{Reconstruction}{Second countability and the usual smooth-section topology are understood; the notes' differential algebra here means a local Lie algebra, not a differential field.}
\GSfield{Fate}{This is the setting of \citet[Sec.~1, Russian pp.~58--59]{kirillov1976local}, published in 1976; received 25 November 1975.}
\GSfield{Present status}{As of 10 November 2026: established in the stated scope; this passage poses no additional open question.}
\GSfield{Assessment}{literal statement: correct; publication match: exact; verification: read the result; development: none; confidence: high}
\end{GSitem}
\begin{GSitem}{1978-459.2}{E}{Locality gives differential expressions; rank one gives first order}
\GSfield{Kind; source}{theorem; \GSpage{P000835}}
\GSfield{The notes}{In a local trivialization the bracket is written \([f,g]=\sum c_{kl}(x)\partial^k f\partial^l g\).}
\GSfield{Statement}{A continuous support-local bilinear operation on smooth sections is locally a bidifferential operator: on every sufficiently small relatively compact chart only finitely many derivatives occur. For a Lie bracket on a line bundle, Jacobi and skew symmetry reduce its order to at most one in each argument. In a local frame it therefore has the form \([f,g]=\Lambda(df,dg)+fE(g)-gE(f)\), where \(\Lambda\) is a bivector and \(E\) a vector field, subject to the Jacobi identity.}
\GSfield{Reconstruction}{The locally finite order and the later first-order restriction are separated. The displayed arbitrary multi-index sum alone does not assert infinite differential order.}
\GSfield{Fate}{\citet[Sec.~2, Lemmas~1--2, Russian pp.~59--60]{kirillov1976local} proves both reductions. The author's \citet[p.~268]{kirillov1977errata} corrects the coordinate Jacobi identities of Lemma~3; it does not withdraw the finite-order and first-order statements of Lemmas~1--2 used here.}
\GSfield{Present status}{As of 10 November 2026: established in the stated scope; this passage poses no additional open question.}
\GSfield{Assessment}{literal statement: correct with missing hypotheses; publication match: exact; verification: read the result; development: none; confidence: high}
\end{GSitem}
\begin{GSitem}{1978-459.3}{E}{The scope of Kirillov's rank-one classification}
\GSfield{Kind; source}{theorem; \GSpage{P000835}}
\GSfield{The notes}{The theorem says rank-one differential Lie algebras are essentially Poisson or contact; the phrase \ru{Все остальные формально изоморфны} follows the two models.}
\GSfield{Statement}{The characteristic distribution generated locally by \(E\) and \(\Lambda^\sharp T^*M\) has integral leaves. On a transitive leaf the local Lie structure is of contact type in odd dimension and of locally conformally symplectic type in positive even dimension; after a local change of line-bundle frame the even model is symplectic Poisson. Rank-zero leaves and singular changes of leaf dimension are permitted for a general local Lie algebra. The theorem is leafwise and local, and does not identify every global bracket on every line bundle with one fixed flat model.}
\GSfield{Reconstruction}{Transitivity, localization and line-bundle changes of frame repair the overly compressed assertion. This is a classification of rank-one bundles, not bundles of arbitrary rank.}
\GSfield{Fate}{See \citet[Sec.~2, Theorem~1; Secs.~3--4, Theorems~2--3]{kirillov1976local}; all precede the session.}
\GSfield{Present status}{As of 10 November 2026: established in the stated scope; this passage poses no additional open question.}
\GSfield{Assessment}{literal statement: false as written; publication match: partial; verification: read the result; development: none; confidence: high}
\end{GSitem}
\begin{GSitem}{1978-459.4}{E}{The standard Poisson model}
\GSfield{Kind; source}{construction; \GSpage{P000835}}
\GSfield{The notes}{Case (a) takes \(M=\mathbb R^{2n}\) and \([f,g]=\sum_i(f_{x_i}g_{y_i}-f_{y_i}g_{x_i})\).}
\GSfield{Statement}{On \(\mathbb R^{2n}\), the inverse bivector to the constant symplectic form gives this Lie bracket on smooth functions. It is a first-order bidifferential operation in each input, obeys the Leibniz rule, and is transitive for \(n\ge1\).}
\GSfield{Reconstruction}{The symplectic interpretation explains the even-dimensional model; its local nature is essential in the preceding classification.}
\GSfield{Fate}{This is the flat model in \citet[Sec.~3]{kirillov1976local}, not a new 1979 result.}
\GSfield{Present status}{As of 10 November 2026: established in the stated scope; this passage poses no additional open question.}
\GSfield{Assessment}{literal statement: correct; publication match: exact; verification: read the result; development: none; confidence: high}
\end{GSitem}
\begin{GSitem}{1978-459.5}{E}{The contact model has a sign error in the notebook}
\GSfield{Kind; source}{construction; \GSpage{P000835}}
\GSfield{The notes}{Case (b) writes \(X_i=\partial_{x_i}\), \(Y_i=\partial_{y_i}+x_i\partial_z\), \(Z=\partial_z\), and \([f,g]=fZg-gZf+\sum_i(X_if\,Y_ig-X_ig\,Y_if)\). Both plus signs are legible.}
\GSfield{Statement}{With the displayed plus sign in \(Y_i\), the operation fails Jacobi already for one pair: \([x,y]=1\), \([y,z]=y\), \([z,x]=-2x\), whence \([x,[y,z]]+[y,[z,x]]+[z,[x,y]]=2\). A consistent contact bracket is obtained by replacing \(Y_i\) with \(\partial_{y_i}-x_i\partial_z\), retaining the rest of the displayed formula. It corresponds to the contact form \(dz+\sum_i x_i\,dy_i\) in these conventions.}
\GSfield{Reconstruction}{The sign replacement is an editorial mathematical emendation, not a claim about what the writer wrote. The explicit Jacobi obstruction and corrected bracket are checked in the critical annotations and the reproducible symbolic calculation.}
\GSfield{Fate}{Contact local Lie algebras themselves are classical; compare \citet[Sec.~4, Theorem~3]{kirillov1976local}. The notebook's literal formula is refuted directly.}
\GSfield{Present status}{As of 10 November 2026: the intended contact model is established; the literal notebook formula is false. No open problem is created by the error.}
\GSfield{Assessment}{literal statement: false as written; publication match: partial; verification: read the result; development: annotated; confidence: high}
\end{GSitem}
\begin{GSitem}{1978-459.6}{F}{The diagram of higher-rank classification gaps}
\GSfield{Kind; source}{remark; \GSpage{P000835}}
\GSfield{The notes}{A diagram has the rank-one row marked Kirillov, the one-dimensional-base/rank-two point marked \ru{10 алгебр Ритта}, and a region marked question marks and \ru{есть только правдоподобная версия}.}
\GSfield{Statement}{The diagram distinguishes base dimension from fiber rank. It supplies no candidate brackets, equivalence relation, transitivity condition, or formal versus smooth category for the region of question marks. Consequently no definite higher-rank conjecture can be reconstructed. Ritt's rank-two point can be identified separately with the formal one-derivation classification.}
\GSfield{Reconstruction}{The count is corrected in item 1978-459.11. The unspecified ``plausible version'' is retained as an unresolved identification rather than turned into a modern classification problem.}
\GSfield{Fate}{Searches of Kirillov's original paper and Ritt/Cassidy material located the two named theories, but no precise conjecture matching this diagram. \citet[Sec.~5, point~6]{kirillov1976local} explicitly distinguishes the greater complexity of higher-rank bundles.}
\GSfield{Present status}{As of 10 November 2026: the intended higher-rank statement is not identifiable from the surviving notes; this is not evidence that a specified problem remains open.}
\GSfield{Assessment}{literal statement: not determinable; publication match: none found; verification: unresolved lead; development: none; confidence: medium}
\end{GSitem}
\begin{GSitem}{1978-459.7}{E}{Formal pseudodifferential symbols as an infinite-component Lie algebra}
\GSfield{Kind; source}{construction; \GSpage{P000835}}
\GSfield{The notes}{The series \(\sigma=\sum_{i\ge1}a_i\xi^{-(i+1)}\) is supplied with a noncommutative product determined by \([\xi,x]=1\); its commutator is said to be a differential Lie algebra on a one-dimensional base with infinite-dimensional fiber.}
\GSfield{Statement}{Let \((F,\partial)\) be a commutative differential algebra of characteristic zero. Formal symbols with orders bounded above have associative multiplication \((a\xi^m)\circ(b\xi^n)=\sum_{j\ge0}\binom mj a\,\partial^jb\,\xi^{m+n-j}\). Each coefficient is a finite sum. The negative-order symbols, and also the order-at-most-minus-two tail actually displayed, are closed under commutator. The relation \([\xi,a]=\partial a\) specializes to \([\xi,x]=1\) when \(\partial x=1\).}
\GSfield{Reconstruction}{We specify formal multiplication and coefficientwise completion. The possibly duplicated extra \(a\) after \(a_i\) is not assigned mathematical significance. The index \(i\ge1\) is preserved; adding a \(\xi^{-1}\) coefficient gives the full usual Volterra Lie algebra. This infinite-rank example lies outside the finite-rank smooth-bundle hypotheses.}
\GSfield{Fate}{The same construction is in \citet[Secs.~3--4, Russian p.~42]{lebedevmanin1979volterra}, received 10 November 1978 and published in 1979. Receipt does not establish public circulation; this passage introduces a formal-symbol construction, not the later Hamiltonian theorem.}
\GSfield{Present status}{As of 10 November 2026: the formal symbol algebra and its commutator are established. The analytic operator realization and convergence are not asserted here.}
\GSfield{Assessment}{literal statement: correct with missing hypotheses; publication match: partial; verification: read the result; development: none; confidence: high}
\end{GSitem}
\begin{GSitem}{1978-459.8}{E}{Formal differential power series and the required completion}
\GSfield{Kind; source}{definition; \GSpage{P000835-P000836}}
\GSfield{The notes}{The notes introduce a characteristic-zero differential field, independent jets \(y_i,y_i',\ldots\), and \(F[\![ y_1,\ldots,y_m]\!]=\bigoplus_{\tau\ge0}F_\tau\). The next page writes \(U=\sum_kP_k\) with each \(P_k\) homogeneous in all jets.}
\GSfield{Statement}{Give every jet \(\partial^j y_i\) total degree one. Let \(F_\tau\) be the space of differential polynomials homogeneous of degree \(\tau\), each containing finitely many monomials. The formal differential power-series algebra used here is the degree completion \(\prod_{\tau\ge0}F_\tau\). Differentiation preserves degree and continuous substitution of series with zero constant term is defined degree by degree.}
\GSfield{Reconstruction}{The direct sum in the displayed definition would give only differential polynomials, so it is replaced by the completed product. The finite-monomial condition at each degree prevents an additional unsupported completion in jet order.}
\GSfield{Fate}{This formal category is explained in \citet[Ch.~III, Sec.~1, pp.~263--264]{cassidy1979differential}, restating Ritt's 1950 formal theory.}
\GSfield{Present status}{As of 10 November 2026: established in the stated scope; this passage poses no additional open question.}
\GSfield{Assessment}{literal statement: false as written; publication match: partial; verification: read the result; development: none; confidence: high}
\end{GSitem}
\begin{GSitem}{1978-459.9}{E}{Formal differential group laws}
\GSfield{Kind; source}{definition; \GSpage{P000835-P000836}}
\GSfield{The notes}{A differential group is given by \(U=(U_1,\ldots,U_m)\) in the jets of \(u,v\), with \(U(u,0)=u\), \(U(0,v)=v\), and \(U(U(u,v),w)=U(u,U(v,w))\).}
\GSfield{Statement}{Over a characteristic-zero differential field, a formal differential group law of dimension \(m\) is an \(m\)-tuple of degree-completed differential series in two \(m\)-tuples, with linear term \(u+v\), satisfying the identities shown. Composition includes all formal derivatives of substituted series. The formal inverse is obtained recursively by total degree. Isomorphisms are invertible formal differential substitutions fixing zero, not arbitrary changes of a smooth infinite-dimensional manifold.}
\GSfield{Reconstruction}{The inverse and equivalence category are made explicit. Here ``dimension'' counts independent differential coordinates, not the infinitely many jet variables as ordinary coordinates.}
\GSfield{Fate}{\citet[Ch.~III, Secs.~1--3]{cassidy1979differential} restates the formal groups and their Lie correspondence from Ritt's 1950--1951 papers.}
\GSfield{Present status}{As of 10 November 2026: established in the stated scope; this passage poses no additional open question.}
\GSfield{Assessment}{literal statement: correct with missing hypotheses; publication match: exact; verification: read the result; development: none; confidence: high}
\end{GSitem}
\begin{GSitem}{1978-459.10}{E}{The substitution group and tensor-density modules}
\GSfield{Kind; source}{construction; \GSpage{P000836}}
\GSfield{The notes}{The basic group \(\Pi\) has \(U(u,v)=u+v(x+u)=u+v+uv'+\tfrac12u^2v''+\cdots\). The notes also introduce the additive group \(A_h\) of expressions \(u(x)(dx)^{-h}\).}
\GSfield{Statement}{Near the identity, write a formal change of variable as \(\phi_u(x)=x+u(x)\). Composition \(\phi_v\circ\phi_u\) has parameter \(u+\sum_{j\ge0}u^j\partial^jv/j!\), a well-defined differential series. The skew bilinear term is \(uv'-vu'\). For scalar \(h\), tensor densities of weight \(-h\) form an additive module: the pullback action of \(\phi\) is \(a(x)(dx)^{-h}\mapsto a(\phi(x))(\phi'(x))^{-h}(dx)^{-h}\), with binomial powers expanded formally near \(\phi'=1\).}
\GSfield{Reconstruction}{The density action is supplied to explain the extension language. These are formal changes near the identity, not transformations which literally fix every value of \(x\). Associativity follows from composition.}
\GSfield{Fate}{The substitution formal group and its Lie algebra occur in \citet[Ch.~IV, Sec.~1, Theorem~5]{cassidy1979differential}; the density action is the explicit interpretation of the notebook's notation.}
\GSfield{Present status}{As of 10 November 2026: established in the stated scope; this passage poses no additional open question.}
\GSfield{Assessment}{literal statement: correct with missing hypotheses; publication match: exact; verification: read the result; development: none; confidence: high}
\end{GSitem}
\begin{GSitem}{1978-459.11}{E}{Ritt's rank-two classification has thirteen types}
\GSfield{Kind; source}{theorem; \GSpage{P000835-P000836}}
\GSfield{The notes}{The diagram says ``10 algebras of Ritt''. Theorem~2 says that when the base has dimension one and the fiber has dimension two, all the groups are extensions of \(\Pi\) by \(A_h\).}
\GSfield{Statement}{Over a universal ordinary differential field of characteristic zero, the differential Lie-algebra structures on its two-dimensional additive vector group, up to differential isomorphism, fall into thirteen types: three finite and ten substitutional types, with parameter families inside some types. The corresponding formal differential-group laws are classified through Ritt's Lie correspondence. They are not all extensions of a substitution group by an abelian density module: the additive two-coordinate group has no nonabelian substitution quotient, while one substitutional type is \(\Pi\times\Pi\). Many remaining substitutional types do have an abelian differential ideal and a substitutional quotient, with nontrivial cocycles possible.}
\GSfield{Reconstruction}{This is a correction of the count and universal quantifier, with the formal category and the universal differential-field setting used by Cassidy restored. Descent to an arbitrary coefficient field is not asserted. No classification of all smooth rank-two bundles over arbitrary curves is inferred. A list of thirteen types is a list of families, not thirteen isomorphism classes.}
\GSfield{Fate}{The original is \citet[pp.~491--519]{ritt1951groups} (metadata checked); the full classification list and exceptions were read in \citet[Introduction, p.~247; Ch.~IV, Sec.~2, pp.~268--269; Sec.~3]{cassidy1979differential}. Cassidy was received 15 September 1976 and published in 1979; Ritt's result was published in 1951.}
\GSfield{Present status}{As of 10 November 2026: the formal one-derivation classification in the stated universal-field setting is established; the unqualified extension statement in the source is false.}
\GSfield{Assessment}{literal statement: false as written; publication match: partial; verification: read the result; development: annotated; confidence: high}
\end{GSitem}

\GStalk{460}{1979-04-02}{G. B. Shabat (catalogue); initial in notes uncertain}{P000834-P000842 (catalogue); actual P000837-P000840}{Web geometry and a Zamolodchikov intervention on vertex models}
The first part reports Chern and Griffiths on algebraic webs, abelian relations, the Castelnuovo rank bound and linearization at maximal rank. Gelfand adds Crofton's formula and the hexagonal three-web problem. The notebook cites the exact 1978 Chern--Griffiths article. The main theorem was then believed proved, but the authors' 1981 correction added normality; Tr\'epreau supplied the unrestricted result in a preprint of 4 August 2005, published in 2006. The heading's last year digit and the speaker's initial are uncertain: the catalogue gives 2 April 1979 and G.~B.~Shabat, while the transcription tentatively reads 1978 and K.~Shabat. Catalogue identity is retained with this qualification.\\
The latter part of P000839 explicitly names \ru{Замолодчиков А. (Саша)}. His matrix contractions, three-line move and statistical-mechanical interpretation continue through P000840, without a second date. Their assignment as an intervention in this session is an editorial reconstruction. These are not assertions by Shabat. The general asymptotic assertions lack enough hypotheses to identify a definite theorem. P000841 begins a different dated talk, R.~E.~Kallosh on 24 September 1979, and P000842 continues it; neither page belongs to this record.

\begin{GSitem}{1978-460.1}{E}{Regular webs and the needed general-position condition}
\GSfield{Kind; source}{definition; \GSpage{P000837}}
\GSfield{The notes}{A web is a finite number of foliations with transverse leaves on a manifold.}
\GSfield{Statement}{For the codimension-one results below, a regular holomorphic \(d\)-web near \(0\in\mathbb C^n\) is a family of submersions \(u_1,\ldots,u_d\), taken up to one-variable reparametrization, whose level sets define the foliations. Every subset of at most \(n\) covectors among \(du_1(0),\ldots,du_d(0)\) is linearly independent. Shrink the neighborhood so that this holds everywhere. Pairwise transversality alone is insufficient when \(n>2\).}
\GSfield{Reconstruction}{The holomorphic germ category and full conormal general position are supplied for the precise theorems, rather than attributed to the notebook's single word ``transverse''.}
\GSfield{Fate}{The contemporary formulation is in \citet[Sec.~1]{beauville1980tissus}, a February 1979 lecture published in 1980; compare the exact hypotheses in \citet[Sec.~1.1]{trepreau2006algebrisation}.}
\GSfield{Present status}{As of 10 November 2026: established in the stated scope; this passage poses no additional open question.}
\GSfield{Assessment}{literal statement: correct with missing hypotheses; publication match: exact; verification: read the result; development: none; confidence: high}
\end{GSitem}
\begin{GSitem}{1978-460.2}{E}{Incidence webs from a projective variety}
\GSfield{Kind; source}{construction; \GSpage{P000837}}
\GSfield{The notes}{For \(X\subset\mathbb P^n\), \(\dim X=m\), \(\deg X=d\), the notes take the open set of complementary planes meeting \(X\) in \(d\) distinct points. They call \(\{H:\exists i,\ P_i\in H\}\) a leaf and say the foliation has codimension one.}
\GSfield{Statement}{Let \(X\) be a reduced pure \(m\)-dimensional projective variety over \(\mathbb C\), and let \(U\subset\operatorname{Gr}(n-m+1,n+1)\) consist of projective \((n-m)\)-planes meeting \(X\) transversely in \(d\) distinct smooth points. Near \(H_0\in U\), label these intersection points holomorphically as \(p_i(H)\). Each map \(p_i:U\to X\) is locally a submersion, and its fibers form a foliation of codimension \(m\). The \(i\)-th leaf through \(H_0\) consists locally of planes containing the fixed point \(p_i(H_0)\). The existential quantifier in the notes gives the union of \(d\) leaves, not a single leaf.}
\GSfield{Reconstruction}{Grassmannian convention, smooth transverse intersections, local labels and the correct codimension are supplied. Global labeling can fail by permutation monodromy. Codimension one applies only after the next page's specialization \(m=1\).}
\GSfield{Fate}{This incidence construction is described in \citet[Sec.~2, pp.~105--106; Sec.~11]{beauville1980tissus}; the erroneous codimension and union are repaired directly from the incidence dimensions.}
\GSfield{Present status}{As of 10 November 2026: established in the stated scope; this passage poses no additional open question.}
\GSfield{Assessment}{literal statement: false as written; publication match: partial; verification: read the result; development: none; confidence: high}
\end{GSitem}
\begin{GSitem}{1978-460.3}{E}{Gelfand's Crofton formula and its normalization}
\GSfield{Kind; source}{theorem; \GSpage{P000837}; I. M. Gelfand}
\GSfield{The notes}{\ru{Крофтона формула}: the length of a curve is the average number of its intersections with a line, averaged using an invariant measure.}
\GSfield{Statement}{For a compact piecewise \(C^1\) immersed plane curve \(\gamma\), counted with parametrization multiplicity, put \(N(\theta,p)=\#\{t:\gamma(t)\cdot(\cos\theta,\sin\theta)=p\}\) for almost all \((\theta,p)\). With unoriented lines represented by \(0\le\theta<\pi\), \(p\in\mathbb R\), one has \(\operatorname{length}(\gamma)=\tfrac12\int_0^\pi\int_{\mathbb R}N(\theta,p)\,dp\,d\theta\). The invariant measure is infinite on all lines, so ``average'' means an integral with a fixed normalization, not a uniform probability distribution on the entire line space.}
\GSfield{Reconstruction}{The coefficient and regularity are supplied. Applying the one-dimensional change-of-variables formula to each projection gives \(\int N(\theta,p)dp=\int|\gamma'(t)\cdot n_\theta|dt\); integrating \(\theta\) gives \(2\int|\gamma'|dt\), independently checking the normalization.}
\GSfield{Fate}{The formula is classical; \citet[Introduction and Sec.~3, pp.~79--80]{linwang1996integral} recalls its nineteenth-century origin and integral-geometric interpretation. That later paper is a checked source for the classical attribution, not the fate of a new 1979 question.}
\GSfield{Present status}{As of 10 November 2026: established in the stated scope; this passage poses no additional open question.}
\GSfield{Assessment}{literal statement: correct with missing hypotheses; publication match: partial; verification: read the result; development: annotated; confidence: high}
\end{GSitem}
\begin{GSitem}{1978-460.4}{E}{When can a planar three-web be made parallel?}
\GSfield{Kind; source}{question; \GSpage{P000837}; I. M. Gelfand}
\GSfield{The notes}{Gelfand asks when three transverse planar foliations can be straightened; two always can. The answer in the notes is hexagonality, dated approximately 1927, with an infinitesimal closed-path diagram and a curvature analogy.}
\GSfield{Statement}{For a regular holomorphic planar three-web, local equivalence to \(x=\mathrm{const}\), \(y=\mathrm{const}\), \(x+y=\mathrm{const}\) is equivalent to closure of every sufficiently small web hexagon. Equivalently, choose defining forms \(\omega_1+\omega_2+\omega_3=0\) and the unique connection form \(\gamma\) satisfying \(d\omega_i=\omega_i\wedge\gamma\); the condition is \(d\gamma=0\). The corresponding statement holds in the regular smooth real category. This is parallelizability, a stronger requirement than merely obtaining three families of nonparallel straight lines.}
\GSfield{Reconstruction}{The standard hexagonal closure, rather than an inferred exact traversal of the faint diagram, is used. The connection formulation makes the notebook's curvature analogy precise. No general planar three-web linearizability conjecture is substituted for this old problem.}
\GSfield{Fate}{Thomsen's 1927 result is explicitly recalled in \citet[Sec.~1, pp.~103--104]{beauville1980tissus}. The notebook itself gives both question and answer.}
\GSfield{Present status}{As of 10 November 2026: solved in the stated parallelizability sense; the hexagon criterion is classical.}
\GSfield{Assessment}{literal statement: correct with missing hypotheses; publication match: exact; verification: read the result; development: none; confidence: high; problem \hyperref[prob:GP-1978-20]{GP-1978-20}}
\end{GSitem}
\begin{GSitem}{1978-460.5}{E}{The dual web of a projective curve}
\GSfield{Kind; source}{construction; \GSpage{P000838}}
\GSfield{The notes}{The notes now put \(m=1\), take a projective curve \(X\), and construct a codimension-one web on the open subset of \(\check{\mathbb P}^{n}\) consisting of transverse hyperplanes.}
\GSfield{Statement}{Let \(C\subset\mathbb P^n\) be a smooth irreducible nondegenerate complex projective curve of degree \(d\). On a sufficiently small open set of hyperplanes with \(d\) distinct transverse intersections in linear general position, label the intersection maps \(p_i\). The hypersurfaces \(p_i(H)=p_i(H_0)\) are pieces of hyperplanes in the dual projective space, giving a linear \(d\)-web. The regularity and general-position open set is part of the construction.}
\GSfield{Reconstruction}{Smoothness, nondegeneracy and the choice of a regular open set avoid unsupported assertions about singular intersections or degenerate curves. The specialization is distinct from the general-codimension construction.}
\GSfield{Fate}{\citet[Sec.~2]{beauville1980tissus} gives exactly this construction; see also \citet[Sec.~1.2]{trepreau2006algebrisation}.}
\GSfield{Present status}{As of 10 November 2026: established in the stated scope; this passage poses no additional open question.}
\GSfield{Assessment}{literal statement: correct with missing hypotheses; publication match: exact; verification: read the result; development: none; confidence: high}
\end{GSitem}
\begin{GSitem}{1978-460.6}{E}{Abelian relations of a web}
\GSfield{Kind; source}{definition; \GSpage{P000838}}
\GSfield{The notes}{The first integrals are displayed as \(u_1,\ldots,u_n\), but the equation is \(\sum_{i=1}^d f_i(u_i)du_i=0\). The coefficients depend on the corresponding first integral.}
\GSfield{Statement}{An abelian relation of a regular holomorphic \(d\)-web is a \(d\)-tuple of closed one-forms \((\eta_i)\), each annihilating the tangent spaces of the \(i\)-th foliation, with \(\sum_i\eta_i=0\). Locally \(\eta_i=f_i(u_i)du_i\); equivalently \(\eta_i=z_i\,du_i\) with \(dz_i\wedge du_i=0\). This definition is invariant under changing the first integral of each foliation.}
\GSfield{Reconstruction}{The terminal subscript in the list of first integrals is corrected from \(n\) to \(d\). The definition is a tuple of forms; distinct coefficient tuples can sum to the same zero form and must not be identified.}
\GSfield{Fate}{Already used in \citet[Sec.~2]{beauville1980tissus}; precise equivalent equations are \citet[Sec.~1.1, Eqs.~(1)--(2)]{trepreau2006algebrisation}.}
\GSfield{Present status}{As of 10 November 2026: established in the stated scope; this passage poses no additional open question.}
\GSfield{Assessment}{literal statement: false as written; publication match: partial; verification: read the result; development: none; confidence: high}
\end{GSitem}
\begin{GSitem}{1978-460.7}{E}{Rank means dimension over the constants}
\GSfield{Kind; source}{definition; \GSpage{P000838}}
\GSfield{The notes}{Rank is the maximum number of linearly independent abelian equations satisfied by the web.}
\GSfield{Statement}{The rank is the complex dimension of the vector space of germs of abelian relations, with linear dependence over constant scalars. It is not the pointwise rank of the cotangent vectors or dimension over the ring of holomorphic functions. In the regular category this dimension is finite, bounded by the number in item 1978-460.9.}
\GSfield{Reconstruction}{The scalar field and the germ interpretation supply the intended meaning of linear independence.}
\GSfield{Fate}{\citet[Sec.~2]{beauville1980tissus} specifies independence over constants, before the April session.}
\GSfield{Present status}{As of 10 November 2026: established in the stated scope; this passage poses no additional open question.}
\GSfield{Assessment}{literal statement: correct with missing hypotheses; publication match: exact; verification: read the result; development: none; confidence: high}
\end{GSitem}
\begin{GSitem}{1978-460.8}{E}{The rank of an algebraic curve web is its genus}
\GSfield{Kind; source}{theorem; \GSpage{P000838}}
\GSfield{The notes}{Proposition 1: the rank of the web constructed from a curve of genus \(g\) equals \(g\).}
\GSfield{Statement}{For the regular dual web of a smooth irreducible nondegenerate complex projective curve \(C\), the trace map \(\omega\mapsto(p_i^*\omega)_{i=1}^d\) identifies \(H^0(C,\Omega_C^1)\) with its abelian-relation space. Abel's theorem gives \(\sum_i p_i^*\omega=0\), and the converse recovers every relation. Consequently the rank is \(g(C)\). For singular curves, a corresponding assertion uses the appropriate regular dualizing differentials and arithmetic genus; the genus of the normalization cannot simply be substituted.}
\GSfield{Reconstruction}{The smooth scope is the complete statement asserted here; the singular caveat explains why the notebook's unrestricted word ``genus'' is insufficient.}
\GSfield{Fate}{\citet[Sec.~2, pp.~105--106]{beauville1980tissus} states the trace identification. \citet[Sec.~1.2]{trepreau2006algebrisation} explicitly identifies the arithmetic-genus version.}
\GSfield{Present status}{As of 10 November 2026: established in the stated scope; this passage poses no additional open question.}
\GSfield{Assessment}{literal statement: correct with missing hypotheses; publication match: exact; verification: read the result; development: none; confidence: high}
\end{GSitem}
\begin{GSitem}{1978-460.9}{E}{The Chern rank bound and Castelnuovo number}
\GSfield{Kind; source}{theorem; \GSpage{P000838}}
\GSfield{The notes}{For an arbitrary \(d\)-web in dimension \(n\), rank \(\le\pi(n,d)\). The notes define \(\pi\) as the maximum genus of a degree-\(d\) curve in \(\mathbb P^n\) and give \(\pi=m(d-1-\tfrac12(m+1)(n-1))\), \(m=\lfloor(d-1)/(n-1)\rfloor\).}
\GSfield{Statement}{For a regular holomorphic codimension-one \(d\)-web in dimension \(n\ge2\), \[\operatorname{rank}\mathcal W\le\pi(n,d)=\sum_{q\ge1}\max\{d-1-q(n-1),0\}.\] This is exactly the notebook's integer formula, including the case where the last summand is zero. For \(d\ge n\), it is also the sharp Castelnuovo bound on the genus of smooth irreducible nondegenerate degree-\(d\) projective curves. The adjective nondegenerate is essential: plane curves embedded in higher projective spaces do not obey the claimed smaller genus bound.}
\GSfield{Reconstruction}{Codimension one, regular general position and nondegeneracy are supplied. The uncertain historical year ``193?'' is not silently made exact. For \(d<n\) the sum still defines the web bound, but no nondegenerate degree-\(d\) curve in \(\mathbb P^n\) is implied.}
\GSfield{Fate}{The bound is classical and is presented in \citet[Sec.~4]{beauville1980tissus}. A proof via independent jets and powers of conormal linear forms is given in \citet[Sec.~2.2, Eqs.~(6)--(8)]{trepreau2006algebrisation}.}
\GSfield{Present status}{As of 10 November 2026: established in the stated scope; this passage poses no additional open question.}
\GSfield{Assessment}{literal statement: correct with missing hypotheses; publication match: exact; verification: read the result; development: none; confidence: high}
\end{GSitem}
\begin{GSitem}{1978-460.10}{B}{Maximal rank implies algebraization: the later repair of the main theorem}
\GSfield{Kind; source}{theorem; \GSpage{P000839}}
\GSfield{The notes}{The main theorem says: if \(n\ge3\), \(d\ge2n\), and rank is \(\pi(n,d)\), the web can be linearized and extends to a web from \(X\subset\mathbb P^n\). The precise reference is Jahresbericht DMV 80 (1978), 13--110.}
\GSfield{Statement}{Let \(\mathcal W\) be a germ of a regular holomorphic codimension-one \(d\)-web on \(\mathbb C^n\), with every set of at most \(n\) defining conormals independent, \(n\ge3\), \(d\ge2n\), and \(\dim_{\mathbb C}\mathcal A(\mathcal W)=\pi(n,d)\). Then a local biholomorphism identifies \(\mathcal W\) with the web dual to a degree-\(d\) projective algebraic curve. In particular its leaves become hyperplane pieces. ``Extends'' means this local algebraic identification, not an assertion about the original web's arbitrary global domain.}
\GSfield{Reconstruction}{The analytic category, general position and local interpretation are restored. The theorem is now true, but its attribution as an already completely proved 1978 result requires correction.}
\GSfield{Fate}{The notebook identifies \citet{cherngriffiths1978abel} exactly. \citet[p.~78]{cherngriffiths1981corrections} acknowledges incorrect divisibility lemmas and a separate gap in Eq.~(4.34), retaining the theorem with an extra normality condition defined on p.~79. The unrestricted result is \citet[Theorem~1.1, preprint p.~3]{trepreau2006algebrisation}: first publicly posted 4 August 2005, published 2006. Theorem~2.4, preprint p.~8, even needs only \(2d-3n+1\) abelian relations with independent first jets. The cases \(n=3\) and \(d=2n\) were classical; the new general range is \(n\ge4,d\ge2n+1\). Beauville's February 1979 lecture repeats the then accepted claim and does not repair the later-discovered gap.}
\GSfield{Present status}{As of 10 November 2026: solved in this exact range by Tr'epreau. No open maximal-rank question is asserted here for the omitted planar or small-degree ranges.}
\GSfield{Assessment}{literal statement: correct with missing hypotheses; publication match: exact; verification: read the result; development: none; confidence: high; problem \hyperref[prob:GP-1978-21]{GP-1978-21}}
\end{GSitem}
\begin{GSitem}{1978-460.11}{E}{The projective-connection strategy and its historical limitation}
\GSfield{Kind; source}{method; \GSpage{P000839}}
\GSfield{The notes}{After the main theorem the notes say that a geometry of projective connections is being developed, with an underlining arrow.}
\GSfield{Statement}{This identifies the connection-based strategy for linearization in the Chern--Griffiths work. The notes specify no connection coefficients, curvature equations, or additional independent theorem. In the later proof, abelian relations construct a projective map whose image must be a surface; proving its differential has rank two is the decisive step.}
\GSfield{Reconstruction}{The description of the later rank-two argument is a literature explanation, not a reconstruction of an unrecorded notebook proof. The vague method line cannot validate the original proof's missing step.}
\GSfield{Fate}{Compare the contemporary exposition \citet[Sec.~6]{beauville1980tissus}, the explicit limitations in \citet[p.~78]{cherngriffiths1981corrections}, and \citet[Sec.~1.3, pp.~3--4]{trepreau2006algebrisation}, who separates the standard strategy from his direct calculation.}
\GSfield{Present status}{As of 10 November 2026: established in the stated scope; this passage poses no additional open question.}
\GSfield{Assessment}{literal statement: not determinable; publication match: partial; verification: read the result; development: none; confidence: medium}
\end{GSitem}
\begin{GSitem}{1978-460.12}{E}{A vertex tensor for an arrangement of lines}
\GSfield{Kind; source}{construction; \GSpage{P000839-P000840}; A. B. Zamolodchikov}
\GSfield{The notes}{A. Zamolodchikov introduces \(S_{ij}^{k\ell}(\lambda)\), colors \(1,\ldots,N\), and \(L\) lines with ordered angles \(\lambda_a\). Put a matrix at each crossing, sum over interior-edge colors and leave endpoint colors free. Pairing equal initial and final colors gives \(Z(\lambda_1,\ldots,\lambda_L)\).}
\GSfield{Statement}{Fix a finite oriented line arrangement without triple crossings, an \(N\)-dimensional vector space \(V\), and a matrix function \(R(u)\in\operatorname{End}(V\otimes V)\). At a crossing of lines \(a,b\), use \(R(\lambda_a-\lambda_b)\) in a fixed input/output convention. Contract all internal indices to obtain a boundary tensor, or a matrix on the endpoint state spaces when their orientations permit. Closing corresponding boundary indices and summing gives a trace partition function. This finite algebraic construction does not itself require positivity, convergence, or the Yang--Baxter identity.}
\GSfield{Reconstruction}{The notes' angle is treated as a spectral parameter; physical relativistic rapidity is a separate interpretation. Endpoint ordering and matrix convention are made explicit because they affect whether permutations are absorbed into \(R\).}
\GSfield{Fate}{The vertex formulation and three-line relation are standard by this session; \citet[Sec.~2, pp.~527--529]{zamolodchikovzamolodchikov1978on} gives the scattering counterpart (received 22 August 1977, revised 14 November 1977, published 27 February 1978). \citet[Secs.~10.1--10.2 and 11.1]{baxter1982exactly} gives a checked later systematic treatment of the lattice contraction.}
\GSfield{Present status}{As of 10 November 2026: established in the stated scope; this passage poses no additional open question.}
\GSfield{Assessment}{literal statement: correct with missing hypotheses; publication match: partial; verification: read the result; development: none; confidence: high}
\end{GSitem}
\begin{GSitem}{1978-460.13}{F}{The unspecified large-system asymptotic and boundary independence}
\GSfield{Kind; source}{problem; \GSpage{P000839-P000840}; A. B. Zamolodchikov}
\GSfield{The notes}{Problem 1 asks for the behavior as \(L\to\infty\), with \(\log Z\) proportional to the number of nodes. P000840 writes \(\log Z_i^j\sim L^2\), even without associativity, and says the leading asymptotic is independent of boundary indices.}
\GSfield{Statement}{As a statement about arbitrary crossing matrices, the claimed nonzero quadratic growth and unrestricted boundary independence are false. Take \(R(u)=I_{V\otimes V}\) and order endpoints by the lines they belong to. Every compatible fixed-boundary matrix entry is \(1\), incompatible entries are \(0\), and the trace is \(N^L\). Thus the compatible logarithm is zero, the other logarithm is undefined, and the trace logarithm is \(L\log N\), despite \(\binom L2\) crossings and the Yang--Baxter identity. A meaningful bulk limit needs a specified family of weights, normalization, rapidity distribution, thermodynamic sequence, and allowed boundary conditions.}
\GSfield{Reconstruction}{The counterexample checks the literal quantifiers, not the truth of a particular integrable model. The notes provide no exact asymptotic function or error estimate that could be indexed as a definite unresolved research problem.}
\GSfield{Fate}{\citet[Secs.~2.2, 10.17, 11.2 and 11.4]{baxter1982exactly} analyzes thermodynamic limits for specified models and qualifications on boundary vectors. Those results do not prove the unrestricted statement. Searches for the precise angle-dependent arbitrary-matrix claim produced no exact match; the identity example settles its literal reading.}
\GSfield{Present status}{As of 10 November 2026: the unrestricted assertion is false; the intended restricted asymptotic problem is not identifiable from these notes.}
\GSfield{Assessment}{literal statement: false as written; publication match: thematic only; verification: read the result; development: annotated; confidence: high}
\end{GSitem}
\begin{GSitem}{1978-460.14}{E}{The three-line move is the Yang--Baxter identity}
\GSfield{Kind; source}{theorem; \GSpage{P000840}; A. B. Zamolodchikov}
\GSfield{The notes}{\ru{Ассоциат.}: the tensor does not depend on the mutual disposition of the lines; it suffices to check a pictured equality for three lines.}
\GSfield{Statement}{In the unpermuted tensor-factor convention, the spectral-parameter relation is \[R_{12}(u)R_{13}(u+v)R_{23}(v)=R_{23}(v)R_{13}(u+v)R_{12}(u).\] This equality is precisely invariance under moving one pairwise crossing past the third line, with the same six external indices. Inserting the equality into a larger contraction preserves its boundary tensor. Hence any sequence of such local moves and exchanges of disjoint crossings leaves that tensor unchanged. Changing endpoints, crossing counts, or allowing a pair creation/deletion requires additional identities.}
\GSfield{Reconstruction}{The displayed equation fixes a convention for the notebook's unlabeled three-line diagram. The conclusion is limited to moves generated by this relation; no unqualified topological invariance is inferred.}
\GSfield{Fate}{\citet[Sec.~2, Eqs.~(2.5)--(2.6)]{zamolodchikovzamolodchikov1978on} derives the factorization equations from associativity; \citet[Sec.~11.1, Eq.~(11.1.6); Sec.~11.2]{baxter1982exactly} proves the lattice move by contraction.}
\GSfield{Present status}{As of 10 November 2026: established in the stated scope; this passage poses no additional open question.}
\GSfield{Assessment}{literal statement: correct with missing hypotheses; publication match: exact; verification: read the result; development: none; confidence: high}
\end{GSitem}
\begin{GSitem}{1978-460.15}{E}{Existence of nontrivial matrices satisfying the three-line identity}
\GSfield{Kind; source}{claim; \GSpage{P000840}; A. B. Zamolodchikov}
\GSfield{The notes}{Assertion 2 says such matrices exist, without displaying a solution at this point.}
\GSfield{Statement}{Nontrivial spectral-parameter solutions exist. For example, on \(V\otimes V\), with \(\dim V=N\ge2\), the rational family \(R(u)=uI+cP\), where \(P(v\otimes w)=w\otimes v\) and \(c\ne0\), satisfies the identity in item 1978-460.14. It is generically invertible and non-scalar. This algebraic existence does not by itself make its entries positive Boltzmann weights or a unitary physical scattering matrix.}
\GSfield{Reconstruction}{The rational example is supplied only as an explicit witness to the existential claim. It is not attributed to the notebook, which gives no formula or normalization for Assertion 2.}
\GSfield{Fate}{The speaker and A.~B.~Zamolodchikov had already given nontrivial \(O(N)\)-invariant relativistic solutions in \citet[Sec.~3]{zamolodchikovzamolodchikov1978on}; the rational permutation witness can also be checked directly in the group algebra of \(S_3\).}
\GSfield{Present status}{As of 10 November 2026: established in the stated scope; this passage poses no additional open question.}
\GSfield{Assessment}{literal statement: correct; publication match: partial; verification: read the result; development: annotated; confidence: high}
\end{GSitem}
\begin{GSitem}{1978-460.16}{F}{The proposed pairwise formula for the leading free energy}
\GSfield{Kind; source}{claim; \GSpage{P000840}; A. B. Zamolodchikov}
\GSfield{The notes}{Assertion 3 reads \(\log Z(\lambda_1,\ldots,\lambda_L)=\sum_{a>b=1}^{L}f(\lambda_a-\lambda_b)+\text{remainder}\), interpreted as a sum of pair contributions to the free energy.}
\GSfield{Statement}{The suggested leading term is a sum over unordered pairs \(1\le b<a\le L\). The source specifies neither \(R\), the function \(f\), the growth of the remainder, nor a limiting empirical distribution of the \(\lambda_a\). Without a remainder bound the displayed equality is not a mathematical asymptotic statement. Moreover, the Yang--Baxter equation alone does not fix \(f\): replacing \(R(u)\) by \(c(u)R(u)\) multiplies the network by \(\prod_{a>b}c(\lambda_a-\lambda_b)\) and shifts \(\log Z\) by its logarithm whenever defined.}
\GSfield{Reconstruction}{The scalar-normalization freedom is a direct obstruction to inferring a universal free-energy density. No specific solvable family is silently inserted.}
\GSfield{Fate}{\citet[Secs.~10.17, 11.2 and 11.4]{baxter1982exactly} obtains additive bulk free energies for specified inhomogeneous/eight-vertex models and cites the 1978 planar-lattice work on p.~281. This is a thematic match, not an exact identification of the notebook's unspecified claim. The 1978 original full text was not successfully retrieved.}
\GSfield{Present status}{As of 10 November 2026: the particular intended pairwise asymptotic remains unidentified; no precise open problem follows from the unspecified remainder.}
\GSfield{Assessment}{literal statement: not determinable; publication match: thematic only; verification: read the result; development: annotated; confidence: medium}
\end{GSitem}
\begin{GSitem}{1978-460.17}{E}{Factorized scattering and Boltzmann-weight interpretations}
\GSfield{Kind; source}{remark; \GSpage{P000840}; A. B. Zamolodchikov}
\GSfield{The notes}{Particles move along straight lines, preserving velocities but changing color at collisions; the whole tensor is an \(S\)-matrix. Alternatively it is a statistical sum, with weights described as exponentials of a potential and followed by \ru{вероятности?}.}
\GSfield{Statement}{A two-dimensional elastic factorized-scattering interpretation assigns a conserved rapidity to each particle and an internal-state space to each line; two-body matrices change internal states, and consistency of different collision orders requires the Yang--Baxter identity. Unitarity, crossing and appropriate analyticity are further physical conditions. A classical lattice interpretation instead assigns nonnegative weights \(w=e^{-\beta E}\) to local configurations and sums products. These weights are not probabilities until a normalization is supplied, and a complex scattering amplitude is not generally a Boltzmann weight.}
\GSfield{Reconstruction}{The two interpretations are separated, and ``exponential of the potential'' is normalized in the usual sign convention without treating the tentative probability gloss as a theorem.}
\GSfield{Fate}{The physical assumptions and their independent roles are explicit in \citet[Sec.~2, Eqs.~(2.1)--(2.6)]{zamolodchikovzamolodchikov1978on}.}
\GSfield{Present status}{As of 10 November 2026: established in the stated scope; this passage poses no additional open question.}
\GSfield{Assessment}{literal statement: correct with missing hypotheses; publication match: partial; verification: read the result; development: none; confidence: high}
\end{GSitem}
\begin{GSitem}{1978-460.18}{E}{Face spins and a nearest-neighbor partition function}
\GSfield{Kind; source}{construction; \GSpage{P000840}; A. B. Zamolodchikov}
\GSfield{The notes}{Spins lie in the polygons cut out by the lines and interact with nearest neighbors. The notes write \(H=\sum h(\sigma_A,\sigma_B)\) and \(Z=\sum_{\{\sigma\}}e^{-\beta H}\).}
\GSfield{Statement}{For a finite planar cell decomposition, assign a spin from a finite set to each face and an interaction \(h_e(\sigma_A,\sigma_B)\) to each edge separating adjacent faces. With specified exterior/boundary spins, \(H(\sigma)=\sum_e h_e(\sigma_A,\sigma_B)\) and \(Z=\sum_{\sigma\text{ allowed}}e^{-\beta H(\sigma)}\). For real finite energies and real \(\beta\), this is a finite positive partition function. A representation as local vertex tensors requires a stated encoding of the spin differences/products and the appropriate global constraints.}
\GSfield{Reconstruction}{Boundary conditions, edge counting and the face adjacency are supplied. The construction alone makes no integrability assertion.}
\GSfield{Fate}{The face-spin/vertex correspondence is exhibited in \citet[Secs.~10.3 and 11.5]{baxter1982exactly}; the Ising case below is checked directly from the notebook's formula.}
\GSfield{Present status}{As of 10 November 2026: established in the stated scope; this passage poses no additional open question.}
\GSfield{Assessment}{literal statement: correct with missing hypotheses; publication match: partial; verification: read the result; development: none; confidence: high}
\end{GSitem}
\begin{GSitem}{1978-460.19}{F}{The unqualified assertion about all exactly computable partition functions}
\GSfield{Kind; source}{claim; \GSpage{P000840}; A. B. Zamolodchikov}
\GSfield{The notes}{After the spin sum the notes add \ru{Точно счит. стат. суммы --- все такие!} (all exactly computable statistical sums are of this kind).}
\GSfield{Statement}{No class of models, meaning of exact computability, or precise meaning of ``of this kind'' is given. It could mean a tensor representation, factorized scattering, or a special integrable vertex model; these are different assertions. Consequently no universal classification theorem is identifiable.}
\GSfield{Reconstruction}{The emphatic aside is retained with its ambiguity. It is not converted into a conjecture that every solvable statistical model satisfies a specified Yang--Baxter equation.}
\GSfield{Fate}{The checked vertex-model literature, including \citet[Chs.~10--11]{baxter1982exactly}, supplies important examples but no theorem matching these undefined universal quantifiers. Searches located the surrounding integrable-model framework, not a statement with this scope.}
\GSfield{Present status}{As of 10 November 2026: the intended universal statement is not identifiable; there is no precise problem to classify as open.}
\GSfield{Assessment}{literal statement: not determinable; publication match: none found; verification: unresolved lead; development: none; confidence: high}
\end{GSitem}
\begin{GSitem}{1978-460.20}{E}{The displayed Ising vertex weight and its parity constraint}
\GSfield{Kind; source}{construction; \GSpage{P000840}; A. B. Zamolodchikov}
\GSfield{The notes}{For Ising, \(S_{ij}^{k\ell}=\tfrac12(ijk\ell+1)\exp[-\tfrac\beta2(i+j+k+\ell)]\). The diagram places spin variables in the faces.}
\GSfield{Statement}{Take \(i,j,k,\ell\in\{\pm1\}\) as successive products of the four face spins around a crossing. Their product is one. The prefactor is precisely the indicator of this even-parity condition; the exponential assigns half of the edge energy \(s_e\) at each incident vertex. Multiplying vertex weights therefore gives \(e^{-\beta\sum_e s_e}\) on interior edges, with boundary half-weights completed separately. On a simply connected connected face complex, the parity constraints recover face spins up to a simultaneous sign flip; on a torus they must be supplemented by trivial cycle products to recover periodic spins. Hence the spin partition function differs by the global twofold multiplicity when all spins, including a reference spin, are summed.}
\GSfield{Reconstruction}{The edge-product convention, boundary factors and topology are reconstructed. With the notebook's minus sign, the interaction is \(+\sum_e\sigma_A\sigma_B\); changing to the usual ferromagnetic \(-J\sum_e\sigma_A\sigma_B\) changes the exponential accordingly. No spectral parameter appears in this displayed weight, so it does not prove a Yang--Baxter family for arbitrary angles.}
\GSfield{Fate}{The parity and face-spin correspondence are discussed in \citet[Secs.~10.3 and 11.5, especially p.~287]{baxter1982exactly}. The exact normalization written here is verified by direct multiplication and exhaustive enumeration of the sixteen four-spin assignments in the supplied check code.}
\GSfield{Present status}{As of 10 November 2026: established in the stated scope; this passage poses no additional open question.}
\GSfield{Assessment}{literal statement: correct with missing hypotheses; publication match: partial; verification: read the result; development: annotated; confidence: high}
\end{GSitem}

\section{Developments}\label{dev:1978}
The subsections below are critical audits of recorded formulas and missing hypotheses. They are not Stage~2 developments of class C or D problems, and no novelty claim is attached to the classical results used in them. Each proof or computation supplied here is marked, and the independent checks and their limits are stated explicitly.

\subsection{The Gaussian-tail example: a finite remainder before surreal interpretation}
\label{audit:gaussian}
Item 1978-446.20 proposes an asymptotic application, but its literal formula is false. This is a correction to a recorded example, not a newly discovered open problem. The following elementary calculation supplies a complete check of the corrected expansion; its content is classical. The separate question of assigning a surreal value to an integral must be answered by a specified extension operator, for example within the precise scope of \citet[Thms.~78 and 85]{costinehrlich2024integration}.

\begin{proposition}[Gaussian tail with signed remainder]\ours{}
For every real \(x>0\) and integer \(N\ge1\), with \((-1)!!=1\),
\begin{align*}
I(x)&:=\int_x^\infty e^{-t^2}\,dt\\
&=\frac{e^{-x^2}}{2x}\sum_{n=0}^{N-1}
 \frac{(-1)^n(2n-1)!!}{2^n x^{2n}}+R_N(x),\\
R_N(x)&=\frac{(-1)^N(2N-1)!!}{2^N}
       \int_x^\infty t^{-2N}e^{-t^2}\,dt,
\end{align*}
and
\[
0<(-1)^N R_N(x)<
\frac{e^{-x^2}}{2x}\frac{(2N-1)!!}{2^N x^{2N}}.
\]
\end{proposition}
\begin{proof}
Let \(J_k(x)=\int_x^\infty t^{-2k}e^{-t^2}\,dt\). Integration by parts, using
\(e^{-t^2}dt=-d(e^{-t^2})/(2t)\), gives the exact identity
\[
J_k(x)=\frac{e^{-x^2}}{2x^{2k+1}}-\frac{2k+1}{2}J_{k+1}(x).
\]
All integrals converge for \(x>0\); the boundary term at infinity vanishes. Iterating from \(k=0\) through \(N-1\) proves the formula. Every \(J_k\) is positive, and the same identity with \(k=N\) gives
\(0<J_N(x)<e^{-x^2}/(2x^{2N+1})\). Multiplication by the positive coefficient gives the strict bound.
\end{proof}

\paragraph{Independent check.}
Writing \(I=e^{-x^2}F/(2x)\) in \(I'=-e^{-x^2}\) gives
\(F'-(2x+x^{-1})F=-2x\). Matching powers in \(F=\sum c_nx^{-2n}\) gives
\(c_0=1\), \(c_n=-(2n-1)c_{n-1}/2\), independently reproducing the coefficients above. The supplied script \texttt{code/check\_independent.py} verifies the truncated differential-equation residual exactly for eight orders and checks signed remainder bounds at four positive real arguments and five truncation orders with 80-digit arithmetic. These finite tests corroborate the proof; they are not used to infer its universal quantifiers.

At any fixed real \(x\), the terms of the infinite asymptotic series eventually increase, since their absolute successive ratio is \((2n+1)/(2x^2)\). At a positive infinite surreal \(x\), its real-coefficient power series in \(x^{-1}\) is summable in Conway's sense. These are different notions of summation. Neither observation makes the notebook's missing exponential prefactor, signs and powers correct. The modern resurgent-extension theory is a later implementation of the broad idea, with hypotheses not present in the notebook. A further research run would need to specify a new function class or extension axiom before there were a bounded open problem to investigate.

\subsection{Omnific fractions and two elementary boundary checks}
The strongest defensible reading of item 1978-446.22 is the known quotient theorem for omnific integers, not the false assertion that every surreal is rational over \(\mathbb Z\).

\begin{proposition}[Normal-form verification of the quotient theorem]\ours{}
Let \(\mathbf{Oz}\) have the standard normal-form characterization: all exponents are nonnegative and the constant coefficient, if present, is an ordinary integer. Then every surreal is a quotient of two members of \(\mathbf{Oz}\).
\end{proposition}
\begin{proof}
For nonzero \(x=\sum_{\alpha<\beta}r_\alpha\omega^{b_\alpha}\), the support is a set. By the simplest-separator theorem choose \(a>0\) with \(a>-b_\alpha\) for every \(\alpha<\beta\). All exponents of \(m=\omega^a x\) are strictly positive. The same holds for \(n=\omega^a\); hence \(m,n\in\mathbf{Oz}\), \(n\ne0\), and \(x=m/n\). For zero take \(0/1\). This proves the quotient assertion from the stated standard characterization; it does not claim a new proof of the equivalence of that characterization with Conway's simplicity definition.
\end{proof}

The normal-form multiplication is independently checked by cancellation of \(\omega^a\omega^{-a}=1\); shifting every exponent back by \(-a\) recovers exactly the original support and coefficients. This check requires set-sized support, which is why that hypothesis is explicit. The historical result belongs to \citet[Chap.~5]{conway1976numbers}; the normal-form theorem was read in \citet[Thm.~2.4]{kuhlmannmatusinski2015classes}.

Two shorter checks resolve additional ambiguities. First, the apparent integer criterion \(|x]-x|=1\) fails for the ordinary integer \(x=3\) and its earlier truncation \(x]=0\). A choice of a particular predecessor would be additional data absent from the notes. Second, the final exponential formula admits \(x=1\), the earlier option zero and \(n=1\), for which its reciprocal denominator is \(E_1(-1)=0\). The missing restrictions cannot be repaired by punctuation alone. These are local counterexamples to literal readings; neither is evidence of an undeveloped research problem.


\subsection{A strict audit of the projected-area formula}

This is a check of the literal minimum-area assertion in occurrence
1978-447.25, not a proposed solution of a newly identified open problem.
The source's phrase \ru{Для не слишком закрученных контуров} supplies no
defined contour class.  The unrestricted equality is false even arbitrarily
close to a planar circle.  The valid Hamilton--Jacobi identity in occurrence
1978-447.24 does not remove this obstruction.

For a closed oriented curve \(C\subset\mathbb R^4\), set
\(\sigma_{ij}=\oint_C x_i\,dx_j\), \(i<j\), regard \(\sigma\) as a
two-form with the independent-component Euclidean norm, and write
\(\sigma_\pm=(\sigma\pm *\sigma)/2\).  Our normalization is
\(A_\pm(C)=\sqrt2\,|\sigma_\pm|\).  Both profiles have value \(\pi\)
on a positively oriented unit circle in the first coordinate plane.
The closed-form/comass inequality giving area lower bounds is classical;
see \citet[Introduction, pp.~48--49]{harveylawson1982calibrated}.

\begin{proposition}[A strict gap near a circle]\ours{}
For \(\epsilon\ne0\), let
\[
C_\epsilon(t)=(\cos t,\sin t,\epsilon\cos 2t,0),\qquad 0\le t\le2\pi.
\]
This is a smooth embedded closed curve, and
\(A_+(C_\epsilon)=A_-(C_\epsilon)=\pi\).
Every oriented Lipschitz spanning disk has area at least
\[
\pi\sqrt{1+\epsilon^2}>\pi.
\]
The same bound holds for compactly supported integral two-currents with
boundary the multiplicity-one oriented curve \(C_\epsilon\), with area
interpreted as mass.  Such disks exist and have finite area.
\end{proposition}

\begin{proof}
Projection to the first two coordinates is the standard unit circle, so
\(C_\epsilon\) is embedded.  Orthogonality of the elementary trigonometric
functions gives \(\sigma_{12}=\pi\) and all other projected areas zero.
Thus \(\sigma_\pm=(\pi/2)(dx_1\wedge dx_2\pm dx_3\wedge dx_4)\),
and both normalized norms equal \(\pi\).

Write \((x,y,z,w)\) for the four coordinates.  Nearest-point projection
\(P\) onto the closed convex cylinder
\(\{x^2+y^2\le1\}\times\mathbb R^2\) is 1-Lipschitz and fixes the
boundary curve.  Composing a spanning disk with \(P\) cannot increase
its two-dimensional area.  Equivalently, push-forward under \(P\)
does not increase the mass of a compactly supported two-current and
preserves its prescribed boundary.

On the cylinder consider the exact two-form
\[
\Omega_\delta=dx\wedge dy-\delta\,d(xy)\wedge dz
             =d\bigl(x\,dy-\delta xy\,dz\bigr).
\]
It is supported in the three-dimensional \((x,y,z)\) directions.  Its
comass is the norm of its three-dimensional Hodge-dual vector:
\[
\|\Omega_\delta\|_{\mathrm{comass}}
 =\sqrt{1+\delta^2(x^2+y^2)}\le\sqrt{1+\delta^2}.
\]
The same comass holds in \(\mathbb R^4\), since adding a direction in
which the form vanishes cannot increase its value on a unit simple
two-vector.  On the boundary,
\[
\oint_{C_\epsilon}x\,dy=\pi,\qquad
\oint_{C_\epsilon}xy\,dz
 =\int_0^{2\pi}\frac{\sin2t}{2}(-2\epsilon\sin2t)\,dt
 =-\epsilon\pi.
\]
Stokes' theorem and the comass bound therefore give, for every projected
spanning disk or current \(S\),
\[
\operatorname{Area}(S)\ge
\frac{\left|\int_S\Omega_\delta\right|}{\sqrt{1+\delta^2}}
=\frac{\pi|1+\delta\epsilon|}{\sqrt{1+\delta^2}}.
\]
Choose \(\delta=\epsilon\).  Projection did not increase area, so the
same lower bound holds for the original filling.

For completeness, the graph
\((x,y)\mapsto(x,y,\epsilon(x^2-y^2),0)\), \(x^2+y^2\le1\),
is a smooth finite-area spanning disk.  No minimizing-disk existence,
regularity, or equality-case theorem is needed for the strict gap.
\end{proof}

\begin{remark}[Scope and independent check]\ours{}
For each fixed finite \(k\), \(C_\epsilon\) tends to the planar circle
in \(C^k\) as \(\epsilon\to0\).  Consequently, an unspecified condition
of being sufficiently close to that circle cannot by itself validate the
projected-area equality.  The flux was checked independently by expanding
\(xy\) and \(dz\) in trigonometric functions; the comass was checked
both as a three-dimensional Hodge-dual norm and as the nonzero singular
value of the corresponding skew matrix.  These checks agree.
This argument does not contradict the source's first-order identity:
the derivative of \(A_\pm\) has a complementary-plane component even
when the boundary lies in \(\mathbb R^3\).
\end{remark}


\subsection{Critical algebraic audit: the printed spin-one density}

This section audits the literal assertion in P000818 (14) and P000819
(16), items 1978-454.17, 1978-454.19 and 1978-454.22. It is a completed
critical check, not a Stage~2 development of a class C or D problem.
The notes give
\[
 h=X^2-2X,\qquad X=\sum_{a=1}^3S^a\otimes S^a,
 \qquad [S^a,S^b]=\varepsilon_{abc}S^c,
\]
and specify the real three-dimensional standard/adjoint representation of
\(SO(3)\). They then ask for parameters fitting this density into the
factorization scheme. We test precisely the analytic, regular,
difference-form construction used by the surrounding transfer-matrix
argument. A coefficient or generator-normalization error changes this
test and cannot be silently repaired.

The necessary condition used below is commonly called the Reshetikhin
condition; see \citet[Sec.~2, eqs.~(1)--(6)]{zhang2025bootstrap}, first posted
in April 2025, in the version dated 6 February 2026 checked here. We give
the coefficient argument completely, so its application does not depend
on an unchecked attribution of sufficiency.

\begin{lemma}[The cubic necessary condition; derivation checked here]\ours{}
Let \(V\) be a finite-dimensional complex vector space and let
\(\check R(u)\in\operatorname{End}(V\otimes V)\) be analytic at zero,
with \(\check R(0)=I\), satisfying the braided difference equation
\[
 \check R_{12}(u)\check R_{23}(u+v)\check R_{12}(v)
 =\check R_{23}(v)\check R_{12}(u+v)\check R_{23}(u).
\]
Put \(h=\check R'(0)\). There exists
\(Q\in\operatorname{End}(V\otimes V)\) such that
\[
 [h_{12}+h_{23},[h_{12},h_{23}]]=Q_{23}-Q_{12}.
\]
\end{lemma}
\begin{proof}
Write \(r_2=\check R''(0)\). Substituting \(v=-u\) and comparing
the coefficient of \(u^2\) gives
\((r_2-h^2)_{12}=(r_2-h^2)_{23}\).
For an operator T, equality \(T_{12}=T_{23}\) implies that T is scalar:
taking the partial trace over the third tensor factor first gives
\(T=I\otimes U\), and substitution gives
\(U\otimes I=I\otimes U\), forcing U scalar. Thus
\(r_2=h^2+cI\). Multiplying \(\check R(u)\) by
\(1-cu^2/2\) preserves the equation, regularity and h, and makes
\(r_2=h^2\). We may therefore write
\[
 \check R(u)=I+uh+\tfrac12u^2h^2+\tfrac16u^3r_3+O(u^4).
\]
Set \(A=h_{12}\), \(B=h_{23}\), \(C=(r_3)_{12}\),
\(D=(r_3)_{23}\). Twice the coefficient of \(u^2v\) on the left is
\[
 D+2AB^2+B^2A+A^2B+A^3+2ABA,
\]
and on the right it is
\[
 C+BA^2+2A^2B+B^3+AB^2+2BAB.
\]
Their equality, using
\[
 [A+B,[A,B]]=A^2B-2ABA+BA^2+2BAB-B^2A-AB^2,
\]
gives the asserted identity with \(Q=r_3-h^3\).
\end{proof}

\begin{proposition}[An exact obstruction in the recorded convention]\oursComputed{}
Let
\[
 S^1=\begin{pmatrix}0&0&0\\0&0&-1\\0&1&0\end{pmatrix},\quad
 S^2=\begin{pmatrix}0&0&1\\0&0&0\\-1&0&0\end{pmatrix},\quad
 S^3=\begin{pmatrix}0&-1&0\\1&0&0\\0&0&0\end{pmatrix}.
\]
These satisfy the recorded commutator relations. For
\(h_t=X^2+tX\) on \(V=\mathbb C^3\), in the basis
\(e_0,e_1,e_2\),
\[
 \langle001|[h_{t,12}+h_{t,23},[h_{t,12},h_{t,23}]]|212\rangle
 =t(t-1)(t+1).
\]
Consequently the printed density, \(t=-2\), fails the lemma.
\end{proposition}
\begin{proof}
Let P be the flip and let
\(K(e_i\otimes e_j)=\delta_{ij}\Omega\),
\(\Omega=\sum_{a=0}^2e_a\otimes e_a\).
Multiplying the displayed matrices gives \(X=K-P\).
Since \(P^2=I\), \(K^2=3K\), and \(PK=KP=K\),
\[
 X^2=I+K,\qquad h_t=I+(t+1)K-tP.
\]
The identity term has no effect on the double commutator. Write
\(a=t+1\), \(c=-t\), and use the two-site operator \(aK+cP\).
Its action is explicitly
\[
 (aK+cP)|ij\rangle=a\delta_{ij}\sum_b|bb\rangle+c|ji\rangle.
\]
For A and B acting on sites 12 and 23 respectively, successive use of
this formula gives the complete six-term calculation
\[
\begin{array}{c|rrrrrr}
 W&A^2B&ABA&BA^2&BAB&B^2A&AB^2\\\hline
 \langle001|W|212\rangle&ac(3a+2c)&a^2c&0&0&0&0.
\end{array}
\]
Inserting the coefficients \(1,-2,1,2,-1,-1\) from the preceding
commutator expansion yields \(ac(a+2c)=t(t-1)(t+1)\).
For any two-site Q,
\[
 \langle001|Q_{23}|212\rangle=0,
 \qquad \langle001|Q_{12}|212\rangle=0,
\]
because the first, respectively third, site would have to be unchanged.
At \(t=-2\) the computed entry is \(-6\), a contradiction.
\end{proof}

\begin{corollary}[Scope of the negative result]\ours{}
No analytic regular difference-form braided R-matrix has first derivative
\(\alpha(X^2-2X)+\beta I\) for \(\alpha\ne0\). In particular, no family
of such matrices can specialize to the printed density with a nonzero
Hamiltonian scale.
\end{corollary}
\begin{proof}
Adding a scalar leaves the double commutator unchanged, while scaling h
by \(\alpha\) scales it by \(\alpha^3\). Its forbidden matrix entry is
therefore \(-6\alpha^3\ne0\).
\end{proof}

\paragraph{Independent checks and the published neighbouring model.}
The supplied \texttt{check\_spin1.py} and output use exact SymPy matrices.
They check the three commutators, the spectra of X and h, and the
polynomial witness. A second calculation in that script uses
Levi--Civita entries and explicit index contractions, without Kronecker
matrix products, and again obtains \(-6\). The independently authored
\texttt{check\_independent.py} constructs permutation and contraction
tensors directly and obtains the same polynomial
\(t(t-1)(t+1)\). As controls, the full necessary-condition residual is
zero at \(t=-1,0,1\), and nonzero at \(t=-2,2\); these passing controls
are not used to assert sufficiency. The six-term calculation above
provides an independently readable finite proof of the witness.

Hermitian generators are \(J^a=iS^a\), so
\(\sum_aJ^a\otimes J^a=-X\). The familiar isotropic spin-one density
\(J\cdot J-(J\cdot J)^2\) is therefore \(-X-X^2\), an overall multiple
of \(h_1\), not of \(h_{-2}\). The anisotropic construction of
\citet{zamolodchikovfateev1980spin} was published in 1980; the explicit
Hamiltonian and this isotropic limit are read in
\citet[Sec.~5, eqs.~(52)--(55)]{bytsko2003xxz}, first circulated as a
preprint in December 2001 and published in 2003. This is the precise
reason for calling that literature a partial match to the notebook's
question. The original coefficient intended by the speaker remains
unidentified. The calculation excludes the stated regular
difference-form construction; it does not classify other possible
meanings of integrability.

There is no residual class C or D problem produced by this audit. A
further bounded source investigation could search an independent record
of the February 1979 talk for a normalization or coefficient correction.
That would resolve the historical identification, not a surviving
anisotropic spin-one existence problem.


\subsection{The contact sign and two vertex-model checks}
These are critical annotations to expository or underspecified items, not developments of new C/D problems.

\begin{proposition}[The literal contact formula fails Jacobi]\oursComputed{}
In one coordinate pair put \(X=\partial_x\), \(Y_\varepsilon=\partial_y+\varepsilon x\partial_z\), \(Z=\partial_z\), and
\[
 B_\varepsilon(f,g)=fZg-gZf+Xf\,Y_\varepsilon g-Xg\,Y_\varepsilon f.
\]
The notebook's choice \(\varepsilon=1\) is not a Lie bracket. In fact, for the convention
\(J(f,g,h)=B(f,B(g,h))+B(g,B(h,f))+B(h,B(f,g))\),
\[
 J_\varepsilon(f,g,h)=(1+\varepsilon)
 \det\begin{pmatrix}f_x&f_y&f_z\\g_x&g_y&g_z\\h_x&h_y&h_z\end{pmatrix}.
\]
\end{proposition}
\begin{proof}
The three brackets at \(\varepsilon=1\) are
\(B(x,y)=1\), \(B(y,z)=y\), \(B(z,x)=-2x\). Thus the three nested terms are \(1,2,-1\), with sum \(2\). This already refutes Jacobi. For the displayed general identity, expand
\[
 B_\varepsilon(f,g)=f_xg_y-f_yg_x+
 \varepsilon x(f_xg_z-f_zg_x)+fg_z-gf_z.
\]
In the cyclic nested expression the second-derivative terms cancel in pairs, as do the undifferentiated-function terms. The six remaining first-derivative products are exactly the determinant above, each with coefficient \(1+\varepsilon\). The supplied symbolic script performs this expansion for arbitrary differentiable functions, not just a finite set of examples.
\end{proof}

\begin{proposition}[A consistent contact convention]\ours{}
On \(\mathbb R^{2n+1}\), the same bracket formula with
\(Y_i=\partial_{y_i}-x_i\partial_z\) is the contact bracket for
\(\alpha=dz+\sum_i x_i\,dy_i\).
\end{proposition}
\begin{proof}
The Reeb field is \(Z=\partial_z\). Define
\[
 V_f=fZ+\sum_i\bigl((X_if)Y_i-(Y_if)X_i\bigr).
\]
Direct substitution gives \(\alpha(V_f)=f\) and
\(\iota_{V_f}d\alpha=(Zf)\alpha-df\), hence
\(\mathcal L_{V_f}\alpha=(Zf)\alpha\). Thus \(V_f\) is a contact vector field, uniquely determined by its \(\alpha\)-component \(f\). The commutator of contact fields is contact, and
\[
 \alpha([V_f,V_g])=V_f(g)-(Zf)g=B_{-1}(f,g).
\]
Uniqueness implies \([V_f,V_g]=V_{B_{-1}(f,g)}\). Injectivity of \(f\mapsto V_f\) transfers the vector-field Jacobi identity to the corrected bracket. This proves the result in every dimension, independently of the three-variable symbolic expansion.
\end{proof}

\paragraph{Second checks.}
The file \texttt{check\_contact.py} verifies the generic-function identity, the coordinate triple, and two unrelated polynomial triples. Its output records the nesting convention explicitly; the alternative convention \([[f,g],h]+\text{cyclic}\) changes the sign. A separate reviewer re-derived the coordinate and polynomial checks and obtained the same obstruction. This is a faithful mathematical error in P000835, where both plus signs are visible, rather than a transcription error. The contact theory itself is classical \citep[Sec.~4]{kirillov1976local}.

\begin{example}[Why the free-energy assertion needs hypotheses]\ours{}
Let \(V\) have dimension \(N\ge2\), let \(L\) generic oriented lines each cross every other line once, and take \(R(u)=I_{V\otimes V}\). Order both sets of endpoints by their lines. Every local crossing preserves both colors, so the boundary operator is \(I_{V^{\otimes L}}\). Compatible matrix entries are \(1\), incompatible ones are \(0\), and the trace is \(N^L\). Therefore \(\log\operatorname{tr}Z=L\log N\), even though there are \(\binom L2\) vertices and the three-line identity holds. This proves that a universal nonzero \(L^2\) law and unrestricted index independence do not follow from the notebook's stated assumptions. More generally, replacing \(R(u)\) by \(c(u)R(u)\) preserves the Yang--Baxter relation but multiplies the tensor by \(\prod_{a>b}c(\lambda_a-\lambda_b)\). Thus even a pairwise leading term requires a specified normalization. These observations do not refute the thermodynamic formulas for particular models in \citet[Secs.~10.17 and 11.4]{baxter1982exactly}.
\end{example}

\begin{proposition}[An explicit witness for the three-line identity]\ours{}
For any vector space \(V\), the family \(R(u)=uI+cP\), where \(P(v\otimes w)=w\otimes v\), satisfies the additive-parameter Yang--Baxter identity used in item 1978-460.14.
\end{proposition}
\begin{proof}
Set \(A=P_{12}\), \(B=P_{13}\), \(C=P_{23}\). In the permutation action on \(V^{\otimes3}\),
\(AB=BC=CA\), \(AC=BA=CB\), and \(ABC=CBA\).
Expand \((uI+cA)((u+v)I+cB)(vI+cC)\) and the product in the reverse order. The coefficients of \(c^0,c^1,c^3\) agree directly. At order \(c^2\), the left side is
\(vAB+(u+v)AC+uBC=(u+v)(AB+AC)\), and the right side is
\(uCB+(u+v)CA+vBA=(u+v)(AC+AB)\). This proves the identity without restriction on the dimension. For \(\dim V\ge2\) and \(c\ne0\) the family is non-scalar and is invertible for generic \(u\).
\end{proof}
The example is a supplied verification, not a claim that this formula appears in the notebook. The speaker's already published nontrivial relativistic solutions are in \citet[Sec.~3]{zamolodchikovzamolodchikov1978on}. The independent check \texttt{check\_vertex.py} verifies the symbolic \(8\times8\) identity for \(V=\mathbb C^2\) and the permutation relations separately.

\begin{proposition}[What the displayed Ising factor encodes]\ours{}
For four successive face spins \(a,b,c,d\in\{\pm1\}\), set
\((i,j,k,\ell)=(ab,bc,cd,da)\). Then \(ijk\ell=1\), and every parity-even quadruple has exactly two such face-spin lifts, related by simultaneous reversal. The notebook's factor
\[
 \frac{1+ijk\ell}{2}\exp\!\left[-\frac\beta2(i+j+k+\ell)\right]
\]
therefore enforces the face-spin compatibility at a vertex and assigns half the energy of each incident bond there.
\end{proposition}
\begin{proof}
Multiplying the four bond products gives \(a^2b^2c^2d^2=1\). Conversely choose \(a\) arbitrarily, put \(b=ai\), \(c=aij\), and \(d=aijk\). The last equation \(da=\ell\) is equivalent to the parity condition, so exactly the two choices of \(a\) give the two lifts. On a connected simply connected face complex, the analogous reconstruction along paths is independent of the path because elementary cycle products are one. Every internal bond meets two vertices, so its two half-exponents add to \(-\beta s_e\). Boundary bonds require the missing half-factor, and a torus requires additional trivial products on noncontractible cycles. Thus no unqualified finite-volume equality with every periodic vertex sector has been assumed.
\end{proof}
The supplied code independently enumerates all sixteen local face-spin assignments and verifies eight admissible bond tuples, each with multiplicity two. The same correspondence is classical \citep[Secs.~10.3 and 11.5]{baxter1982exactly}. This local encoding does not supply an angle-dependent integrable family or prove the notebook's sweeping asymptotic claims.

\section{Problem index}\label{problems:1978}
Earlier-year occurrence identifiers are retained as text; current-year identifiers link to their entries.

\subsection{GP-1978-01: A spectral representation for the nonabelian Toda chain}\label{prob:GP-1978-01}
For the matrix chain \(\dot A_s=B_s-B_{s-1}\), \(\dot B_s=B_sA_{s+1}-A_sB_s\), construct a spectral auxiliary linear system whose compatibility is exactly these equations. Finite periodic chains give a concrete algebraic setting; infinite-chain spectral statements require additional decay and domain conditions.
\par\noindent\textit{Occurrences:} \hyperref[item:1978-447.4]{1978-447.4}.
\par\noindent\textit{Status (2026-11):} As of 10 November 2026: resolved. The difference-operator pair \(\mathcal L=E^{-1}+A+BE\), \(\mathcal M=BE\) verifies the compatibility directly. Subsequent nonabelian Toda solution theory is traced to Krichever 1981 by \citet[\S1 and \S9]{takebezabrodin2025multicomponent}. The earlier scalar periodic-Toda involution problem GP-1975-05 concerns a different assertion and is not merged with this matrix Lax-representation question.

\subsection{GP-1978-10: Anisotropic regular factorization for the spin-one density: literal and corrected versions}\label{prob:GP-1978-10}
Find a parameter family of analytic regular difference-form R-matrices whose Hamiltonian specializes to the rotationally invariant spin-one density printed on P000818. The literal density is \(X^2-2X\), with \([S^a,S^b]=\varepsilon_{abc}S^c\) and \(X=\sum_aS^a\otimes S^a\). Separately, the identified corrected historical problem concerns anisotropic deformations of \(J\cdot J-(J\cdot J)^2\) for Hermitian spin-one generators.
\par\noindent\textit{Occurrences:} \hyperref[item:1978-454.17]{1978-454.17}, \hyperref[item:1978-454.19]{1978-454.19}, \hyperref[item:1978-454.22]{1978-454.22}.
\par\noindent\textit{Status (2026-11):} As of 10 November 2026: the literal regular difference-form extension is excluded by the exact cubic obstruction proved in this chapter. The nearby correctly normalized anisotropic model was published by \citet{zamolodchikovfateev1980spin} in 1980; its explicit relation to the isotropic density is checked in \citet[Sec.~5]{bytsko2003xxz}. This is a closed negative literal result plus a solved corrected construction, not a manufactured open remainder.

\subsection{GP-1978-11: The finite periodic XYZ spectral problem}\label{prob:GP-1978-11}
For real \(J_1,J_2,J_3\) and \(N\ge2\), diagonalize \(H=\sum_{n=1}^N\sum_{a=1}^3J_a\sigma_n^a\sigma_{n+1}^a\) on \((\mathbb C^2)^{\otimes N}\), with periodic boundaries, by an integrable description uniform in N. Multiplicities and exceptional parameter values require explicit handling.
\par\noindent\textit{Occurrences:} \hyperref[item:1978-453.8]{1978-453.8}.
\par\noindent\textit{Status (2026-11):} As of 10 November 2026: an established exactly solvable family, already called solved by Baxter in the notebook. The 1971 announcement is verified \citep{baxter1971eight}; the 1973 eigenvector article is bibliographically identified but its direct URL remained inaccessible \citep{baxter1973eigenvectors}. No all-exceptional-parameter completeness theorem is asserted from inaccessible text.

\subsection{GP-1978-20: Hexagonal closure and parallelizability of planar three-webs}\label{prob:GP-1978-20}
For a regular holomorphic planar three-web, is local equivalence to the three parallel foliations \(x,y,x+y=\mathrm{const}\) equivalent to closure of every sufficiently small web hexagon? Yes: Thomsen's theorem; equivalently the web curvature vanishes.
\par\noindent\textit{Occurrences:} \hyperref[item:1978-460.4]{1978-460.4}.
\par\noindent\textit{Status (2026-11):} As of 10 November 2026: solved before the seminar (Thomsen, 1927), as recalled in \citet[Sec.~1]{beauville1980tissus}. General three-web linearizability is a different question.

\subsection{GP-1978-21: Algebraization of maximal-rank codimension-one webs in dimension at least three}\label{prob:GP-1978-21}
For a germ of a regular holomorphic codimension-one \(d\)-web on \(\mathbb C^n\), with every at-most-\(n\) subset of conormals independent, \(n\ge3\), \(d\ge2n\), and rank \(\pi(n,d)\), must the web be locally biholomorphic to the dual web of a projective degree-\(d\) curve?
\par\noindent\textit{Occurrences:} \hyperref[item:1978-460.10]{1978-460.10}.
\par\noindent\textit{Status (2026-11):} As of 10 November 2026: yes. The 1978 Chern--Griffiths proof had two defects; their 1981 correction added normality. Tr\'epreau, Theorem~1.1, removes that extra condition: preprint first posted 4 August 2005, published 2006. His Theorem~2.4 gives a stronger first-jet criterion. The notebook's theorem is thus solved, not an open problem.

\section{Corrections to the transcription and translation}\label{corr:1978}
The evidence distinguishes a transcription correction from a mathematical emendation of wording faithfully preserved in the notebook.
\subsection*{P000786 (both; provenance): Speaker and catalogue initials}
\noindent\textit{Transcription reads.} The catalogue gives V.~S.~Kruskal; the manuscript only has \ru{Крускал}.\par
\noindent\textit{Critical reading.} M.~D.~Kruskal is the probable identification; retain the catalogue initials as discrepant metadata.\par
\noindent\textit{Evidence.} The scan contains no initials. \citet[Introduction]{gonshor1986surreal} identifies M.~D.~Kruskal as the speaker who introduced this theory in January 1977. This supports identification, but is not an independent attendance record for Moscow.\par
\subsection*{P000786-P000787 (both; provenance): Earlier-index overlap}
\noindent\textit{Transcription reads.} 1977-445.1--4 are indexed as earlier-year occurrences.\par
\noindent\textit{Critical reading.} They overlap the opening of record 446 dated 30 October 1978.\par
\noindent\textit{Evidence.} The supplied item titles match sign sequences, simplicity, the separator theorem and recursive addition on these scans. The date is visible on P000786; a shared source must not count as recurrence.\par
\subsection*{P000787 (both; ambiguity): Associative-law sign}
\noindent\textit{Transcription reads.} \((x+y)+z=x+(y\mathbin{\text{[?]}}z)\).\par
\noindent\textit{Critical reading.} The mathematical reading is \((x+y)+z=x+(y+z)\).\par
\noindent\textit{Evidence.} The scan has an overwritten plus-like stroke. The neighboring text asserts ordinary additive laws; this establishes the intended formula, not a certain letter-by-letter reading.\par
\subsection*{P000788 (both; ambiguity): Product inequality}
\noindent\textit{Transcription reads.} \(ru+sv\mathrel{\text{[?]}}rv+su\).\par
\noindent\textit{Critical reading.} \(ru+sv>rv+su\), if \(r<s\) and \(u<v\).\par
\noindent\textit{Evidence.} The scan supports a retraced greater-than sign; independently the difference is \((s-r)(v-u)>0\).\par
\subsection*{P000789 (both; ambiguity): Historical surname}
\noindent\textit{Transcription reads.} Hessenburg\par
\noindent\textit{Critical reading.} Hessenberg (retain the manuscript spelling in diplomatic text).\par
\noindent\textit{Evidence.} The scan does resemble Hessenburg. The uniquely matching natural ordinal operations identify Gerhard Hessenberg; this is an editorial name normalization, not a claim that the transcription miscopied the final letters.\par
\subsection*{P000790 (both; faithful error): Gaussian-tail example}
\noindent\textit{Transcription reads.} \(\int_x^{\text{class-long }\uparrow}e^{-\xi^2}d\xi=\sum n!/x^n\).\par
\noindent\textit{Critical reading.} For real \(x\to+\infty\), \(\int_x^\infty e^{-t^2}dt\sim e^{-x^2}(2x)^{-1}\sum_{n\ge0}(-1)^n(2n-1)!!/(2^n x^{2n})\).\par
\noindent\textit{Evidence.} Both the exponential and the factorial series are visible on the scan. Integration by parts gives the stated correction with a signed remainder, proved and checked in the critical annotations. A surreal integral needs an additional definition.\par
\subsection*{P000790 (both; faithful error): Integer criterion}
\noindent\textit{Transcription reads.} \(|x]-x|=1\) is the definition of an integer.\par
\noindent\textit{Critical reading.} The intended integers are omnific integers. Equality of distance to an arbitrary truncation is not their defining property.\par
\noindent\textit{Evidence.} For the ordinary integer \(3\), the zero truncation has distance three. The scan contains no quantifier or choice of a special predecessor, so the literal formula cannot be certified as an equivalent definition.\par
\subsection*{P000790 (both; ambiguity): Power-series coefficient qualifier}
\noindent\textit{Transcription reads.} \(x\) infinitesimal, \(n\) real; \(\sum x^n a_n\) converges for all \(a_n\).\par
\noindent\textit{Critical reading.} Intended theorem: \(n\in\mathbb N_0\), \(a_n\in\mathbb R\); summability is in the normal-form sense.\par
\noindent\textit{Evidence.} The scan does not securely identify the abbreviated qualifier. The precise theorem is \citet[Prop.~4]{costinehrlich2024integration}. Unrestricted surreal coefficients fail for \(a_n=x^{-n}\).\par
\subsection*{P000790 (both; faithful error): Recursive exponential}
\noindent\textit{Transcription reads.} The left Taylor sum begins at \(m=1\); right sums have unrestricted denominators; left/right options are not separated.\par
\noindent\textit{Critical reading.} Record the formula as incomplete. The modern recursion uses constant Taylor terms, separate option families and positive-denominator restrictions.\par
\noindent\textit{Evidence.} The omissions are visible in the scan. At \(x=1\), earlier option \(0\), the displayed denominator at \(n=1\) is zero. See \citet[Chap.~10]{gonshor1986surreal}.\par
\subsection*{P000791 (both; ambiguity): Toda kinetic term}
\noindent\textit{Transcription reads.} \(\dot g^{-1}\dot g\), with no factor \(1/2\)\par
\noindent\textit{Critical reading.} Distinguish derivative of \(g^{-1}\) from inverse of \(\dot g\); a normalized compatible Lagrangian is \(\operatorname{Tr}(\tfrac12 A^2-B)\)\par
\noindent\textit{Evidence.} The scan preserves the ambiguous notation. Differentiating \(g^{-1}\) gives \(-g^{-1}\dot g\,g^{-1}\); direct variation shows the factor needed for the displayed equations.\par
\subsection*{P000791 (en; translation): word about gauge structure}
\noindent\textit{Transcription reads.} connectedness\par
\noindent\textit{Critical reading.} connection\par
\noindent\textit{Evidence.} Russian \ru{связность} and matrix-link/continuum context identify a connection; the word is marked provisionally/cancelled in the diplomatic notes.\par
\subsection*{P000791-P000798 (both; provenance): record 447 boundaries}
\noindent\textit{Transcription reads.} Catalogue range starts P789 and overlaps following talk\par
\noindent\textit{Critical reading.} Polyakov begins at his heading on P791 and ends before Schwarz on P799\par
\noindent\textit{Evidence.} Both transcripts and source headings on P791 and P799; catalogue range is an inclusive locator, not exclusive ownership.\par
\subsection*{P000792 (both; faithful error): derivative of transfer product}
\noindent\textit{Transcription reads.} \(\frac{d}{dt}\prod_s L_s=0\)\par
\noindent\textit{Critical reading.} Finite product derivative is \(T M_{\rm right}-M_{\rm left}T\); zero requires endpoint normalization\par
\noindent\textit{Evidence.} Direct telescoping of the displayed zero-curvature equation. Infinite product also requires convergence; periodicity conserves conjugacy invariants, not generally \(T\).\par
\subsection*{P000792 (both; faithful error): power-series expansion}
\noindent\textit{Transcription reads.} Truncated expansion written as equality\par
\noindent\textit{Critical reading.} Insert the appropriate \(O(\lambda^2)\) or \(O(\lambda^3)\) remainder, or state formal coefficient extraction\par
\noindent\textit{Evidence.} Scan and both transcriptions show truncated terms; expansion of the given recursion verifies the first two coefficients.\par
\subsection*{P000793 (both; faithful error): expanded backward plaquette}
\noindent\textit{Transcription reads.} \(B^{-1}_{x,\beta}\) at the first backward edge\par
\noindent\textit{Critical reading.} \(B^{-1}_{x-\beta,\beta}\)\par
\noindent\textit{Evidence.} The scanned square diagram requires an edge from \(x\) to \(x-\beta\). The literal index fails to close the four-link path.\par
\subsection*{P000793 (both; faithful error): plaquette exponential expansion}
\noindent\textit{Transcription reads.} Finite polynomial equals exponential\par
\noindent\textit{Critical reading.} Add the higher-order remainder\par
\noindent\textit{Evidence.} Ordinary exponential series; the scan contains no remainder.\par
\subsection*{P000795 (both; ambiguity): inverting the cross-product relation}
\noindent\textit{Transcription reads.} \(F=t\times d\chi\) without tangent normalization\par
\noindent\textit{Critical reading.} \(F=(t\times d\chi)/|t|^2\) for transverse \(F\), or impose \(|t|=1\)\par
\noindent\textit{Evidence.} Vector triple-product calculation from \(d\chi=-t\times F\); the source does not state unit-speed parametrization.\par
\subsection*{P000796 (both; faithful error): flat closed-loop holonomy}
\noindent\textit{Transcription reads.} Flatness implies every closed-loop ordered exponential equals 1\par
\noindent\textit{Critical reading.} Restrict to contractible loops, or add trivial global monodromy\par
\noindent\textit{Evidence.} Flat connections can have nontrivial holonomy on noncontractible loops; the notebook states no topology.\par
\subsection*{P000796 (both; faithful error): massive Green asymptotic}
\noindent\textit{Transcription reads.} \(G\sim\exp(-mr)\)\par
\noindent\textit{Critical reading.} \(G\sim c_d r^{-(d-1)/2}\exp(-mr)\), or state only its logarithmic decay rate\par
\noindent\textit{Evidence.} Modified-Bessel asymptotic of the Euclidean massive fundamental solution; exponential rate is correct, literal equivalence misses the prefactor.\par
\subsection*{P000796 (both; ambiguity): spectral denominator}
\noindent\textit{Transcription reads.} \(1-z^2\) with no signature stated\par
\noindent\textit{Critical reading.} Specify Lorentzian \({*}^2=1\) convention; Euclidean \({*}^2=-1\) changes the algebraic denominator\par
\noindent\textit{Evidence.} The displayed inversion depends on the square of the two-dimensional Hodge operator.\par
\subsection*{P000797 (both; ambiguity): Hodge-star convention}
\noindent\textit{Transcription reads.} \(\varepsilon_{\alpha\beta\gamma\delta}\pi_{\gamma\delta}\) without stated factor\par
\noindent\textit{Critical reading.} Use \(\tfrac12\varepsilon_{\alpha\beta\gamma\delta}\pi_{\gamma\delta}\) when both antisymmetric indices are summed\par
\noindent\textit{Evidence.} Standard full-index Hodge convention; a sum only over independent \(\gamma<\delta\) would remove the \(1/2\). Convention must be explicit.\par
\subsection*{P000797 (both; faithful error): contracted antisymmetric derivatives}
\noindent\textit{Transcription reads.} \(\pi_{\alpha\beta}\pi_{\beta\gamma}\) proportional to \(\delta_{\alpha\beta}\)\par
\noindent\textit{Critical reading.} \(\delta_{\alpha\gamma}\)\par
\noindent\textit{Evidence.} The free indices on the left are \(\alpha,\gamma\). Source and diplomatic transcription preserve the mismatch.\par
\subsection*{P000797-P000798 (both; faithful error): minimum-area assertion}
\noindent\textit{Transcription reads.} \(A_{\min}\) equals the self-dual projected-area norm for contours not too twisted\par
\noindent\textit{Critical reading.} The normalized projected-area expression is a lower bound; equality requires a common calibration. The quoted contour class is undefined.\par
\noindent\textit{Evidence.} Explicit arbitrarily small nonplanar perturbations of a circle have projected value \(\pi\) but area at least \(\pi\sqrt{1+\epsilon^2}\); see audit.tex. The separate Hamilton--Jacobi identity is valid.\par
\subsection*{P000799-P000803 (both; provenance): record 448 boundaries}
\noindent\textit{Transcription reads.} Catalogue overlap with adjacent records\par
\noindent\textit{Critical reading.} Schwarz meeting starts on P799 and ends before dated P804 meeting\par
\noindent\textit{Evidence.} P799 begins with the palindromes and Krichever remark, then Schwarz; P804 begins 27.11.78.\par
\subsection*{P000800 (both; ambiguity): quotient determinant}
\noindent\textit{Transcription reads.} \(\det(TT^*)\) for the orbit tangent map \(T:\operatorname{Lie}(G)\to T_xM\)\par
\noindent\textit{Critical reading.} Use the determinant on the nonzero orbit image, or \(\det(T^*T)\) on the injective Lie-algebra domain, with consistent volume conventions\par
\noindent\textit{Evidence.} The dimensions of the rectangular map and its kernel make the unrestricted determinant singular; the notebook omits the nonzero-subspace convention.\par
\subsection*{P000801 (both; ambiguity): instanton dimension}
\noindent\textit{Transcription reads.} \(8q-3\) instantons\par
\noindent\textit{Critical reading.} An \((8q-3)\)-dimensional unframed moduli family for \(SU(2)\) charge \(q\) on \(S^4\)\par
\noindent\textit{Evidence.} Atiyah--Hitchin--Singer primary abstract gives dimension, not a finite count. Framing adds gauge-orientation dimensions.\par
\subsection*{P000801 (both; ambiguity): rational-map extremals}
\noindent\textit{Transcription reads.} All extremals are rational \(w(z)\)\par
\noindent\textit{Critical reading.} In the compactified finite-action sphere-to-sphere setting, nonconstant solutions are holomorphic or antiholomorphic; specify the degree sign\par
\noindent\textit{Evidence.} The displayed holomorphic formula describes the positive-charge minimizing sector. Reversing orientation gives the other sign; statements on general domains need boundary conditions.\par
\subsection*{P000802 (both; faithful error): heat integral for determinant}
\noindent\textit{Transcription reads.} Positive heat-trace integral equals log determinant\par
\noindent\textit{Critical reading.} The regularized log determinant uses a minus sign, removes the kernel, and specifies finite-part subtraction\par
\noindent\textit{Evidence.} Fateev--Frolov--Schwarz equation (2.10), zeta derivative definition, and scalar Frullani identity agree. Determinants of \(A^*A\) versus \(|A|\) also differ by \(1/2\).\par
\subsection*{P000803 (both; faithful error): heat-trace variation}
\noindent\textit{Transcription reads.} Derivative of heat trace written as a trace difference without \(t\) derivative\par
\noindent\textit{Critical reading.} Restore the \(t\) derivative and factor produced by Duhamel differentiation before integrating in \(t\)\par
\noindent\textit{Evidence.} At finite-dimensional level \(\delta\operatorname{Tr}\exp(-tL)=-t\operatorname{Tr}(\delta L\exp(-tL))\); the notebook abbreviation drops this essential operation.\par
\subsection*{P000803 (both; faithful error): instanton-measure product}
\noindent\textit{Transcription reads.} Denominator contains a free \(j\)\par
\noindent\textit{Critical reading.} Product over all \(i,j\) for unlike-species pairs\par
\noindent\textit{Evidence.} Fateev--Frolov--Schwarz equations (1.9),(3.10) and \(q\)-instanton zero/pole symmetry. The \(c\) factor \(|c^2|\) is already equal to \(|c|^2\) and is not an error.\par
\subsection*{P000804 (both; faithful error): equation (3)}
\noindent\textit{Transcription reads.} \(U_\xi-V_\eta+[U,V]=0\)\par
\noindent\textit{Critical reading.} \(U_\eta-V_\xi+[U,V]=0\)\par
\noindent\textit{Evidence.} Differentiate \(\Psi_\xi=U\Psi\) and \(\Psi_\eta=V\Psi\). The scan confirms the reversed derivative labels, so this is not an OCR correction.\par
\subsection*{P000804-P000807 (both; provenance): record 449 speaker and boundaries}
\noindent\textit{Transcription reads.} Catalogue names only Zakharov and range continues beyond his contribution\par
\noindent\textit{Critical reading.} Zakharov P804-upper P806; Lipatov P806-P807, with Zakharov intervention at end\par
\noindent\textit{Evidence.} P804 explicitly names Lev Nikolaevich Lipatov; P806 has a second underlined Lipatov heading and ZhETF 72,411,1977 citation; P808 starts Guichardet.\par
\subsection*{P000805 (both; faithful error): equation (10), sums and final equality}
\noindent\textit{Transcription reads.} Sums use \(m\) in bounds and \(n\) in factors; final interaction indices are interchanged; equals sign breaks the expression\par
\noindent\textit{Critical reading.} Use sums \(n=1,\ldots,N_1\), \(m=1,\ldots,N_2\) and interaction \(U_nV_m/(\lambda_n-\mu_m)\); interpret last line as continuation\par
\noindent\textit{Evidence.} Both transcriptions flag the same source indices. The corrected gauge-fixed action is verified by compactly supported variation.\par
\subsection*{P000805 (both; faithful error): equation (4), residue commutators}
\noindent\textit{Transcription reads.} \(U_{n\eta}=[U_n,A_n]\), \(V_{m\xi}=[V_m,B_m]\)\par
\noindent\textit{Critical reading.} \(U_{n\eta}=[A_n,U_n]\), \(V_{m\xi}=[B_m,V_m]\)\par
\noindent\textit{Evidence.} Residue extraction from the corrected compatibility equation with \(A_n=V(\lambda_n)\), \(B_m=U(\mu_m)\). The frame equations (9) use this corrected convention.\par
\subsection*{P000805 (both; faithful error): equation (7), second Jordan representative}
\noindent\textit{Transcription reads.} \(V_0^{(n)}\)\par
\noindent\textit{Critical reading.} \(V_n^{(0)}\)\par
\noindent\textit{Evidence.} The preceding equation (6) and the parallel \(U\) formula fix the intended representative; the scan confirms reversed placement in the notebook.\par
\subsection*{P000806 (both; ambiguity): mass pole}
\noindent\textit{Transcription reads.} Physical mass squared is the pole while the free denominator is \(p^2+m^2\)\par
\noindent\textit{Critical reading.} Euclidean pole is at \(p_E^2=-m_{\rm phys}^2\); positive \(m_{\rm phys}^2\) refers to a Minkowski convention\par
\noindent\textit{Evidence.} The source switches between Lorentzian and Euclidean notation without explaining analytic continuation.\par
\subsection*{P000806-P000807 (both; ambiguity): amputated vertex and invariant charge}
\noindent\textit{Transcription reads.} \(\Gamma\) defined by amputation, then multiplied by \(\lambda_0\) again in invariant charge\par
\noindent\textit{Critical reading.} Either normalize \(\Gamma\) to \(1\) at tree level before multiplying by \(\lambda_0\), or omit that separate factor\par
\noindent\textit{Evidence.} Tree-level coupling count and the preceding vertex definition. The source supplies no consistent convention.\par
\subsection*{P000807 (both; faithful error): Landau limit}
\noindent\textit{Transcription reads.} \(\lambda_\mu\) tends to infinity at fixed \(\lambda_0>0\) as \(\Lambda\) tends to infinity\par
\noindent\textit{Critical reading.} For the displayed plus denominator it tends to zero\par
\noindent\textit{Evidence.} Source PNG explicitly confirms both signs; algebra of the displayed expression is decisive. Holding renormalized coupling fixed is a different limit.\par
\subsection*{P000807 (both; ambiguity): beta derivative convention}
\noindent\textit{Transcription reads.} \(d\lambda_\mu/d\log(\Lambda/\mu^2)\) called the series function without held-fixed variables\par
\noindent\textit{Critical reading.} Specify fixed bare coupling/cutoff convention; its sign differs from differentiation with respect to physical scale\par
\noindent\textit{Evidence.} Differentiating the displayed leading-log relation gives a negative coefficient at fixed \(\lambda_0\). No unseen coefficient is corrected.\par
\subsection*{P000808-P000811 (both; provenance): record 450 boundaries}
\noindent\textit{Transcription reads.} Catalogue includes a broader page range\par
\noindent\textit{Critical reading.} Guichardet begins P808 and ends P811; P812 already belongs to 1979 material\par
\noindent\textit{Evidence.} Dated Guichardet heading P808; semidirect-product setup ends P811; following heading P812.\par
\subsection*{P000809-P000811 (both; provenance): previous\_index.md, 1977-451.1--10}
\noindent\textit{Transcription reads.} Guichardet material assigned to academic year 1977 and record 451\par
\noindent\textit{Critical reading.} Treat as misassigned occurrences of the present 1978-450 material, not a genuine previous-year recurrence\par
\noindent\textit{Evidence.} Same infinitesimal-character lemma, relative \(H^1\), affine-group example, contradictory \(SO\) fragment, \(SL_2\) lead and final semidirect-product setup; the present scan date is 4.12.78.\par
\subsection*{P000810 (both; ambiguity): additive-group example and following sentence}
\noindent\textit{Transcription reads.} No corresponding group and Lie-algebra representations after an explicit \(\mathbb R\) unipotent example\par
\noindent\textit{Critical reading.} The displayed nilpotent Lie-algebra representation integrates to \(\mathbb R\); a compact-period obstruction arises for a circle quotient\par
\noindent\textit{Evidence.} Matrix exponentiation checks the \(\mathbb R\) example. The circle interpretation is an explanatory reconstruction, not a corrected reading of the source.\par
\subsection*{P000810 (both; faithful error): relative and absolute \(H^1\)}
\noindent\textit{Transcription reads.} Absolute \(H^1\) is relative \(H^1\); later relative \(H^1\) is an element of absolute \(H^1\)\par
\noindent\textit{Critical reading.} Natural degree-one injection, whose image is the kernel of restriction to \(H^1(\mathfrak h,M)\)\par
\noindent\textit{Evidence.} Relative cocycles vanish on \(\mathfrak h\); relative coboundaries are \(du\) with \(\mathfrak h\)-invariant \(u\). The source has the literal element-of sign.\par
\subsection*{P000810-P000812 (both; provenance): Year-heading source range}
\noindent\textit{Transcription reads.} Record 451 links pages containing Guichardet and Belavin mathematics.\par
\noindent\textit{Critical reading.} Retain 451 only as a calendar-year heading; assign actual mathematics by handwritten speaker/date headings.\par
\noindent\textit{Evidence.} The new dated session starts on P000812; the supplied previous-index items 1977-451.1--13 cover overlapping source content, not a distinct 1977 lecture.\par
\subsection*{P000811 (both; ambiguity): orthogonal-group fragment}
\noindent\textit{Transcription reads.} For nonsemisimple groups, \(SO_0(m,n)\); no nonunitary representations with nonzero cohomology\par
\noindent\textit{Critical reading.} Preserve as an unresolved contradictory statement; do not silently substitute unitary or semisimple\par
\noindent\textit{Evidence.} PNG and Russian both contain the conflicting heading and abbreviation. No dimensions or representations identify a theorem.\par
\subsection*{P000812 (both; provenance): Boundary between reading programme and Belavin}
\noindent\textit{Transcription reads.} Overlapping catalogue pages for records 452 and 453\par
\noindent\textit{Critical reading.} The reading programme above the separate underlined Belavin heading is record 452; the mathematical talk from that heading is record 453.\par
\noindent\textit{Evidence.} Both transcript and scan show the separate heading; do not manufacture mathematical assertions from book titles.\par
\subsection*{P000812 (both; provenance): Prior-index alias}
\noindent\textit{Transcription reads.} 1977-451.11--13 treated as earlier-year occurrences\par
\noindent\textit{Critical reading.} Same source segments: the ADHM and gauge-averaging paragraphs belong to 12 February 1979, record 453.\par
\noindent\textit{Evidence.} The same pages and wording are present here. Prior-year IDs are preserved as provenance aliases, not evidence of historical recurrence.\par
\subsection*{P000812 (both; ambiguity): Quaternionic matrix action}
\noindent\textit{Transcription reads.} \(M(x)=xA+B\), \(M\mapsto GMV\); groups unstated\par
\noindent\textit{Critical reading.} In the row convention take \(G\in GL(N,\mathbb R)\), \(V\in Sp(N+1)\), left multiplication \(xA\), and nondegenerate data.\par
\noindent\textit{Evidence.} ADHM, pp.~185--187. The identity \(G(xA+B)V=x(GAV)+GBV\) follows from real G and associativity; no commuting x through V is assumed. Full rank is omitted from the notebook.\par
\subsection*{P000812 (both; faithful error): Reading list, Feynman coauthor}
\noindent\textit{Transcription reads.} Feynman with Gibbs / \ru{Гиббсом}\par
\noindent\textit{Critical reading.} Feynman and Hibbs / \ru{Хиббсом}\par
\noindent\textit{Evidence.} The scan reads \ru{Гиббсом}; the 1965 book and its emended-edition editor identify Albert Hibbs. This is a notebook name error.\par
\subsection*{P000813 (both; faithful error): Equation (2), Pauli commutator}
\noindent\textit{Transcription reads.} \([\sigma^a,\sigma^b]=\varepsilon_{abc}\sigma^c\)\par
\noindent\textit{Critical reading.} \([\sigma^a,\sigma^b]=2i\varepsilon_{abc}\sigma^c\)\par
\noindent\textit{Evidence.} The scan displays the standard Pauli matrices; their direct product gives \(2i\). The omission is faithful in both languages.\par
\subsection*{P000813 (both; faithful error): Equations (3)--(4), Jordan--Wigner normalization}
\noindent\textit{Transcription reads.} No factor \(1/2\) in the fermion operator, but canonical anticommutator 1\par
\noindent\textit{Critical reading.} Insert \(1/2\) in each fermion operator, or change the anticommutator to 4.\par
\noindent\textit{Evidence.} Direct multiplication of the displayed Pauli matrices; periodic fermions also require a parity-sector boundary term.\par
\subsection*{P000814 (both; ambiguity): Equation (10), continuum quartic interaction}
\noindent\textit{Transcription reads.} \(\psi_1^+\psi_2\psi_1^+\psi_2^+\)\par
\noindent\textit{Critical reading.} A regularized Thirring current-current term; exact intended order/indices unidentified\par
\noindent\textit{Evidence.} The repeated creation component and missing regularization do not specify the current-current operator. No guessed unique emendation is made.\par
\subsection*{P000814 (both; faithful error): Equation (12), bosonic CCR}
\noindent\textit{Transcription reads.} \([\dot\varphi(x),\varphi(y)]=\delta(x-y)\)\par
\noindent\textit{Critical reading.} \([\Pi(x),\varphi(y)]=-i\delta(x-y)\) for canonical Hermitian fields.\par
\noindent\textit{Evidence.} Canonical commutator and Hermitian-adjoint consistency; the notebook has no imaginary factor.\par
\subsection*{P000814 (both; faithful error): Equation (7), time evolution}
\noindent\textit{Transcription reads.} Commutator order/sign as printed\par
\noindent\textit{Critical reading.} For \(U(t)=e^{-itH}\), use \(\dot O=i[H,O]\).\par
\noindent\textit{Evidence.} This correction states the time-evolution convention explicitly; no convention is supplied that would justify silently switching the sign.\par
\subsection*{P000814 (both; faithful error): Equation (8), lattice fermion evolution}
\noindent\textit{Transcription reads.} Linear equation without \(\Delta\) or adjoint terms\par
\noindent\textit{Critical reading.} The generic transverse-Ising equation includes the transverse field and creation/annihilation mixing\par
\noindent\textit{Evidence.} Direct commutation with displayed Hamiltonian (6); no unique corrected scalar equation follows without a component convention.\par
\subsection*{P000814 (both; ambiguity): Equations (13)--(14), vertex operator}
\noindent\textit{Transcription reads.} Unqualified exponential and \(\alpha\beta=4\pi\)\par
\noindent\textit{Critical reading.} A regulated normally ordered vertex/disorder operator with explicit imaginary factors and component/Klein conventions\par
\noindent\textit{Evidence.} The notebook formula omits conventions essential to its claimed anticommutation. The exact parameter relation cannot be certified without them.\par
\subsection*{P000814-P000815 (both; faithful error): Exchange-generator index range}
\noindent\textit{Transcription reads.} \(1\le i<M\)\par
\noindent\textit{Critical reading.} \(1\le i\le M\)\par
\noindent\textit{Evidence.} The following \(M=2\) example uses both \(A_1\) and \(A_2\). The scan confirms that the strict upper bound is a notebook inconsistency.\par
\subsection*{P000815 (both; faithful error): All solutions assertion}
\noindent\textit{Transcription reads.} All solutions up to a scalar factor\par
\noindent\textit{Critical reading.} Only the restricted eight-vertex ansatz and its qualified generic solutions\par
\noindent\textit{Evidence.} An unrestricted classification of all \(4\times4\) Yang--Baxter matrices does not follow from four specified weights.\par
\subsection*{P000815 (both; ambiguity): Elliptic weights and spectral origin}
\noindent\textit{Transcription reads.} \(a=\operatorname{sn}(\theta+\eta)\), \(b=\operatorname{sn}(\theta-\eta)\)\par
\noindent\textit{Critical reading.} Regular parameter \(u=\theta-\eta\), with \(\eta_0=2\eta\); use the corresponding addition law.\par
\noindent\textit{Evidence.} The weights themselves can be shifted to a regular Baxter parametrization, but the shift cannot be ignored in the difference/additive functional equation.\par
\subsection*{P000816 (both; faithful error): Last factor of the triple-reordering chain}
\noindent\textit{Transcription reads.} \(A_i(\theta)\)\par
\noindent\textit{Critical reading.} \(A_{i_2}(\theta_1)\)\par
\noindent\textit{Evidence.} The previous line ends with that contracted index and rapidity. The scan and both transcriptions agree on the dropped subscripts.\par
\subsection*{P000816 (both; ambiguity): Physical parameter \(\theta\)}
\noindent\textit{Transcription reads.} \(\theta_i\) are velocities\par
\noindent\textit{Critical reading.} For relativistic difference-form amplitudes, \(\theta_i\) are rapidities and \(v_i=\tanh\theta_i\).\par
\noindent\textit{Evidence.} The difference form is Lorentz invariant for rapidity, not ordinary velocity.\par
\subsection*{P000817 (both; ambiguity): Equation (2), missing last argument}
\noindent\textit{Transcription reads.} The final S has no argument\par
\noindent\textit{Critical reading.} Supply the spectral argument required by the fixed leg/permutation convention; use the invariant Yang--Baxter form in the item\par
\noindent\textit{Evidence.} Component indices also require consistent contraction. The scan omits the argument; no isolated guess is substituted for the complete identity.\par
\subsection*{P000817 (both; faithful error): Equation (6) with initial condition (7)}
\noindent\textit{Transcription reads.} \(a(u)a(u+v)d(v)=d(u+v)a(v)+b(u)b(u+v)d(v)\)\par
\noindent\textit{Critical reading.} Literal equation inconsistent with a regular elliptic family; exact intended component equation unidentified\par
\noindent\textit{Evidence.} Set \(v=0\) and use \(a(0)=1\), \(b(0)=d(0)=0\) from \(R(0)=P\): the equation forces \(d(u)=0\).\par
\subsection*{P000817-P000818 (both; ambiguity): Equation (8) and omitted constraints}
\noindent\textit{Transcription reads.} Four ODEs, two omitted equations and unprinted algebraic constraints\par
\noindent\textit{Critical reading.} An incomplete differential list cannot certify the original functional Yang--Baxter equation\par
\noindent\textit{Evidence.} The preceding sample functional equation already fails the zero-parameter test. No omitted equation is reconstructed from memory.\par
\subsection*{P000818 (both; faithful error): After equation (12)}
\noindent\textit{Transcription reads.} \(I_2\) contains three spins\par
\noindent\textit{Critical reading.} Local three-site charge comes from the appropriate logarithmic derivative, not the raw Taylor coefficient\par
\noindent\textit{Evidence.} The second logarithmic derivative subtracts the square of the first; raw coefficients retain disconnected products.\par
\subsection*{P000818-P000819 (both; faithful error): Equations (14) and (16), spin-one normalization}
\noindent\textit{Transcription reads.} \(h=X^2-2X\); \([S^a,S^b]=\varepsilon_{abc}S^c\)\par
\noindent\textit{Critical reading.} The literal regular difference-form R-matrix assertion is false; intended coefficient unidentified\par
\noindent\textit{Evidence.} For real skew generators, \(X=K-P\). The cubic-condition entry \(\langle001|D|212\rangle=-6\), whereas every \(Q_{23}-Q_{12}\) has zero there. Exact code and an independent proof are supplied.\par
\subsection*{P000819 (both; faithful error): Equation (18), internal product indices}
\noindent\textit{Transcription reads.} \(S_{ij_1}S_{j_2j_3}\cdots S_{j_Nk}\)\par
\noindent\textit{Critical reading.} Use one uninterrupted internal-index chain, or the ordered monodromy product\par
\noindent\textit{Evidence.} The scan explicitly jumps from \(j_1\) to \(j_2\). A consistent product has successive internal indices and terminal k.\par
\subsection*{P000819 (both; faithful error): General nearest-neighbour Hamiltonian claim}
\noindent\textit{Transcription reads.} Always has a nontrivial integral because S second derivative reconstructs\par
\noindent\textit{Critical reading.} Second-order extension is insufficient; a cubic necessary obstruction remains\par
\noindent\textit{Evidence.} Complete expansion and explicit spin-one counterexample in the critical algebraic audit.\par
\subsection*{P000820 (both; ambiguity): Historical line after Onsager 1944}
\noindent\textit{Transcription reads.} Onsager 1965 introduced a commuting transfer matrix\par
\noindent\textit{Critical reading.} Retain as an unverified historical attribution; no definite corrected year inferred\par
\noindent\textit{Evidence.} The verb is itself uncertain; Onsager 1944 and Baxter 1971 are checked, but this 1965 priority assertion was not identified.\par
\subsection*{P000820 (both; ambiguity): Last sentence}
\noindent\textit{Transcription reads.} \(T(J_1,J_2)\) commutes \ru{пр}\ldots\par
\noindent\textit{Critical reading.} Unfinished conditional assertion; no condition supplied\par
\noindent\textit{Evidence.} The scan stops after \ru{пр}. Do not fill in a commuting curve or spectral parametrization as source text.\par
\subsection*{P000821 (both; faithful error): Equation (1), real homogeneous quotient}
\noindent\textit{Transcription reads.} \(\widehat M=SU(2,2)/SO(3,1)\)\par
\noindent\textit{Critical reading.} \(\widehat M=SU(2,2)/P\), where P is the eleven-dimensional null-two-plane parabolic stabilizer.\par
\noindent\textit{Evidence.} The literal quotient has real dimension nine. The stabilizer Lie algebra is \((\mathfrak{so}(3,1)\oplus\mathbb R)\ltimes\mathbb R^{3,1}\), dimension eleven.\par
\subsection*{P000821 (both; faithful error): Equation (2), conformal groups}
\noindent\textit{Transcription reads.} \(SU(2,2)=SO(4,2)\)\par
\noindent\textit{Critical reading.} \(SU(2,2)\cong\operatorname{Spin}_0(4,2)\), a double cover of \(SO_0(4,2)\).\par
\noindent\textit{Evidence.} Penrose--MacCallum, Section 1.1, gives the covering relations; the projective effective action has further central identifications.\par
\subsection*{P000821 (both; ambiguity): Theorem 1, manifold of null twistor lines}
\noindent\textit{Transcription reads.} The manifold of all lines in R is Minkowski space\par
\noindent\textit{Critical reading.} It is conformally compactified Minkowski space, isomorphic to \(U(2)\); affine Minkowski space is a chart.\par
\noindent\textit{Evidence.} The graph formula on the same page gives compact \(U(2)\). The opening abbreviation remains uncertain.\par
\subsection*{P000822 (both; faithful error): Equation (6), Yang--Mills derivative}
\noindent\textit{Transcription reads.} \(+[F_{\mu\nu},A_\mu]\)\par
\noindent\textit{Critical reading.} \(+[A^\mu,F_{\mu\nu}]\), with metric contraction\par
\noindent\textit{Evidence.} With the explicitly printed \(F=dA+A\wedge A\), the derivative is \(D=d+[A,\cdot]\). The scan confirms the reversed order, hence a faithful error.\par
\subsection*{P000822 (both; ambiguity): Equation (8), dual space and dimension glyph}
\noindent\textit{Transcription reads.} \(\mathbb P^3\times\mathbb P^3\); \(\dim_{\mathbb C}G=5\)\par
\noindent\textit{Critical reading.} \(\mathbb P^3\times\mathbb P^{3*}\); \(\dim_{\mathbb C}Q=5\)\par
\noindent\textit{Evidence.} Hyperplanes form the dual projective space. P823 explicitly restores the star. The handwritten dimension variable may already be Q, so its letter is an ambiguity.\par
\subsection*{P000822 (both; faithful error): Tentative stabilizer}
\noindent\textit{Transcription reads.} \(GL(2,\mathbb C)\times\mathbb R^3\) (?)\par
\noindent\textit{Critical reading.} The parabolic stabilizer has Lie algebra \((\mathfrak{so}(3,1)\oplus\mathbb R)\ltimes\mathbb R^{3,1}\).\par
\noindent\textit{Evidence.} The guessed direct product has wrong Lie-group structure; retain the notebook question mark.\par
\subsection*{P000822-P000823 (both; faithful error): All connections/all bundles correspondence}
\noindent\textit{Transcription reads.} All bundles on Q\par
\noindent\textit{Critical reading.} Holomorphic bundles trivial on each incidence quadric \(Q_x\cong\mathbb P^1\times\mathbb P^1\), on admissible local domains.\par
\noindent\textit{Evidence.} P823 explicitly supplies this condition. Topological or Chern-class triviality alone is insufficient.\par
\subsection*{P000823 (both; ambiguity): Canonical trivialization along a line}
\noindent\textit{Transcription reads.} Canonically trivialized because curvature vanishes and \(\pi_1=0\)\par
\noindent\textit{Critical reading.} Canonically identified with the horizontal-section space, without a chosen basis; local curves require simply connected intersections\par
\noindent\textit{Evidence.} A holomorphic two-form vanishes on a complex curve. The full projective-line argument must be adjusted to the open domain stated on P822.\par
\subsection*{P000823-P000829 (both; provenance): Misassigned prior-year IDs}
\noindent\textit{Transcription reads.} 1977-457.1--19 as one earlier occurrence\par
\noindent\textit{Critical reading.} P823 is Beilinson in 1979; P824--P825/P829 are Fuchs in 1979; unavailable intermediate pages are not inferred.\par
\noindent\textit{Evidence.} Prior-index names match exact available source segments but academic-year assignment conflicts with dated headings.\par
\subsection*{P000824 (both; provenance): Dated new lecture inside catalogue record 456 page range}
\noindent\textit{Transcription reads.} Absorbed into Beilinson catalogue range P822--P825\par
\noindent\textit{Critical reading.} Separate Fuchs lecture, 5 March 1979; source-established added record 457\par
\noindent\textit{Evidence.} The scan visibly reads 5/III-79 and \ru{Д. Б. Фукс}. The attached fourteen-record list omits it.\par
\subsection*{P000824 (both; faithful error): Four-family classification}
\noindent\textit{Transcription reads.} Lower bound and transitivity imply four series\par
\noindent\textit{Critical reading.} Additional simplicity/irreducibility/finite-piece and generation hypotheses required; general transitive versions include extensions\par
\noindent\textit{Evidence.} Kac 1968, Theorems 4--5 and explicit conventions. Merely transitive bounded-below examples need not be simple or infinite-dimensional.\par
\subsection*{P000824-P000825 (both; faithful error): Recover everything from \(\mathfrak g_{-1},\mathfrak g_0\)}
\noindent\textit{Transcription reads.} Every graded algebra uniquely recovered\par
\noindent\textit{Critical reading.} The tensor formula gives the maximal Cartan prolongation; an arbitrary transitive depth-one algebra only embeds in it\par
\noindent\textit{Evidence.} The successive inclusion signs in the notes become equality only when choosing maximal prolongation. A faithful \(\mathfrak g_0\)-action is necessary.\par
\subsection*{P000829 (both; provenance): First evidence wrapper}
\noindent\textit{Transcription reads.} \texttt{GelfandAdjacentEvidence\{GELFAND-P000826\}\{1\}}\par
\noindent\textit{Critical reading.} \texttt{GelfandAdjacentEvidence\{GELFAND-P000829\}\{1\}}\par
\noindent\textit{Evidence.} SourcePage and the actual attached scan are P000829, handwritten page 12. P000826 is not attached.\par
\subsection*{P000829 (both; faithful error): Simplicity of the period-three example}
\noindent\textit{Transcription reads.} Simple, with no qualification\par
\noindent\textit{Critical reading.} Graded-simple with respect to homogeneous ideals; not simple as an ungraded algebra\par
\noindent\textit{Evidence.} The explicit map identifies \(\mathfrak{sl}_2\) Laurent loops. Evaluation at \(z=1\) gives a proper nonzero ideal; degrees \(-1,0,1\) generate only constant \(\mathfrak{sl}_2\).\par
\subsection*{P000829 (both; ambiguity): Three simple algebras and exhaustion}
\noindent\textit{Transcription reads.} Three simple ones; two constructions exhaust all\par
\noindent\textit{Critical reading.} Incomplete list with missing classification quantifiers; no unrestricted theorem reconstructed\par
\noindent\textit{Evidence.} Only one bracket is displayed, and P826--P828 are unavailable. Kac finite-growth theorems require explicit extra assumptions.\par
\subsection*{P000830 (both; faithful error): Vector-spinor conversion}
\noindent\textit{Transcription reads.} \(P_m=\sigma^m_{\alpha\dot\beta}P_{\alpha\dot\beta}\).\par
\noindent\textit{Critical reading.} With \(P_{\alpha\dot\beta}=\sigma^m_{\alpha\dot\beta}P_m\), inversion uses \(P_m=\tfrac12\bar\sigma_m^{\dot\beta\alpha}P_{\alpha\dot\beta}\).\par
\noindent\textit{Evidence.} The displayed Hermitian matrix fixes the normalization: \(\operatorname{tr}(\bar\sigma_m\sigma^n)=2\delta_m^n\). The scan preserves the unqualified shorthand.\par
\subsection*{P000830-P000831 (both; provenance): Academic-year attribution}
\noindent\textit{Transcription reads.} Previous index: 1977-457.20--25.\par
\noindent\textit{Critical reading.} Dated 19 March 1979: record 458; previous identifiers describe the same source pages.\par
\noindent\textit{Evidence.} The scan P000830 explicitly reads 19/III--79 and names Ogievetskii. Parallel formulas identify duplication, not a repeated talk.\par
\subsection*{P000831 (both; ambiguity): Mixed quadratic superfield term}
\noindent\textit{Transcription reads.} \(\theta\sigma^\alpha\bar\theta h_\alpha(x)\).\par
\noindent\textit{Critical reading.} \(\theta\sigma^m\bar\theta h_m(x)\), with a vector index \(m=0,1,2,3\).\par
\noindent\textit{Evidence.} The handwritten letter resembles alpha or Latin a. Lorentz covariance requires a four-vector index, unlike the two-valued spinor index; preserve the diplomatic ambiguity.\par
\subsection*{P000832 (both; faithful error): Equation (7)}
\noindent\textit{Transcription reads.} Coefficient \(i\) in \([\delta_\lambda,\delta_\mu]x\).\par
\noindent\textit{Critical reading.} Coefficient \(2i\); with the displayed \(\delta\) convention the right side is \(2i(\lambda\sigma\bar\mu-\mu\sigma\bar\lambda)\).\par
\noindent\textit{Evidence.} Visual scan confirms absence of 2. Direct differentiation of (6) gives two equal cross terms; reversing commutator order accounts only for a sign.\par
\subsection*{P000832 (both; faithful error): Sentence before equation (7)}
\noindent\textit{Transcription reads.} Two successive supertranslations give a pure translation.\par
\noindent\textit{Critical reading.} The commutator of two supertranslations gives a pure translation.\par
\noindent\textit{Evidence.} Formula (6) makes the successive odd shifts add to \(\lambda+\mu\), whereas the commutator cancels them.\par
\subsection*{P000835 (both; faithful error): After the two rank-one models}
\noindent\textit{Transcription reads.} All others are formally isomorphic.\par
\noindent\textit{Critical reading.} The nonzero transitive leafwise models are contact or locally conformally symplectic, with local frame changes; arbitrary global and singular structures are not all identical.\par
\noindent\textit{Evidence.} Kirillov 1976, Sec.2 Theorem1 and Secs.3--4, Theorems2--3; nontransitive examples and global line bundles prevent the unqualified conclusion.\par
\subsection*{P000835 (both; faithful error): Classification diagram}
\noindent\textit{Transcription reads.} 10 algebras of Ritt.\par
\noindent\textit{Critical reading.} Thirteen types in total: three finite and ten substitutional types; some are parameter families.\par
\noindent\textit{Evidence.} Cassidy 1979, Ch.IV Sec.2 pp.268--269 explicitly enumerates the thirteen types and credits Ritt 1951.\par
\subsection*{P000835 (both; faithful error): Contact-model vector field and bracket}
\noindent\textit{Transcription reads.} \(Y_i=\partial_{y_i}+x_i\partial_z\), \([f,g]=fZg-gZf+\sum(X_ifY_ig-X_igY_if)\).\par
\noindent\textit{Critical reading.} A consistent convention retaining this bracket is \(Y_i=\partial_{y_i}-x_i\partial_z\).\par
\noindent\textit{Evidence.} Both plus signs are visible. In one pair the literal bracket has cyclic Jacobiator \([x,[y,z]]+[y,[z,x]]+[z,[x,y]]=2\); symbolic expansion with general functions and an independent hand check agree.\par
\subsection*{P000835-P000836 (both; faithful error): Formal-series completion}
\noindent\textit{Transcription reads.} \(F[\![ y]\!]=\bigoplus_{\tau\ge0}F_\tau\).\par
\noindent\textit{Critical reading.} \(F[\![ y]\!]=\prod_{\tau\ge0}F_\tau\) for the formal degree completion.\par
\noindent\textit{Evidence.} The next page explicitly allows \(\sum_{k\ge0}P_k\); a direct sum admits only finitely many nonzero homogeneous components.\par
\subsection*{P000836 (both; faithful error): Theorem2 universal extension assertion}
\noindent\textit{Transcription reads.} All two-coordinate groups are extensions of \(\Pi\) by \(A_h\).\par
\noindent\textit{Critical reading.} This describes only a part of Ritt\'s classification; include the finite types and the \(\Pi\times\Pi\) exception and specify the formal differential category.\par
\noindent\textit{Evidence.} Cassidy 1979, Ch.IV Secs.2--3, especially first finite and second substitutional types. The abelian group alone refutes the unrestricted quotient assertion.\par
\subsection*{P000837 (both; faithful error): Codimension before specialization}
\noindent\textit{Transcription reads.} This is a foliation of codimension1 for \(\dim X=m\).\par
\noindent\textit{Critical reading.} The incidence foliations have codimension \(m\); codimension1 is obtained when \(m=1\), imposed on P000838.\par
\noindent\textit{Evidence.} The Grassmannian has dimension \(m(n-m+1)\); imposing containment of one fixed point reduces dimension by \(m\).\par
\subsection*{P000837 (both; ambiguity): Dated heading and speaker initial}
\noindent\textit{Transcription reads.} 2/IV--78; K. Shabat (both tentative).\par
\noindent\textit{Critical reading.} Retain the diplomatic uncertainty; catalogue record460 gives 2 April1979, G.B.Shabat.\par
\noindent\textit{Evidence.} The scan does not securely settle the last year digit or initial; the catalogue supplies a separate source, not an invisible correction to handwriting.\par
\subsection*{P000837 (both; faithful error): Leaf definition, Eq.(2)}
\noindent\textit{Transcription reads.} \(\{H:\exists i,\ P_i\in H\}\) is a leaf.\par
\noindent\textit{Critical reading.} The \(i\)-th leaf is \(\{H:P_i\in H\}\) locally for fixed \(i\); the existential set is a union of leaves.\par
\noindent\textit{Evidence.} The distinct branches of the incidence covering give distinct foliations; the scan clearly includes the existential quantifier.\par
\subsection*{P000837 (both; ambiguity): Three-web straightening criterion}
\noindent\textit{Transcription reads.} Three foliations can be straightened iff hexagonal.\par
\noindent\textit{Critical reading.} Require local coordinates in which the three foliations are \(x=\mathrm{const}\), \(y=\mathrm{const}\), and \(x+y=\mathrm{const}\). This is stronger than general straight-line foliations.\par
\noindent\textit{Evidence.} Beauville1980 Sec.1 states the stronger precise target; the hexagon diagram supports this interpretation.\par
\subsection*{P000838 (both; faithful error): Castelnuovo interpretation Eq.(6)}
\noindent\textit{Transcription reads.} Maximum genus of a degree-d curve in projective n-space.\par
\noindent\textit{Critical reading.} Maximum genus of a nondegenerate curve (in the smooth irreducible scope stated), with d at least n.\par
\noindent\textit{Evidence.} A plane curve embedded in higher-dimensional projective space violates the smaller claimed bound if nondegeneracy is omitted.\par
\subsection*{P000838 (both; faithful error): First-integral display Eq.(4)}
\noindent\textit{Transcription reads.} \(u_1,\ldots,u_n\).\par
\noindent\textit{Critical reading.} \(u_1,\ldots,u_d\) for a \(d\)-web.\par
\noindent\textit{Evidence.} The adjacent Eq.(5) sums from1 to d; the scan faithfully shows n in Eq.(4).\par
\subsection*{P000839 (both; faithful error): Historical status of main theorem}
\noindent\textit{Transcription reads.} Main theorem attributed without caveat to Chern--Griffiths1978.\par
\noindent\textit{Critical reading.} The original announcement was corrected in1981 to normal webs; the unrestricted theorem was proved by Trepreau, preprint2005/publication2006.\par
\noindent\textit{Evidence.} Primary1981 correction pp.78--79 explicitly identifies both errors and adds normality; Trepreau Theorem1.1 proves the full stated range.\par
\subsection*{P000839-P000840 (both; provenance): Speaker transition at Zamolodchikov heading}
\noindent\textit{Transcription reads.} Catalogue groups these pages under Shabat.\par
\noindent\textit{Critical reading.} A.B.Zamolodchikov intervention starts in the latter part of P000839 and continues through P000840; same-session assignment is reconstructed.\par
\noindent\textit{Evidence.} Explicit handwritten Zamolodchikov A.(Sasha), repeated A.Z. on P000840, and no intervening date.\par
\subsection*{P000840 (both; faithful error): Quadratic logarithmic growth and index independence}
\noindent\textit{Transcription reads.} \(\log Z_i^j\sim L^2\), independent of boundary indices.\par
\noindent\textit{Critical reading.} This needs specified nondegenerate weights and thermodynamic/boundary hypotheses; it is false for arbitrary matrices.\par
\noindent\textit{Evidence.} With identity crossing matrices compatible entries are1, incompatible entries0, and the trace is \(N^L\). This example satisfies the three-line identity.\par
\subsection*{P000841-P000842 (both; provenance): Extent assigned to record460}
\noindent\textit{Transcription reads.} Pages included in the Shabat catalogue extent; P000842 bears handwritten66.\par
\noindent\textit{Critical reading.} P000841 starts R.E.Kallosh,24September1979; P000842 continues that talk and belongs to1979--80.\par
\noindent\textit{Evidence.} Date and speaker heading are explicit on P000841. Formula labels(1)--(2) continue as(3)--(5), with the same A,B,S,T,zero/nonzero modes and supergravity discussion. The66 is a notebook page number, not evidence of Belavin affiliation.\par
